\documentclass[11pt,a4paper]{article}
\usepackage[utf8]{inputenc}
\usepackage[T1]{fontenc}
\usepackage[french]{babel}
\usepackage{amsmath, amssymb, amsthm}
\usepackage{geometry}
\geometry{margin=2.5cm}
\usepackage{listings}
\usepackage{hyperref}

\newtheorem{theorem}{Théorème}
\newtheorem{lemma}{Lemme}
\newtheorem{definition}{Définition}

\title{Sur la Conjecture Diophantienne d'Erd\H{o}s-Sierpi\'{n}ski}
\author{Charles EDOU NZE\thanks{Chercheur indépendant / Independent Researcher}}
\date{\today}

\begin{document}

\maketitle

\begin{abstract}
Ce document présente une analyse formelle rigoureuse et une démonstration constructive partielle de la conjecture d'Erd\H{o}s-Sierpi\'{n}ski concernant les fractions égyptiennes. Nous y détaillons l'axiomatisation, la décomposition en classes de congruences, et la génération de solutions explicites pour les premières valeurs entières.
\end{abstract}

\section{Analyse et Décomposition}

\begin{definition}
Soit $n \in \mathbb{N}$ avec $n \geq 2$. Le prédicat d'Erd\H{o}s-Sierpi\'{n}ski $P(n)$ est défini par :
$$ P(n) \iff \exists x, y, z \in \mathbb{Z}^{+}, \quad 5xyz = n(xy + yz + zx) $$
\end{definition}

L'équation se reformule en géométrie algébrique comme la recherche de points rationnels sur une surface de Châtelet affine. Les symétries inhérentes au groupe symétrique $S_3$ opérant sur $(x,y,z)$ permettent d'imposer $x \leq y \leq z$ sans perte de généralité.

\section{Littérature Contextuelle et Analogies}

Le problème d'Erd\H{o}s-Sierpi\'{n}ski est la généralisation naturelle de l'équation d'Erd\H{o}s-Straus ($4/n = 1/x + 1/y + 1/z$). Les travaux asymptotiques de Vaughan (1970) et de Tao sur les densités de fractions égyptiennes fournissent un cadre analytique. L'analogie directe s'établit avec le crible de grand crible et la méthode du cercle de Hardy-Littlewood, qui bornent le nombre d'exceptions potentielles.

\section{Architecture d'Autoformalisation (Lean 4)}

Le squelette de preuve formelle est donné ci-dessous.

\begin{lstlisting}[language=Caml]
import Mathlib.Data.Nat.Basic
import Mathlib.Data.Nat.Parity
import Mathlib.Tactic.Ring
import Mathlib.Tactic.Linarith
import Mathlib.Tactic.Omega

-- Definition stricte du predicat de Sierpinski
def ErdosSierpinskiPredicate (n : Nat) : Prop :=
  exists x y z : Nat, x > 0 /\ y > 0 /\ z > 0 /\ 5 * x * y * z = n * (x * y + y * z + z * x)

-- Theoreme principal
theorem erdos_sierpinski_conjecture : forall n : Nat, n >= 2 -> ErdosSierpinskiPredicate n := by
  -- Il s'agit d'une esquisse de preuve incomplete destinee a une autoformalisation future.
  intro n hn
  sorry

-- Lemme de reduction modulo 5
lemma erdos_sierpinski_mod5_0 (k : Nat) (hk : k >= 1) : ErdosSierpinskiPredicate (5 * k) := by
  unfold ErdosSierpinskiPredicate
  use 2 * k, 3 * k, 6 * k
  exact And.intro (by omega) (And.intro (by omega) (And.intro (by omega) (by ring_nf)))

lemma erdos_sierpinski_constructive (n x y z : Nat) (hx : x > 0) (hy : y > 0) (hz : z > 0)
  (h1 : 5*x*y*z = n*(x*y + y*z + z*x)) : ErdosSierpinskiPredicate n := by
  unfold ErdosSierpinskiPredicate
  use x, y, z
  exact ⟨hx, hy, hz, h1⟩
\end{lstlisting}

\section{Lemmes Stratégiques et Preuve Informelle}

\subsection{Lemme 1 : Facteurs de congruence triviaux}

\begin{lemma}
Pour tout entier $k \in \mathbb{Z}^{+}$, l'équation $\frac{5}{n} = \frac{1}{x} + \frac{1}{y} + \frac{1}{z}$ admet une solution triviale pour $n = 5k$.
\end{lemma}

\begin{proof}
Soit $n = 5k$. Nous substituons cette valeur dans l'équation initiale, ce qui donne $\frac{5}{5k} = \frac{1}{k}$.
L'objectif est de trouver trois entiers $x, y, z$ strictement positifs tels que $\frac{1}{x} + \frac{1}{y} + \frac{1}{z} = \frac{1}{k}$.
Nous utilisons l'identité classique de la somme de fractions unitaires : $\frac{1}{k} = \frac{1}{2k} + \frac{1}{3k} + \frac{1}{6k}$.
Nous posons explicitement $x = 2k$, $y = 3k$, et $z = 6k$.
Évaluons la somme $\frac{1}{2k} + \frac{1}{3k} + \frac{1}{6k}$. Le dénominateur commun est $6k$.
$$ \frac{1}{2k} = \frac{3}{6k} $$
$$ \frac{1}{3k} = \frac{2}{6k} $$
$$ \frac{1}{6k} = \frac{1}{6k} $$
En sommant ces trois termes, nous obtenons :
$$ \frac{3 + 2 + 1}{6k} = \frac{6}{6k} = \frac{1}{k} $$
Puisque $k \geq 1$, les entiers $2k$, $3k$ et $6k$ sont strictement positifs. La solution est donc exacte et le lemme est démontré.
\end{proof}

\subsection{Lemme 2 : Construction algorithmique pour les classes complexes}

\begin{lemma}
S'il existe un entier $x \in [\lceil n/5 \rceil, \infty)$ tel que la fraction $\frac{5x-n}{nx}$ puisse être décomposée en deux fractions unitaires, alors $P(n)$ est satisfait.
\end{lemma}

\begin{proof}
L'équation initiale s'écrit $\frac{5}{n} - \frac{1}{x} = \frac{1}{y} + \frac{1}{z}$.
La mise au même dénominateur du terme de gauche donne $\frac{5x-n}{nx} = \frac{1}{y} + \frac{1}{z}$.
En posant $a = 5x-n$ et $b = nx$, il faut résoudre $a y z = b(y + z)$, ce qui équivaut à $(ay-b)(az-b) = b^2$.
Soit $D$ un diviseur de $b^2$. Si $D \equiv -b \pmod a$ et $\frac{b^2}{D} \equiv -b \pmod a$, alors les entiers $y = \frac{b+D}{a}$ et $z = \frac{b + b^2/D}{a}$ sont des entiers strictement positifs. Par construction, ces entiers valident l'équation.
\end{proof}

\section{Démonstrations Constructives Explicites pour les Entiers Initiaux}

\subsection{Démonstration pour $n = 2$}
Soit $n = 2$. Nous cherchons $x, y, z \in \mathbb{Z}^{+}$ tels que $\frac{5}{2} = \frac{1}{x} + \frac{1}{y} + \frac{1}{z}$.
Posons $x = 1$, $y = 1$, $z = 2$.
Les conditions $x > 0$, $y > 0$ et $z > 0$ sont satisfaites car $x, y, z \geq 1$.
Le PPCM des dénominateurs est $\text{PPCM}(1, 1, 2) = 2$.
En réduisant au même dénominateur :
\begin{itemize}
    \item $\frac{1}{1} = \frac{2}{2}$
    \item $\frac{1}{1} = \frac{2}{2}$
    \item $\frac{1}{2} = \frac{1}{2}$
\end{itemize}
La somme des numérateurs est :
$$ \frac{1}{1} + \frac{1}{1} + \frac{1}{2} = \frac{2 + 2 + 1}{2} = \frac{5}{2} $$
Le PGCD du numérateur et du dénominateur est $\text{PGCD}(5, 2) = 1$.
La fraction irréductible est :
$$ \frac{5}{2} = \frac{5 \div 1}{2 \div 1} = \frac{5}{2} $$
Cette fraction est bien égale à $\frac{5}{2}$.

\subsection{Démonstration pour $n = 3$}
Soit $n = 3$. Nous cherchons $x, y, z \in \mathbb{Z}^{+}$ tels que $\frac{5}{3} = \frac{1}{x} + \frac{1}{y} + \frac{1}{z}$.
Posons $x = 1$, $y = 2$, $z = 6$.
Les conditions $x > 0$, $y > 0$ et $z > 0$ sont satisfaites car $x, y, z \geq 1$.
Le PPCM des dénominateurs est $\text{PPCM}(1, 2, 6) = 6$.
En réduisant au même dénominateur :
\begin{itemize}
    \item $\frac{1}{1} = \frac{6}{6}$
    \item $\frac{1}{2} = \frac{3}{6}$
    \item $\frac{1}{6} = \frac{1}{6}$
\end{itemize}
La somme des numérateurs est :
$$ \frac{1}{1} + \frac{1}{2} + \frac{1}{6} = \frac{6 + 3 + 1}{6} = \frac{10}{6} $$
Le PGCD du numérateur et du dénominateur est $\text{PGCD}(10, 6) = 2$.
La fraction irréductible est :
$$ \frac{10}{6} = \frac{10 \div 2}{6 \div 2} = \frac{5}{3} $$
Cette fraction est bien égale à $\frac{5}{3}$.

\subsection{Démonstration pour $n = 4$}
Soit $n = 4$. Nous cherchons $x, y, z \in \mathbb{Z}^{+}$ tels que $\frac{5}{4} = \frac{1}{x} + \frac{1}{y} + \frac{1}{z}$.
Posons $x = 1$, $y = 5$, $z = 20$.
Les conditions $x > 0$, $y > 0$ et $z > 0$ sont satisfaites car $x, y, z \geq 1$.
Le PPCM des dénominateurs est $\text{PPCM}(1, 5, 20) = 20$.
En réduisant au même dénominateur :
\begin{itemize}
    \item $\frac{1}{1} = \frac{20}{20}$
    \item $\frac{1}{5} = \frac{4}{20}$
    \item $\frac{1}{20} = \frac{1}{20}$
\end{itemize}
La somme des numérateurs est :
$$ \frac{1}{1} + \frac{1}{5} + \frac{1}{20} = \frac{20 + 4 + 1}{20} = \frac{25}{20} $$
Le PGCD du numérateur et du dénominateur est $\text{PGCD}(25, 20) = 5$.
La fraction irréductible est :
$$ \frac{25}{20} = \frac{25 \div 5}{20 \div 5} = \frac{5}{4} $$
Cette fraction est bien égale à $\frac{5}{4}$.

\subsection{Démonstration pour $n = 5$}
Soit $n = 5$. Nous cherchons $x, y, z \in \mathbb{Z}^{+}$ tels que $\frac{5}{5} = \frac{1}{x} + \frac{1}{y} + \frac{1}{z}$.
Posons $x = 2$, $y = 3$, $z = 6$.
Les conditions $x > 0$, $y > 0$ et $z > 0$ sont satisfaites car $x, y, z \geq 1$.
Le PPCM des dénominateurs est $\text{PPCM}(2, 3, 6) = 6$.
En réduisant au même dénominateur :
\begin{itemize}
    \item $\frac{1}{2} = \frac{3}{6}$
    \item $\frac{1}{3} = \frac{2}{6}$
    \item $\frac{1}{6} = \frac{1}{6}$
\end{itemize}
La somme des numérateurs est :
$$ \frac{1}{2} + \frac{1}{3} + \frac{1}{6} = \frac{3 + 2 + 1}{6} = \frac{6}{6} $$
Le PGCD du numérateur et du dénominateur est $\text{PGCD}(6, 6) = 6$.
La fraction irréductible est :
$$ \frac{6}{6} = \frac{6 \div 6}{6 \div 6} = \frac{1}{1} $$
Cette fraction est bien égale à $\frac{5}{5}$.

\subsection{Démonstration pour $n = 6$}
Soit $n = 6$. Nous cherchons $x, y, z \in \mathbb{Z}^{+}$ tels que $\frac{5}{6} = \frac{1}{x} + \frac{1}{y} + \frac{1}{z}$.
Posons $x = 2$, $y = 4$, $z = 12$.
Les conditions $x > 0$, $y > 0$ et $z > 0$ sont satisfaites car $x, y, z \geq 1$.
Le PPCM des dénominateurs est $\text{PPCM}(2, 4, 12) = 12$.
En réduisant au même dénominateur :
\begin{itemize}
    \item $\frac{1}{2} = \frac{6}{12}$
    \item $\frac{1}{4} = \frac{3}{12}$
    \item $\frac{1}{12} = \frac{1}{12}$
\end{itemize}
La somme des numérateurs est :
$$ \frac{1}{2} + \frac{1}{4} + \frac{1}{12} = \frac{6 + 3 + 1}{12} = \frac{10}{12} $$
Le PGCD du numérateur et du dénominateur est $\text{PGCD}(10, 12) = 2$.
La fraction irréductible est :
$$ \frac{10}{12} = \frac{10 \div 2}{12 \div 2} = \frac{5}{6} $$
Cette fraction est bien égale à $\frac{5}{6}$.

\subsection{Démonstration pour $n = 7$}
Soit $n = 7$. Nous cherchons $x, y, z \in \mathbb{Z}^{+}$ tels que $\frac{5}{7} = \frac{1}{x} + \frac{1}{y} + \frac{1}{z}$.
Posons $x = 2$, $y = 5$, $z = 70$.
Les conditions $x > 0$, $y > 0$ et $z > 0$ sont satisfaites car $x, y, z \geq 1$.
Le PPCM des dénominateurs est $\text{PPCM}(2, 5, 70) = 70$.
En réduisant au même dénominateur :
\begin{itemize}
    \item $\frac{1}{2} = \frac{35}{70}$
    \item $\frac{1}{5} = \frac{14}{70}$
    \item $\frac{1}{70} = \frac{1}{70}$
\end{itemize}
La somme des numérateurs est :
$$ \frac{1}{2} + \frac{1}{5} + \frac{1}{70} = \frac{35 + 14 + 1}{70} = \frac{50}{70} $$
Le PGCD du numérateur et du dénominateur est $\text{PGCD}(50, 70) = 10$.
La fraction irréductible est :
$$ \frac{50}{70} = \frac{50 \div 10}{70 \div 10} = \frac{5}{7} $$
Cette fraction est bien égale à $\frac{5}{7}$.

\subsection{Démonstration pour $n = 8$}
Soit $n = 8$. Nous cherchons $x, y, z \in \mathbb{Z}^{+}$ tels que $\frac{5}{8} = \frac{1}{x} + \frac{1}{y} + \frac{1}{z}$.
Posons $x = 2$, $y = 9$, $z = 72$.
Les conditions $x > 0$, $y > 0$ et $z > 0$ sont satisfaites car $x, y, z \geq 1$.
Le PPCM des dénominateurs est $\text{PPCM}(2, 9, 72) = 72$.
En réduisant au même dénominateur :
\begin{itemize}
    \item $\frac{1}{2} = \frac{36}{72}$
    \item $\frac{1}{9} = \frac{8}{72}$
    \item $\frac{1}{72} = \frac{1}{72}$
\end{itemize}
La somme des numérateurs est :
$$ \frac{1}{2} + \frac{1}{9} + \frac{1}{72} = \frac{36 + 8 + 1}{72} = \frac{45}{72} $$
Le PGCD du numérateur et du dénominateur est $\text{PGCD}(45, 72) = 9$.
La fraction irréductible est :
$$ \frac{45}{72} = \frac{45 \div 9}{72 \div 9} = \frac{5}{8} $$
Cette fraction est bien égale à $\frac{5}{8}$.

\subsection{Démonstration pour $n = 9$}
Soit $n = 9$. Nous cherchons $x, y, z \in \mathbb{Z}^{+}$ tels que $\frac{5}{9} = \frac{1}{x} + \frac{1}{y} + \frac{1}{z}$.
Posons $x = 2$, $y = 19$, $z = 342$.
Les conditions $x > 0$, $y > 0$ et $z > 0$ sont satisfaites car $x, y, z \geq 1$.
Le PPCM des dénominateurs est $\text{PPCM}(2, 19, 342) = 342$.
En réduisant au même dénominateur :
\begin{itemize}
    \item $\frac{1}{2} = \frac{171}{342}$
    \item $\frac{1}{19} = \frac{18}{342}$
    \item $\frac{1}{342} = \frac{1}{342}$
\end{itemize}
La somme des numérateurs est :
$$ \frac{1}{2} + \frac{1}{19} + \frac{1}{342} = \frac{171 + 18 + 1}{342} = \frac{190}{342} $$
Le PGCD du numérateur et du dénominateur est $\text{PGCD}(190, 342) = 38$.
La fraction irréductible est :
$$ \frac{190}{342} = \frac{190 \div 38}{342 \div 38} = \frac{5}{9} $$
Cette fraction est bien égale à $\frac{5}{9}$.

\subsection{Démonstration pour $n = 10$}
Soit $n = 10$. Nous cherchons $x, y, z \in \mathbb{Z}^{+}$ tels que $\frac{5}{10} = \frac{1}{x} + \frac{1}{y} + \frac{1}{z}$.
Posons $x = 3$, $y = 7$, $z = 42$.
Les conditions $x > 0$, $y > 0$ et $z > 0$ sont satisfaites car $x, y, z \geq 1$.
Le PPCM des dénominateurs est $\text{PPCM}(3, 7, 42) = 42$.
En réduisant au même dénominateur :
\begin{itemize}
    \item $\frac{1}{3} = \frac{14}{42}$
    \item $\frac{1}{7} = \frac{6}{42}$
    \item $\frac{1}{42} = \frac{1}{42}$
\end{itemize}
La somme des numérateurs est :
$$ \frac{1}{3} + \frac{1}{7} + \frac{1}{42} = \frac{14 + 6 + 1}{42} = \frac{21}{42} $$
Le PGCD du numérateur et du dénominateur est $\text{PGCD}(21, 42) = 21$.
La fraction irréductible est :
$$ \frac{21}{42} = \frac{21 \div 21}{42 \div 21} = \frac{1}{2} $$
Cette fraction est bien égale à $\frac{5}{10}$.

\subsection{Démonstration pour $n = 11$}
Soit $n = 11$. Nous cherchons $x, y, z \in \mathbb{Z}^{+}$ tels que $\frac{5}{11} = \frac{1}{x} + \frac{1}{y} + \frac{1}{z}$.
Posons $x = 3$, $y = 9$, $z = 99$.
Les conditions $x > 0$, $y > 0$ et $z > 0$ sont satisfaites car $x, y, z \geq 1$.
Le PPCM des dénominateurs est $\text{PPCM}(3, 9, 99) = 99$.
En réduisant au même dénominateur :
\begin{itemize}
    \item $\frac{1}{3} = \frac{33}{99}$
    \item $\frac{1}{9} = \frac{11}{99}$
    \item $\frac{1}{99} = \frac{1}{99}$
\end{itemize}
La somme des numérateurs est :
$$ \frac{1}{3} + \frac{1}{9} + \frac{1}{99} = \frac{33 + 11 + 1}{99} = \frac{45}{99} $$
Le PGCD du numérateur et du dénominateur est $\text{PGCD}(45, 99) = 9$.
La fraction irréductible est :
$$ \frac{45}{99} = \frac{45 \div 9}{99 \div 9} = \frac{5}{11} $$
Cette fraction est bien égale à $\frac{5}{11}$.

\subsection{Démonstration pour $n = 12$}
Soit $n = 12$. Nous cherchons $x, y, z \in \mathbb{Z}^{+}$ tels que $\frac{5}{12} = \frac{1}{x} + \frac{1}{y} + \frac{1}{z}$.
Posons $x = 3$, $y = 13$, $z = 156$.
Les conditions $x > 0$, $y > 0$ et $z > 0$ sont satisfaites car $x, y, z \geq 1$.
Le PPCM des dénominateurs est $\text{PPCM}(3, 13, 156) = 156$.
En réduisant au même dénominateur :
\begin{itemize}
    \item $\frac{1}{3} = \frac{52}{156}$
    \item $\frac{1}{13} = \frac{12}{156}$
    \item $\frac{1}{156} = \frac{1}{156}$
\end{itemize}
La somme des numérateurs est :
$$ \frac{1}{3} + \frac{1}{13} + \frac{1}{156} = \frac{52 + 12 + 1}{156} = \frac{65}{156} $$
Le PGCD du numérateur et du dénominateur est $\text{PGCD}(65, 156) = 13$.
La fraction irréductible est :
$$ \frac{65}{156} = \frac{65 \div 13}{156 \div 13} = \frac{5}{12} $$
Cette fraction est bien égale à $\frac{5}{12}$.

\subsection{Démonstration pour $n = 13$}
Soit $n = 13$. Nous cherchons $x, y, z \in \mathbb{Z}^{+}$ tels que $\frac{5}{13} = \frac{1}{x} + \frac{1}{y} + \frac{1}{z}$.
Posons $x = 3$, $y = 20$, $z = 780$.
Les conditions $x > 0$, $y > 0$ et $z > 0$ sont satisfaites car $x, y, z \geq 1$.
Le PPCM des dénominateurs est $\text{PPCM}(3, 20, 780) = 780$.
En réduisant au même dénominateur :
\begin{itemize}
    \item $\frac{1}{3} = \frac{260}{780}$
    \item $\frac{1}{20} = \frac{39}{780}$
    \item $\frac{1}{780} = \frac{1}{780}$
\end{itemize}
La somme des numérateurs est :
$$ \frac{1}{3} + \frac{1}{20} + \frac{1}{780} = \frac{260 + 39 + 1}{780} = \frac{300}{780} $$
Le PGCD du numérateur et du dénominateur est $\text{PGCD}(300, 780) = 60$.
La fraction irréductible est :
$$ \frac{300}{780} = \frac{300 \div 60}{780 \div 60} = \frac{5}{13} $$
Cette fraction est bien égale à $\frac{5}{13}$.

\subsection{Démonstration pour $n = 14$}
Soit $n = 14$. Nous cherchons $x, y, z \in \mathbb{Z}^{+}$ tels que $\frac{5}{14} = \frac{1}{x} + \frac{1}{y} + \frac{1}{z}$.
Posons $x = 3$, $y = 43$, $z = 1806$.
Les conditions $x > 0$, $y > 0$ et $z > 0$ sont satisfaites car $x, y, z \geq 1$.
Le PPCM des dénominateurs est $\text{PPCM}(3, 43, 1806) = 1806$.
En réduisant au même dénominateur :
\begin{itemize}
    \item $\frac{1}{3} = \frac{602}{1806}$
    \item $\frac{1}{43} = \frac{42}{1806}$
    \item $\frac{1}{1806} = \frac{1}{1806}$
\end{itemize}
La somme des numérateurs est :
$$ \frac{1}{3} + \frac{1}{43} + \frac{1}{1806} = \frac{602 + 42 + 1}{1806} = \frac{645}{1806} $$
Le PGCD du numérateur et du dénominateur est $\text{PGCD}(645, 1806) = 129$.
La fraction irréductible est :
$$ \frac{645}{1806} = \frac{645 \div 129}{1806 \div 129} = \frac{5}{14} $$
Cette fraction est bien égale à $\frac{5}{14}$.

\subsection{Démonstration pour $n = 15$}
Soit $n = 15$. Nous cherchons $x, y, z \in \mathbb{Z}^{+}$ tels que $\frac{5}{15} = \frac{1}{x} + \frac{1}{y} + \frac{1}{z}$.
Posons $x = 4$, $y = 13$, $z = 156$.
Les conditions $x > 0$, $y > 0$ et $z > 0$ sont satisfaites car $x, y, z \geq 1$.
Le PPCM des dénominateurs est $\text{PPCM}(4, 13, 156) = 156$.
En réduisant au même dénominateur :
\begin{itemize}
    \item $\frac{1}{4} = \frac{39}{156}$
    \item $\frac{1}{13} = \frac{12}{156}$
    \item $\frac{1}{156} = \frac{1}{156}$
\end{itemize}
La somme des numérateurs est :
$$ \frac{1}{4} + \frac{1}{13} + \frac{1}{156} = \frac{39 + 12 + 1}{156} = \frac{52}{156} $$
Le PGCD du numérateur et du dénominateur est $\text{PGCD}(52, 156) = 52$.
La fraction irréductible est :
$$ \frac{52}{156} = \frac{52 \div 52}{156 \div 52} = \frac{1}{3} $$
Cette fraction est bien égale à $\frac{5}{15}$.

\subsection{Démonstration pour $n = 16$}
Soit $n = 16$. Nous cherchons $x, y, z \in \mathbb{Z}^{+}$ tels que $\frac{5}{16} = \frac{1}{x} + \frac{1}{y} + \frac{1}{z}$.
Posons $x = 4$, $y = 17$, $z = 272$.
Les conditions $x > 0$, $y > 0$ et $z > 0$ sont satisfaites car $x, y, z \geq 1$.
Le PPCM des dénominateurs est $\text{PPCM}(4, 17, 272) = 272$.
En réduisant au même dénominateur :
\begin{itemize}
    \item $\frac{1}{4} = \frac{68}{272}$
    \item $\frac{1}{17} = \frac{16}{272}$
    \item $\frac{1}{272} = \frac{1}{272}$
\end{itemize}
La somme des numérateurs est :
$$ \frac{1}{4} + \frac{1}{17} + \frac{1}{272} = \frac{68 + 16 + 1}{272} = \frac{85}{272} $$
Le PGCD du numérateur et du dénominateur est $\text{PGCD}(85, 272) = 17$.
La fraction irréductible est :
$$ \frac{85}{272} = \frac{85 \div 17}{272 \div 17} = \frac{5}{16} $$
Cette fraction est bien égale à $\frac{5}{16}$.

\subsection{Démonstration pour $n = 17$}
Soit $n = 17$. Nous cherchons $x, y, z \in \mathbb{Z}^{+}$ tels que $\frac{5}{17} = \frac{1}{x} + \frac{1}{y} + \frac{1}{z}$.
Posons $x = 4$, $y = 23$, $z = 1564$.
Les conditions $x > 0$, $y > 0$ et $z > 0$ sont satisfaites car $x, y, z \geq 1$.
Le PPCM des dénominateurs est $\text{PPCM}(4, 23, 1564) = 1564$.
En réduisant au même dénominateur :
\begin{itemize}
    \item $\frac{1}{4} = \frac{391}{1564}$
    \item $\frac{1}{23} = \frac{68}{1564}$
    \item $\frac{1}{1564} = \frac{1}{1564}$
\end{itemize}
La somme des numérateurs est :
$$ \frac{1}{4} + \frac{1}{23} + \frac{1}{1564} = \frac{391 + 68 + 1}{1564} = \frac{460}{1564} $$
Le PGCD du numérateur et du dénominateur est $\text{PGCD}(460, 1564) = 92$.
La fraction irréductible est :
$$ \frac{460}{1564} = \frac{460 \div 92}{1564 \div 92} = \frac{5}{17} $$
Cette fraction est bien égale à $\frac{5}{17}$.

\subsection{Démonstration pour $n = 18$}
Soit $n = 18$. Nous cherchons $x, y, z \in \mathbb{Z}^{+}$ tels que $\frac{5}{18} = \frac{1}{x} + \frac{1}{y} + \frac{1}{z}$.
Posons $x = 4$, $y = 37$, $z = 1332$.
Les conditions $x > 0$, $y > 0$ et $z > 0$ sont satisfaites car $x, y, z \geq 1$.
Le PPCM des dénominateurs est $\text{PPCM}(4, 37, 1332) = 1332$.
En réduisant au même dénominateur :
\begin{itemize}
    \item $\frac{1}{4} = \frac{333}{1332}$
    \item $\frac{1}{37} = \frac{36}{1332}$
    \item $\frac{1}{1332} = \frac{1}{1332}$
\end{itemize}
La somme des numérateurs est :
$$ \frac{1}{4} + \frac{1}{37} + \frac{1}{1332} = \frac{333 + 36 + 1}{1332} = \frac{370}{1332} $$
Le PGCD du numérateur et du dénominateur est $\text{PGCD}(370, 1332) = 74$.
La fraction irréductible est :
$$ \frac{370}{1332} = \frac{370 \div 74}{1332 \div 74} = \frac{5}{18} $$
Cette fraction est bien égale à $\frac{5}{18}$.

\subsection{Démonstration pour $n = 19$}
Soit $n = 19$. Nous cherchons $x, y, z \in \mathbb{Z}^{+}$ tels que $\frac{5}{19} = \frac{1}{x} + \frac{1}{y} + \frac{1}{z}$.
Posons $x = 4$, $y = 77$, $z = 5852$.
Les conditions $x > 0$, $y > 0$ et $z > 0$ sont satisfaites car $x, y, z \geq 1$.
Le PPCM des dénominateurs est $\text{PPCM}(4, 77, 5852) = 5852$.
En réduisant au même dénominateur :
\begin{itemize}
    \item $\frac{1}{4} = \frac{1463}{5852}$
    \item $\frac{1}{77} = \frac{76}{5852}$
    \item $\frac{1}{5852} = \frac{1}{5852}$
\end{itemize}
La somme des numérateurs est :
$$ \frac{1}{4} + \frac{1}{77} + \frac{1}{5852} = \frac{1463 + 76 + 1}{5852} = \frac{1540}{5852} $$
Le PGCD du numérateur et du dénominateur est $\text{PGCD}(1540, 5852) = 308$.
La fraction irréductible est :
$$ \frac{1540}{5852} = \frac{1540 \div 308}{5852 \div 308} = \frac{5}{19} $$
Cette fraction est bien égale à $\frac{5}{19}$.

\subsection{Démonstration pour $n = 20$}
Soit $n = 20$. Nous cherchons $x, y, z \in \mathbb{Z}^{+}$ tels que $\frac{5}{20} = \frac{1}{x} + \frac{1}{y} + \frac{1}{z}$.
Posons $x = 5$, $y = 21$, $z = 420$.
Les conditions $x > 0$, $y > 0$ et $z > 0$ sont satisfaites car $x, y, z \geq 1$.
Le PPCM des dénominateurs est $\text{PPCM}(5, 21, 420) = 420$.
En réduisant au même dénominateur :
\begin{itemize}
    \item $\frac{1}{5} = \frac{84}{420}$
    \item $\frac{1}{21} = \frac{20}{420}$
    \item $\frac{1}{420} = \frac{1}{420}$
\end{itemize}
La somme des numérateurs est :
$$ \frac{1}{5} + \frac{1}{21} + \frac{1}{420} = \frac{84 + 20 + 1}{420} = \frac{105}{420} $$
Le PGCD du numérateur et du dénominateur est $\text{PGCD}(105, 420) = 105$.
La fraction irréductible est :
$$ \frac{105}{420} = \frac{105 \div 105}{420 \div 105} = \frac{1}{4} $$
Cette fraction est bien égale à $\frac{5}{20}$.

\subsection{Démonstration pour $n = 21$}
Soit $n = 21$. Nous cherchons $x, y, z \in \mathbb{Z}^{+}$ tels que $\frac{5}{21} = \frac{1}{x} + \frac{1}{y} + \frac{1}{z}$.
Posons $x = 5$, $y = 27$, $z = 945$.
Les conditions $x > 0$, $y > 0$ et $z > 0$ sont satisfaites car $x, y, z \geq 1$.
Le PPCM des dénominateurs est $\text{PPCM}(5, 27, 945) = 945$.
En réduisant au même dénominateur :
\begin{itemize}
    \item $\frac{1}{5} = \frac{189}{945}$
    \item $\frac{1}{27} = \frac{35}{945}$
    \item $\frac{1}{945} = \frac{1}{945}$
\end{itemize}
La somme des numérateurs est :
$$ \frac{1}{5} + \frac{1}{27} + \frac{1}{945} = \frac{189 + 35 + 1}{945} = \frac{225}{945} $$
Le PGCD du numérateur et du dénominateur est $\text{PGCD}(225, 945) = 45$.
La fraction irréductible est :
$$ \frac{225}{945} = \frac{225 \div 45}{945 \div 45} = \frac{5}{21} $$
Cette fraction est bien égale à $\frac{5}{21}$.

\subsection{Démonstration pour $n = 22$}
Soit $n = 22$. Nous cherchons $x, y, z \in \mathbb{Z}^{+}$ tels que $\frac{5}{22} = \frac{1}{x} + \frac{1}{y} + \frac{1}{z}$.
Posons $x = 5$, $y = 37$, $z = 4070$.
Les conditions $x > 0$, $y > 0$ et $z > 0$ sont satisfaites car $x, y, z \geq 1$.
Le PPCM des dénominateurs est $\text{PPCM}(5, 37, 4070) = 4070$.
En réduisant au même dénominateur :
\begin{itemize}
    \item $\frac{1}{5} = \frac{814}{4070}$
    \item $\frac{1}{37} = \frac{110}{4070}$
    \item $\frac{1}{4070} = \frac{1}{4070}$
\end{itemize}
La somme des numérateurs est :
$$ \frac{1}{5} + \frac{1}{37} + \frac{1}{4070} = \frac{814 + 110 + 1}{4070} = \frac{925}{4070} $$
Le PGCD du numérateur et du dénominateur est $\text{PGCD}(925, 4070) = 185$.
La fraction irréductible est :
$$ \frac{925}{4070} = \frac{925 \div 185}{4070 \div 185} = \frac{5}{22} $$
Cette fraction est bien égale à $\frac{5}{22}$.

\subsection{Démonstration pour $n = 23$}
Soit $n = 23$. Nous cherchons $x, y, z \in \mathbb{Z}^{+}$ tels que $\frac{5}{23} = \frac{1}{x} + \frac{1}{y} + \frac{1}{z}$.
Posons $x = 5$, $y = 58$, $z = 6670$.
Les conditions $x > 0$, $y > 0$ et $z > 0$ sont satisfaites car $x, y, z \geq 1$.
Le PPCM des dénominateurs est $\text{PPCM}(5, 58, 6670) = 6670$.
En réduisant au même dénominateur :
\begin{itemize}
    \item $\frac{1}{5} = \frac{1334}{6670}$
    \item $\frac{1}{58} = \frac{115}{6670}$
    \item $\frac{1}{6670} = \frac{1}{6670}$
\end{itemize}
La somme des numérateurs est :
$$ \frac{1}{5} + \frac{1}{58} + \frac{1}{6670} = \frac{1334 + 115 + 1}{6670} = \frac{1450}{6670} $$
Le PGCD du numérateur et du dénominateur est $\text{PGCD}(1450, 6670) = 290$.
La fraction irréductible est :
$$ \frac{1450}{6670} = \frac{1450 \div 290}{6670 \div 290} = \frac{5}{23} $$
Cette fraction est bien égale à $\frac{5}{23}$.

\subsection{Démonstration pour $n = 24$}
Soit $n = 24$. Nous cherchons $x, y, z \in \mathbb{Z}^{+}$ tels que $\frac{5}{24} = \frac{1}{x} + \frac{1}{y} + \frac{1}{z}$.
Posons $x = 5$, $y = 121$, $z = 14520$.
Les conditions $x > 0$, $y > 0$ et $z > 0$ sont satisfaites car $x, y, z \geq 1$.
Le PPCM des dénominateurs est $\text{PPCM}(5, 121, 14520) = 14520$.
En réduisant au même dénominateur :
\begin{itemize}
    \item $\frac{1}{5} = \frac{2904}{14520}$
    \item $\frac{1}{121} = \frac{120}{14520}$
    \item $\frac{1}{14520} = \frac{1}{14520}$
\end{itemize}
La somme des numérateurs est :
$$ \frac{1}{5} + \frac{1}{121} + \frac{1}{14520} = \frac{2904 + 120 + 1}{14520} = \frac{3025}{14520} $$
Le PGCD du numérateur et du dénominateur est $\text{PGCD}(3025, 14520) = 605$.
La fraction irréductible est :
$$ \frac{3025}{14520} = \frac{3025 \div 605}{14520 \div 605} = \frac{5}{24} $$
Cette fraction est bien égale à $\frac{5}{24}$.

\subsection{Démonstration pour $n = 25$}
Soit $n = 25$. Nous cherchons $x, y, z \in \mathbb{Z}^{+}$ tels que $\frac{5}{25} = \frac{1}{x} + \frac{1}{y} + \frac{1}{z}$.
Posons $x = 6$, $y = 31$, $z = 930$.
Les conditions $x > 0$, $y > 0$ et $z > 0$ sont satisfaites car $x, y, z \geq 1$.
Le PPCM des dénominateurs est $\text{PPCM}(6, 31, 930) = 930$.
En réduisant au même dénominateur :
\begin{itemize}
    \item $\frac{1}{6} = \frac{155}{930}$
    \item $\frac{1}{31} = \frac{30}{930}$
    \item $\frac{1}{930} = \frac{1}{930}$
\end{itemize}
La somme des numérateurs est :
$$ \frac{1}{6} + \frac{1}{31} + \frac{1}{930} = \frac{155 + 30 + 1}{930} = \frac{186}{930} $$
Le PGCD du numérateur et du dénominateur est $\text{PGCD}(186, 930) = 186$.
La fraction irréductible est :
$$ \frac{186}{930} = \frac{186 \div 186}{930 \div 186} = \frac{1}{5} $$
Cette fraction est bien égale à $\frac{5}{25}$.

\subsection{Démonstration pour $n = 26$}
Soit $n = 26$. Nous cherchons $x, y, z \in \mathbb{Z}^{+}$ tels que $\frac{5}{26} = \frac{1}{x} + \frac{1}{y} + \frac{1}{z}$.
Posons $x = 6$, $y = 40$, $z = 1560$.
Les conditions $x > 0$, $y > 0$ et $z > 0$ sont satisfaites car $x, y, z \geq 1$.
Le PPCM des dénominateurs est $\text{PPCM}(6, 40, 1560) = 1560$.
En réduisant au même dénominateur :
\begin{itemize}
    \item $\frac{1}{6} = \frac{260}{1560}$
    \item $\frac{1}{40} = \frac{39}{1560}$
    \item $\frac{1}{1560} = \frac{1}{1560}$
\end{itemize}
La somme des numérateurs est :
$$ \frac{1}{6} + \frac{1}{40} + \frac{1}{1560} = \frac{260 + 39 + 1}{1560} = \frac{300}{1560} $$
Le PGCD du numérateur et du dénominateur est $\text{PGCD}(300, 1560) = 60$.
La fraction irréductible est :
$$ \frac{300}{1560} = \frac{300 \div 60}{1560 \div 60} = \frac{5}{26} $$
Cette fraction est bien égale à $\frac{5}{26}$.

\subsection{Démonstration pour $n = 27$}
Soit $n = 27$. Nous cherchons $x, y, z \in \mathbb{Z}^{+}$ tels que $\frac{5}{27} = \frac{1}{x} + \frac{1}{y} + \frac{1}{z}$.
Posons $x = 6$, $y = 55$, $z = 2970$.
Les conditions $x > 0$, $y > 0$ et $z > 0$ sont satisfaites car $x, y, z \geq 1$.
Le PPCM des dénominateurs est $\text{PPCM}(6, 55, 2970) = 2970$.
En réduisant au même dénominateur :
\begin{itemize}
    \item $\frac{1}{6} = \frac{495}{2970}$
    \item $\frac{1}{55} = \frac{54}{2970}$
    \item $\frac{1}{2970} = \frac{1}{2970}$
\end{itemize}
La somme des numérateurs est :
$$ \frac{1}{6} + \frac{1}{55} + \frac{1}{2970} = \frac{495 + 54 + 1}{2970} = \frac{550}{2970} $$
Le PGCD du numérateur et du dénominateur est $\text{PGCD}(550, 2970) = 110$.
La fraction irréductible est :
$$ \frac{550}{2970} = \frac{550 \div 110}{2970 \div 110} = \frac{5}{27} $$
Cette fraction est bien égale à $\frac{5}{27}$.

\subsection{Démonstration pour $n = 28$}
Soit $n = 28$. Nous cherchons $x, y, z \in \mathbb{Z}^{+}$ tels que $\frac{5}{28} = \frac{1}{x} + \frac{1}{y} + \frac{1}{z}$.
Posons $x = 6$, $y = 85$, $z = 7140$.
Les conditions $x > 0$, $y > 0$ et $z > 0$ sont satisfaites car $x, y, z \geq 1$.
Le PPCM des dénominateurs est $\text{PPCM}(6, 85, 7140) = 7140$.
En réduisant au même dénominateur :
\begin{itemize}
    \item $\frac{1}{6} = \frac{1190}{7140}$
    \item $\frac{1}{85} = \frac{84}{7140}$
    \item $\frac{1}{7140} = \frac{1}{7140}$
\end{itemize}
La somme des numérateurs est :
$$ \frac{1}{6} + \frac{1}{85} + \frac{1}{7140} = \frac{1190 + 84 + 1}{7140} = \frac{1275}{7140} $$
Le PGCD du numérateur et du dénominateur est $\text{PGCD}(1275, 7140) = 255$.
La fraction irréductible est :
$$ \frac{1275}{7140} = \frac{1275 \div 255}{7140 \div 255} = \frac{5}{28} $$
Cette fraction est bien égale à $\frac{5}{28}$.

\subsection{Démonstration pour $n = 29$}
Soit $n = 29$. Nous cherchons $x, y, z \in \mathbb{Z}^{+}$ tels que $\frac{5}{29} = \frac{1}{x} + \frac{1}{y} + \frac{1}{z}$.
Posons $x = 6$, $y = 175$, $z = 30450$.
Les conditions $x > 0$, $y > 0$ et $z > 0$ sont satisfaites car $x, y, z \geq 1$.
Le PPCM des dénominateurs est $\text{PPCM}(6, 175, 30450) = 30450$.
En réduisant au même dénominateur :
\begin{itemize}
    \item $\frac{1}{6} = \frac{5075}{30450}$
    \item $\frac{1}{175} = \frac{174}{30450}$
    \item $\frac{1}{30450} = \frac{1}{30450}$
\end{itemize}
La somme des numérateurs est :
$$ \frac{1}{6} + \frac{1}{175} + \frac{1}{30450} = \frac{5075 + 174 + 1}{30450} = \frac{5250}{30450} $$
Le PGCD du numérateur et du dénominateur est $\text{PGCD}(5250, 30450) = 1050$.
La fraction irréductible est :
$$ \frac{5250}{30450} = \frac{5250 \div 1050}{30450 \div 1050} = \frac{5}{29} $$
Cette fraction est bien égale à $\frac{5}{29}$.

\subsection{Démonstration pour $n = 30$}
Soit $n = 30$. Nous cherchons $x, y, z \in \mathbb{Z}^{+}$ tels que $\frac{5}{30} = \frac{1}{x} + \frac{1}{y} + \frac{1}{z}$.
Posons $x = 7$, $y = 43$, $z = 1806$.
Les conditions $x > 0$, $y > 0$ et $z > 0$ sont satisfaites car $x, y, z \geq 1$.
Le PPCM des dénominateurs est $\text{PPCM}(7, 43, 1806) = 1806$.
En réduisant au même dénominateur :
\begin{itemize}
    \item $\frac{1}{7} = \frac{258}{1806}$
    \item $\frac{1}{43} = \frac{42}{1806}$
    \item $\frac{1}{1806} = \frac{1}{1806}$
\end{itemize}
La somme des numérateurs est :
$$ \frac{1}{7} + \frac{1}{43} + \frac{1}{1806} = \frac{258 + 42 + 1}{1806} = \frac{301}{1806} $$
Le PGCD du numérateur et du dénominateur est $\text{PGCD}(301, 1806) = 301$.
La fraction irréductible est :
$$ \frac{301}{1806} = \frac{301 \div 301}{1806 \div 301} = \frac{1}{6} $$
Cette fraction est bien égale à $\frac{5}{30}$.

\subsection{Démonstration pour $n = 31$}
Soit $n = 31$. Nous cherchons $x, y, z \in \mathbb{Z}^{+}$ tels que $\frac{5}{31} = \frac{1}{x} + \frac{1}{y} + \frac{1}{z}$.
Posons $x = 7$, $y = 56$, $z = 1736$.
Les conditions $x > 0$, $y > 0$ et $z > 0$ sont satisfaites car $x, y, z \geq 1$.
Le PPCM des dénominateurs est $\text{PPCM}(7, 56, 1736) = 1736$.
En réduisant au même dénominateur :
\begin{itemize}
    \item $\frac{1}{7} = \frac{248}{1736}$
    \item $\frac{1}{56} = \frac{31}{1736}$
    \item $\frac{1}{1736} = \frac{1}{1736}$
\end{itemize}
La somme des numérateurs est :
$$ \frac{1}{7} + \frac{1}{56} + \frac{1}{1736} = \frac{248 + 31 + 1}{1736} = \frac{280}{1736} $$
Le PGCD du numérateur et du dénominateur est $\text{PGCD}(280, 1736) = 56$.
La fraction irréductible est :
$$ \frac{280}{1736} = \frac{280 \div 56}{1736 \div 56} = \frac{5}{31} $$
Cette fraction est bien égale à $\frac{5}{31}$.

\subsection{Démonstration pour $n = 32$}
Soit $n = 32$. Nous cherchons $x, y, z \in \mathbb{Z}^{+}$ tels que $\frac{5}{32} = \frac{1}{x} + \frac{1}{y} + \frac{1}{z}$.
Posons $x = 7$, $y = 75$, $z = 16800$.
Les conditions $x > 0$, $y > 0$ et $z > 0$ sont satisfaites car $x, y, z \geq 1$.
Le PPCM des dénominateurs est $\text{PPCM}(7, 75, 16800) = 16800$.
En réduisant au même dénominateur :
\begin{itemize}
    \item $\frac{1}{7} = \frac{2400}{16800}$
    \item $\frac{1}{75} = \frac{224}{16800}$
    \item $\frac{1}{16800} = \frac{1}{16800}$
\end{itemize}
La somme des numérateurs est :
$$ \frac{1}{7} + \frac{1}{75} + \frac{1}{16800} = \frac{2400 + 224 + 1}{16800} = \frac{2625}{16800} $$
Le PGCD du numérateur et du dénominateur est $\text{PGCD}(2625, 16800) = 525$.
La fraction irréductible est :
$$ \frac{2625}{16800} = \frac{2625 \div 525}{16800 \div 525} = \frac{5}{32} $$
Cette fraction est bien égale à $\frac{5}{32}$.

\subsection{Démonstration pour $n = 33$}
Soit $n = 33$. Nous cherchons $x, y, z \in \mathbb{Z}^{+}$ tels que $\frac{5}{33} = \frac{1}{x} + \frac{1}{y} + \frac{1}{z}$.
Posons $x = 7$, $y = 116$, $z = 26796$.
Les conditions $x > 0$, $y > 0$ et $z > 0$ sont satisfaites car $x, y, z \geq 1$.
Le PPCM des dénominateurs est $\text{PPCM}(7, 116, 26796) = 26796$.
En réduisant au même dénominateur :
\begin{itemize}
    \item $\frac{1}{7} = \frac{3828}{26796}$
    \item $\frac{1}{116} = \frac{231}{26796}$
    \item $\frac{1}{26796} = \frac{1}{26796}$
\end{itemize}
La somme des numérateurs est :
$$ \frac{1}{7} + \frac{1}{116} + \frac{1}{26796} = \frac{3828 + 231 + 1}{26796} = \frac{4060}{26796} $$
Le PGCD du numérateur et du dénominateur est $\text{PGCD}(4060, 26796) = 812$.
La fraction irréductible est :
$$ \frac{4060}{26796} = \frac{4060 \div 812}{26796 \div 812} = \frac{5}{33} $$
Cette fraction est bien égale à $\frac{5}{33}$.

\subsection{Démonstration pour $n = 34$}
Soit $n = 34$. Nous cherchons $x, y, z \in \mathbb{Z}^{+}$ tels que $\frac{5}{34} = \frac{1}{x} + \frac{1}{y} + \frac{1}{z}$.
Posons $x = 7$, $y = 239$, $z = 56882$.
Les conditions $x > 0$, $y > 0$ et $z > 0$ sont satisfaites car $x, y, z \geq 1$.
Le PPCM des dénominateurs est $\text{PPCM}(7, 239, 56882) = 56882$.
En réduisant au même dénominateur :
\begin{itemize}
    \item $\frac{1}{7} = \frac{8126}{56882}$
    \item $\frac{1}{239} = \frac{238}{56882}$
    \item $\frac{1}{56882} = \frac{1}{56882}$
\end{itemize}
La somme des numérateurs est :
$$ \frac{1}{7} + \frac{1}{239} + \frac{1}{56882} = \frac{8126 + 238 + 1}{56882} = \frac{8365}{56882} $$
Le PGCD du numérateur et du dénominateur est $\text{PGCD}(8365, 56882) = 1673$.
La fraction irréductible est :
$$ \frac{8365}{56882} = \frac{8365 \div 1673}{56882 \div 1673} = \frac{5}{34} $$
Cette fraction est bien égale à $\frac{5}{34}$.

\subsection{Démonstration pour $n = 35$}
Soit $n = 35$. Nous cherchons $x, y, z \in \mathbb{Z}^{+}$ tels que $\frac{5}{35} = \frac{1}{x} + \frac{1}{y} + \frac{1}{z}$.
Posons $x = 8$, $y = 57$, $z = 3192$.
Les conditions $x > 0$, $y > 0$ et $z > 0$ sont satisfaites car $x, y, z \geq 1$.
Le PPCM des dénominateurs est $\text{PPCM}(8, 57, 3192) = 3192$.
En réduisant au même dénominateur :
\begin{itemize}
    \item $\frac{1}{8} = \frac{399}{3192}$
    \item $\frac{1}{57} = \frac{56}{3192}$
    \item $\frac{1}{3192} = \frac{1}{3192}$
\end{itemize}
La somme des numérateurs est :
$$ \frac{1}{8} + \frac{1}{57} + \frac{1}{3192} = \frac{399 + 56 + 1}{3192} = \frac{456}{3192} $$
Le PGCD du numérateur et du dénominateur est $\text{PGCD}(456, 3192) = 456$.
La fraction irréductible est :
$$ \frac{456}{3192} = \frac{456 \div 456}{3192 \div 456} = \frac{1}{7} $$
Cette fraction est bien égale à $\frac{5}{35}$.

\subsection{Démonstration pour $n = 36$}
Soit $n = 36$. Nous cherchons $x, y, z \in \mathbb{Z}^{+}$ tels que $\frac{5}{36} = \frac{1}{x} + \frac{1}{y} + \frac{1}{z}$.
Posons $x = 8$, $y = 73$, $z = 5256$.
Les conditions $x > 0$, $y > 0$ et $z > 0$ sont satisfaites car $x, y, z \geq 1$.
Le PPCM des dénominateurs est $\text{PPCM}(8, 73, 5256) = 5256$.
En réduisant au même dénominateur :
\begin{itemize}
    \item $\frac{1}{8} = \frac{657}{5256}$
    \item $\frac{1}{73} = \frac{72}{5256}$
    \item $\frac{1}{5256} = \frac{1}{5256}$
\end{itemize}
La somme des numérateurs est :
$$ \frac{1}{8} + \frac{1}{73} + \frac{1}{5256} = \frac{657 + 72 + 1}{5256} = \frac{730}{5256} $$
Le PGCD du numérateur et du dénominateur est $\text{PGCD}(730, 5256) = 146$.
La fraction irréductible est :
$$ \frac{730}{5256} = \frac{730 \div 146}{5256 \div 146} = \frac{5}{36} $$
Cette fraction est bien égale à $\frac{5}{36}$.

\subsection{Démonstration pour $n = 37$}
Soit $n = 37$. Nous cherchons $x, y, z \in \mathbb{Z}^{+}$ tels que $\frac{5}{37} = \frac{1}{x} + \frac{1}{y} + \frac{1}{z}$.
Posons $x = 8$, $y = 99$, $z = 29304$.
Les conditions $x > 0$, $y > 0$ et $z > 0$ sont satisfaites car $x, y, z \geq 1$.
Le PPCM des dénominateurs est $\text{PPCM}(8, 99, 29304) = 29304$.
En réduisant au même dénominateur :
\begin{itemize}
    \item $\frac{1}{8} = \frac{3663}{29304}$
    \item $\frac{1}{99} = \frac{296}{29304}$
    \item $\frac{1}{29304} = \frac{1}{29304}$
\end{itemize}
La somme des numérateurs est :
$$ \frac{1}{8} + \frac{1}{99} + \frac{1}{29304} = \frac{3663 + 296 + 1}{29304} = \frac{3960}{29304} $$
Le PGCD du numérateur et du dénominateur est $\text{PGCD}(3960, 29304) = 792$.
La fraction irréductible est :
$$ \frac{3960}{29304} = \frac{3960 \div 792}{29304 \div 792} = \frac{5}{37} $$
Cette fraction est bien égale à $\frac{5}{37}$.

\subsection{Démonstration pour $n = 38$}
Soit $n = 38$. Nous cherchons $x, y, z \in \mathbb{Z}^{+}$ tels que $\frac{5}{38} = \frac{1}{x} + \frac{1}{y} + \frac{1}{z}$.
Posons $x = 8$, $y = 153$, $z = 23256$.
Les conditions $x > 0$, $y > 0$ et $z > 0$ sont satisfaites car $x, y, z \geq 1$.
Le PPCM des dénominateurs est $\text{PPCM}(8, 153, 23256) = 23256$.
En réduisant au même dénominateur :
\begin{itemize}
    \item $\frac{1}{8} = \frac{2907}{23256}$
    \item $\frac{1}{153} = \frac{152}{23256}$
    \item $\frac{1}{23256} = \frac{1}{23256}$
\end{itemize}
La somme des numérateurs est :
$$ \frac{1}{8} + \frac{1}{153} + \frac{1}{23256} = \frac{2907 + 152 + 1}{23256} = \frac{3060}{23256} $$
Le PGCD du numérateur et du dénominateur est $\text{PGCD}(3060, 23256) = 612$.
La fraction irréductible est :
$$ \frac{3060}{23256} = \frac{3060 \div 612}{23256 \div 612} = \frac{5}{38} $$
Cette fraction est bien égale à $\frac{5}{38}$.

\subsection{Démonstration pour $n = 39$}
Soit $n = 39$. Nous cherchons $x, y, z \in \mathbb{Z}^{+}$ tels que $\frac{5}{39} = \frac{1}{x} + \frac{1}{y} + \frac{1}{z}$.
Posons $x = 8$, $y = 313$, $z = 97656$.
Les conditions $x > 0$, $y > 0$ et $z > 0$ sont satisfaites car $x, y, z \geq 1$.
Le PPCM des dénominateurs est $\text{PPCM}(8, 313, 97656) = 97656$.
En réduisant au même dénominateur :
\begin{itemize}
    \item $\frac{1}{8} = \frac{12207}{97656}$
    \item $\frac{1}{313} = \frac{312}{97656}$
    \item $\frac{1}{97656} = \frac{1}{97656}$
\end{itemize}
La somme des numérateurs est :
$$ \frac{1}{8} + \frac{1}{313} + \frac{1}{97656} = \frac{12207 + 312 + 1}{97656} = \frac{12520}{97656} $$
Le PGCD du numérateur et du dénominateur est $\text{PGCD}(12520, 97656) = 2504$.
La fraction irréductible est :
$$ \frac{12520}{97656} = \frac{12520 \div 2504}{97656 \div 2504} = \frac{5}{39} $$
Cette fraction est bien égale à $\frac{5}{39}$.

\subsection{Démonstration pour $n = 40$}
Soit $n = 40$. Nous cherchons $x, y, z \in \mathbb{Z}^{+}$ tels que $\frac{5}{40} = \frac{1}{x} + \frac{1}{y} + \frac{1}{z}$.
Posons $x = 9$, $y = 73$, $z = 5256$.
Les conditions $x > 0$, $y > 0$ et $z > 0$ sont satisfaites car $x, y, z \geq 1$.
Le PPCM des dénominateurs est $\text{PPCM}(9, 73, 5256) = 5256$.
En réduisant au même dénominateur :
\begin{itemize}
    \item $\frac{1}{9} = \frac{584}{5256}$
    \item $\frac{1}{73} = \frac{72}{5256}$
    \item $\frac{1}{5256} = \frac{1}{5256}$
\end{itemize}
La somme des numérateurs est :
$$ \frac{1}{9} + \frac{1}{73} + \frac{1}{5256} = \frac{584 + 72 + 1}{5256} = \frac{657}{5256} $$
Le PGCD du numérateur et du dénominateur est $\text{PGCD}(657, 5256) = 657$.
La fraction irréductible est :
$$ \frac{657}{5256} = \frac{657 \div 657}{5256 \div 657} = \frac{1}{8} $$
Cette fraction est bien égale à $\frac{5}{40}$.

\subsection{Démonstration pour $n = 41$}
Soit $n = 41$. Nous cherchons $x, y, z \in \mathbb{Z}^{+}$ tels que $\frac{5}{41} = \frac{1}{x} + \frac{1}{y} + \frac{1}{z}$.
Posons $x = 9$, $y = 93$, $z = 11439$.
Les conditions $x > 0$, $y > 0$ et $z > 0$ sont satisfaites car $x, y, z \geq 1$.
Le PPCM des dénominateurs est $\text{PPCM}(9, 93, 11439) = 11439$.
En réduisant au même dénominateur :
\begin{itemize}
    \item $\frac{1}{9} = \frac{1271}{11439}$
    \item $\frac{1}{93} = \frac{123}{11439}$
    \item $\frac{1}{11439} = \frac{1}{11439}$
\end{itemize}
La somme des numérateurs est :
$$ \frac{1}{9} + \frac{1}{93} + \frac{1}{11439} = \frac{1271 + 123 + 1}{11439} = \frac{1395}{11439} $$
Le PGCD du numérateur et du dénominateur est $\text{PGCD}(1395, 11439) = 279$.
La fraction irréductible est :
$$ \frac{1395}{11439} = \frac{1395 \div 279}{11439 \div 279} = \frac{5}{41} $$
Cette fraction est bien égale à $\frac{5}{41}$.

\subsection{Démonstration pour $n = 42$}
Soit $n = 42$. Nous cherchons $x, y, z \in \mathbb{Z}^{+}$ tels que $\frac{5}{42} = \frac{1}{x} + \frac{1}{y} + \frac{1}{z}$.
Posons $x = 9$, $y = 127$, $z = 16002$.
Les conditions $x > 0$, $y > 0$ et $z > 0$ sont satisfaites car $x, y, z \geq 1$.
Le PPCM des dénominateurs est $\text{PPCM}(9, 127, 16002) = 16002$.
En réduisant au même dénominateur :
\begin{itemize}
    \item $\frac{1}{9} = \frac{1778}{16002}$
    \item $\frac{1}{127} = \frac{126}{16002}$
    \item $\frac{1}{16002} = \frac{1}{16002}$
\end{itemize}
La somme des numérateurs est :
$$ \frac{1}{9} + \frac{1}{127} + \frac{1}{16002} = \frac{1778 + 126 + 1}{16002} = \frac{1905}{16002} $$
Le PGCD du numérateur et du dénominateur est $\text{PGCD}(1905, 16002) = 381$.
La fraction irréductible est :
$$ \frac{1905}{16002} = \frac{1905 \div 381}{16002 \div 381} = \frac{5}{42} $$
Cette fraction est bien égale à $\frac{5}{42}$.

\subsection{Démonstration pour $n = 43$}
Soit $n = 43$. Nous cherchons $x, y, z \in \mathbb{Z}^{+}$ tels que $\frac{5}{43} = \frac{1}{x} + \frac{1}{y} + \frac{1}{z}$.
Posons $x = 9$, $y = 194$, $z = 75078$.
Les conditions $x > 0$, $y > 0$ et $z > 0$ sont satisfaites car $x, y, z \geq 1$.
Le PPCM des dénominateurs est $\text{PPCM}(9, 194, 75078) = 75078$.
En réduisant au même dénominateur :
\begin{itemize}
    \item $\frac{1}{9} = \frac{8342}{75078}$
    \item $\frac{1}{194} = \frac{387}{75078}$
    \item $\frac{1}{75078} = \frac{1}{75078}$
\end{itemize}
La somme des numérateurs est :
$$ \frac{1}{9} + \frac{1}{194} + \frac{1}{75078} = \frac{8342 + 387 + 1}{75078} = \frac{8730}{75078} $$
Le PGCD du numérateur et du dénominateur est $\text{PGCD}(8730, 75078) = 1746$.
La fraction irréductible est :
$$ \frac{8730}{75078} = \frac{8730 \div 1746}{75078 \div 1746} = \frac{5}{43} $$
Cette fraction est bien égale à $\frac{5}{43}$.

\subsection{Démonstration pour $n = 44$}
Soit $n = 44$. Nous cherchons $x, y, z \in \mathbb{Z}^{+}$ tels que $\frac{5}{44} = \frac{1}{x} + \frac{1}{y} + \frac{1}{z}$.
Posons $x = 9$, $y = 397$, $z = 157212$.
Les conditions $x > 0$, $y > 0$ et $z > 0$ sont satisfaites car $x, y, z \geq 1$.
Le PPCM des dénominateurs est $\text{PPCM}(9, 397, 157212) = 157212$.
En réduisant au même dénominateur :
\begin{itemize}
    \item $\frac{1}{9} = \frac{17468}{157212}$
    \item $\frac{1}{397} = \frac{396}{157212}$
    \item $\frac{1}{157212} = \frac{1}{157212}$
\end{itemize}
La somme des numérateurs est :
$$ \frac{1}{9} + \frac{1}{397} + \frac{1}{157212} = \frac{17468 + 396 + 1}{157212} = \frac{17865}{157212} $$
Le PGCD du numérateur et du dénominateur est $\text{PGCD}(17865, 157212) = 3573$.
La fraction irréductible est :
$$ \frac{17865}{157212} = \frac{17865 \div 3573}{157212 \div 3573} = \frac{5}{44} $$
Cette fraction est bien égale à $\frac{5}{44}$.

\subsection{Démonstration pour $n = 45$}
Soit $n = 45$. Nous cherchons $x, y, z \in \mathbb{Z}^{+}$ tels que $\frac{5}{45} = \frac{1}{x} + \frac{1}{y} + \frac{1}{z}$.
Posons $x = 10$, $y = 91$, $z = 8190$.
Les conditions $x > 0$, $y > 0$ et $z > 0$ sont satisfaites car $x, y, z \geq 1$.
Le PPCM des dénominateurs est $\text{PPCM}(10, 91, 8190) = 8190$.
En réduisant au même dénominateur :
\begin{itemize}
    \item $\frac{1}{10} = \frac{819}{8190}$
    \item $\frac{1}{91} = \frac{90}{8190}$
    \item $\frac{1}{8190} = \frac{1}{8190}$
\end{itemize}
La somme des numérateurs est :
$$ \frac{1}{10} + \frac{1}{91} + \frac{1}{8190} = \frac{819 + 90 + 1}{8190} = \frac{910}{8190} $$
Le PGCD du numérateur et du dénominateur est $\text{PGCD}(910, 8190) = 910$.
La fraction irréductible est :
$$ \frac{910}{8190} = \frac{910 \div 910}{8190 \div 910} = \frac{1}{9} $$
Cette fraction est bien égale à $\frac{5}{45}$.

\subsection{Démonstration pour $n = 46$}
Soit $n = 46$. Nous cherchons $x, y, z \in \mathbb{Z}^{+}$ tels que $\frac{5}{46} = \frac{1}{x} + \frac{1}{y} + \frac{1}{z}$.
Posons $x = 10$, $y = 116$, $z = 13340$.
Les conditions $x > 0$, $y > 0$ et $z > 0$ sont satisfaites car $x, y, z \geq 1$.
Le PPCM des dénominateurs est $\text{PPCM}(10, 116, 13340) = 13340$.
En réduisant au même dénominateur :
\begin{itemize}
    \item $\frac{1}{10} = \frac{1334}{13340}$
    \item $\frac{1}{116} = \frac{115}{13340}$
    \item $\frac{1}{13340} = \frac{1}{13340}$
\end{itemize}
La somme des numérateurs est :
$$ \frac{1}{10} + \frac{1}{116} + \frac{1}{13340} = \frac{1334 + 115 + 1}{13340} = \frac{1450}{13340} $$
Le PGCD du numérateur et du dénominateur est $\text{PGCD}(1450, 13340) = 290$.
La fraction irréductible est :
$$ \frac{1450}{13340} = \frac{1450 \div 290}{13340 \div 290} = \frac{5}{46} $$
Cette fraction est bien égale à $\frac{5}{46}$.

\subsection{Démonstration pour $n = 47$}
Soit $n = 47$. Nous cherchons $x, y, z \in \mathbb{Z}^{+}$ tels que $\frac{5}{47} = \frac{1}{x} + \frac{1}{y} + \frac{1}{z}$.
Posons $x = 10$, $y = 157$, $z = 73790$.
Les conditions $x > 0$, $y > 0$ et $z > 0$ sont satisfaites car $x, y, z \geq 1$.
Le PPCM des dénominateurs est $\text{PPCM}(10, 157, 73790) = 73790$.
En réduisant au même dénominateur :
\begin{itemize}
    \item $\frac{1}{10} = \frac{7379}{73790}$
    \item $\frac{1}{157} = \frac{470}{73790}$
    \item $\frac{1}{73790} = \frac{1}{73790}$
\end{itemize}
La somme des numérateurs est :
$$ \frac{1}{10} + \frac{1}{157} + \frac{1}{73790} = \frac{7379 + 470 + 1}{73790} = \frac{7850}{73790} $$
Le PGCD du numérateur et du dénominateur est $\text{PGCD}(7850, 73790) = 1570$.
La fraction irréductible est :
$$ \frac{7850}{73790} = \frac{7850 \div 1570}{73790 \div 1570} = \frac{5}{47} $$
Cette fraction est bien égale à $\frac{5}{47}$.

\subsection{Démonstration pour $n = 48$}
Soit $n = 48$. Nous cherchons $x, y, z \in \mathbb{Z}^{+}$ tels que $\frac{5}{48} = \frac{1}{x} + \frac{1}{y} + \frac{1}{z}$.
Posons $x = 10$, $y = 241$, $z = 57840$.
Les conditions $x > 0$, $y > 0$ et $z > 0$ sont satisfaites car $x, y, z \geq 1$.
Le PPCM des dénominateurs est $\text{PPCM}(10, 241, 57840) = 57840$.
En réduisant au même dénominateur :
\begin{itemize}
    \item $\frac{1}{10} = \frac{5784}{57840}$
    \item $\frac{1}{241} = \frac{240}{57840}$
    \item $\frac{1}{57840} = \frac{1}{57840}$
\end{itemize}
La somme des numérateurs est :
$$ \frac{1}{10} + \frac{1}{241} + \frac{1}{57840} = \frac{5784 + 240 + 1}{57840} = \frac{6025}{57840} $$
Le PGCD du numérateur et du dénominateur est $\text{PGCD}(6025, 57840) = 1205$.
La fraction irréductible est :
$$ \frac{6025}{57840} = \frac{6025 \div 1205}{57840 \div 1205} = \frac{5}{48} $$
Cette fraction est bien égale à $\frac{5}{48}$.

\subsection{Démonstration pour $n = 49$}
Soit $n = 49$. Nous cherchons $x, y, z \in \mathbb{Z}^{+}$ tels que $\frac{5}{49} = \frac{1}{x} + \frac{1}{y} + \frac{1}{z}$.
Posons $x = 10$, $y = 491$, $z = 240590$.
Les conditions $x > 0$, $y > 0$ et $z > 0$ sont satisfaites car $x, y, z \geq 1$.
Le PPCM des dénominateurs est $\text{PPCM}(10, 491, 240590) = 240590$.
En réduisant au même dénominateur :
\begin{itemize}
    \item $\frac{1}{10} = \frac{24059}{240590}$
    \item $\frac{1}{491} = \frac{490}{240590}$
    \item $\frac{1}{240590} = \frac{1}{240590}$
\end{itemize}
La somme des numérateurs est :
$$ \frac{1}{10} + \frac{1}{491} + \frac{1}{240590} = \frac{24059 + 490 + 1}{240590} = \frac{24550}{240590} $$
Le PGCD du numérateur et du dénominateur est $\text{PGCD}(24550, 240590) = 4910$.
La fraction irréductible est :
$$ \frac{24550}{240590} = \frac{24550 \div 4910}{240590 \div 4910} = \frac{5}{49} $$
Cette fraction est bien égale à $\frac{5}{49}$.

\subsection{Démonstration pour $n = 50$}
Soit $n = 50$. Nous cherchons $x, y, z \in \mathbb{Z}^{+}$ tels que $\frac{5}{50} = \frac{1}{x} + \frac{1}{y} + \frac{1}{z}$.
Posons $x = 11$, $y = 111$, $z = 12210$.
Les conditions $x > 0$, $y > 0$ et $z > 0$ sont satisfaites car $x, y, z \geq 1$.
Le PPCM des dénominateurs est $\text{PPCM}(11, 111, 12210) = 12210$.
En réduisant au même dénominateur :
\begin{itemize}
    \item $\frac{1}{11} = \frac{1110}{12210}$
    \item $\frac{1}{111} = \frac{110}{12210}$
    \item $\frac{1}{12210} = \frac{1}{12210}$
\end{itemize}
La somme des numérateurs est :
$$ \frac{1}{11} + \frac{1}{111} + \frac{1}{12210} = \frac{1110 + 110 + 1}{12210} = \frac{1221}{12210} $$
Le PGCD du numérateur et du dénominateur est $\text{PGCD}(1221, 12210) = 1221$.
La fraction irréductible est :
$$ \frac{1221}{12210} = \frac{1221 \div 1221}{12210 \div 1221} = \frac{1}{10} $$
Cette fraction est bien égale à $\frac{5}{50}$.

\subsection{Démonstration pour $n = 51$}
Soit $n = 51$. Nous cherchons $x, y, z \in \mathbb{Z}^{+}$ tels que $\frac{5}{51} = \frac{1}{x} + \frac{1}{y} + \frac{1}{z}$.
Posons $x = 11$, $y = 141$, $z = 26367$.
Les conditions $x > 0$, $y > 0$ et $z > 0$ sont satisfaites car $x, y, z \geq 1$.
Le PPCM des dénominateurs est $\text{PPCM}(11, 141, 26367) = 26367$.
En réduisant au même dénominateur :
\begin{itemize}
    \item $\frac{1}{11} = \frac{2397}{26367}$
    \item $\frac{1}{141} = \frac{187}{26367}$
    \item $\frac{1}{26367} = \frac{1}{26367}$
\end{itemize}
La somme des numérateurs est :
$$ \frac{1}{11} + \frac{1}{141} + \frac{1}{26367} = \frac{2397 + 187 + 1}{26367} = \frac{2585}{26367} $$
Le PGCD du numérateur et du dénominateur est $\text{PGCD}(2585, 26367) = 517$.
La fraction irréductible est :
$$ \frac{2585}{26367} = \frac{2585 \div 517}{26367 \div 517} = \frac{5}{51} $$
Cette fraction est bien égale à $\frac{5}{51}$.

\subsection{Démonstration pour $n = 52$}
Soit $n = 52$. Nous cherchons $x, y, z \in \mathbb{Z}^{+}$ tels que $\frac{5}{52} = \frac{1}{x} + \frac{1}{y} + \frac{1}{z}$.
Posons $x = 11$, $y = 191$, $z = 109252$.
Les conditions $x > 0$, $y > 0$ et $z > 0$ sont satisfaites car $x, y, z \geq 1$.
Le PPCM des dénominateurs est $\text{PPCM}(11, 191, 109252) = 109252$.
En réduisant au même dénominateur :
\begin{itemize}
    \item $\frac{1}{11} = \frac{9932}{109252}$
    \item $\frac{1}{191} = \frac{572}{109252}$
    \item $\frac{1}{109252} = \frac{1}{109252}$
\end{itemize}
La somme des numérateurs est :
$$ \frac{1}{11} + \frac{1}{191} + \frac{1}{109252} = \frac{9932 + 572 + 1}{109252} = \frac{10505}{109252} $$
Le PGCD du numérateur et du dénominateur est $\text{PGCD}(10505, 109252) = 2101$.
La fraction irréductible est :
$$ \frac{10505}{109252} = \frac{10505 \div 2101}{109252 \div 2101} = \frac{5}{52} $$
Cette fraction est bien égale à $\frac{5}{52}$.

\subsection{Démonstration pour $n = 53$}
Soit $n = 53$. Nous cherchons $x, y, z \in \mathbb{Z}^{+}$ tels que $\frac{5}{53} = \frac{1}{x} + \frac{1}{y} + \frac{1}{z}$.
Posons $x = 11$, $y = 292$, $z = 170236$.
Les conditions $x > 0$, $y > 0$ et $z > 0$ sont satisfaites car $x, y, z \geq 1$.
Le PPCM des dénominateurs est $\text{PPCM}(11, 292, 170236) = 170236$.
En réduisant au même dénominateur :
\begin{itemize}
    \item $\frac{1}{11} = \frac{15476}{170236}$
    \item $\frac{1}{292} = \frac{583}{170236}$
    \item $\frac{1}{170236} = \frac{1}{170236}$
\end{itemize}
La somme des numérateurs est :
$$ \frac{1}{11} + \frac{1}{292} + \frac{1}{170236} = \frac{15476 + 583 + 1}{170236} = \frac{16060}{170236} $$
Le PGCD du numérateur et du dénominateur est $\text{PGCD}(16060, 170236) = 3212$.
La fraction irréductible est :
$$ \frac{16060}{170236} = \frac{16060 \div 3212}{170236 \div 3212} = \frac{5}{53} $$
Cette fraction est bien égale à $\frac{5}{53}$.

\subsection{Démonstration pour $n = 54$}
Soit $n = 54$. Nous cherchons $x, y, z \in \mathbb{Z}^{+}$ tels que $\frac{5}{54} = \frac{1}{x} + \frac{1}{y} + \frac{1}{z}$.
Posons $x = 11$, $y = 595$, $z = 353430$.
Les conditions $x > 0$, $y > 0$ et $z > 0$ sont satisfaites car $x, y, z \geq 1$.
Le PPCM des dénominateurs est $\text{PPCM}(11, 595, 353430) = 353430$.
En réduisant au même dénominateur :
\begin{itemize}
    \item $\frac{1}{11} = \frac{32130}{353430}$
    \item $\frac{1}{595} = \frac{594}{353430}$
    \item $\frac{1}{353430} = \frac{1}{353430}$
\end{itemize}
La somme des numérateurs est :
$$ \frac{1}{11} + \frac{1}{595} + \frac{1}{353430} = \frac{32130 + 594 + 1}{353430} = \frac{32725}{353430} $$
Le PGCD du numérateur et du dénominateur est $\text{PGCD}(32725, 353430) = 6545$.
La fraction irréductible est :
$$ \frac{32725}{353430} = \frac{32725 \div 6545}{353430 \div 6545} = \frac{5}{54} $$
Cette fraction est bien égale à $\frac{5}{54}$.

\subsection{Démonstration pour $n = 55$}
Soit $n = 55$. Nous cherchons $x, y, z \in \mathbb{Z}^{+}$ tels que $\frac{5}{55} = \frac{1}{x} + \frac{1}{y} + \frac{1}{z}$.
Posons $x = 12$, $y = 133$, $z = 17556$.
Les conditions $x > 0$, $y > 0$ et $z > 0$ sont satisfaites car $x, y, z \geq 1$.
Le PPCM des dénominateurs est $\text{PPCM}(12, 133, 17556) = 17556$.
En réduisant au même dénominateur :
\begin{itemize}
    \item $\frac{1}{12} = \frac{1463}{17556}$
    \item $\frac{1}{133} = \frac{132}{17556}$
    \item $\frac{1}{17556} = \frac{1}{17556}$
\end{itemize}
La somme des numérateurs est :
$$ \frac{1}{12} + \frac{1}{133} + \frac{1}{17556} = \frac{1463 + 132 + 1}{17556} = \frac{1596}{17556} $$
Le PGCD du numérateur et du dénominateur est $\text{PGCD}(1596, 17556) = 1596$.
La fraction irréductible est :
$$ \frac{1596}{17556} = \frac{1596 \div 1596}{17556 \div 1596} = \frac{1}{11} $$
Cette fraction est bien égale à $\frac{5}{55}$.

\subsection{Démonstration pour $n = 56$}
Soit $n = 56$. Nous cherchons $x, y, z \in \mathbb{Z}^{+}$ tels que $\frac{5}{56} = \frac{1}{x} + \frac{1}{y} + \frac{1}{z}$.
Posons $x = 12$, $y = 169$, $z = 28392$.
Les conditions $x > 0$, $y > 0$ et $z > 0$ sont satisfaites car $x, y, z \geq 1$.
Le PPCM des dénominateurs est $\text{PPCM}(12, 169, 28392) = 28392$.
En réduisant au même dénominateur :
\begin{itemize}
    \item $\frac{1}{12} = \frac{2366}{28392}$
    \item $\frac{1}{169} = \frac{168}{28392}$
    \item $\frac{1}{28392} = \frac{1}{28392}$
\end{itemize}
La somme des numérateurs est :
$$ \frac{1}{12} + \frac{1}{169} + \frac{1}{28392} = \frac{2366 + 168 + 1}{28392} = \frac{2535}{28392} $$
Le PGCD du numérateur et du dénominateur est $\text{PGCD}(2535, 28392) = 507$.
La fraction irréductible est :
$$ \frac{2535}{28392} = \frac{2535 \div 507}{28392 \div 507} = \frac{5}{56} $$
Cette fraction est bien égale à $\frac{5}{56}$.

\subsection{Démonstration pour $n = 57$}
Soit $n = 57$. Nous cherchons $x, y, z \in \mathbb{Z}^{+}$ tels que $\frac{5}{57} = \frac{1}{x} + \frac{1}{y} + \frac{1}{z}$.
Posons $x = 12$, $y = 229$, $z = 52212$.
Les conditions $x > 0$, $y > 0$ et $z > 0$ sont satisfaites car $x, y, z \geq 1$.
Le PPCM des dénominateurs est $\text{PPCM}(12, 229, 52212) = 52212$.
En réduisant au même dénominateur :
\begin{itemize}
    \item $\frac{1}{12} = \frac{4351}{52212}$
    \item $\frac{1}{229} = \frac{228}{52212}$
    \item $\frac{1}{52212} = \frac{1}{52212}$
\end{itemize}
La somme des numérateurs est :
$$ \frac{1}{12} + \frac{1}{229} + \frac{1}{52212} = \frac{4351 + 228 + 1}{52212} = \frac{4580}{52212} $$
Le PGCD du numérateur et du dénominateur est $\text{PGCD}(4580, 52212) = 916$.
La fraction irréductible est :
$$ \frac{4580}{52212} = \frac{4580 \div 916}{52212 \div 916} = \frac{5}{57} $$
Cette fraction est bien égale à $\frac{5}{57}$.

\subsection{Démonstration pour $n = 58$}
Soit $n = 58$. Nous cherchons $x, y, z \in \mathbb{Z}^{+}$ tels que $\frac{5}{58} = \frac{1}{x} + \frac{1}{y} + \frac{1}{z}$.
Posons $x = 12$, $y = 349$, $z = 121452$.
Les conditions $x > 0$, $y > 0$ et $z > 0$ sont satisfaites car $x, y, z \geq 1$.
Le PPCM des dénominateurs est $\text{PPCM}(12, 349, 121452) = 121452$.
En réduisant au même dénominateur :
\begin{itemize}
    \item $\frac{1}{12} = \frac{10121}{121452}$
    \item $\frac{1}{349} = \frac{348}{121452}$
    \item $\frac{1}{121452} = \frac{1}{121452}$
\end{itemize}
La somme des numérateurs est :
$$ \frac{1}{12} + \frac{1}{349} + \frac{1}{121452} = \frac{10121 + 348 + 1}{121452} = \frac{10470}{121452} $$
Le PGCD du numérateur et du dénominateur est $\text{PGCD}(10470, 121452) = 2094$.
La fraction irréductible est :
$$ \frac{10470}{121452} = \frac{10470 \div 2094}{121452 \div 2094} = \frac{5}{58} $$
Cette fraction est bien égale à $\frac{5}{58}$.

\subsection{Démonstration pour $n = 59$}
Soit $n = 59$. Nous cherchons $x, y, z \in \mathbb{Z}^{+}$ tels que $\frac{5}{59} = \frac{1}{x} + \frac{1}{y} + \frac{1}{z}$.
Posons $x = 12$, $y = 709$, $z = 501972$.
Les conditions $x > 0$, $y > 0$ et $z > 0$ sont satisfaites car $x, y, z \geq 1$.
Le PPCM des dénominateurs est $\text{PPCM}(12, 709, 501972) = 501972$.
En réduisant au même dénominateur :
\begin{itemize}
    \item $\frac{1}{12} = \frac{41831}{501972}$
    \item $\frac{1}{709} = \frac{708}{501972}$
    \item $\frac{1}{501972} = \frac{1}{501972}$
\end{itemize}
La somme des numérateurs est :
$$ \frac{1}{12} + \frac{1}{709} + \frac{1}{501972} = \frac{41831 + 708 + 1}{501972} = \frac{42540}{501972} $$
Le PGCD du numérateur et du dénominateur est $\text{PGCD}(42540, 501972) = 8508$.
La fraction irréductible est :
$$ \frac{42540}{501972} = \frac{42540 \div 8508}{501972 \div 8508} = \frac{5}{59} $$
Cette fraction est bien égale à $\frac{5}{59}$.

\subsection{Démonstration pour $n = 60$}
Soit $n = 60$. Nous cherchons $x, y, z \in \mathbb{Z}^{+}$ tels que $\frac{5}{60} = \frac{1}{x} + \frac{1}{y} + \frac{1}{z}$.
Posons $x = 13$, $y = 157$, $z = 24492$.
Les conditions $x > 0$, $y > 0$ et $z > 0$ sont satisfaites car $x, y, z \geq 1$.
Le PPCM des dénominateurs est $\text{PPCM}(13, 157, 24492) = 24492$.
En réduisant au même dénominateur :
\begin{itemize}
    \item $\frac{1}{13} = \frac{1884}{24492}$
    \item $\frac{1}{157} = \frac{156}{24492}$
    \item $\frac{1}{24492} = \frac{1}{24492}$
\end{itemize}
La somme des numérateurs est :
$$ \frac{1}{13} + \frac{1}{157} + \frac{1}{24492} = \frac{1884 + 156 + 1}{24492} = \frac{2041}{24492} $$
Le PGCD du numérateur et du dénominateur est $\text{PGCD}(2041, 24492) = 2041$.
La fraction irréductible est :
$$ \frac{2041}{24492} = \frac{2041 \div 2041}{24492 \div 2041} = \frac{1}{12} $$
Cette fraction est bien égale à $\frac{5}{60}$.

\subsection{Démonstration pour $n = 61$}
Soit $n = 61$. Nous cherchons $x, y, z \in \mathbb{Z}^{+}$ tels que $\frac{5}{61} = \frac{1}{x} + \frac{1}{y} + \frac{1}{z}$.
Posons $x = 14$, $y = 95$, $z = 81130$.
Les conditions $x > 0$, $y > 0$ et $z > 0$ sont satisfaites car $x, y, z \geq 1$.
Le PPCM des dénominateurs est $\text{PPCM}(14, 95, 81130) = 81130$.
En réduisant au même dénominateur :
\begin{itemize}
    \item $\frac{1}{14} = \frac{5795}{81130}$
    \item $\frac{1}{95} = \frac{854}{81130}$
    \item $\frac{1}{81130} = \frac{1}{81130}$
\end{itemize}
La somme des numérateurs est :
$$ \frac{1}{14} + \frac{1}{95} + \frac{1}{81130} = \frac{5795 + 854 + 1}{81130} = \frac{6650}{81130} $$
Le PGCD du numérateur et du dénominateur est $\text{PGCD}(6650, 81130) = 1330$.
La fraction irréductible est :
$$ \frac{6650}{81130} = \frac{6650 \div 1330}{81130 \div 1330} = \frac{5}{61} $$
Cette fraction est bien égale à $\frac{5}{61}$.

\subsection{Démonstration pour $n = 62$}
Soit $n = 62$. Nous cherchons $x, y, z \in \mathbb{Z}^{+}$ tels que $\frac{5}{62} = \frac{1}{x} + \frac{1}{y} + \frac{1}{z}$.
Posons $x = 13$, $y = 269$, $z = 216814$.
Les conditions $x > 0$, $y > 0$ et $z > 0$ sont satisfaites car $x, y, z \geq 1$.
Le PPCM des dénominateurs est $\text{PPCM}(13, 269, 216814) = 216814$.
En réduisant au même dénominateur :
\begin{itemize}
    \item $\frac{1}{13} = \frac{16678}{216814}$
    \item $\frac{1}{269} = \frac{806}{216814}$
    \item $\frac{1}{216814} = \frac{1}{216814}$
\end{itemize}
La somme des numérateurs est :
$$ \frac{1}{13} + \frac{1}{269} + \frac{1}{216814} = \frac{16678 + 806 + 1}{216814} = \frac{17485}{216814} $$
Le PGCD du numérateur et du dénominateur est $\text{PGCD}(17485, 216814) = 3497$.
La fraction irréductible est :
$$ \frac{17485}{216814} = \frac{17485 \div 3497}{216814 \div 3497} = \frac{5}{62} $$
Cette fraction est bien égale à $\frac{5}{62}$.

\subsection{Démonstration pour $n = 63$}
Soit $n = 63$. Nous cherchons $x, y, z \in \mathbb{Z}^{+}$ tels que $\frac{5}{63} = \frac{1}{x} + \frac{1}{y} + \frac{1}{z}$.
Posons $x = 13$, $y = 410$, $z = 335790$.
Les conditions $x > 0$, $y > 0$ et $z > 0$ sont satisfaites car $x, y, z \geq 1$.
Le PPCM des dénominateurs est $\text{PPCM}(13, 410, 335790) = 335790$.
En réduisant au même dénominateur :
\begin{itemize}
    \item $\frac{1}{13} = \frac{25830}{335790}$
    \item $\frac{1}{410} = \frac{819}{335790}$
    \item $\frac{1}{335790} = \frac{1}{335790}$
\end{itemize}
La somme des numérateurs est :
$$ \frac{1}{13} + \frac{1}{410} + \frac{1}{335790} = \frac{25830 + 819 + 1}{335790} = \frac{26650}{335790} $$
Le PGCD du numérateur et du dénominateur est $\text{PGCD}(26650, 335790) = 5330$.
La fraction irréductible est :
$$ \frac{26650}{335790} = \frac{26650 \div 5330}{335790 \div 5330} = \frac{5}{63} $$
Cette fraction est bien égale à $\frac{5}{63}$.

\subsection{Démonstration pour $n = 64$}
Soit $n = 64$. Nous cherchons $x, y, z \in \mathbb{Z}^{+}$ tels que $\frac{5}{64} = \frac{1}{x} + \frac{1}{y} + \frac{1}{z}$.
Posons $x = 13$, $y = 833$, $z = 693056$.
Les conditions $x > 0$, $y > 0$ et $z > 0$ sont satisfaites car $x, y, z \geq 1$.
Le PPCM des dénominateurs est $\text{PPCM}(13, 833, 693056) = 693056$.
En réduisant au même dénominateur :
\begin{itemize}
    \item $\frac{1}{13} = \frac{53312}{693056}$
    \item $\frac{1}{833} = \frac{832}{693056}$
    \item $\frac{1}{693056} = \frac{1}{693056}$
\end{itemize}
La somme des numérateurs est :
$$ \frac{1}{13} + \frac{1}{833} + \frac{1}{693056} = \frac{53312 + 832 + 1}{693056} = \frac{54145}{693056} $$
Le PGCD du numérateur et du dénominateur est $\text{PGCD}(54145, 693056) = 10829$.
La fraction irréductible est :
$$ \frac{54145}{693056} = \frac{54145 \div 10829}{693056 \div 10829} = \frac{5}{64} $$
Cette fraction est bien égale à $\frac{5}{64}$.

\subsection{Démonstration pour $n = 65$}
Soit $n = 65$. Nous cherchons $x, y, z \in \mathbb{Z}^{+}$ tels que $\frac{5}{65} = \frac{1}{x} + \frac{1}{y} + \frac{1}{z}$.
Posons $x = 14$, $y = 183$, $z = 33306$.
Les conditions $x > 0$, $y > 0$ et $z > 0$ sont satisfaites car $x, y, z \geq 1$.
Le PPCM des dénominateurs est $\text{PPCM}(14, 183, 33306) = 33306$.
En réduisant au même dénominateur :
\begin{itemize}
    \item $\frac{1}{14} = \frac{2379}{33306}$
    \item $\frac{1}{183} = \frac{182}{33306}$
    \item $\frac{1}{33306} = \frac{1}{33306}$
\end{itemize}
La somme des numérateurs est :
$$ \frac{1}{14} + \frac{1}{183} + \frac{1}{33306} = \frac{2379 + 182 + 1}{33306} = \frac{2562}{33306} $$
Le PGCD du numérateur et du dénominateur est $\text{PGCD}(2562, 33306) = 2562$.
La fraction irréductible est :
$$ \frac{2562}{33306} = \frac{2562 \div 2562}{33306 \div 2562} = \frac{1}{13} $$
Cette fraction est bien égale à $\frac{5}{65}$.

\subsection{Démonstration pour $n = 66$}
Soit $n = 66$. Nous cherchons $x, y, z \in \mathbb{Z}^{+}$ tels que $\frac{5}{66} = \frac{1}{x} + \frac{1}{y} + \frac{1}{z}$.
Posons $x = 14$, $y = 232$, $z = 53592$.
Les conditions $x > 0$, $y > 0$ et $z > 0$ sont satisfaites car $x, y, z \geq 1$.
Le PPCM des dénominateurs est $\text{PPCM}(14, 232, 53592) = 53592$.
En réduisant au même dénominateur :
\begin{itemize}
    \item $\frac{1}{14} = \frac{3828}{53592}$
    \item $\frac{1}{232} = \frac{231}{53592}$
    \item $\frac{1}{53592} = \frac{1}{53592}$
\end{itemize}
La somme des numérateurs est :
$$ \frac{1}{14} + \frac{1}{232} + \frac{1}{53592} = \frac{3828 + 231 + 1}{53592} = \frac{4060}{53592} $$
Le PGCD du numérateur et du dénominateur est $\text{PGCD}(4060, 53592) = 812$.
La fraction irréductible est :
$$ \frac{4060}{53592} = \frac{4060 \div 812}{53592 \div 812} = \frac{5}{66} $$
Cette fraction est bien égale à $\frac{5}{66}$.

\subsection{Démonstration pour $n = 67$}
Soit $n = 67$. Nous cherchons $x, y, z \in \mathbb{Z}^{+}$ tels que $\frac{5}{67} = \frac{1}{x} + \frac{1}{y} + \frac{1}{z}$.
Posons $x = 14$, $y = 313$, $z = 293594$.
Les conditions $x > 0$, $y > 0$ et $z > 0$ sont satisfaites car $x, y, z \geq 1$.
Le PPCM des dénominateurs est $\text{PPCM}(14, 313, 293594) = 293594$.
En réduisant au même dénominateur :
\begin{itemize}
    \item $\frac{1}{14} = \frac{20971}{293594}$
    \item $\frac{1}{313} = \frac{938}{293594}$
    \item $\frac{1}{293594} = \frac{1}{293594}$
\end{itemize}
La somme des numérateurs est :
$$ \frac{1}{14} + \frac{1}{313} + \frac{1}{293594} = \frac{20971 + 938 + 1}{293594} = \frac{21910}{293594} $$
Le PGCD du numérateur et du dénominateur est $\text{PGCD}(21910, 293594) = 4382$.
La fraction irréductible est :
$$ \frac{21910}{293594} = \frac{21910 \div 4382}{293594 \div 4382} = \frac{5}{67} $$
Cette fraction est bien égale à $\frac{5}{67}$.

\subsection{Démonstration pour $n = 68$}
Soit $n = 68$. Nous cherchons $x, y, z \in \mathbb{Z}^{+}$ tels que $\frac{5}{68} = \frac{1}{x} + \frac{1}{y} + \frac{1}{z}$.
Posons $x = 14$, $y = 477$, $z = 227052$.
Les conditions $x > 0$, $y > 0$ et $z > 0$ sont satisfaites car $x, y, z \geq 1$.
Le PPCM des dénominateurs est $\text{PPCM}(14, 477, 227052) = 227052$.
En réduisant au même dénominateur :
\begin{itemize}
    \item $\frac{1}{14} = \frac{16218}{227052}$
    \item $\frac{1}{477} = \frac{476}{227052}$
    \item $\frac{1}{227052} = \frac{1}{227052}$
\end{itemize}
La somme des numérateurs est :
$$ \frac{1}{14} + \frac{1}{477} + \frac{1}{227052} = \frac{16218 + 476 + 1}{227052} = \frac{16695}{227052} $$
Le PGCD du numérateur et du dénominateur est $\text{PGCD}(16695, 227052) = 3339$.
La fraction irréductible est :
$$ \frac{16695}{227052} = \frac{16695 \div 3339}{227052 \div 3339} = \frac{5}{68} $$
Cette fraction est bien égale à $\frac{5}{68}$.

\subsection{Démonstration pour $n = 69$}
Soit $n = 69$. Nous cherchons $x, y, z \in \mathbb{Z}^{+}$ tels que $\frac{5}{69} = \frac{1}{x} + \frac{1}{y} + \frac{1}{z}$.
Posons $x = 14$, $y = 967$, $z = 934122$.
Les conditions $x > 0$, $y > 0$ et $z > 0$ sont satisfaites car $x, y, z \geq 1$.
Le PPCM des dénominateurs est $\text{PPCM}(14, 967, 934122) = 934122$.
En réduisant au même dénominateur :
\begin{itemize}
    \item $\frac{1}{14} = \frac{66723}{934122}$
    \item $\frac{1}{967} = \frac{966}{934122}$
    \item $\frac{1}{934122} = \frac{1}{934122}$
\end{itemize}
La somme des numérateurs est :
$$ \frac{1}{14} + \frac{1}{967} + \frac{1}{934122} = \frac{66723 + 966 + 1}{934122} = \frac{67690}{934122} $$
Le PGCD du numérateur et du dénominateur est $\text{PGCD}(67690, 934122) = 13538$.
La fraction irréductible est :
$$ \frac{67690}{934122} = \frac{67690 \div 13538}{934122 \div 13538} = \frac{5}{69} $$
Cette fraction est bien égale à $\frac{5}{69}$.

\subsection{Démonstration pour $n = 70$}
Soit $n = 70$. Nous cherchons $x, y, z \in \mathbb{Z}^{+}$ tels que $\frac{5}{70} = \frac{1}{x} + \frac{1}{y} + \frac{1}{z}$.
Posons $x = 15$, $y = 211$, $z = 44310$.
Les conditions $x > 0$, $y > 0$ et $z > 0$ sont satisfaites car $x, y, z \geq 1$.
Le PPCM des dénominateurs est $\text{PPCM}(15, 211, 44310) = 44310$.
En réduisant au même dénominateur :
\begin{itemize}
    \item $\frac{1}{15} = \frac{2954}{44310}$
    \item $\frac{1}{211} = \frac{210}{44310}$
    \item $\frac{1}{44310} = \frac{1}{44310}$
\end{itemize}
La somme des numérateurs est :
$$ \frac{1}{15} + \frac{1}{211} + \frac{1}{44310} = \frac{2954 + 210 + 1}{44310} = \frac{3165}{44310} $$
Le PGCD du numérateur et du dénominateur est $\text{PGCD}(3165, 44310) = 3165$.
La fraction irréductible est :
$$ \frac{3165}{44310} = \frac{3165 \div 3165}{44310 \div 3165} = \frac{1}{14} $$
Cette fraction est bien égale à $\frac{5}{70}$.

\subsection{Démonstration pour $n = 71$}
Soit $n = 71$. Nous cherchons $x, y, z \in \mathbb{Z}^{+}$ tels que $\frac{5}{71} = \frac{1}{x} + \frac{1}{y} + \frac{1}{z}$.
Posons $x = 15$, $y = 267$, $z = 94785$.
Les conditions $x > 0$, $y > 0$ et $z > 0$ sont satisfaites car $x, y, z \geq 1$.
Le PPCM des dénominateurs est $\text{PPCM}(15, 267, 94785) = 94785$.
En réduisant au même dénominateur :
\begin{itemize}
    \item $\frac{1}{15} = \frac{6319}{94785}$
    \item $\frac{1}{267} = \frac{355}{94785}$
    \item $\frac{1}{94785} = \frac{1}{94785}$
\end{itemize}
La somme des numérateurs est :
$$ \frac{1}{15} + \frac{1}{267} + \frac{1}{94785} = \frac{6319 + 355 + 1}{94785} = \frac{6675}{94785} $$
Le PGCD du numérateur et du dénominateur est $\text{PGCD}(6675, 94785) = 1335$.
La fraction irréductible est :
$$ \frac{6675}{94785} = \frac{6675 \div 1335}{94785 \div 1335} = \frac{5}{71} $$
Cette fraction est bien égale à $\frac{5}{71}$.

\subsection{Démonstration pour $n = 72$}
Soit $n = 72$. Nous cherchons $x, y, z \in \mathbb{Z}^{+}$ tels que $\frac{5}{72} = \frac{1}{x} + \frac{1}{y} + \frac{1}{z}$.
Posons $x = 15$, $y = 361$, $z = 129960$.
Les conditions $x > 0$, $y > 0$ et $z > 0$ sont satisfaites car $x, y, z \geq 1$.
Le PPCM des dénominateurs est $\text{PPCM}(15, 361, 129960) = 129960$.
En réduisant au même dénominateur :
\begin{itemize}
    \item $\frac{1}{15} = \frac{8664}{129960}$
    \item $\frac{1}{361} = \frac{360}{129960}$
    \item $\frac{1}{129960} = \frac{1}{129960}$
\end{itemize}
La somme des numérateurs est :
$$ \frac{1}{15} + \frac{1}{361} + \frac{1}{129960} = \frac{8664 + 360 + 1}{129960} = \frac{9025}{129960} $$
Le PGCD du numérateur et du dénominateur est $\text{PGCD}(9025, 129960) = 1805$.
La fraction irréductible est :
$$ \frac{9025}{129960} = \frac{9025 \div 1805}{129960 \div 1805} = \frac{5}{72} $$
Cette fraction est bien égale à $\frac{5}{72}$.

\subsection{Démonstration pour $n = 73$}
Soit $n = 73$. Nous cherchons $x, y, z \in \mathbb{Z}^{+}$ tels que $\frac{5}{73} = \frac{1}{x} + \frac{1}{y} + \frac{1}{z}$.
Posons $x = 15$, $y = 548$, $z = 600060$.
Les conditions $x > 0$, $y > 0$ et $z > 0$ sont satisfaites car $x, y, z \geq 1$.
Le PPCM des dénominateurs est $\text{PPCM}(15, 548, 600060) = 600060$.
En réduisant au même dénominateur :
\begin{itemize}
    \item $\frac{1}{15} = \frac{40004}{600060}$
    \item $\frac{1}{548} = \frac{1095}{600060}$
    \item $\frac{1}{600060} = \frac{1}{600060}$
\end{itemize}
La somme des numérateurs est :
$$ \frac{1}{15} + \frac{1}{548} + \frac{1}{600060} = \frac{40004 + 1095 + 1}{600060} = \frac{41100}{600060} $$
Le PGCD du numérateur et du dénominateur est $\text{PGCD}(41100, 600060) = 8220$.
La fraction irréductible est :
$$ \frac{41100}{600060} = \frac{41100 \div 8220}{600060 \div 8220} = \frac{5}{73} $$
Cette fraction est bien égale à $\frac{5}{73}$.

\subsection{Démonstration pour $n = 74$}
Soit $n = 74$. Nous cherchons $x, y, z \in \mathbb{Z}^{+}$ tels que $\frac{5}{74} = \frac{1}{x} + \frac{1}{y} + \frac{1}{z}$.
Posons $x = 15$, $y = 1111$, $z = 1233210$.
Les conditions $x > 0$, $y > 0$ et $z > 0$ sont satisfaites car $x, y, z \geq 1$.
Le PPCM des dénominateurs est $\text{PPCM}(15, 1111, 1233210) = 1233210$.
En réduisant au même dénominateur :
\begin{itemize}
    \item $\frac{1}{15} = \frac{82214}{1233210}$
    \item $\frac{1}{1111} = \frac{1110}{1233210}$
    \item $\frac{1}{1233210} = \frac{1}{1233210}$
\end{itemize}
La somme des numérateurs est :
$$ \frac{1}{15} + \frac{1}{1111} + \frac{1}{1233210} = \frac{82214 + 1110 + 1}{1233210} = \frac{83325}{1233210} $$
Le PGCD du numérateur et du dénominateur est $\text{PGCD}(83325, 1233210) = 16665$.
La fraction irréductible est :
$$ \frac{83325}{1233210} = \frac{83325 \div 16665}{1233210 \div 16665} = \frac{5}{74} $$
Cette fraction est bien égale à $\frac{5}{74}$.

\subsection{Démonstration pour $n = 75$}
Soit $n = 75$. Nous cherchons $x, y, z \in \mathbb{Z}^{+}$ tels que $\frac{5}{75} = \frac{1}{x} + \frac{1}{y} + \frac{1}{z}$.
Posons $x = 16$, $y = 241$, $z = 57840$.
Les conditions $x > 0$, $y > 0$ et $z > 0$ sont satisfaites car $x, y, z \geq 1$.
Le PPCM des dénominateurs est $\text{PPCM}(16, 241, 57840) = 57840$.
En réduisant au même dénominateur :
\begin{itemize}
    \item $\frac{1}{16} = \frac{3615}{57840}$
    \item $\frac{1}{241} = \frac{240}{57840}$
    \item $\frac{1}{57840} = \frac{1}{57840}$
\end{itemize}
La somme des numérateurs est :
$$ \frac{1}{16} + \frac{1}{241} + \frac{1}{57840} = \frac{3615 + 240 + 1}{57840} = \frac{3856}{57840} $$
Le PGCD du numérateur et du dénominateur est $\text{PGCD}(3856, 57840) = 3856$.
La fraction irréductible est :
$$ \frac{3856}{57840} = \frac{3856 \div 3856}{57840 \div 3856} = \frac{1}{15} $$
Cette fraction est bien égale à $\frac{5}{75}$.

\subsection{Démonstration pour $n = 76$}
Soit $n = 76$. Nous cherchons $x, y, z \in \mathbb{Z}^{+}$ tels que $\frac{5}{76} = \frac{1}{x} + \frac{1}{y} + \frac{1}{z}$.
Posons $x = 16$, $y = 305$, $z = 92720$.
Les conditions $x > 0$, $y > 0$ et $z > 0$ sont satisfaites car $x, y, z \geq 1$.
Le PPCM des dénominateurs est $\text{PPCM}(16, 305, 92720) = 92720$.
En réduisant au même dénominateur :
\begin{itemize}
    \item $\frac{1}{16} = \frac{5795}{92720}$
    \item $\frac{1}{305} = \frac{304}{92720}$
    \item $\frac{1}{92720} = \frac{1}{92720}$
\end{itemize}
La somme des numérateurs est :
$$ \frac{1}{16} + \frac{1}{305} + \frac{1}{92720} = \frac{5795 + 304 + 1}{92720} = \frac{6100}{92720} $$
Le PGCD du numérateur et du dénominateur est $\text{PGCD}(6100, 92720) = 1220$.
La fraction irréductible est :
$$ \frac{6100}{92720} = \frac{6100 \div 1220}{92720 \div 1220} = \frac{5}{76} $$
Cette fraction est bien égale à $\frac{5}{76}$.

\subsection{Démonstration pour $n = 77$}
Soit $n = 77$. Nous cherchons $x, y, z \in \mathbb{Z}^{+}$ tels que $\frac{5}{77} = \frac{1}{x} + \frac{1}{y} + \frac{1}{z}$.
Posons $x = 16$, $y = 411$, $z = 506352$.
Les conditions $x > 0$, $y > 0$ et $z > 0$ sont satisfaites car $x, y, z \geq 1$.
Le PPCM des dénominateurs est $\text{PPCM}(16, 411, 506352) = 506352$.
En réduisant au même dénominateur :
\begin{itemize}
    \item $\frac{1}{16} = \frac{31647}{506352}$
    \item $\frac{1}{411} = \frac{1232}{506352}$
    \item $\frac{1}{506352} = \frac{1}{506352}$
\end{itemize}
La somme des numérateurs est :
$$ \frac{1}{16} + \frac{1}{411} + \frac{1}{506352} = \frac{31647 + 1232 + 1}{506352} = \frac{32880}{506352} $$
Le PGCD du numérateur et du dénominateur est $\text{PGCD}(32880, 506352) = 6576$.
La fraction irréductible est :
$$ \frac{32880}{506352} = \frac{32880 \div 6576}{506352 \div 6576} = \frac{5}{77} $$
Cette fraction est bien égale à $\frac{5}{77}$.

\subsection{Démonstration pour $n = 78$}
Soit $n = 78$. Nous cherchons $x, y, z \in \mathbb{Z}^{+}$ tels que $\frac{5}{78} = \frac{1}{x} + \frac{1}{y} + \frac{1}{z}$.
Posons $x = 16$, $y = 625$, $z = 390000$.
Les conditions $x > 0$, $y > 0$ et $z > 0$ sont satisfaites car $x, y, z \geq 1$.
Le PPCM des dénominateurs est $\text{PPCM}(16, 625, 390000) = 390000$.
En réduisant au même dénominateur :
\begin{itemize}
    \item $\frac{1}{16} = \frac{24375}{390000}$
    \item $\frac{1}{625} = \frac{624}{390000}$
    \item $\frac{1}{390000} = \frac{1}{390000}$
\end{itemize}
La somme des numérateurs est :
$$ \frac{1}{16} + \frac{1}{625} + \frac{1}{390000} = \frac{24375 + 624 + 1}{390000} = \frac{25000}{390000} $$
Le PGCD du numérateur et du dénominateur est $\text{PGCD}(25000, 390000) = 5000$.
La fraction irréductible est :
$$ \frac{25000}{390000} = \frac{25000 \div 5000}{390000 \div 5000} = \frac{5}{78} $$
Cette fraction est bien égale à $\frac{5}{78}$.

\subsection{Démonstration pour $n = 79$}
Soit $n = 79$. Nous cherchons $x, y, z \in \mathbb{Z}^{+}$ tels que $\frac{5}{79} = \frac{1}{x} + \frac{1}{y} + \frac{1}{z}$.
Posons $x = 16$, $y = 1265$, $z = 1598960$.
Les conditions $x > 0$, $y > 0$ et $z > 0$ sont satisfaites car $x, y, z \geq 1$.
Le PPCM des dénominateurs est $\text{PPCM}(16, 1265, 1598960) = 1598960$.
En réduisant au même dénominateur :
\begin{itemize}
    \item $\frac{1}{16} = \frac{99935}{1598960}$
    \item $\frac{1}{1265} = \frac{1264}{1598960}$
    \item $\frac{1}{1598960} = \frac{1}{1598960}$
\end{itemize}
La somme des numérateurs est :
$$ \frac{1}{16} + \frac{1}{1265} + \frac{1}{1598960} = \frac{99935 + 1264 + 1}{1598960} = \frac{101200}{1598960} $$
Le PGCD du numérateur et du dénominateur est $\text{PGCD}(101200, 1598960) = 20240$.
La fraction irréductible est :
$$ \frac{101200}{1598960} = \frac{101200 \div 20240}{1598960 \div 20240} = \frac{5}{79} $$
Cette fraction est bien égale à $\frac{5}{79}$.

\subsection{Démonstration pour $n = 80$}
Soit $n = 80$. Nous cherchons $x, y, z \in \mathbb{Z}^{+}$ tels que $\frac{5}{80} = \frac{1}{x} + \frac{1}{y} + \frac{1}{z}$.
Posons $x = 17$, $y = 273$, $z = 74256$.
Les conditions $x > 0$, $y > 0$ et $z > 0$ sont satisfaites car $x, y, z \geq 1$.
Le PPCM des dénominateurs est $\text{PPCM}(17, 273, 74256) = 74256$.
En réduisant au même dénominateur :
\begin{itemize}
    \item $\frac{1}{17} = \frac{4368}{74256}$
    \item $\frac{1}{273} = \frac{272}{74256}$
    \item $\frac{1}{74256} = \frac{1}{74256}$
\end{itemize}
La somme des numérateurs est :
$$ \frac{1}{17} + \frac{1}{273} + \frac{1}{74256} = \frac{4368 + 272 + 1}{74256} = \frac{4641}{74256} $$
Le PGCD du numérateur et du dénominateur est $\text{PGCD}(4641, 74256) = 4641$.
La fraction irréductible est :
$$ \frac{4641}{74256} = \frac{4641 \div 4641}{74256 \div 4641} = \frac{1}{16} $$
Cette fraction est bien égale à $\frac{5}{80}$.

\subsection{Démonstration pour $n = 81$}
Soit $n = 81$. Nous cherchons $x, y, z \in \mathbb{Z}^{+}$ tels que $\frac{5}{81} = \frac{1}{x} + \frac{1}{y} + \frac{1}{z}$.
Posons $x = 17$, $y = 345$, $z = 158355$.
Les conditions $x > 0$, $y > 0$ et $z > 0$ sont satisfaites car $x, y, z \geq 1$.
Le PPCM des dénominateurs est $\text{PPCM}(17, 345, 158355) = 158355$.
En réduisant au même dénominateur :
\begin{itemize}
    \item $\frac{1}{17} = \frac{9315}{158355}$
    \item $\frac{1}{345} = \frac{459}{158355}$
    \item $\frac{1}{158355} = \frac{1}{158355}$
\end{itemize}
La somme des numérateurs est :
$$ \frac{1}{17} + \frac{1}{345} + \frac{1}{158355} = \frac{9315 + 459 + 1}{158355} = \frac{9775}{158355} $$
Le PGCD du numérateur et du dénominateur est $\text{PGCD}(9775, 158355) = 1955$.
La fraction irréductible est :
$$ \frac{9775}{158355} = \frac{9775 \div 1955}{158355 \div 1955} = \frac{5}{81} $$
Cette fraction est bien égale à $\frac{5}{81}$.

\subsection{Démonstration pour $n = 82$}
Soit $n = 82$. Nous cherchons $x, y, z \in \mathbb{Z}^{+}$ tels que $\frac{5}{82} = \frac{1}{x} + \frac{1}{y} + \frac{1}{z}$.
Posons $x = 17$, $y = 465$, $z = 648210$.
Les conditions $x > 0$, $y > 0$ et $z > 0$ sont satisfaites car $x, y, z \geq 1$.
Le PPCM des dénominateurs est $\text{PPCM}(17, 465, 648210) = 648210$.
En réduisant au même dénominateur :
\begin{itemize}
    \item $\frac{1}{17} = \frac{38130}{648210}$
    \item $\frac{1}{465} = \frac{1394}{648210}$
    \item $\frac{1}{648210} = \frac{1}{648210}$
\end{itemize}
La somme des numérateurs est :
$$ \frac{1}{17} + \frac{1}{465} + \frac{1}{648210} = \frac{38130 + 1394 + 1}{648210} = \frac{39525}{648210} $$
Le PGCD du numérateur et du dénominateur est $\text{PGCD}(39525, 648210) = 7905$.
La fraction irréductible est :
$$ \frac{39525}{648210} = \frac{39525 \div 7905}{648210 \div 7905} = \frac{5}{82} $$
Cette fraction est bien égale à $\frac{5}{82}$.

\subsection{Démonstration pour $n = 83$}
Soit $n = 83$. Nous cherchons $x, y, z \in \mathbb{Z}^{+}$ tels que $\frac{5}{83} = \frac{1}{x} + \frac{1}{y} + \frac{1}{z}$.
Posons $x = 17$, $y = 706$, $z = 996166$.
Les conditions $x > 0$, $y > 0$ et $z > 0$ sont satisfaites car $x, y, z \geq 1$.
Le PPCM des dénominateurs est $\text{PPCM}(17, 706, 996166) = 996166$.
En réduisant au même dénominateur :
\begin{itemize}
    \item $\frac{1}{17} = \frac{58598}{996166}$
    \item $\frac{1}{706} = \frac{1411}{996166}$
    \item $\frac{1}{996166} = \frac{1}{996166}$
\end{itemize}
La somme des numérateurs est :
$$ \frac{1}{17} + \frac{1}{706} + \frac{1}{996166} = \frac{58598 + 1411 + 1}{996166} = \frac{60010}{996166} $$
Le PGCD du numérateur et du dénominateur est $\text{PGCD}(60010, 996166) = 12002$.
La fraction irréductible est :
$$ \frac{60010}{996166} = \frac{60010 \div 12002}{996166 \div 12002} = \frac{5}{83} $$
Cette fraction est bien égale à $\frac{5}{83}$.

\subsection{Démonstration pour $n = 84$}
Soit $n = 84$. Nous cherchons $x, y, z \in \mathbb{Z}^{+}$ tels que $\frac{5}{84} = \frac{1}{x} + \frac{1}{y} + \frac{1}{z}$.
Posons $x = 17$, $y = 1429$, $z = 2040612$.
Les conditions $x > 0$, $y > 0$ et $z > 0$ sont satisfaites car $x, y, z \geq 1$.
Le PPCM des dénominateurs est $\text{PPCM}(17, 1429, 2040612) = 2040612$.
En réduisant au même dénominateur :
\begin{itemize}
    \item $\frac{1}{17} = \frac{120036}{2040612}$
    \item $\frac{1}{1429} = \frac{1428}{2040612}$
    \item $\frac{1}{2040612} = \frac{1}{2040612}$
\end{itemize}
La somme des numérateurs est :
$$ \frac{1}{17} + \frac{1}{1429} + \frac{1}{2040612} = \frac{120036 + 1428 + 1}{2040612} = \frac{121465}{2040612} $$
Le PGCD du numérateur et du dénominateur est $\text{PGCD}(121465, 2040612) = 24293$.
La fraction irréductible est :
$$ \frac{121465}{2040612} = \frac{121465 \div 24293}{2040612 \div 24293} = \frac{5}{84} $$
Cette fraction est bien égale à $\frac{5}{84}$.

\subsection{Démonstration pour $n = 85$}
Soit $n = 85$. Nous cherchons $x, y, z \in \mathbb{Z}^{+}$ tels que $\frac{5}{85} = \frac{1}{x} + \frac{1}{y} + \frac{1}{z}$.
Posons $x = 18$, $y = 307$, $z = 93942$.
Les conditions $x > 0$, $y > 0$ et $z > 0$ sont satisfaites car $x, y, z \geq 1$.
Le PPCM des dénominateurs est $\text{PPCM}(18, 307, 93942) = 93942$.
En réduisant au même dénominateur :
\begin{itemize}
    \item $\frac{1}{18} = \frac{5219}{93942}$
    \item $\frac{1}{307} = \frac{306}{93942}$
    \item $\frac{1}{93942} = \frac{1}{93942}$
\end{itemize}
La somme des numérateurs est :
$$ \frac{1}{18} + \frac{1}{307} + \frac{1}{93942} = \frac{5219 + 306 + 1}{93942} = \frac{5526}{93942} $$
Le PGCD du numérateur et du dénominateur est $\text{PGCD}(5526, 93942) = 5526$.
La fraction irréductible est :
$$ \frac{5526}{93942} = \frac{5526 \div 5526}{93942 \div 5526} = \frac{1}{17} $$
Cette fraction est bien égale à $\frac{5}{85}$.

\subsection{Démonstration pour $n = 86$}
Soit $n = 86$. Nous cherchons $x, y, z \in \mathbb{Z}^{+}$ tels que $\frac{5}{86} = \frac{1}{x} + \frac{1}{y} + \frac{1}{z}$.
Posons $x = 18$, $y = 388$, $z = 150156$.
Les conditions $x > 0$, $y > 0$ et $z > 0$ sont satisfaites car $x, y, z \geq 1$.
Le PPCM des dénominateurs est $\text{PPCM}(18, 388, 150156) = 150156$.
En réduisant au même dénominateur :
\begin{itemize}
    \item $\frac{1}{18} = \frac{8342}{150156}$
    \item $\frac{1}{388} = \frac{387}{150156}$
    \item $\frac{1}{150156} = \frac{1}{150156}$
\end{itemize}
La somme des numérateurs est :
$$ \frac{1}{18} + \frac{1}{388} + \frac{1}{150156} = \frac{8342 + 387 + 1}{150156} = \frac{8730}{150156} $$
Le PGCD du numérateur et du dénominateur est $\text{PGCD}(8730, 150156) = 1746$.
La fraction irréductible est :
$$ \frac{8730}{150156} = \frac{8730 \div 1746}{150156 \div 1746} = \frac{5}{86} $$
Cette fraction est bien égale à $\frac{5}{86}$.

\subsection{Démonstration pour $n = 87$}
Soit $n = 87$. Nous cherchons $x, y, z \in \mathbb{Z}^{+}$ tels que $\frac{5}{87} = \frac{1}{x} + \frac{1}{y} + \frac{1}{z}$.
Posons $x = 18$, $y = 523$, $z = 273006$.
Les conditions $x > 0$, $y > 0$ et $z > 0$ sont satisfaites car $x, y, z \geq 1$.
Le PPCM des dénominateurs est $\text{PPCM}(18, 523, 273006) = 273006$.
En réduisant au même dénominateur :
\begin{itemize}
    \item $\frac{1}{18} = \frac{15167}{273006}$
    \item $\frac{1}{523} = \frac{522}{273006}$
    \item $\frac{1}{273006} = \frac{1}{273006}$
\end{itemize}
La somme des numérateurs est :
$$ \frac{1}{18} + \frac{1}{523} + \frac{1}{273006} = \frac{15167 + 522 + 1}{273006} = \frac{15690}{273006} $$
Le PGCD du numérateur et du dénominateur est $\text{PGCD}(15690, 273006) = 3138$.
La fraction irréductible est :
$$ \frac{15690}{273006} = \frac{15690 \div 3138}{273006 \div 3138} = \frac{5}{87} $$
Cette fraction est bien égale à $\frac{5}{87}$.

\subsection{Démonstration pour $n = 88$}
Soit $n = 88$. Nous cherchons $x, y, z \in \mathbb{Z}^{+}$ tels que $\frac{5}{88} = \frac{1}{x} + \frac{1}{y} + \frac{1}{z}$.
Posons $x = 18$, $y = 793$, $z = 628056$.
Les conditions $x > 0$, $y > 0$ et $z > 0$ sont satisfaites car $x, y, z \geq 1$.
Le PPCM des dénominateurs est $\text{PPCM}(18, 793, 628056) = 628056$.
En réduisant au même dénominateur :
\begin{itemize}
    \item $\frac{1}{18} = \frac{34892}{628056}$
    \item $\frac{1}{793} = \frac{792}{628056}$
    \item $\frac{1}{628056} = \frac{1}{628056}$
\end{itemize}
La somme des numérateurs est :
$$ \frac{1}{18} + \frac{1}{793} + \frac{1}{628056} = \frac{34892 + 792 + 1}{628056} = \frac{35685}{628056} $$
Le PGCD du numérateur et du dénominateur est $\text{PGCD}(35685, 628056) = 7137$.
La fraction irréductible est :
$$ \frac{35685}{628056} = \frac{35685 \div 7137}{628056 \div 7137} = \frac{5}{88} $$
Cette fraction est bien égale à $\frac{5}{88}$.

\subsection{Démonstration pour $n = 89$}
Soit $n = 89$. Nous cherchons $x, y, z \in \mathbb{Z}^{+}$ tels que $\frac{5}{89} = \frac{1}{x} + \frac{1}{y} + \frac{1}{z}$.
Posons $x = 18$, $y = 1603$, $z = 2568006$.
Les conditions $x > 0$, $y > 0$ et $z > 0$ sont satisfaites car $x, y, z \geq 1$.
Le PPCM des dénominateurs est $\text{PPCM}(18, 1603, 2568006) = 2568006$.
En réduisant au même dénominateur :
\begin{itemize}
    \item $\frac{1}{18} = \frac{142667}{2568006}$
    \item $\frac{1}{1603} = \frac{1602}{2568006}$
    \item $\frac{1}{2568006} = \frac{1}{2568006}$
\end{itemize}
La somme des numérateurs est :
$$ \frac{1}{18} + \frac{1}{1603} + \frac{1}{2568006} = \frac{142667 + 1602 + 1}{2568006} = \frac{144270}{2568006} $$
Le PGCD du numérateur et du dénominateur est $\text{PGCD}(144270, 2568006) = 28854$.
La fraction irréductible est :
$$ \frac{144270}{2568006} = \frac{144270 \div 28854}{2568006 \div 28854} = \frac{5}{89} $$
Cette fraction est bien égale à $\frac{5}{89}$.

\subsection{Démonstration pour $n = 90$}
Soit $n = 90$. Nous cherchons $x, y, z \in \mathbb{Z}^{+}$ tels que $\frac{5}{90} = \frac{1}{x} + \frac{1}{y} + \frac{1}{z}$.
Posons $x = 19$, $y = 343$, $z = 117306$.
Les conditions $x > 0$, $y > 0$ et $z > 0$ sont satisfaites car $x, y, z \geq 1$.
Le PPCM des dénominateurs est $\text{PPCM}(19, 343, 117306) = 117306$.
En réduisant au même dénominateur :
\begin{itemize}
    \item $\frac{1}{19} = \frac{6174}{117306}$
    \item $\frac{1}{343} = \frac{342}{117306}$
    \item $\frac{1}{117306} = \frac{1}{117306}$
\end{itemize}
La somme des numérateurs est :
$$ \frac{1}{19} + \frac{1}{343} + \frac{1}{117306} = \frac{6174 + 342 + 1}{117306} = \frac{6517}{117306} $$
Le PGCD du numérateur et du dénominateur est $\text{PGCD}(6517, 117306) = 6517$.
La fraction irréductible est :
$$ \frac{6517}{117306} = \frac{6517 \div 6517}{117306 \div 6517} = \frac{1}{18} $$
Cette fraction est bien égale à $\frac{5}{90}$.

\subsection{Démonstration pour $n = 91$}
Soit $n = 91$. Nous cherchons $x, y, z \in \mathbb{Z}^{+}$ tels que $\frac{5}{91} = \frac{1}{x} + \frac{1}{y} + \frac{1}{z}$.
Posons $x = 19$, $y = 434$, $z = 107198$.
Les conditions $x > 0$, $y > 0$ et $z > 0$ sont satisfaites car $x, y, z \geq 1$.
Le PPCM des dénominateurs est $\text{PPCM}(19, 434, 107198) = 107198$.
En réduisant au même dénominateur :
\begin{itemize}
    \item $\frac{1}{19} = \frac{5642}{107198}$
    \item $\frac{1}{434} = \frac{247}{107198}$
    \item $\frac{1}{107198} = \frac{1}{107198}$
\end{itemize}
La somme des numérateurs est :
$$ \frac{1}{19} + \frac{1}{434} + \frac{1}{107198} = \frac{5642 + 247 + 1}{107198} = \frac{5890}{107198} $$
Le PGCD du numérateur et du dénominateur est $\text{PGCD}(5890, 107198) = 1178$.
La fraction irréductible est :
$$ \frac{5890}{107198} = \frac{5890 \div 1178}{107198 \div 1178} = \frac{5}{91} $$
Cette fraction est bien égale à $\frac{5}{91}$.

\subsection{Démonstration pour $n = 92$}
Soit $n = 92$. Nous cherchons $x, y, z \in \mathbb{Z}^{+}$ tels que $\frac{5}{92} = \frac{1}{x} + \frac{1}{y} + \frac{1}{z}$.
Posons $x = 19$, $y = 583$, $z = 1019084$.
Les conditions $x > 0$, $y > 0$ et $z > 0$ sont satisfaites car $x, y, z \geq 1$.
Le PPCM des dénominateurs est $\text{PPCM}(19, 583, 1019084) = 1019084$.
En réduisant au même dénominateur :
\begin{itemize}
    \item $\frac{1}{19} = \frac{53636}{1019084}$
    \item $\frac{1}{583} = \frac{1748}{1019084}$
    \item $\frac{1}{1019084} = \frac{1}{1019084}$
\end{itemize}
La somme des numérateurs est :
$$ \frac{1}{19} + \frac{1}{583} + \frac{1}{1019084} = \frac{53636 + 1748 + 1}{1019084} = \frac{55385}{1019084} $$
Le PGCD du numérateur et du dénominateur est $\text{PGCD}(55385, 1019084) = 11077$.
La fraction irréductible est :
$$ \frac{55385}{1019084} = \frac{55385 \div 11077}{1019084 \div 11077} = \frac{5}{92} $$
Cette fraction est bien égale à $\frac{5}{92}$.

\subsection{Démonstration pour $n = 93$}
Soit $n = 93$. Nous cherchons $x, y, z \in \mathbb{Z}^{+}$ tels que $\frac{5}{93} = \frac{1}{x} + \frac{1}{y} + \frac{1}{z}$.
Posons $x = 19$, $y = 884$, $z = 1562028$.
Les conditions $x > 0$, $y > 0$ et $z > 0$ sont satisfaites car $x, y, z \geq 1$.
Le PPCM des dénominateurs est $\text{PPCM}(19, 884, 1562028) = 1562028$.
En réduisant au même dénominateur :
\begin{itemize}
    \item $\frac{1}{19} = \frac{82212}{1562028}$
    \item $\frac{1}{884} = \frac{1767}{1562028}$
    \item $\frac{1}{1562028} = \frac{1}{1562028}$
\end{itemize}
La somme des numérateurs est :
$$ \frac{1}{19} + \frac{1}{884} + \frac{1}{1562028} = \frac{82212 + 1767 + 1}{1562028} = \frac{83980}{1562028} $$
Le PGCD du numérateur et du dénominateur est $\text{PGCD}(83980, 1562028) = 16796$.
La fraction irréductible est :
$$ \frac{83980}{1562028} = \frac{83980 \div 16796}{1562028 \div 16796} = \frac{5}{93} $$
Cette fraction est bien égale à $\frac{5}{93}$.

\subsection{Démonstration pour $n = 94$}
Soit $n = 94$. Nous cherchons $x, y, z \in \mathbb{Z}^{+}$ tels que $\frac{5}{94} = \frac{1}{x} + \frac{1}{y} + \frac{1}{z}$.
Posons $x = 19$, $y = 1787$, $z = 3191582$.
Les conditions $x > 0$, $y > 0$ et $z > 0$ sont satisfaites car $x, y, z \geq 1$.
Le PPCM des dénominateurs est $\text{PPCM}(19, 1787, 3191582) = 3191582$.
En réduisant au même dénominateur :
\begin{itemize}
    \item $\frac{1}{19} = \frac{167978}{3191582}$
    \item $\frac{1}{1787} = \frac{1786}{3191582}$
    \item $\frac{1}{3191582} = \frac{1}{3191582}$
\end{itemize}
La somme des numérateurs est :
$$ \frac{1}{19} + \frac{1}{1787} + \frac{1}{3191582} = \frac{167978 + 1786 + 1}{3191582} = \frac{169765}{3191582} $$
Le PGCD du numérateur et du dénominateur est $\text{PGCD}(169765, 3191582) = 33953$.
La fraction irréductible est :
$$ \frac{169765}{3191582} = \frac{169765 \div 33953}{3191582 \div 33953} = \frac{5}{94} $$
Cette fraction est bien égale à $\frac{5}{94}$.

\subsection{Démonstration pour $n = 95$}
Soit $n = 95$. Nous cherchons $x, y, z \in \mathbb{Z}^{+}$ tels que $\frac{5}{95} = \frac{1}{x} + \frac{1}{y} + \frac{1}{z}$.
Posons $x = 20$, $y = 381$, $z = 144780$.
Les conditions $x > 0$, $y > 0$ et $z > 0$ sont satisfaites car $x, y, z \geq 1$.
Le PPCM des dénominateurs est $\text{PPCM}(20, 381, 144780) = 144780$.
En réduisant au même dénominateur :
\begin{itemize}
    \item $\frac{1}{20} = \frac{7239}{144780}$
    \item $\frac{1}{381} = \frac{380}{144780}$
    \item $\frac{1}{144780} = \frac{1}{144780}$
\end{itemize}
La somme des numérateurs est :
$$ \frac{1}{20} + \frac{1}{381} + \frac{1}{144780} = \frac{7239 + 380 + 1}{144780} = \frac{7620}{144780} $$
Le PGCD du numérateur et du dénominateur est $\text{PGCD}(7620, 144780) = 7620$.
La fraction irréductible est :
$$ \frac{7620}{144780} = \frac{7620 \div 7620}{144780 \div 7620} = \frac{1}{19} $$
Cette fraction est bien égale à $\frac{5}{95}$.

\subsection{Démonstration pour $n = 96$}
Soit $n = 96$. Nous cherchons $x, y, z \in \mathbb{Z}^{+}$ tels que $\frac{5}{96} = \frac{1}{x} + \frac{1}{y} + \frac{1}{z}$.
Posons $x = 20$, $y = 481$, $z = 230880$.
Les conditions $x > 0$, $y > 0$ et $z > 0$ sont satisfaites car $x, y, z \geq 1$.
Le PPCM des dénominateurs est $\text{PPCM}(20, 481, 230880) = 230880$.
En réduisant au même dénominateur :
\begin{itemize}
    \item $\frac{1}{20} = \frac{11544}{230880}$
    \item $\frac{1}{481} = \frac{480}{230880}$
    \item $\frac{1}{230880} = \frac{1}{230880}$
\end{itemize}
La somme des numérateurs est :
$$ \frac{1}{20} + \frac{1}{481} + \frac{1}{230880} = \frac{11544 + 480 + 1}{230880} = \frac{12025}{230880} $$
Le PGCD du numérateur et du dénominateur est $\text{PGCD}(12025, 230880) = 2405$.
La fraction irréductible est :
$$ \frac{12025}{230880} = \frac{12025 \div 2405}{230880 \div 2405} = \frac{5}{96} $$
Cette fraction est bien égale à $\frac{5}{96}$.

\subsection{Démonstration pour $n = 97$}
Soit $n = 97$. Nous cherchons $x, y, z \in \mathbb{Z}^{+}$ tels que $\frac{5}{97} = \frac{1}{x} + \frac{1}{y} + \frac{1}{z}$.
Posons $x = 20$, $y = 647$, $z = 1255180$.
Les conditions $x > 0$, $y > 0$ et $z > 0$ sont satisfaites car $x, y, z \geq 1$.
Le PPCM des dénominateurs est $\text{PPCM}(20, 647, 1255180) = 1255180$.
En réduisant au même dénominateur :
\begin{itemize}
    \item $\frac{1}{20} = \frac{62759}{1255180}$
    \item $\frac{1}{647} = \frac{1940}{1255180}$
    \item $\frac{1}{1255180} = \frac{1}{1255180}$
\end{itemize}
La somme des numérateurs est :
$$ \frac{1}{20} + \frac{1}{647} + \frac{1}{1255180} = \frac{62759 + 1940 + 1}{1255180} = \frac{64700}{1255180} $$
Le PGCD du numérateur et du dénominateur est $\text{PGCD}(64700, 1255180) = 12940$.
La fraction irréductible est :
$$ \frac{64700}{1255180} = \frac{64700 \div 12940}{1255180 \div 12940} = \frac{5}{97} $$
Cette fraction est bien égale à $\frac{5}{97}$.

\subsection{Démonstration pour $n = 98$}
Soit $n = 98$. Nous cherchons $x, y, z \in \mathbb{Z}^{+}$ tels que $\frac{5}{98} = \frac{1}{x} + \frac{1}{y} + \frac{1}{z}$.
Posons $x = 20$, $y = 981$, $z = 961380$.
Les conditions $x > 0$, $y > 0$ et $z > 0$ sont satisfaites car $x, y, z \geq 1$.
Le PPCM des dénominateurs est $\text{PPCM}(20, 981, 961380) = 961380$.
En réduisant au même dénominateur :
\begin{itemize}
    \item $\frac{1}{20} = \frac{48069}{961380}$
    \item $\frac{1}{981} = \frac{980}{961380}$
    \item $\frac{1}{961380} = \frac{1}{961380}$
\end{itemize}
La somme des numérateurs est :
$$ \frac{1}{20} + \frac{1}{981} + \frac{1}{961380} = \frac{48069 + 980 + 1}{961380} = \frac{49050}{961380} $$
Le PGCD du numérateur et du dénominateur est $\text{PGCD}(49050, 961380) = 9810$.
La fraction irréductible est :
$$ \frac{49050}{961380} = \frac{49050 \div 9810}{961380 \div 9810} = \frac{5}{98} $$
Cette fraction est bien égale à $\frac{5}{98}$.

\subsection{Démonstration pour $n = 99$}
Soit $n = 99$. Nous cherchons $x, y, z \in \mathbb{Z}^{+}$ tels que $\frac{5}{99} = \frac{1}{x} + \frac{1}{y} + \frac{1}{z}$.
Posons $x = 20$, $y = 1981$, $z = 3922380$.
Les conditions $x > 0$, $y > 0$ et $z > 0$ sont satisfaites car $x, y, z \geq 1$.
Le PPCM des dénominateurs est $\text{PPCM}(20, 1981, 3922380) = 3922380$.
En réduisant au même dénominateur :
\begin{itemize}
    \item $\frac{1}{20} = \frac{196119}{3922380}$
    \item $\frac{1}{1981} = \frac{1980}{3922380}$
    \item $\frac{1}{3922380} = \frac{1}{3922380}$
\end{itemize}
La somme des numérateurs est :
$$ \frac{1}{20} + \frac{1}{1981} + \frac{1}{3922380} = \frac{196119 + 1980 + 1}{3922380} = \frac{198100}{3922380} $$
Le PGCD du numérateur et du dénominateur est $\text{PGCD}(198100, 3922380) = 39620$.
La fraction irréductible est :
$$ \frac{198100}{3922380} = \frac{198100 \div 39620}{3922380 \div 39620} = \frac{5}{99} $$
Cette fraction est bien égale à $\frac{5}{99}$.

\subsection{Démonstration pour $n = 100$}
Soit $n = 100$. Nous cherchons $x, y, z \in \mathbb{Z}^{+}$ tels que $\frac{5}{100} = \frac{1}{x} + \frac{1}{y} + \frac{1}{z}$.
Posons $x = 21$, $y = 421$, $z = 176820$.
Les conditions $x > 0$, $y > 0$ et $z > 0$ sont satisfaites car $x, y, z \geq 1$.
Le PPCM des dénominateurs est $\text{PPCM}(21, 421, 176820) = 176820$.
En réduisant au même dénominateur :
\begin{itemize}
    \item $\frac{1}{21} = \frac{8420}{176820}$
    \item $\frac{1}{421} = \frac{420}{176820}$
    \item $\frac{1}{176820} = \frac{1}{176820}$
\end{itemize}
La somme des numérateurs est :
$$ \frac{1}{21} + \frac{1}{421} + \frac{1}{176820} = \frac{8420 + 420 + 1}{176820} = \frac{8841}{176820} $$
Le PGCD du numérateur et du dénominateur est $\text{PGCD}(8841, 176820) = 8841$.
La fraction irréductible est :
$$ \frac{8841}{176820} = \frac{8841 \div 8841}{176820 \div 8841} = \frac{1}{20} $$
Cette fraction est bien égale à $\frac{5}{100}$.

\subsection{Démonstration pour $n = 101$}
Soit $n = 101$. Nous cherchons $x, y, z \in \mathbb{Z}^{+}$ tels que $\frac{5}{101} = \frac{1}{x} + \frac{1}{y} + \frac{1}{z}$.
Posons $x = 21$, $y = 531$, $z = 375417$.
Les conditions $x > 0$, $y > 0$ et $z > 0$ sont satisfaites car $x, y, z \geq 1$.
Le PPCM des dénominateurs est $\text{PPCM}(21, 531, 375417) = 375417$.
En réduisant au même dénominateur :
\begin{itemize}
    \item $\frac{1}{21} = \frac{17877}{375417}$
    \item $\frac{1}{531} = \frac{707}{375417}$
    \item $\frac{1}{375417} = \frac{1}{375417}$
\end{itemize}
La somme des numérateurs est :
$$ \frac{1}{21} + \frac{1}{531} + \frac{1}{375417} = \frac{17877 + 707 + 1}{375417} = \frac{18585}{375417} $$
Le PGCD du numérateur et du dénominateur est $\text{PGCD}(18585, 375417) = 3717$.
La fraction irréductible est :
$$ \frac{18585}{375417} = \frac{18585 \div 3717}{375417 \div 3717} = \frac{5}{101} $$
Cette fraction est bien égale à $\frac{5}{101}$.

\subsection{Démonstration pour $n = 102$}
Soit $n = 102$. Nous cherchons $x, y, z \in \mathbb{Z}^{+}$ tels que $\frac{5}{102} = \frac{1}{x} + \frac{1}{y} + \frac{1}{z}$.
Posons $x = 21$, $y = 715$, $z = 510510$.
Les conditions $x > 0$, $y > 0$ et $z > 0$ sont satisfaites car $x, y, z \geq 1$.
Le PPCM des dénominateurs est $\text{PPCM}(21, 715, 510510) = 510510$.
En réduisant au même dénominateur :
\begin{itemize}
    \item $\frac{1}{21} = \frac{24310}{510510}$
    \item $\frac{1}{715} = \frac{714}{510510}$
    \item $\frac{1}{510510} = \frac{1}{510510}$
\end{itemize}
La somme des numérateurs est :
$$ \frac{1}{21} + \frac{1}{715} + \frac{1}{510510} = \frac{24310 + 714 + 1}{510510} = \frac{25025}{510510} $$
Le PGCD du numérateur et du dénominateur est $\text{PGCD}(25025, 510510) = 5005$.
La fraction irréductible est :
$$ \frac{25025}{510510} = \frac{25025 \div 5005}{510510 \div 5005} = \frac{5}{102} $$
Cette fraction est bien égale à $\frac{5}{102}$.

\subsection{Démonstration pour $n = 103$}
Soit $n = 103$. Nous cherchons $x, y, z \in \mathbb{Z}^{+}$ tels que $\frac{5}{103} = \frac{1}{x} + \frac{1}{y} + \frac{1}{z}$.
Posons $x = 21$, $y = 1082$, $z = 2340366$.
Les conditions $x > 0$, $y > 0$ et $z > 0$ sont satisfaites car $x, y, z \geq 1$.
Le PPCM des dénominateurs est $\text{PPCM}(21, 1082, 2340366) = 2340366$.
En réduisant au même dénominateur :
\begin{itemize}
    \item $\frac{1}{21} = \frac{111446}{2340366}$
    \item $\frac{1}{1082} = \frac{2163}{2340366}$
    \item $\frac{1}{2340366} = \frac{1}{2340366}$
\end{itemize}
La somme des numérateurs est :
$$ \frac{1}{21} + \frac{1}{1082} + \frac{1}{2340366} = \frac{111446 + 2163 + 1}{2340366} = \frac{113610}{2340366} $$
Le PGCD du numérateur et du dénominateur est $\text{PGCD}(113610, 2340366) = 22722$.
La fraction irréductible est :
$$ \frac{113610}{2340366} = \frac{113610 \div 22722}{2340366 \div 22722} = \frac{5}{103} $$
Cette fraction est bien égale à $\frac{5}{103}$.

\subsection{Démonstration pour $n = 104$}
Soit $n = 104$. Nous cherchons $x, y, z \in \mathbb{Z}^{+}$ tels que $\frac{5}{104} = \frac{1}{x} + \frac{1}{y} + \frac{1}{z}$.
Posons $x = 21$, $y = 2185$, $z = 4772040$.
Les conditions $x > 0$, $y > 0$ et $z > 0$ sont satisfaites car $x, y, z \geq 1$.
Le PPCM des dénominateurs est $\text{PPCM}(21, 2185, 4772040) = 4772040$.
En réduisant au même dénominateur :
\begin{itemize}
    \item $\frac{1}{21} = \frac{227240}{4772040}$
    \item $\frac{1}{2185} = \frac{2184}{4772040}$
    \item $\frac{1}{4772040} = \frac{1}{4772040}$
\end{itemize}
La somme des numérateurs est :
$$ \frac{1}{21} + \frac{1}{2185} + \frac{1}{4772040} = \frac{227240 + 2184 + 1}{4772040} = \frac{229425}{4772040} $$
Le PGCD du numérateur et du dénominateur est $\text{PGCD}(229425, 4772040) = 45885$.
La fraction irréductible est :
$$ \frac{229425}{4772040} = \frac{229425 \div 45885}{4772040 \div 45885} = \frac{5}{104} $$
Cette fraction est bien égale à $\frac{5}{104}$.

\subsection{Démonstration pour $n = 105$}
Soit $n = 105$. Nous cherchons $x, y, z \in \mathbb{Z}^{+}$ tels que $\frac{5}{105} = \frac{1}{x} + \frac{1}{y} + \frac{1}{z}$.
Posons $x = 22$, $y = 463$, $z = 213906$.
Les conditions $x > 0$, $y > 0$ et $z > 0$ sont satisfaites car $x, y, z \geq 1$.
Le PPCM des dénominateurs est $\text{PPCM}(22, 463, 213906) = 213906$.
En réduisant au même dénominateur :
\begin{itemize}
    \item $\frac{1}{22} = \frac{9723}{213906}$
    \item $\frac{1}{463} = \frac{462}{213906}$
    \item $\frac{1}{213906} = \frac{1}{213906}$
\end{itemize}
La somme des numérateurs est :
$$ \frac{1}{22} + \frac{1}{463} + \frac{1}{213906} = \frac{9723 + 462 + 1}{213906} = \frac{10186}{213906} $$
Le PGCD du numérateur et du dénominateur est $\text{PGCD}(10186, 213906) = 10186$.
La fraction irréductible est :
$$ \frac{10186}{213906} = \frac{10186 \div 10186}{213906 \div 10186} = \frac{1}{21} $$
Cette fraction est bien égale à $\frac{5}{105}$.

\subsection{Démonstration pour $n = 106$}
Soit $n = 106$. Nous cherchons $x, y, z \in \mathbb{Z}^{+}$ tels que $\frac{5}{106} = \frac{1}{x} + \frac{1}{y} + \frac{1}{z}$.
Posons $x = 22$, $y = 584$, $z = 340472$.
Les conditions $x > 0$, $y > 0$ et $z > 0$ sont satisfaites car $x, y, z \geq 1$.
Le PPCM des dénominateurs est $\text{PPCM}(22, 584, 340472) = 340472$.
En réduisant au même dénominateur :
\begin{itemize}
    \item $\frac{1}{22} = \frac{15476}{340472}$
    \item $\frac{1}{584} = \frac{583}{340472}$
    \item $\frac{1}{340472} = \frac{1}{340472}$
\end{itemize}
La somme des numérateurs est :
$$ \frac{1}{22} + \frac{1}{584} + \frac{1}{340472} = \frac{15476 + 583 + 1}{340472} = \frac{16060}{340472} $$
Le PGCD du numérateur et du dénominateur est $\text{PGCD}(16060, 340472) = 3212$.
La fraction irréductible est :
$$ \frac{16060}{340472} = \frac{16060 \div 3212}{340472 \div 3212} = \frac{5}{106} $$
Cette fraction est bien égale à $\frac{5}{106}$.

\subsection{Démonstration pour $n = 107$}
Soit $n = 107$. Nous cherchons $x, y, z \in \mathbb{Z}^{+}$ tels que $\frac{5}{107} = \frac{1}{x} + \frac{1}{y} + \frac{1}{z}$.
Posons $x = 22$, $y = 785$, $z = 1847890$.
Les conditions $x > 0$, $y > 0$ et $z > 0$ sont satisfaites car $x, y, z \geq 1$.
Le PPCM des dénominateurs est $\text{PPCM}(22, 785, 1847890) = 1847890$.
En réduisant au même dénominateur :
\begin{itemize}
    \item $\frac{1}{22} = \frac{83995}{1847890}$
    \item $\frac{1}{785} = \frac{2354}{1847890}$
    \item $\frac{1}{1847890} = \frac{1}{1847890}$
\end{itemize}
La somme des numérateurs est :
$$ \frac{1}{22} + \frac{1}{785} + \frac{1}{1847890} = \frac{83995 + 2354 + 1}{1847890} = \frac{86350}{1847890} $$
Le PGCD du numérateur et du dénominateur est $\text{PGCD}(86350, 1847890) = 17270$.
La fraction irréductible est :
$$ \frac{86350}{1847890} = \frac{86350 \div 17270}{1847890 \div 17270} = \frac{5}{107} $$
Cette fraction est bien égale à $\frac{5}{107}$.

\subsection{Démonstration pour $n = 108$}
Soit $n = 108$. Nous cherchons $x, y, z \in \mathbb{Z}^{+}$ tels que $\frac{5}{108} = \frac{1}{x} + \frac{1}{y} + \frac{1}{z}$.
Posons $x = 22$, $y = 1189$, $z = 1412532$.
Les conditions $x > 0$, $y > 0$ et $z > 0$ sont satisfaites car $x, y, z \geq 1$.
Le PPCM des dénominateurs est $\text{PPCM}(22, 1189, 1412532) = 1412532$.
En réduisant au même dénominateur :
\begin{itemize}
    \item $\frac{1}{22} = \frac{64206}{1412532}$
    \item $\frac{1}{1189} = \frac{1188}{1412532}$
    \item $\frac{1}{1412532} = \frac{1}{1412532}$
\end{itemize}
La somme des numérateurs est :
$$ \frac{1}{22} + \frac{1}{1189} + \frac{1}{1412532} = \frac{64206 + 1188 + 1}{1412532} = \frac{65395}{1412532} $$
Le PGCD du numérateur et du dénominateur est $\text{PGCD}(65395, 1412532) = 13079$.
La fraction irréductible est :
$$ \frac{65395}{1412532} = \frac{65395 \div 13079}{1412532 \div 13079} = \frac{5}{108} $$
Cette fraction est bien égale à $\frac{5}{108}$.

\subsection{Démonstration pour $n = 109$}
Soit $n = 109$. Nous cherchons $x, y, z \in \mathbb{Z}^{+}$ tels que $\frac{5}{109} = \frac{1}{x} + \frac{1}{y} + \frac{1}{z}$.
Posons $x = 22$, $y = 2399$, $z = 5752802$.
Les conditions $x > 0$, $y > 0$ et $z > 0$ sont satisfaites car $x, y, z \geq 1$.
Le PPCM des dénominateurs est $\text{PPCM}(22, 2399, 5752802) = 5752802$.
En réduisant au même dénominateur :
\begin{itemize}
    \item $\frac{1}{22} = \frac{261491}{5752802}$
    \item $\frac{1}{2399} = \frac{2398}{5752802}$
    \item $\frac{1}{5752802} = \frac{1}{5752802}$
\end{itemize}
La somme des numérateurs est :
$$ \frac{1}{22} + \frac{1}{2399} + \frac{1}{5752802} = \frac{261491 + 2398 + 1}{5752802} = \frac{263890}{5752802} $$
Le PGCD du numérateur et du dénominateur est $\text{PGCD}(263890, 5752802) = 52778$.
La fraction irréductible est :
$$ \frac{263890}{5752802} = \frac{263890 \div 52778}{5752802 \div 52778} = \frac{5}{109} $$
Cette fraction est bien égale à $\frac{5}{109}$.

\subsection{Démonstration pour $n = 110$}
Soit $n = 110$. Nous cherchons $x, y, z \in \mathbb{Z}^{+}$ tels que $\frac{5}{110} = \frac{1}{x} + \frac{1}{y} + \frac{1}{z}$.
Posons $x = 23$, $y = 507$, $z = 256542$.
Les conditions $x > 0$, $y > 0$ et $z > 0$ sont satisfaites car $x, y, z \geq 1$.
Le PPCM des dénominateurs est $\text{PPCM}(23, 507, 256542) = 256542$.
En réduisant au même dénominateur :
\begin{itemize}
    \item $\frac{1}{23} = \frac{11154}{256542}$
    \item $\frac{1}{507} = \frac{506}{256542}$
    \item $\frac{1}{256542} = \frac{1}{256542}$
\end{itemize}
La somme des numérateurs est :
$$ \frac{1}{23} + \frac{1}{507} + \frac{1}{256542} = \frac{11154 + 506 + 1}{256542} = \frac{11661}{256542} $$
Le PGCD du numérateur et du dénominateur est $\text{PGCD}(11661, 256542) = 11661$.
La fraction irréductible est :
$$ \frac{11661}{256542} = \frac{11661 \div 11661}{256542 \div 11661} = \frac{1}{22} $$
Cette fraction est bien égale à $\frac{5}{110}$.

\subsection{Démonstration pour $n = 111$}
Soit $n = 111$. Nous cherchons $x, y, z \in \mathbb{Z}^{+}$ tels que $\frac{5}{111} = \frac{1}{x} + \frac{1}{y} + \frac{1}{z}$.
Posons $x = 23$, $y = 639$, $z = 543789$.
Les conditions $x > 0$, $y > 0$ et $z > 0$ sont satisfaites car $x, y, z \geq 1$.
Le PPCM des dénominateurs est $\text{PPCM}(23, 639, 543789) = 543789$.
En réduisant au même dénominateur :
\begin{itemize}
    \item $\frac{1}{23} = \frac{23643}{543789}$
    \item $\frac{1}{639} = \frac{851}{543789}$
    \item $\frac{1}{543789} = \frac{1}{543789}$
\end{itemize}
La somme des numérateurs est :
$$ \frac{1}{23} + \frac{1}{639} + \frac{1}{543789} = \frac{23643 + 851 + 1}{543789} = \frac{24495}{543789} $$
Le PGCD du numérateur et du dénominateur est $\text{PGCD}(24495, 543789) = 4899$.
La fraction irréductible est :
$$ \frac{24495}{543789} = \frac{24495 \div 4899}{543789 \div 4899} = \frac{5}{111} $$
Cette fraction est bien égale à $\frac{5}{111}$.

\subsection{Démonstration pour $n = 112$}
Soit $n = 112$. Nous cherchons $x, y, z \in \mathbb{Z}^{+}$ tels que $\frac{5}{112} = \frac{1}{x} + \frac{1}{y} + \frac{1}{z}$.
Posons $x = 23$, $y = 859$, $z = 2212784$.
Les conditions $x > 0$, $y > 0$ et $z > 0$ sont satisfaites car $x, y, z \geq 1$.
Le PPCM des dénominateurs est $\text{PPCM}(23, 859, 2212784) = 2212784$.
En réduisant au même dénominateur :
\begin{itemize}
    \item $\frac{1}{23} = \frac{96208}{2212784}$
    \item $\frac{1}{859} = \frac{2576}{2212784}$
    \item $\frac{1}{2212784} = \frac{1}{2212784}$
\end{itemize}
La somme des numérateurs est :
$$ \frac{1}{23} + \frac{1}{859} + \frac{1}{2212784} = \frac{96208 + 2576 + 1}{2212784} = \frac{98785}{2212784} $$
Le PGCD du numérateur et du dénominateur est $\text{PGCD}(98785, 2212784) = 19757$.
La fraction irréductible est :
$$ \frac{98785}{2212784} = \frac{98785 \div 19757}{2212784 \div 19757} = \frac{5}{112} $$
Cette fraction est bien égale à $\frac{5}{112}$.

\subsection{Démonstration pour $n = 113$}
Soit $n = 113$. Nous cherchons $x, y, z \in \mathbb{Z}^{+}$ tels que $\frac{5}{113} = \frac{1}{x} + \frac{1}{y} + \frac{1}{z}$.
Posons $x = 23$, $y = 1300$, $z = 3378700$.
Les conditions $x > 0$, $y > 0$ et $z > 0$ sont satisfaites car $x, y, z \geq 1$.
Le PPCM des dénominateurs est $\text{PPCM}(23, 1300, 3378700) = 3378700$.
En réduisant au même dénominateur :
\begin{itemize}
    \item $\frac{1}{23} = \frac{146900}{3378700}$
    \item $\frac{1}{1300} = \frac{2599}{3378700}$
    \item $\frac{1}{3378700} = \frac{1}{3378700}$
\end{itemize}
La somme des numérateurs est :
$$ \frac{1}{23} + \frac{1}{1300} + \frac{1}{3378700} = \frac{146900 + 2599 + 1}{3378700} = \frac{149500}{3378700} $$
Le PGCD du numérateur et du dénominateur est $\text{PGCD}(149500, 3378700) = 29900$.
La fraction irréductible est :
$$ \frac{149500}{3378700} = \frac{149500 \div 29900}{3378700 \div 29900} = \frac{5}{113} $$
Cette fraction est bien égale à $\frac{5}{113}$.

\subsection{Démonstration pour $n = 114$}
Soit $n = 114$. Nous cherchons $x, y, z \in \mathbb{Z}^{+}$ tels que $\frac{5}{114} = \frac{1}{x} + \frac{1}{y} + \frac{1}{z}$.
Posons $x = 23$, $y = 2623$, $z = 6877506$.
Les conditions $x > 0$, $y > 0$ et $z > 0$ sont satisfaites car $x, y, z \geq 1$.
Le PPCM des dénominateurs est $\text{PPCM}(23, 2623, 6877506) = 6877506$.
En réduisant au même dénominateur :
\begin{itemize}
    \item $\frac{1}{23} = \frac{299022}{6877506}$
    \item $\frac{1}{2623} = \frac{2622}{6877506}$
    \item $\frac{1}{6877506} = \frac{1}{6877506}$
\end{itemize}
La somme des numérateurs est :
$$ \frac{1}{23} + \frac{1}{2623} + \frac{1}{6877506} = \frac{299022 + 2622 + 1}{6877506} = \frac{301645}{6877506} $$
Le PGCD du numérateur et du dénominateur est $\text{PGCD}(301645, 6877506) = 60329$.
La fraction irréductible est :
$$ \frac{301645}{6877506} = \frac{301645 \div 60329}{6877506 \div 60329} = \frac{5}{114} $$
Cette fraction est bien égale à $\frac{5}{114}$.

\subsection{Démonstration pour $n = 115$}
Soit $n = 115$. Nous cherchons $x, y, z \in \mathbb{Z}^{+}$ tels que $\frac{5}{115} = \frac{1}{x} + \frac{1}{y} + \frac{1}{z}$.
Posons $x = 24$, $y = 553$, $z = 305256$.
Les conditions $x > 0$, $y > 0$ et $z > 0$ sont satisfaites car $x, y, z \geq 1$.
Le PPCM des dénominateurs est $\text{PPCM}(24, 553, 305256) = 305256$.
En réduisant au même dénominateur :
\begin{itemize}
    \item $\frac{1}{24} = \frac{12719}{305256}$
    \item $\frac{1}{553} = \frac{552}{305256}$
    \item $\frac{1}{305256} = \frac{1}{305256}$
\end{itemize}
La somme des numérateurs est :
$$ \frac{1}{24} + \frac{1}{553} + \frac{1}{305256} = \frac{12719 + 552 + 1}{305256} = \frac{13272}{305256} $$
Le PGCD du numérateur et du dénominateur est $\text{PGCD}(13272, 305256) = 13272$.
La fraction irréductible est :
$$ \frac{13272}{305256} = \frac{13272 \div 13272}{305256 \div 13272} = \frac{1}{23} $$
Cette fraction est bien égale à $\frac{5}{115}$.

\subsection{Démonstration pour $n = 116$}
Soit $n = 116$. Nous cherchons $x, y, z \in \mathbb{Z}^{+}$ tels que $\frac{5}{116} = \frac{1}{x} + \frac{1}{y} + \frac{1}{z}$.
Posons $x = 24$, $y = 697$, $z = 485112$.
Les conditions $x > 0$, $y > 0$ et $z > 0$ sont satisfaites car $x, y, z \geq 1$.
Le PPCM des dénominateurs est $\text{PPCM}(24, 697, 485112) = 485112$.
En réduisant au même dénominateur :
\begin{itemize}
    \item $\frac{1}{24} = \frac{20213}{485112}$
    \item $\frac{1}{697} = \frac{696}{485112}$
    \item $\frac{1}{485112} = \frac{1}{485112}$
\end{itemize}
La somme des numérateurs est :
$$ \frac{1}{24} + \frac{1}{697} + \frac{1}{485112} = \frac{20213 + 696 + 1}{485112} = \frac{20910}{485112} $$
Le PGCD du numérateur et du dénominateur est $\text{PGCD}(20910, 485112) = 4182$.
La fraction irréductible est :
$$ \frac{20910}{485112} = \frac{20910 \div 4182}{485112 \div 4182} = \frac{5}{116} $$
Cette fraction est bien égale à $\frac{5}{116}$.

\subsection{Démonstration pour $n = 117$}
Soit $n = 117$. Nous cherchons $x, y, z \in \mathbb{Z}^{+}$ tels que $\frac{5}{117} = \frac{1}{x} + \frac{1}{y} + \frac{1}{z}$.
Posons $x = 24$, $y = 937$, $z = 877032$.
Les conditions $x > 0$, $y > 0$ et $z > 0$ sont satisfaites car $x, y, z \geq 1$.
Le PPCM des dénominateurs est $\text{PPCM}(24, 937, 877032) = 877032$.
En réduisant au même dénominateur :
\begin{itemize}
    \item $\frac{1}{24} = \frac{36543}{877032}$
    \item $\frac{1}{937} = \frac{936}{877032}$
    \item $\frac{1}{877032} = \frac{1}{877032}$
\end{itemize}
La somme des numérateurs est :
$$ \frac{1}{24} + \frac{1}{937} + \frac{1}{877032} = \frac{36543 + 936 + 1}{877032} = \frac{37480}{877032} $$
Le PGCD du numérateur et du dénominateur est $\text{PGCD}(37480, 877032) = 7496$.
La fraction irréductible est :
$$ \frac{37480}{877032} = \frac{37480 \div 7496}{877032 \div 7496} = \frac{5}{117} $$
Cette fraction est bien égale à $\frac{5}{117}$.

\subsection{Démonstration pour $n = 118$}
Soit $n = 118$. Nous cherchons $x, y, z \in \mathbb{Z}^{+}$ tels que $\frac{5}{118} = \frac{1}{x} + \frac{1}{y} + \frac{1}{z}$.
Posons $x = 24$, $y = 1417$, $z = 2006472$.
Les conditions $x > 0$, $y > 0$ et $z > 0$ sont satisfaites car $x, y, z \geq 1$.
Le PPCM des dénominateurs est $\text{PPCM}(24, 1417, 2006472) = 2006472$.
En réduisant au même dénominateur :
\begin{itemize}
    \item $\frac{1}{24} = \frac{83603}{2006472}$
    \item $\frac{1}{1417} = \frac{1416}{2006472}$
    \item $\frac{1}{2006472} = \frac{1}{2006472}$
\end{itemize}
La somme des numérateurs est :
$$ \frac{1}{24} + \frac{1}{1417} + \frac{1}{2006472} = \frac{83603 + 1416 + 1}{2006472} = \frac{85020}{2006472} $$
Le PGCD du numérateur et du dénominateur est $\text{PGCD}(85020, 2006472) = 17004$.
La fraction irréductible est :
$$ \frac{85020}{2006472} = \frac{85020 \div 17004}{2006472 \div 17004} = \frac{5}{118} $$
Cette fraction est bien égale à $\frac{5}{118}$.

\subsection{Démonstration pour $n = 119$}
Soit $n = 119$. Nous cherchons $x, y, z \in \mathbb{Z}^{+}$ tels que $\frac{5}{119} = \frac{1}{x} + \frac{1}{y} + \frac{1}{z}$.
Posons $x = 24$, $y = 2857$, $z = 8159592$.
Les conditions $x > 0$, $y > 0$ et $z > 0$ sont satisfaites car $x, y, z \geq 1$.
Le PPCM des dénominateurs est $\text{PPCM}(24, 2857, 8159592) = 8159592$.
En réduisant au même dénominateur :
\begin{itemize}
    \item $\frac{1}{24} = \frac{339983}{8159592}$
    \item $\frac{1}{2857} = \frac{2856}{8159592}$
    \item $\frac{1}{8159592} = \frac{1}{8159592}$
\end{itemize}
La somme des numérateurs est :
$$ \frac{1}{24} + \frac{1}{2857} + \frac{1}{8159592} = \frac{339983 + 2856 + 1}{8159592} = \frac{342840}{8159592} $$
Le PGCD du numérateur et du dénominateur est $\text{PGCD}(342840, 8159592) = 68568$.
La fraction irréductible est :
$$ \frac{342840}{8159592} = \frac{342840 \div 68568}{8159592 \div 68568} = \frac{5}{119} $$
Cette fraction est bien égale à $\frac{5}{119}$.

\subsection{Démonstration pour $n = 120$}
Soit $n = 120$. Nous cherchons $x, y, z \in \mathbb{Z}^{+}$ tels que $\frac{5}{120} = \frac{1}{x} + \frac{1}{y} + \frac{1}{z}$.
Posons $x = 25$, $y = 601$, $z = 360600$.
Les conditions $x > 0$, $y > 0$ et $z > 0$ sont satisfaites car $x, y, z \geq 1$.
Le PPCM des dénominateurs est $\text{PPCM}(25, 601, 360600) = 360600$.
En réduisant au même dénominateur :
\begin{itemize}
    \item $\frac{1}{25} = \frac{14424}{360600}$
    \item $\frac{1}{601} = \frac{600}{360600}$
    \item $\frac{1}{360600} = \frac{1}{360600}$
\end{itemize}
La somme des numérateurs est :
$$ \frac{1}{25} + \frac{1}{601} + \frac{1}{360600} = \frac{14424 + 600 + 1}{360600} = \frac{15025}{360600} $$
Le PGCD du numérateur et du dénominateur est $\text{PGCD}(15025, 360600) = 15025$.
La fraction irréductible est :
$$ \frac{15025}{360600} = \frac{15025 \div 15025}{360600 \div 15025} = \frac{1}{24} $$
Cette fraction est bien égale à $\frac{5}{120}$.

\subsection{Démonstration pour $n = 121$}
Soit $n = 121$. Nous cherchons $x, y, z \in \mathbb{Z}^{+}$ tels que $\frac{5}{121} = \frac{1}{x} + \frac{1}{y} + \frac{1}{z}$.
Posons $x = 25$, $y = 759$, $z = 208725$.
Les conditions $x > 0$, $y > 0$ et $z > 0$ sont satisfaites car $x, y, z \geq 1$.
Le PPCM des dénominateurs est $\text{PPCM}(25, 759, 208725) = 208725$.
En réduisant au même dénominateur :
\begin{itemize}
    \item $\frac{1}{25} = \frac{8349}{208725}$
    \item $\frac{1}{759} = \frac{275}{208725}$
    \item $\frac{1}{208725} = \frac{1}{208725}$
\end{itemize}
La somme des numérateurs est :
$$ \frac{1}{25} + \frac{1}{759} + \frac{1}{208725} = \frac{8349 + 275 + 1}{208725} = \frac{8625}{208725} $$
Le PGCD du numérateur et du dénominateur est $\text{PGCD}(8625, 208725) = 1725$.
La fraction irréductible est :
$$ \frac{8625}{208725} = \frac{8625 \div 1725}{208725 \div 1725} = \frac{5}{121} $$
Cette fraction est bien égale à $\frac{5}{121}$.

\subsection{Démonstration pour $n = 122$}
Soit $n = 122$. Nous cherchons $x, y, z \in \mathbb{Z}^{+}$ tels que $\frac{5}{122} = \frac{1}{x} + \frac{1}{y} + \frac{1}{z}$.
Posons $x = 25$, $y = 1017$, $z = 3101850$.
Les conditions $x > 0$, $y > 0$ et $z > 0$ sont satisfaites car $x, y, z \geq 1$.
Le PPCM des dénominateurs est $\text{PPCM}(25, 1017, 3101850) = 3101850$.
En réduisant au même dénominateur :
\begin{itemize}
    \item $\frac{1}{25} = \frac{124074}{3101850}$
    \item $\frac{1}{1017} = \frac{3050}{3101850}$
    \item $\frac{1}{3101850} = \frac{1}{3101850}$
\end{itemize}
La somme des numérateurs est :
$$ \frac{1}{25} + \frac{1}{1017} + \frac{1}{3101850} = \frac{124074 + 3050 + 1}{3101850} = \frac{127125}{3101850} $$
Le PGCD du numérateur et du dénominateur est $\text{PGCD}(127125, 3101850) = 25425$.
La fraction irréductible est :
$$ \frac{127125}{3101850} = \frac{127125 \div 25425}{3101850 \div 25425} = \frac{5}{122} $$
Cette fraction est bien égale à $\frac{5}{122}$.

\subsection{Démonstration pour $n = 123$}
Soit $n = 123$. Nous cherchons $x, y, z \in \mathbb{Z}^{+}$ tels que $\frac{5}{123} = \frac{1}{x} + \frac{1}{y} + \frac{1}{z}$.
Posons $x = 25$, $y = 1538$, $z = 4729350$.
Les conditions $x > 0$, $y > 0$ et $z > 0$ sont satisfaites car $x, y, z \geq 1$.
Le PPCM des dénominateurs est $\text{PPCM}(25, 1538, 4729350) = 4729350$.
En réduisant au même dénominateur :
\begin{itemize}
    \item $\frac{1}{25} = \frac{189174}{4729350}$
    \item $\frac{1}{1538} = \frac{3075}{4729350}$
    \item $\frac{1}{4729350} = \frac{1}{4729350}$
\end{itemize}
La somme des numérateurs est :
$$ \frac{1}{25} + \frac{1}{1538} + \frac{1}{4729350} = \frac{189174 + 3075 + 1}{4729350} = \frac{192250}{4729350} $$
Le PGCD du numérateur et du dénominateur est $\text{PGCD}(192250, 4729350) = 38450$.
La fraction irréductible est :
$$ \frac{192250}{4729350} = \frac{192250 \div 38450}{4729350 \div 38450} = \frac{5}{123} $$
Cette fraction est bien égale à $\frac{5}{123}$.

\subsection{Démonstration pour $n = 124$}
Soit $n = 124$. Nous cherchons $x, y, z \in \mathbb{Z}^{+}$ tels que $\frac{5}{124} = \frac{1}{x} + \frac{1}{y} + \frac{1}{z}$.
Posons $x = 25$, $y = 3101$, $z = 9613100$.
Les conditions $x > 0$, $y > 0$ et $z > 0$ sont satisfaites car $x, y, z \geq 1$.
Le PPCM des dénominateurs est $\text{PPCM}(25, 3101, 9613100) = 9613100$.
En réduisant au même dénominateur :
\begin{itemize}
    \item $\frac{1}{25} = \frac{384524}{9613100}$
    \item $\frac{1}{3101} = \frac{3100}{9613100}$
    \item $\frac{1}{9613100} = \frac{1}{9613100}$
\end{itemize}
La somme des numérateurs est :
$$ \frac{1}{25} + \frac{1}{3101} + \frac{1}{9613100} = \frac{384524 + 3100 + 1}{9613100} = \frac{387625}{9613100} $$
Le PGCD du numérateur et du dénominateur est $\text{PGCD}(387625, 9613100) = 77525$.
La fraction irréductible est :
$$ \frac{387625}{9613100} = \frac{387625 \div 77525}{9613100 \div 77525} = \frac{5}{124} $$
Cette fraction est bien égale à $\frac{5}{124}$.

\subsection{Démonstration pour $n = 125$}
Soit $n = 125$. Nous cherchons $x, y, z \in \mathbb{Z}^{+}$ tels que $\frac{5}{125} = \frac{1}{x} + \frac{1}{y} + \frac{1}{z}$.
Posons $x = 26$, $y = 651$, $z = 423150$.
Les conditions $x > 0$, $y > 0$ et $z > 0$ sont satisfaites car $x, y, z \geq 1$.
Le PPCM des dénominateurs est $\text{PPCM}(26, 651, 423150) = 423150$.
En réduisant au même dénominateur :
\begin{itemize}
    \item $\frac{1}{26} = \frac{16275}{423150}$
    \item $\frac{1}{651} = \frac{650}{423150}$
    \item $\frac{1}{423150} = \frac{1}{423150}$
\end{itemize}
La somme des numérateurs est :
$$ \frac{1}{26} + \frac{1}{651} + \frac{1}{423150} = \frac{16275 + 650 + 1}{423150} = \frac{16926}{423150} $$
Le PGCD du numérateur et du dénominateur est $\text{PGCD}(16926, 423150) = 16926$.
La fraction irréductible est :
$$ \frac{16926}{423150} = \frac{16926 \div 16926}{423150 \div 16926} = \frac{1}{25} $$
Cette fraction est bien égale à $\frac{5}{125}$.

\subsection{Démonstration pour $n = 126$}
Soit $n = 126$. Nous cherchons $x, y, z \in \mathbb{Z}^{+}$ tels que $\frac{5}{126} = \frac{1}{x} + \frac{1}{y} + \frac{1}{z}$.
Posons $x = 26$, $y = 820$, $z = 671580$.
Les conditions $x > 0$, $y > 0$ et $z > 0$ sont satisfaites car $x, y, z \geq 1$.
Le PPCM des dénominateurs est $\text{PPCM}(26, 820, 671580) = 671580$.
En réduisant au même dénominateur :
\begin{itemize}
    \item $\frac{1}{26} = \frac{25830}{671580}$
    \item $\frac{1}{820} = \frac{819}{671580}$
    \item $\frac{1}{671580} = \frac{1}{671580}$
\end{itemize}
La somme des numérateurs est :
$$ \frac{1}{26} + \frac{1}{820} + \frac{1}{671580} = \frac{25830 + 819 + 1}{671580} = \frac{26650}{671580} $$
Le PGCD du numérateur et du dénominateur est $\text{PGCD}(26650, 671580) = 5330$.
La fraction irréductible est :
$$ \frac{26650}{671580} = \frac{26650 \div 5330}{671580 \div 5330} = \frac{5}{126} $$
Cette fraction est bien égale à $\frac{5}{126}$.

\subsection{Démonstration pour $n = 127$}
Soit $n = 127$. Nous cherchons $x, y, z \in \mathbb{Z}^{+}$ tels que $\frac{5}{127} = \frac{1}{x} + \frac{1}{y} + \frac{1}{z}$.
Posons $x = 26$, $y = 1101$, $z = 3635502$.
Les conditions $x > 0$, $y > 0$ et $z > 0$ sont satisfaites car $x, y, z \geq 1$.
Le PPCM des dénominateurs est $\text{PPCM}(26, 1101, 3635502) = 3635502$.
En réduisant au même dénominateur :
\begin{itemize}
    \item $\frac{1}{26} = \frac{139827}{3635502}$
    \item $\frac{1}{1101} = \frac{3302}{3635502}$
    \item $\frac{1}{3635502} = \frac{1}{3635502}$
\end{itemize}
La somme des numérateurs est :
$$ \frac{1}{26} + \frac{1}{1101} + \frac{1}{3635502} = \frac{139827 + 3302 + 1}{3635502} = \frac{143130}{3635502} $$
Le PGCD du numérateur et du dénominateur est $\text{PGCD}(143130, 3635502) = 28626$.
La fraction irréductible est :
$$ \frac{143130}{3635502} = \frac{143130 \div 28626}{3635502 \div 28626} = \frac{5}{127} $$
Cette fraction est bien égale à $\frac{5}{127}$.

\subsection{Démonstration pour $n = 128$}
Soit $n = 128$. Nous cherchons $x, y, z \in \mathbb{Z}^{+}$ tels que $\frac{5}{128} = \frac{1}{x} + \frac{1}{y} + \frac{1}{z}$.
Posons $x = 26$, $y = 1665$, $z = 2770560$.
Les conditions $x > 0$, $y > 0$ et $z > 0$ sont satisfaites car $x, y, z \geq 1$.
Le PPCM des dénominateurs est $\text{PPCM}(26, 1665, 2770560) = 2770560$.
En réduisant au même dénominateur :
\begin{itemize}
    \item $\frac{1}{26} = \frac{106560}{2770560}$
    \item $\frac{1}{1665} = \frac{1664}{2770560}$
    \item $\frac{1}{2770560} = \frac{1}{2770560}$
\end{itemize}
La somme des numérateurs est :
$$ \frac{1}{26} + \frac{1}{1665} + \frac{1}{2770560} = \frac{106560 + 1664 + 1}{2770560} = \frac{108225}{2770560} $$
Le PGCD du numérateur et du dénominateur est $\text{PGCD}(108225, 2770560) = 21645$.
La fraction irréductible est :
$$ \frac{108225}{2770560} = \frac{108225 \div 21645}{2770560 \div 21645} = \frac{5}{128} $$
Cette fraction est bien égale à $\frac{5}{128}$.

\subsection{Démonstration pour $n = 129$}
Soit $n = 129$. Nous cherchons $x, y, z \in \mathbb{Z}^{+}$ tels que $\frac{5}{129} = \frac{1}{x} + \frac{1}{y} + \frac{1}{z}$.
Posons $x = 26$, $y = 3355$, $z = 11252670$.
Les conditions $x > 0$, $y > 0$ et $z > 0$ sont satisfaites car $x, y, z \geq 1$.
Le PPCM des dénominateurs est $\text{PPCM}(26, 3355, 11252670) = 11252670$.
En réduisant au même dénominateur :
\begin{itemize}
    \item $\frac{1}{26} = \frac{432795}{11252670}$
    \item $\frac{1}{3355} = \frac{3354}{11252670}$
    \item $\frac{1}{11252670} = \frac{1}{11252670}$
\end{itemize}
La somme des numérateurs est :
$$ \frac{1}{26} + \frac{1}{3355} + \frac{1}{11252670} = \frac{432795 + 3354 + 1}{11252670} = \frac{436150}{11252670} $$
Le PGCD du numérateur et du dénominateur est $\text{PGCD}(436150, 11252670) = 87230$.
La fraction irréductible est :
$$ \frac{436150}{11252670} = \frac{436150 \div 87230}{11252670 \div 87230} = \frac{5}{129} $$
Cette fraction est bien égale à $\frac{5}{129}$.

\subsection{Démonstration pour $n = 130$}
Soit $n = 130$. Nous cherchons $x, y, z \in \mathbb{Z}^{+}$ tels que $\frac{5}{130} = \frac{1}{x} + \frac{1}{y} + \frac{1}{z}$.
Posons $x = 27$, $y = 703$, $z = 493506$.
Les conditions $x > 0$, $y > 0$ et $z > 0$ sont satisfaites car $x, y, z \geq 1$.
Le PPCM des dénominateurs est $\text{PPCM}(27, 703, 493506) = 493506$.
En réduisant au même dénominateur :
\begin{itemize}
    \item $\frac{1}{27} = \frac{18278}{493506}$
    \item $\frac{1}{703} = \frac{702}{493506}$
    \item $\frac{1}{493506} = \frac{1}{493506}$
\end{itemize}
La somme des numérateurs est :
$$ \frac{1}{27} + \frac{1}{703} + \frac{1}{493506} = \frac{18278 + 702 + 1}{493506} = \frac{18981}{493506} $$
Le PGCD du numérateur et du dénominateur est $\text{PGCD}(18981, 493506) = 18981$.
La fraction irréductible est :
$$ \frac{18981}{493506} = \frac{18981 \div 18981}{493506 \div 18981} = \frac{1}{26} $$
Cette fraction est bien égale à $\frac{5}{130}$.

\subsection{Démonstration pour $n = 131$}
Soit $n = 131$. Nous cherchons $x, y, z \in \mathbb{Z}^{+}$ tels que $\frac{5}{131} = \frac{1}{x} + \frac{1}{y} + \frac{1}{z}$.
Posons $x = 27$, $y = 885$, $z = 1043415$.
Les conditions $x > 0$, $y > 0$ et $z > 0$ sont satisfaites car $x, y, z \geq 1$.
Le PPCM des dénominateurs est $\text{PPCM}(27, 885, 1043415) = 1043415$.
En réduisant au même dénominateur :
\begin{itemize}
    \item $\frac{1}{27} = \frac{38645}{1043415}$
    \item $\frac{1}{885} = \frac{1179}{1043415}$
    \item $\frac{1}{1043415} = \frac{1}{1043415}$
\end{itemize}
La somme des numérateurs est :
$$ \frac{1}{27} + \frac{1}{885} + \frac{1}{1043415} = \frac{38645 + 1179 + 1}{1043415} = \frac{39825}{1043415} $$
Le PGCD du numérateur et du dénominateur est $\text{PGCD}(39825, 1043415) = 7965$.
La fraction irréductible est :
$$ \frac{39825}{1043415} = \frac{39825 \div 7965}{1043415 \div 7965} = \frac{5}{131} $$
Cette fraction est bien égale à $\frac{5}{131}$.

\subsection{Démonstration pour $n = 132$}
Soit $n = 132$. Nous cherchons $x, y, z \in \mathbb{Z}^{+}$ tels que $\frac{5}{132} = \frac{1}{x} + \frac{1}{y} + \frac{1}{z}$.
Posons $x = 27$, $y = 1189$, $z = 1412532$.
Les conditions $x > 0$, $y > 0$ et $z > 0$ sont satisfaites car $x, y, z \geq 1$.
Le PPCM des dénominateurs est $\text{PPCM}(27, 1189, 1412532) = 1412532$.
En réduisant au même dénominateur :
\begin{itemize}
    \item $\frac{1}{27} = \frac{52316}{1412532}$
    \item $\frac{1}{1189} = \frac{1188}{1412532}$
    \item $\frac{1}{1412532} = \frac{1}{1412532}$
\end{itemize}
La somme des numérateurs est :
$$ \frac{1}{27} + \frac{1}{1189} + \frac{1}{1412532} = \frac{52316 + 1188 + 1}{1412532} = \frac{53505}{1412532} $$
Le PGCD du numérateur et du dénominateur est $\text{PGCD}(53505, 1412532) = 10701$.
La fraction irréductible est :
$$ \frac{53505}{1412532} = \frac{53505 \div 10701}{1412532 \div 10701} = \frac{5}{132} $$
Cette fraction est bien égale à $\frac{5}{132}$.

\subsection{Démonstration pour $n = 133$}
Soit $n = 133$. Nous cherchons $x, y, z \in \mathbb{Z}^{+}$ tels que $\frac{5}{133} = \frac{1}{x} + \frac{1}{y} + \frac{1}{z}$.
Posons $x = 27$, $y = 1796$, $z = 6449436$.
Les conditions $x > 0$, $y > 0$ et $z > 0$ sont satisfaites car $x, y, z \geq 1$.
Le PPCM des dénominateurs est $\text{PPCM}(27, 1796, 6449436) = 6449436$.
En réduisant au même dénominateur :
\begin{itemize}
    \item $\frac{1}{27} = \frac{238868}{6449436}$
    \item $\frac{1}{1796} = \frac{3591}{6449436}$
    \item $\frac{1}{6449436} = \frac{1}{6449436}$
\end{itemize}
La somme des numérateurs est :
$$ \frac{1}{27} + \frac{1}{1796} + \frac{1}{6449436} = \frac{238868 + 3591 + 1}{6449436} = \frac{242460}{6449436} $$
Le PGCD du numérateur et du dénominateur est $\text{PGCD}(242460, 6449436) = 48492$.
La fraction irréductible est :
$$ \frac{242460}{6449436} = \frac{242460 \div 48492}{6449436 \div 48492} = \frac{5}{133} $$
Cette fraction est bien égale à $\frac{5}{133}$.

\subsection{Démonstration pour $n = 134$}
Soit $n = 134$. Nous cherchons $x, y, z \in \mathbb{Z}^{+}$ tels que $\frac{5}{134} = \frac{1}{x} + \frac{1}{y} + \frac{1}{z}$.
Posons $x = 27$, $y = 3619$, $z = 13093542$.
Les conditions $x > 0$, $y > 0$ et $z > 0$ sont satisfaites car $x, y, z \geq 1$.
Le PPCM des dénominateurs est $\text{PPCM}(27, 3619, 13093542) = 13093542$.
En réduisant au même dénominateur :
\begin{itemize}
    \item $\frac{1}{27} = \frac{484946}{13093542}$
    \item $\frac{1}{3619} = \frac{3618}{13093542}$
    \item $\frac{1}{13093542} = \frac{1}{13093542}$
\end{itemize}
La somme des numérateurs est :
$$ \frac{1}{27} + \frac{1}{3619} + \frac{1}{13093542} = \frac{484946 + 3618 + 1}{13093542} = \frac{488565}{13093542} $$
Le PGCD du numérateur et du dénominateur est $\text{PGCD}(488565, 13093542) = 97713$.
La fraction irréductible est :
$$ \frac{488565}{13093542} = \frac{488565 \div 97713}{13093542 \div 97713} = \frac{5}{134} $$
Cette fraction est bien égale à $\frac{5}{134}$.

\subsection{Démonstration pour $n = 135$}
Soit $n = 135$. Nous cherchons $x, y, z \in \mathbb{Z}^{+}$ tels que $\frac{5}{135} = \frac{1}{x} + \frac{1}{y} + \frac{1}{z}$.
Posons $x = 28$, $y = 757$, $z = 572292$.
Les conditions $x > 0$, $y > 0$ et $z > 0$ sont satisfaites car $x, y, z \geq 1$.
Le PPCM des dénominateurs est $\text{PPCM}(28, 757, 572292) = 572292$.
En réduisant au même dénominateur :
\begin{itemize}
    \item $\frac{1}{28} = \frac{20439}{572292}$
    \item $\frac{1}{757} = \frac{756}{572292}$
    \item $\frac{1}{572292} = \frac{1}{572292}$
\end{itemize}
La somme des numérateurs est :
$$ \frac{1}{28} + \frac{1}{757} + \frac{1}{572292} = \frac{20439 + 756 + 1}{572292} = \frac{21196}{572292} $$
Le PGCD du numérateur et du dénominateur est $\text{PGCD}(21196, 572292) = 21196$.
La fraction irréductible est :
$$ \frac{21196}{572292} = \frac{21196 \div 21196}{572292 \div 21196} = \frac{1}{27} $$
Cette fraction est bien égale à $\frac{5}{135}$.

\subsection{Démonstration pour $n = 136$}
Soit $n = 136$. Nous cherchons $x, y, z \in \mathbb{Z}^{+}$ tels que $\frac{5}{136} = \frac{1}{x} + \frac{1}{y} + \frac{1}{z}$.
Posons $x = 28$, $y = 953$, $z = 907256$.
Les conditions $x > 0$, $y > 0$ et $z > 0$ sont satisfaites car $x, y, z \geq 1$.
Le PPCM des dénominateurs est $\text{PPCM}(28, 953, 907256) = 907256$.
En réduisant au même dénominateur :
\begin{itemize}
    \item $\frac{1}{28} = \frac{32402}{907256}$
    \item $\frac{1}{953} = \frac{952}{907256}$
    \item $\frac{1}{907256} = \frac{1}{907256}$
\end{itemize}
La somme des numérateurs est :
$$ \frac{1}{28} + \frac{1}{953} + \frac{1}{907256} = \frac{32402 + 952 + 1}{907256} = \frac{33355}{907256} $$
Le PGCD du numérateur et du dénominateur est $\text{PGCD}(33355, 907256) = 6671$.
La fraction irréductible est :
$$ \frac{33355}{907256} = \frac{33355 \div 6671}{907256 \div 6671} = \frac{5}{136} $$
Cette fraction est bien égale à $\frac{5}{136}$.

\subsection{Démonstration pour $n = 137$}
Soit $n = 137$. Nous cherchons $x, y, z \in \mathbb{Z}^{+}$ tels que $\frac{5}{137} = \frac{1}{x} + \frac{1}{y} + \frac{1}{z}$.
Posons $x = 28$, $y = 1279$, $z = 4906244$.
Les conditions $x > 0$, $y > 0$ et $z > 0$ sont satisfaites car $x, y, z \geq 1$.
Le PPCM des dénominateurs est $\text{PPCM}(28, 1279, 4906244) = 4906244$.
En réduisant au même dénominateur :
\begin{itemize}
    \item $\frac{1}{28} = \frac{175223}{4906244}$
    \item $\frac{1}{1279} = \frac{3836}{4906244}$
    \item $\frac{1}{4906244} = \frac{1}{4906244}$
\end{itemize}
La somme des numérateurs est :
$$ \frac{1}{28} + \frac{1}{1279} + \frac{1}{4906244} = \frac{175223 + 3836 + 1}{4906244} = \frac{179060}{4906244} $$
Le PGCD du numérateur et du dénominateur est $\text{PGCD}(179060, 4906244) = 35812$.
La fraction irréductible est :
$$ \frac{179060}{4906244} = \frac{179060 \div 35812}{4906244 \div 35812} = \frac{5}{137} $$
Cette fraction est bien égale à $\frac{5}{137}$.

\subsection{Démonstration pour $n = 138$}
Soit $n = 138$. Nous cherchons $x, y, z \in \mathbb{Z}^{+}$ tels que $\frac{5}{138} = \frac{1}{x} + \frac{1}{y} + \frac{1}{z}$.
Posons $x = 28$, $y = 1933$, $z = 3734556$.
Les conditions $x > 0$, $y > 0$ et $z > 0$ sont satisfaites car $x, y, z \geq 1$.
Le PPCM des dénominateurs est $\text{PPCM}(28, 1933, 3734556) = 3734556$.
En réduisant au même dénominateur :
\begin{itemize}
    \item $\frac{1}{28} = \frac{133377}{3734556}$
    \item $\frac{1}{1933} = \frac{1932}{3734556}$
    \item $\frac{1}{3734556} = \frac{1}{3734556}$
\end{itemize}
La somme des numérateurs est :
$$ \frac{1}{28} + \frac{1}{1933} + \frac{1}{3734556} = \frac{133377 + 1932 + 1}{3734556} = \frac{135310}{3734556} $$
Le PGCD du numérateur et du dénominateur est $\text{PGCD}(135310, 3734556) = 27062$.
La fraction irréductible est :
$$ \frac{135310}{3734556} = \frac{135310 \div 27062}{3734556 \div 27062} = \frac{5}{138} $$
Cette fraction est bien égale à $\frac{5}{138}$.

\subsection{Démonstration pour $n = 139$}
Soit $n = 139$. Nous cherchons $x, y, z \in \mathbb{Z}^{+}$ tels que $\frac{5}{139} = \frac{1}{x} + \frac{1}{y} + \frac{1}{z}$.
Posons $x = 28$, $y = 3893$, $z = 15151556$.
Les conditions $x > 0$, $y > 0$ et $z > 0$ sont satisfaites car $x, y, z \geq 1$.
Le PPCM des dénominateurs est $\text{PPCM}(28, 3893, 15151556) = 15151556$.
En réduisant au même dénominateur :
\begin{itemize}
    \item $\frac{1}{28} = \frac{541127}{15151556}$
    \item $\frac{1}{3893} = \frac{3892}{15151556}$
    \item $\frac{1}{15151556} = \frac{1}{15151556}$
\end{itemize}
La somme des numérateurs est :
$$ \frac{1}{28} + \frac{1}{3893} + \frac{1}{15151556} = \frac{541127 + 3892 + 1}{15151556} = \frac{545020}{15151556} $$
Le PGCD du numérateur et du dénominateur est $\text{PGCD}(545020, 15151556) = 109004$.
La fraction irréductible est :
$$ \frac{545020}{15151556} = \frac{545020 \div 109004}{15151556 \div 109004} = \frac{5}{139} $$
Cette fraction est bien égale à $\frac{5}{139}$.

\subsection{Démonstration pour $n = 140$}
Soit $n = 140$. Nous cherchons $x, y, z \in \mathbb{Z}^{+}$ tels que $\frac{5}{140} = \frac{1}{x} + \frac{1}{y} + \frac{1}{z}$.
Posons $x = 29$, $y = 813$, $z = 660156$.
Les conditions $x > 0$, $y > 0$ et $z > 0$ sont satisfaites car $x, y, z \geq 1$.
Le PPCM des dénominateurs est $\text{PPCM}(29, 813, 660156) = 660156$.
En réduisant au même dénominateur :
\begin{itemize}
    \item $\frac{1}{29} = \frac{22764}{660156}$
    \item $\frac{1}{813} = \frac{812}{660156}$
    \item $\frac{1}{660156} = \frac{1}{660156}$
\end{itemize}
La somme des numérateurs est :
$$ \frac{1}{29} + \frac{1}{813} + \frac{1}{660156} = \frac{22764 + 812 + 1}{660156} = \frac{23577}{660156} $$
Le PGCD du numérateur et du dénominateur est $\text{PGCD}(23577, 660156) = 23577$.
La fraction irréductible est :
$$ \frac{23577}{660156} = \frac{23577 \div 23577}{660156 \div 23577} = \frac{1}{28} $$
Cette fraction est bien égale à $\frac{5}{140}$.

\subsection{Démonstration pour $n = 141$}
Soit $n = 141$. Nous cherchons $x, y, z \in \mathbb{Z}^{+}$ tels que $\frac{5}{141} = \frac{1}{x} + \frac{1}{y} + \frac{1}{z}$.
Posons $x = 29$, $y = 1023$, $z = 1394349$.
Les conditions $x > 0$, $y > 0$ et $z > 0$ sont satisfaites car $x, y, z \geq 1$.
Le PPCM des dénominateurs est $\text{PPCM}(29, 1023, 1394349) = 1394349$.
En réduisant au même dénominateur :
\begin{itemize}
    \item $\frac{1}{29} = \frac{48081}{1394349}$
    \item $\frac{1}{1023} = \frac{1363}{1394349}$
    \item $\frac{1}{1394349} = \frac{1}{1394349}$
\end{itemize}
La somme des numérateurs est :
$$ \frac{1}{29} + \frac{1}{1023} + \frac{1}{1394349} = \frac{48081 + 1363 + 1}{1394349} = \frac{49445}{1394349} $$
Le PGCD du numérateur et du dénominateur est $\text{PGCD}(49445, 1394349) = 9889$.
La fraction irréductible est :
$$ \frac{49445}{1394349} = \frac{49445 \div 9889}{1394349 \div 9889} = \frac{5}{141} $$
Cette fraction est bien égale à $\frac{5}{141}$.

\subsection{Démonstration pour $n = 142$}
Soit $n = 142$. Nous cherchons $x, y, z \in \mathbb{Z}^{+}$ tels que $\frac{5}{142} = \frac{1}{x} + \frac{1}{y} + \frac{1}{z}$.
Posons $x = 29$, $y = 1373$, $z = 5654014$.
Les conditions $x > 0$, $y > 0$ et $z > 0$ sont satisfaites car $x, y, z \geq 1$.
Le PPCM des dénominateurs est $\text{PPCM}(29, 1373, 5654014) = 5654014$.
En réduisant au même dénominateur :
\begin{itemize}
    \item $\frac{1}{29} = \frac{194966}{5654014}$
    \item $\frac{1}{1373} = \frac{4118}{5654014}$
    \item $\frac{1}{5654014} = \frac{1}{5654014}$
\end{itemize}
La somme des numérateurs est :
$$ \frac{1}{29} + \frac{1}{1373} + \frac{1}{5654014} = \frac{194966 + 4118 + 1}{5654014} = \frac{199085}{5654014} $$
Le PGCD du numérateur et du dénominateur est $\text{PGCD}(199085, 5654014) = 39817$.
La fraction irréductible est :
$$ \frac{199085}{5654014} = \frac{199085 \div 39817}{5654014 \div 39817} = \frac{5}{142} $$
Cette fraction est bien égale à $\frac{5}{142}$.

\subsection{Démonstration pour $n = 143$}
Soit $n = 143$. Nous cherchons $x, y, z \in \mathbb{Z}^{+}$ tels que $\frac{5}{143} = \frac{1}{x} + \frac{1}{y} + \frac{1}{z}$.
Posons $x = 29$, $y = 2074$, $z = 8600878$.
Les conditions $x > 0$, $y > 0$ et $z > 0$ sont satisfaites car $x, y, z \geq 1$.
Le PPCM des dénominateurs est $\text{PPCM}(29, 2074, 8600878) = 8600878$.
En réduisant au même dénominateur :
\begin{itemize}
    \item $\frac{1}{29} = \frac{296582}{8600878}$
    \item $\frac{1}{2074} = \frac{4147}{8600878}$
    \item $\frac{1}{8600878} = \frac{1}{8600878}$
\end{itemize}
La somme des numérateurs est :
$$ \frac{1}{29} + \frac{1}{2074} + \frac{1}{8600878} = \frac{296582 + 4147 + 1}{8600878} = \frac{300730}{8600878} $$
Le PGCD du numérateur et du dénominateur est $\text{PGCD}(300730, 8600878) = 60146$.
La fraction irréductible est :
$$ \frac{300730}{8600878} = \frac{300730 \div 60146}{8600878 \div 60146} = \frac{5}{143} $$
Cette fraction est bien égale à $\frac{5}{143}$.

\subsection{Démonstration pour $n = 144$}
Soit $n = 144$. Nous cherchons $x, y, z \in \mathbb{Z}^{+}$ tels que $\frac{5}{144} = \frac{1}{x} + \frac{1}{y} + \frac{1}{z}$.
Posons $x = 29$, $y = 4177$, $z = 17443152$.
Les conditions $x > 0$, $y > 0$ et $z > 0$ sont satisfaites car $x, y, z \geq 1$.
Le PPCM des dénominateurs est $\text{PPCM}(29, 4177, 17443152) = 17443152$.
En réduisant au même dénominateur :
\begin{itemize}
    \item $\frac{1}{29} = \frac{601488}{17443152}$
    \item $\frac{1}{4177} = \frac{4176}{17443152}$
    \item $\frac{1}{17443152} = \frac{1}{17443152}$
\end{itemize}
La somme des numérateurs est :
$$ \frac{1}{29} + \frac{1}{4177} + \frac{1}{17443152} = \frac{601488 + 4176 + 1}{17443152} = \frac{605665}{17443152} $$
Le PGCD du numérateur et du dénominateur est $\text{PGCD}(605665, 17443152) = 121133$.
La fraction irréductible est :
$$ \frac{605665}{17443152} = \frac{605665 \div 121133}{17443152 \div 121133} = \frac{5}{144} $$
Cette fraction est bien égale à $\frac{5}{144}$.

\subsection{Démonstration pour $n = 145$}
Soit $n = 145$. Nous cherchons $x, y, z \in \mathbb{Z}^{+}$ tels que $\frac{5}{145} = \frac{1}{x} + \frac{1}{y} + \frac{1}{z}$.
Posons $x = 30$, $y = 871$, $z = 757770$.
Les conditions $x > 0$, $y > 0$ et $z > 0$ sont satisfaites car $x, y, z \geq 1$.
Le PPCM des dénominateurs est $\text{PPCM}(30, 871, 757770) = 757770$.
En réduisant au même dénominateur :
\begin{itemize}
    \item $\frac{1}{30} = \frac{25259}{757770}$
    \item $\frac{1}{871} = \frac{870}{757770}$
    \item $\frac{1}{757770} = \frac{1}{757770}$
\end{itemize}
La somme des numérateurs est :
$$ \frac{1}{30} + \frac{1}{871} + \frac{1}{757770} = \frac{25259 + 870 + 1}{757770} = \frac{26130}{757770} $$
Le PGCD du numérateur et du dénominateur est $\text{PGCD}(26130, 757770) = 26130$.
La fraction irréductible est :
$$ \frac{26130}{757770} = \frac{26130 \div 26130}{757770 \div 26130} = \frac{1}{29} $$
Cette fraction est bien égale à $\frac{5}{145}$.

\subsection{Démonstration pour $n = 146$}
Soit $n = 146$. Nous cherchons $x, y, z \in \mathbb{Z}^{+}$ tels que $\frac{5}{146} = \frac{1}{x} + \frac{1}{y} + \frac{1}{z}$.
Posons $x = 30$, $y = 1096$, $z = 1200120$.
Les conditions $x > 0$, $y > 0$ et $z > 0$ sont satisfaites car $x, y, z \geq 1$.
Le PPCM des dénominateurs est $\text{PPCM}(30, 1096, 1200120) = 1200120$.
En réduisant au même dénominateur :
\begin{itemize}
    \item $\frac{1}{30} = \frac{40004}{1200120}$
    \item $\frac{1}{1096} = \frac{1095}{1200120}$
    \item $\frac{1}{1200120} = \frac{1}{1200120}$
\end{itemize}
La somme des numérateurs est :
$$ \frac{1}{30} + \frac{1}{1096} + \frac{1}{1200120} = \frac{40004 + 1095 + 1}{1200120} = \frac{41100}{1200120} $$
Le PGCD du numérateur et du dénominateur est $\text{PGCD}(41100, 1200120) = 8220$.
La fraction irréductible est :
$$ \frac{41100}{1200120} = \frac{41100 \div 8220}{1200120 \div 8220} = \frac{5}{146} $$
Cette fraction est bien égale à $\frac{5}{146}$.

\subsection{Démonstration pour $n = 147$}
Soit $n = 147$. Nous cherchons $x, y, z \in \mathbb{Z}^{+}$ tels que $\frac{5}{147} = \frac{1}{x} + \frac{1}{y} + \frac{1}{z}$.
Posons $x = 30$, $y = 1471$, $z = 2162370$.
Les conditions $x > 0$, $y > 0$ et $z > 0$ sont satisfaites car $x, y, z \geq 1$.
Le PPCM des dénominateurs est $\text{PPCM}(30, 1471, 2162370) = 2162370$.
En réduisant au même dénominateur :
\begin{itemize}
    \item $\frac{1}{30} = \frac{72079}{2162370}$
    \item $\frac{1}{1471} = \frac{1470}{2162370}$
    \item $\frac{1}{2162370} = \frac{1}{2162370}$
\end{itemize}
La somme des numérateurs est :
$$ \frac{1}{30} + \frac{1}{1471} + \frac{1}{2162370} = \frac{72079 + 1470 + 1}{2162370} = \frac{73550}{2162370} $$
Le PGCD du numérateur et du dénominateur est $\text{PGCD}(73550, 2162370) = 14710$.
La fraction irréductible est :
$$ \frac{73550}{2162370} = \frac{73550 \div 14710}{2162370 \div 14710} = \frac{5}{147} $$
Cette fraction est bien égale à $\frac{5}{147}$.

\subsection{Démonstration pour $n = 148$}
Soit $n = 148$. Nous cherchons $x, y, z \in \mathbb{Z}^{+}$ tels que $\frac{5}{148} = \frac{1}{x} + \frac{1}{y} + \frac{1}{z}$.
Posons $x = 30$, $y = 2221$, $z = 4930620$.
Les conditions $x > 0$, $y > 0$ et $z > 0$ sont satisfaites car $x, y, z \geq 1$.
Le PPCM des dénominateurs est $\text{PPCM}(30, 2221, 4930620) = 4930620$.
En réduisant au même dénominateur :
\begin{itemize}
    \item $\frac{1}{30} = \frac{164354}{4930620}$
    \item $\frac{1}{2221} = \frac{2220}{4930620}$
    \item $\frac{1}{4930620} = \frac{1}{4930620}$
\end{itemize}
La somme des numérateurs est :
$$ \frac{1}{30} + \frac{1}{2221} + \frac{1}{4930620} = \frac{164354 + 2220 + 1}{4930620} = \frac{166575}{4930620} $$
Le PGCD du numérateur et du dénominateur est $\text{PGCD}(166575, 4930620) = 33315$.
La fraction irréductible est :
$$ \frac{166575}{4930620} = \frac{166575 \div 33315}{4930620 \div 33315} = \frac{5}{148} $$
Cette fraction est bien égale à $\frac{5}{148}$.

\subsection{Démonstration pour $n = 149$}
Soit $n = 149$. Nous cherchons $x, y, z \in \mathbb{Z}^{+}$ tels que $\frac{5}{149} = \frac{1}{x} + \frac{1}{y} + \frac{1}{z}$.
Posons $x = 30$, $y = 4471$, $z = 19985370$.
Les conditions $x > 0$, $y > 0$ et $z > 0$ sont satisfaites car $x, y, z \geq 1$.
Le PPCM des dénominateurs est $\text{PPCM}(30, 4471, 19985370) = 19985370$.
En réduisant au même dénominateur :
\begin{itemize}
    \item $\frac{1}{30} = \frac{666179}{19985370}$
    \item $\frac{1}{4471} = \frac{4470}{19985370}$
    \item $\frac{1}{19985370} = \frac{1}{19985370}$
\end{itemize}
La somme des numérateurs est :
$$ \frac{1}{30} + \frac{1}{4471} + \frac{1}{19985370} = \frac{666179 + 4470 + 1}{19985370} = \frac{670650}{19985370} $$
Le PGCD du numérateur et du dénominateur est $\text{PGCD}(670650, 19985370) = 134130$.
La fraction irréductible est :
$$ \frac{670650}{19985370} = \frac{670650 \div 134130}{19985370 \div 134130} = \frac{5}{149} $$
Cette fraction est bien égale à $\frac{5}{149}$.

\subsection{Démonstration pour $n = 150$}
Soit $n = 150$. Nous cherchons $x, y, z \in \mathbb{Z}^{+}$ tels que $\frac{5}{150} = \frac{1}{x} + \frac{1}{y} + \frac{1}{z}$.
Posons $x = 31$, $y = 931$, $z = 865830$.
Les conditions $x > 0$, $y > 0$ et $z > 0$ sont satisfaites car $x, y, z \geq 1$.
Le PPCM des dénominateurs est $\text{PPCM}(31, 931, 865830) = 865830$.
En réduisant au même dénominateur :
\begin{itemize}
    \item $\frac{1}{31} = \frac{27930}{865830}$
    \item $\frac{1}{931} = \frac{930}{865830}$
    \item $\frac{1}{865830} = \frac{1}{865830}$
\end{itemize}
La somme des numérateurs est :
$$ \frac{1}{31} + \frac{1}{931} + \frac{1}{865830} = \frac{27930 + 930 + 1}{865830} = \frac{28861}{865830} $$
Le PGCD du numérateur et du dénominateur est $\text{PGCD}(28861, 865830) = 28861$.
La fraction irréductible est :
$$ \frac{28861}{865830} = \frac{28861 \div 28861}{865830 \div 28861} = \frac{1}{30} $$
Cette fraction est bien égale à $\frac{5}{150}$.

\subsection{Démonstration pour $n = 151$}
Soit $n = 151$. Nous cherchons $x, y, z \in \mathbb{Z}^{+}$ tels que $\frac{5}{151} = \frac{1}{x} + \frac{1}{y} + \frac{1}{z}$.
Posons $x = 31$, $y = 1178$, $z = 177878$.
Les conditions $x > 0$, $y > 0$ et $z > 0$ sont satisfaites car $x, y, z \geq 1$.
Le PPCM des dénominateurs est $\text{PPCM}(31, 1178, 177878) = 177878$.
En réduisant au même dénominateur :
\begin{itemize}
    \item $\frac{1}{31} = \frac{5738}{177878}$
    \item $\frac{1}{1178} = \frac{151}{177878}$
    \item $\frac{1}{177878} = \frac{1}{177878}$
\end{itemize}
La somme des numérateurs est :
$$ \frac{1}{31} + \frac{1}{1178} + \frac{1}{177878} = \frac{5738 + 151 + 1}{177878} = \frac{5890}{177878} $$
Le PGCD du numérateur et du dénominateur est $\text{PGCD}(5890, 177878) = 1178$.
La fraction irréductible est :
$$ \frac{5890}{177878} = \frac{5890 \div 1178}{177878 \div 1178} = \frac{5}{151} $$
Cette fraction est bien égale à $\frac{5}{151}$.

\subsection{Démonstration pour $n = 152$}
Soit $n = 152$. Nous cherchons $x, y, z \in \mathbb{Z}^{+}$ tels que $\frac{5}{152} = \frac{1}{x} + \frac{1}{y} + \frac{1}{z}$.
Posons $x = 31$, $y = 1571$, $z = 7402552$.
Les conditions $x > 0$, $y > 0$ et $z > 0$ sont satisfaites car $x, y, z \geq 1$.
Le PPCM des dénominateurs est $\text{PPCM}(31, 1571, 7402552) = 7402552$.
En réduisant au même dénominateur :
\begin{itemize}
    \item $\frac{1}{31} = \frac{238792}{7402552}$
    \item $\frac{1}{1571} = \frac{4712}{7402552}$
    \item $\frac{1}{7402552} = \frac{1}{7402552}$
\end{itemize}
La somme des numérateurs est :
$$ \frac{1}{31} + \frac{1}{1571} + \frac{1}{7402552} = \frac{238792 + 4712 + 1}{7402552} = \frac{243505}{7402552} $$
Le PGCD du numérateur et du dénominateur est $\text{PGCD}(243505, 7402552) = 48701$.
La fraction irréductible est :
$$ \frac{243505}{7402552} = \frac{243505 \div 48701}{7402552 \div 48701} = \frac{5}{152} $$
Cette fraction est bien égale à $\frac{5}{152}$.

\subsection{Démonstration pour $n = 153$}
Soit $n = 153$. Nous cherchons $x, y, z \in \mathbb{Z}^{+}$ tels que $\frac{5}{153} = \frac{1}{x} + \frac{1}{y} + \frac{1}{z}$.
Posons $x = 31$, $y = 2372$, $z = 11250396$.
Les conditions $x > 0$, $y > 0$ et $z > 0$ sont satisfaites car $x, y, z \geq 1$.
Le PPCM des dénominateurs est $\text{PPCM}(31, 2372, 11250396) = 11250396$.
En réduisant au même dénominateur :
\begin{itemize}
    \item $\frac{1}{31} = \frac{362916}{11250396}$
    \item $\frac{1}{2372} = \frac{4743}{11250396}$
    \item $\frac{1}{11250396} = \frac{1}{11250396}$
\end{itemize}
La somme des numérateurs est :
$$ \frac{1}{31} + \frac{1}{2372} + \frac{1}{11250396} = \frac{362916 + 4743 + 1}{11250396} = \frac{367660}{11250396} $$
Le PGCD du numérateur et du dénominateur est $\text{PGCD}(367660, 11250396) = 73532$.
La fraction irréductible est :
$$ \frac{367660}{11250396} = \frac{367660 \div 73532}{11250396 \div 73532} = \frac{5}{153} $$
Cette fraction est bien égale à $\frac{5}{153}$.

\subsection{Démonstration pour $n = 154$}
Soit $n = 154$. Nous cherchons $x, y, z \in \mathbb{Z}^{+}$ tels que $\frac{5}{154} = \frac{1}{x} + \frac{1}{y} + \frac{1}{z}$.
Posons $x = 31$, $y = 4775$, $z = 22795850$.
Les conditions $x > 0$, $y > 0$ et $z > 0$ sont satisfaites car $x, y, z \geq 1$.
Le PPCM des dénominateurs est $\text{PPCM}(31, 4775, 22795850) = 22795850$.
En réduisant au même dénominateur :
\begin{itemize}
    \item $\frac{1}{31} = \frac{735350}{22795850}$
    \item $\frac{1}{4775} = \frac{4774}{22795850}$
    \item $\frac{1}{22795850} = \frac{1}{22795850}$
\end{itemize}
La somme des numérateurs est :
$$ \frac{1}{31} + \frac{1}{4775} + \frac{1}{22795850} = \frac{735350 + 4774 + 1}{22795850} = \frac{740125}{22795850} $$
Le PGCD du numérateur et du dénominateur est $\text{PGCD}(740125, 22795850) = 148025$.
La fraction irréductible est :
$$ \frac{740125}{22795850} = \frac{740125 \div 148025}{22795850 \div 148025} = \frac{5}{154} $$
Cette fraction est bien égale à $\frac{5}{154}$.

\subsection{Démonstration pour $n = 155$}
Soit $n = 155$. Nous cherchons $x, y, z \in \mathbb{Z}^{+}$ tels que $\frac{5}{155} = \frac{1}{x} + \frac{1}{y} + \frac{1}{z}$.
Posons $x = 32$, $y = 993$, $z = 985056$.
Les conditions $x > 0$, $y > 0$ et $z > 0$ sont satisfaites car $x, y, z \geq 1$.
Le PPCM des dénominateurs est $\text{PPCM}(32, 993, 985056) = 985056$.
En réduisant au même dénominateur :
\begin{itemize}
    \item $\frac{1}{32} = \frac{30783}{985056}$
    \item $\frac{1}{993} = \frac{992}{985056}$
    \item $\frac{1}{985056} = \frac{1}{985056}$
\end{itemize}
La somme des numérateurs est :
$$ \frac{1}{32} + \frac{1}{993} + \frac{1}{985056} = \frac{30783 + 992 + 1}{985056} = \frac{31776}{985056} $$
Le PGCD du numérateur et du dénominateur est $\text{PGCD}(31776, 985056) = 31776$.
La fraction irréductible est :
$$ \frac{31776}{985056} = \frac{31776 \div 31776}{985056 \div 31776} = \frac{1}{31} $$
Cette fraction est bien égale à $\frac{5}{155}$.

\subsection{Démonstration pour $n = 156$}
Soit $n = 156$. Nous cherchons $x, y, z \in \mathbb{Z}^{+}$ tels que $\frac{5}{156} = \frac{1}{x} + \frac{1}{y} + \frac{1}{z}$.
Posons $x = 32$, $y = 1249$, $z = 1558752$.
Les conditions $x > 0$, $y > 0$ et $z > 0$ sont satisfaites car $x, y, z \geq 1$.
Le PPCM des dénominateurs est $\text{PPCM}(32, 1249, 1558752) = 1558752$.
En réduisant au même dénominateur :
\begin{itemize}
    \item $\frac{1}{32} = \frac{48711}{1558752}$
    \item $\frac{1}{1249} = \frac{1248}{1558752}$
    \item $\frac{1}{1558752} = \frac{1}{1558752}$
\end{itemize}
La somme des numérateurs est :
$$ \frac{1}{32} + \frac{1}{1249} + \frac{1}{1558752} = \frac{48711 + 1248 + 1}{1558752} = \frac{49960}{1558752} $$
Le PGCD du numérateur et du dénominateur est $\text{PGCD}(49960, 1558752) = 9992$.
La fraction irréductible est :
$$ \frac{49960}{1558752} = \frac{49960 \div 9992}{1558752 \div 9992} = \frac{5}{156} $$
Cette fraction est bien égale à $\frac{5}{156}$.

\subsection{Démonstration pour $n = 157$}
Soit $n = 157$. Nous cherchons $x, y, z \in \mathbb{Z}^{+}$ tels que $\frac{5}{157} = \frac{1}{x} + \frac{1}{y} + \frac{1}{z}$.
Posons $x = 32$, $y = 1675$, $z = 8415200$.
Les conditions $x > 0$, $y > 0$ et $z > 0$ sont satisfaites car $x, y, z \geq 1$.
Le PPCM des dénominateurs est $\text{PPCM}(32, 1675, 8415200) = 8415200$.
En réduisant au même dénominateur :
\begin{itemize}
    \item $\frac{1}{32} = \frac{262975}{8415200}$
    \item $\frac{1}{1675} = \frac{5024}{8415200}$
    \item $\frac{1}{8415200} = \frac{1}{8415200}$
\end{itemize}
La somme des numérateurs est :
$$ \frac{1}{32} + \frac{1}{1675} + \frac{1}{8415200} = \frac{262975 + 5024 + 1}{8415200} = \frac{268000}{8415200} $$
Le PGCD du numérateur et du dénominateur est $\text{PGCD}(268000, 8415200) = 53600$.
La fraction irréductible est :
$$ \frac{268000}{8415200} = \frac{268000 \div 53600}{8415200 \div 53600} = \frac{5}{157} $$
Cette fraction est bien égale à $\frac{5}{157}$.

\subsection{Démonstration pour $n = 158$}
Soit $n = 158$. Nous cherchons $x, y, z \in \mathbb{Z}^{+}$ tels que $\frac{5}{158} = \frac{1}{x} + \frac{1}{y} + \frac{1}{z}$.
Posons $x = 32$, $y = 2529$, $z = 6393312$.
Les conditions $x > 0$, $y > 0$ et $z > 0$ sont satisfaites car $x, y, z \geq 1$.
Le PPCM des dénominateurs est $\text{PPCM}(32, 2529, 6393312) = 6393312$.
En réduisant au même dénominateur :
\begin{itemize}
    \item $\frac{1}{32} = \frac{199791}{6393312}$
    \item $\frac{1}{2529} = \frac{2528}{6393312}$
    \item $\frac{1}{6393312} = \frac{1}{6393312}$
\end{itemize}
La somme des numérateurs est :
$$ \frac{1}{32} + \frac{1}{2529} + \frac{1}{6393312} = \frac{199791 + 2528 + 1}{6393312} = \frac{202320}{6393312} $$
Le PGCD du numérateur et du dénominateur est $\text{PGCD}(202320, 6393312) = 40464$.
La fraction irréductible est :
$$ \frac{202320}{6393312} = \frac{202320 \div 40464}{6393312 \div 40464} = \frac{5}{158} $$
Cette fraction est bien égale à $\frac{5}{158}$.

\subsection{Démonstration pour $n = 159$}
Soit $n = 159$. Nous cherchons $x, y, z \in \mathbb{Z}^{+}$ tels que $\frac{5}{159} = \frac{1}{x} + \frac{1}{y} + \frac{1}{z}$.
Posons $x = 32$, $y = 5089$, $z = 25892832$.
Les conditions $x > 0$, $y > 0$ et $z > 0$ sont satisfaites car $x, y, z \geq 1$.
Le PPCM des dénominateurs est $\text{PPCM}(32, 5089, 25892832) = 25892832$.
En réduisant au même dénominateur :
\begin{itemize}
    \item $\frac{1}{32} = \frac{809151}{25892832}$
    \item $\frac{1}{5089} = \frac{5088}{25892832}$
    \item $\frac{1}{25892832} = \frac{1}{25892832}$
\end{itemize}
La somme des numérateurs est :
$$ \frac{1}{32} + \frac{1}{5089} + \frac{1}{25892832} = \frac{809151 + 5088 + 1}{25892832} = \frac{814240}{25892832} $$
Le PGCD du numérateur et du dénominateur est $\text{PGCD}(814240, 25892832) = 162848$.
La fraction irréductible est :
$$ \frac{814240}{25892832} = \frac{814240 \div 162848}{25892832 \div 162848} = \frac{5}{159} $$
Cette fraction est bien égale à $\frac{5}{159}$.

\subsection{Démonstration pour $n = 160$}
Soit $n = 160$. Nous cherchons $x, y, z \in \mathbb{Z}^{+}$ tels que $\frac{5}{160} = \frac{1}{x} + \frac{1}{y} + \frac{1}{z}$.
Posons $x = 33$, $y = 1057$, $z = 1116192$.
Les conditions $x > 0$, $y > 0$ et $z > 0$ sont satisfaites car $x, y, z \geq 1$.
Le PPCM des dénominateurs est $\text{PPCM}(33, 1057, 1116192) = 1116192$.
En réduisant au même dénominateur :
\begin{itemize}
    \item $\frac{1}{33} = \frac{33824}{1116192}$
    \item $\frac{1}{1057} = \frac{1056}{1116192}$
    \item $\frac{1}{1116192} = \frac{1}{1116192}$
\end{itemize}
La somme des numérateurs est :
$$ \frac{1}{33} + \frac{1}{1057} + \frac{1}{1116192} = \frac{33824 + 1056 + 1}{1116192} = \frac{34881}{1116192} $$
Le PGCD du numérateur et du dénominateur est $\text{PGCD}(34881, 1116192) = 34881$.
La fraction irréductible est :
$$ \frac{34881}{1116192} = \frac{34881 \div 34881}{1116192 \div 34881} = \frac{1}{32} $$
Cette fraction est bien égale à $\frac{5}{160}$.

\subsection{Démonstration pour $n = 161$}
Soit $n = 161$. Nous cherchons $x, y, z \in \mathbb{Z}^{+}$ tels que $\frac{5}{161} = \frac{1}{x} + \frac{1}{y} + \frac{1}{z}$.
Posons $x = 33$, $y = 1329$, $z = 2353659$.
Les conditions $x > 0$, $y > 0$ et $z > 0$ sont satisfaites car $x, y, z \geq 1$.
Le PPCM des dénominateurs est $\text{PPCM}(33, 1329, 2353659) = 2353659$.
En réduisant au même dénominateur :
\begin{itemize}
    \item $\frac{1}{33} = \frac{71323}{2353659}$
    \item $\frac{1}{1329} = \frac{1771}{2353659}$
    \item $\frac{1}{2353659} = \frac{1}{2353659}$
\end{itemize}
La somme des numérateurs est :
$$ \frac{1}{33} + \frac{1}{1329} + \frac{1}{2353659} = \frac{71323 + 1771 + 1}{2353659} = \frac{73095}{2353659} $$
Le PGCD du numérateur et du dénominateur est $\text{PGCD}(73095, 2353659) = 14619$.
La fraction irréductible est :
$$ \frac{73095}{2353659} = \frac{73095 \div 14619}{2353659 \div 14619} = \frac{5}{161} $$
Cette fraction est bien égale à $\frac{5}{161}$.

\subsection{Démonstration pour $n = 162$}
Soit $n = 162$. Nous cherchons $x, y, z \in \mathbb{Z}^{+}$ tels que $\frac{5}{162} = \frac{1}{x} + \frac{1}{y} + \frac{1}{z}$.
Posons $x = 33$, $y = 1783$, $z = 3177306$.
Les conditions $x > 0$, $y > 0$ et $z > 0$ sont satisfaites car $x, y, z \geq 1$.
Le PPCM des dénominateurs est $\text{PPCM}(33, 1783, 3177306) = 3177306$.
En réduisant au même dénominateur :
\begin{itemize}
    \item $\frac{1}{33} = \frac{96282}{3177306}$
    \item $\frac{1}{1783} = \frac{1782}{3177306}$
    \item $\frac{1}{3177306} = \frac{1}{3177306}$
\end{itemize}
La somme des numérateurs est :
$$ \frac{1}{33} + \frac{1}{1783} + \frac{1}{3177306} = \frac{96282 + 1782 + 1}{3177306} = \frac{98065}{3177306} $$
Le PGCD du numérateur et du dénominateur est $\text{PGCD}(98065, 3177306) = 19613$.
La fraction irréductible est :
$$ \frac{98065}{3177306} = \frac{98065 \div 19613}{3177306 \div 19613} = \frac{5}{162} $$
Cette fraction est bien égale à $\frac{5}{162}$.

\subsection{Démonstration pour $n = 163$}
Soit $n = 163$. Nous cherchons $x, y, z \in \mathbb{Z}^{+}$ tels que $\frac{5}{163} = \frac{1}{x} + \frac{1}{y} + \frac{1}{z}$.
Posons $x = 33$, $y = 2690$, $z = 14469510$.
Les conditions $x > 0$, $y > 0$ et $z > 0$ sont satisfaites car $x, y, z \geq 1$.
Le PPCM des dénominateurs est $\text{PPCM}(33, 2690, 14469510) = 14469510$.
En réduisant au même dénominateur :
\begin{itemize}
    \item $\frac{1}{33} = \frac{438470}{14469510}$
    \item $\frac{1}{2690} = \frac{5379}{14469510}$
    \item $\frac{1}{14469510} = \frac{1}{14469510}$
\end{itemize}
La somme des numérateurs est :
$$ \frac{1}{33} + \frac{1}{2690} + \frac{1}{14469510} = \frac{438470 + 5379 + 1}{14469510} = \frac{443850}{14469510} $$
Le PGCD du numérateur et du dénominateur est $\text{PGCD}(443850, 14469510) = 88770$.
La fraction irréductible est :
$$ \frac{443850}{14469510} = \frac{443850 \div 88770}{14469510 \div 88770} = \frac{5}{163} $$
Cette fraction est bien égale à $\frac{5}{163}$.

\subsection{Démonstration pour $n = 164$}
Soit $n = 164$. Nous cherchons $x, y, z \in \mathbb{Z}^{+}$ tels que $\frac{5}{164} = \frac{1}{x} + \frac{1}{y} + \frac{1}{z}$.
Posons $x = 33$, $y = 5413$, $z = 29295156$.
Les conditions $x > 0$, $y > 0$ et $z > 0$ sont satisfaites car $x, y, z \geq 1$.
Le PPCM des dénominateurs est $\text{PPCM}(33, 5413, 29295156) = 29295156$.
En réduisant au même dénominateur :
\begin{itemize}
    \item $\frac{1}{33} = \frac{887732}{29295156}$
    \item $\frac{1}{5413} = \frac{5412}{29295156}$
    \item $\frac{1}{29295156} = \frac{1}{29295156}$
\end{itemize}
La somme des numérateurs est :
$$ \frac{1}{33} + \frac{1}{5413} + \frac{1}{29295156} = \frac{887732 + 5412 + 1}{29295156} = \frac{893145}{29295156} $$
Le PGCD du numérateur et du dénominateur est $\text{PGCD}(893145, 29295156) = 178629$.
La fraction irréductible est :
$$ \frac{893145}{29295156} = \frac{893145 \div 178629}{29295156 \div 178629} = \frac{5}{164} $$
Cette fraction est bien égale à $\frac{5}{164}$.

\subsection{Démonstration pour $n = 165$}
Soit $n = 165$. Nous cherchons $x, y, z \in \mathbb{Z}^{+}$ tels que $\frac{5}{165} = \frac{1}{x} + \frac{1}{y} + \frac{1}{z}$.
Posons $x = 34$, $y = 1123$, $z = 1260006$.
Les conditions $x > 0$, $y > 0$ et $z > 0$ sont satisfaites car $x, y, z \geq 1$.
Le PPCM des dénominateurs est $\text{PPCM}(34, 1123, 1260006) = 1260006$.
En réduisant au même dénominateur :
\begin{itemize}
    \item $\frac{1}{34} = \frac{37059}{1260006}$
    \item $\frac{1}{1123} = \frac{1122}{1260006}$
    \item $\frac{1}{1260006} = \frac{1}{1260006}$
\end{itemize}
La somme des numérateurs est :
$$ \frac{1}{34} + \frac{1}{1123} + \frac{1}{1260006} = \frac{37059 + 1122 + 1}{1260006} = \frac{38182}{1260006} $$
Le PGCD du numérateur et du dénominateur est $\text{PGCD}(38182, 1260006) = 38182$.
La fraction irréductible est :
$$ \frac{38182}{1260006} = \frac{38182 \div 38182}{1260006 \div 38182} = \frac{1}{33} $$
Cette fraction est bien égale à $\frac{5}{165}$.

\subsection{Démonstration pour $n = 166$}
Soit $n = 166$. Nous cherchons $x, y, z \in \mathbb{Z}^{+}$ tels que $\frac{5}{166} = \frac{1}{x} + \frac{1}{y} + \frac{1}{z}$.
Posons $x = 34$, $y = 1412$, $z = 1992332$.
Les conditions $x > 0$, $y > 0$ et $z > 0$ sont satisfaites car $x, y, z \geq 1$.
Le PPCM des dénominateurs est $\text{PPCM}(34, 1412, 1992332) = 1992332$.
En réduisant au même dénominateur :
\begin{itemize}
    \item $\frac{1}{34} = \frac{58598}{1992332}$
    \item $\frac{1}{1412} = \frac{1411}{1992332}$
    \item $\frac{1}{1992332} = \frac{1}{1992332}$
\end{itemize}
La somme des numérateurs est :
$$ \frac{1}{34} + \frac{1}{1412} + \frac{1}{1992332} = \frac{58598 + 1411 + 1}{1992332} = \frac{60010}{1992332} $$
Le PGCD du numérateur et du dénominateur est $\text{PGCD}(60010, 1992332) = 12002$.
La fraction irréductible est :
$$ \frac{60010}{1992332} = \frac{60010 \div 12002}{1992332 \div 12002} = \frac{5}{166} $$
Cette fraction est bien égale à $\frac{5}{166}$.

\subsection{Démonstration pour $n = 167$}
Soit $n = 167$. Nous cherchons $x, y, z \in \mathbb{Z}^{+}$ tels que $\frac{5}{167} = \frac{1}{x} + \frac{1}{y} + \frac{1}{z}$.
Posons $x = 34$, $y = 1893$, $z = 10748454$.
Les conditions $x > 0$, $y > 0$ et $z > 0$ sont satisfaites car $x, y, z \geq 1$.
Le PPCM des dénominateurs est $\text{PPCM}(34, 1893, 10748454) = 10748454$.
En réduisant au même dénominateur :
\begin{itemize}
    \item $\frac{1}{34} = \frac{316131}{10748454}$
    \item $\frac{1}{1893} = \frac{5678}{10748454}$
    \item $\frac{1}{10748454} = \frac{1}{10748454}$
\end{itemize}
La somme des numérateurs est :
$$ \frac{1}{34} + \frac{1}{1893} + \frac{1}{10748454} = \frac{316131 + 5678 + 1}{10748454} = \frac{321810}{10748454} $$
Le PGCD du numérateur et du dénominateur est $\text{PGCD}(321810, 10748454) = 64362$.
La fraction irréductible est :
$$ \frac{321810}{10748454} = \frac{321810 \div 64362}{10748454 \div 64362} = \frac{5}{167} $$
Cette fraction est bien égale à $\frac{5}{167}$.

\subsection{Démonstration pour $n = 168$}
Soit $n = 168$. Nous cherchons $x, y, z \in \mathbb{Z}^{+}$ tels que $\frac{5}{168} = \frac{1}{x} + \frac{1}{y} + \frac{1}{z}$.
Posons $x = 34$, $y = 2857$, $z = 8159592$.
Les conditions $x > 0$, $y > 0$ et $z > 0$ sont satisfaites car $x, y, z \geq 1$.
Le PPCM des dénominateurs est $\text{PPCM}(34, 2857, 8159592) = 8159592$.
En réduisant au même dénominateur :
\begin{itemize}
    \item $\frac{1}{34} = \frac{239988}{8159592}$
    \item $\frac{1}{2857} = \frac{2856}{8159592}$
    \item $\frac{1}{8159592} = \frac{1}{8159592}$
\end{itemize}
La somme des numérateurs est :
$$ \frac{1}{34} + \frac{1}{2857} + \frac{1}{8159592} = \frac{239988 + 2856 + 1}{8159592} = \frac{242845}{8159592} $$
Le PGCD du numérateur et du dénominateur est $\text{PGCD}(242845, 8159592) = 48569$.
La fraction irréductible est :
$$ \frac{242845}{8159592} = \frac{242845 \div 48569}{8159592 \div 48569} = \frac{5}{168} $$
Cette fraction est bien égale à $\frac{5}{168}$.

\subsection{Démonstration pour $n = 169$}
Soit $n = 169$. Nous cherchons $x, y, z \in \mathbb{Z}^{+}$ tels que $\frac{5}{169} = \frac{1}{x} + \frac{1}{y} + \frac{1}{z}$.
Posons $x = 34$, $y = 5747$, $z = 33022262$.
Les conditions $x > 0$, $y > 0$ et $z > 0$ sont satisfaites car $x, y, z \geq 1$.
Le PPCM des dénominateurs est $\text{PPCM}(34, 5747, 33022262) = 33022262$.
En réduisant au même dénominateur :
\begin{itemize}
    \item $\frac{1}{34} = \frac{971243}{33022262}$
    \item $\frac{1}{5747} = \frac{5746}{33022262}$
    \item $\frac{1}{33022262} = \frac{1}{33022262}$
\end{itemize}
La somme des numérateurs est :
$$ \frac{1}{34} + \frac{1}{5747} + \frac{1}{33022262} = \frac{971243 + 5746 + 1}{33022262} = \frac{976990}{33022262} $$
Le PGCD du numérateur et du dénominateur est $\text{PGCD}(976990, 33022262) = 195398$.
La fraction irréductible est :
$$ \frac{976990}{33022262} = \frac{976990 \div 195398}{33022262 \div 195398} = \frac{5}{169} $$
Cette fraction est bien égale à $\frac{5}{169}$.

\subsection{Démonstration pour $n = 170$}
Soit $n = 170$. Nous cherchons $x, y, z \in \mathbb{Z}^{+}$ tels que $\frac{5}{170} = \frac{1}{x} + \frac{1}{y} + \frac{1}{z}$.
Posons $x = 35$, $y = 1191$, $z = 1417290$.
Les conditions $x > 0$, $y > 0$ et $z > 0$ sont satisfaites car $x, y, z \geq 1$.
Le PPCM des dénominateurs est $\text{PPCM}(35, 1191, 1417290) = 1417290$.
En réduisant au même dénominateur :
\begin{itemize}
    \item $\frac{1}{35} = \frac{40494}{1417290}$
    \item $\frac{1}{1191} = \frac{1190}{1417290}$
    \item $\frac{1}{1417290} = \frac{1}{1417290}$
\end{itemize}
La somme des numérateurs est :
$$ \frac{1}{35} + \frac{1}{1191} + \frac{1}{1417290} = \frac{40494 + 1190 + 1}{1417290} = \frac{41685}{1417290} $$
Le PGCD du numérateur et du dénominateur est $\text{PGCD}(41685, 1417290) = 41685$.
La fraction irréductible est :
$$ \frac{41685}{1417290} = \frac{41685 \div 41685}{1417290 \div 41685} = \frac{1}{34} $$
Cette fraction est bien égale à $\frac{5}{170}$.

\subsection{Démonstration pour $n = 171$}
Soit $n = 171$. Nous cherchons $x, y, z \in \mathbb{Z}^{+}$ tels que $\frac{5}{171} = \frac{1}{x} + \frac{1}{y} + \frac{1}{z}$.
Posons $x = 35$, $y = 1497$, $z = 2986515$.
Les conditions $x > 0$, $y > 0$ et $z > 0$ sont satisfaites car $x, y, z \geq 1$.
Le PPCM des dénominateurs est $\text{PPCM}(35, 1497, 2986515) = 2986515$.
En réduisant au même dénominateur :
\begin{itemize}
    \item $\frac{1}{35} = \frac{85329}{2986515}$
    \item $\frac{1}{1497} = \frac{1995}{2986515}$
    \item $\frac{1}{2986515} = \frac{1}{2986515}$
\end{itemize}
La somme des numérateurs est :
$$ \frac{1}{35} + \frac{1}{1497} + \frac{1}{2986515} = \frac{85329 + 1995 + 1}{2986515} = \frac{87325}{2986515} $$
Le PGCD du numérateur et du dénominateur est $\text{PGCD}(87325, 2986515) = 17465$.
La fraction irréductible est :
$$ \frac{87325}{2986515} = \frac{87325 \div 17465}{2986515 \div 17465} = \frac{5}{171} $$
Cette fraction est bien égale à $\frac{5}{171}$.

\subsection{Démonstration pour $n = 172$}
Soit $n = 172$. Nous cherchons $x, y, z \in \mathbb{Z}^{+}$ tels que $\frac{5}{172} = \frac{1}{x} + \frac{1}{y} + \frac{1}{z}$.
Posons $x = 35$, $y = 2007$, $z = 12082140$.
Les conditions $x > 0$, $y > 0$ et $z > 0$ sont satisfaites car $x, y, z \geq 1$.
Le PPCM des dénominateurs est $\text{PPCM}(35, 2007, 12082140) = 12082140$.
En réduisant au même dénominateur :
\begin{itemize}
    \item $\frac{1}{35} = \frac{345204}{12082140}$
    \item $\frac{1}{2007} = \frac{6020}{12082140}$
    \item $\frac{1}{12082140} = \frac{1}{12082140}$
\end{itemize}
La somme des numérateurs est :
$$ \frac{1}{35} + \frac{1}{2007} + \frac{1}{12082140} = \frac{345204 + 6020 + 1}{12082140} = \frac{351225}{12082140} $$
Le PGCD du numérateur et du dénominateur est $\text{PGCD}(351225, 12082140) = 70245$.
La fraction irréductible est :
$$ \frac{351225}{12082140} = \frac{351225 \div 70245}{12082140 \div 70245} = \frac{5}{172} $$
Cette fraction est bien égale à $\frac{5}{172}$.

\subsection{Démonstration pour $n = 173$}
Soit $n = 173$. Nous cherchons $x, y, z \in \mathbb{Z}^{+}$ tels que $\frac{5}{173} = \frac{1}{x} + \frac{1}{y} + \frac{1}{z}$.
Posons $x = 35$, $y = 3028$, $z = 18334540$.
Les conditions $x > 0$, $y > 0$ et $z > 0$ sont satisfaites car $x, y, z \geq 1$.
Le PPCM des dénominateurs est $\text{PPCM}(35, 3028, 18334540) = 18334540$.
En réduisant au même dénominateur :
\begin{itemize}
    \item $\frac{1}{35} = \frac{523844}{18334540}$
    \item $\frac{1}{3028} = \frac{6055}{18334540}$
    \item $\frac{1}{18334540} = \frac{1}{18334540}$
\end{itemize}
La somme des numérateurs est :
$$ \frac{1}{35} + \frac{1}{3028} + \frac{1}{18334540} = \frac{523844 + 6055 + 1}{18334540} = \frac{529900}{18334540} $$
Le PGCD du numérateur et du dénominateur est $\text{PGCD}(529900, 18334540) = 105980$.
La fraction irréductible est :
$$ \frac{529900}{18334540} = \frac{529900 \div 105980}{18334540 \div 105980} = \frac{5}{173} $$
Cette fraction est bien égale à $\frac{5}{173}$.

\subsection{Démonstration pour $n = 174$}
Soit $n = 174$. Nous cherchons $x, y, z \in \mathbb{Z}^{+}$ tels que $\frac{5}{174} = \frac{1}{x} + \frac{1}{y} + \frac{1}{z}$.
Posons $x = 35$, $y = 6091$, $z = 37094190$.
Les conditions $x > 0$, $y > 0$ et $z > 0$ sont satisfaites car $x, y, z \geq 1$.
Le PPCM des dénominateurs est $\text{PPCM}(35, 6091, 37094190) = 37094190$.
En réduisant au même dénominateur :
\begin{itemize}
    \item $\frac{1}{35} = \frac{1059834}{37094190}$
    \item $\frac{1}{6091} = \frac{6090}{37094190}$
    \item $\frac{1}{37094190} = \frac{1}{37094190}$
\end{itemize}
La somme des numérateurs est :
$$ \frac{1}{35} + \frac{1}{6091} + \frac{1}{37094190} = \frac{1059834 + 6090 + 1}{37094190} = \frac{1065925}{37094190} $$
Le PGCD du numérateur et du dénominateur est $\text{PGCD}(1065925, 37094190) = 213185$.
La fraction irréductible est :
$$ \frac{1065925}{37094190} = \frac{1065925 \div 213185}{37094190 \div 213185} = \frac{5}{174} $$
Cette fraction est bien égale à $\frac{5}{174}$.

\subsection{Démonstration pour $n = 175$}
Soit $n = 175$. Nous cherchons $x, y, z \in \mathbb{Z}^{+}$ tels que $\frac{5}{175} = \frac{1}{x} + \frac{1}{y} + \frac{1}{z}$.
Posons $x = 36$, $y = 1261$, $z = 1588860$.
Les conditions $x > 0$, $y > 0$ et $z > 0$ sont satisfaites car $x, y, z \geq 1$.
Le PPCM des dénominateurs est $\text{PPCM}(36, 1261, 1588860) = 1588860$.
En réduisant au même dénominateur :
\begin{itemize}
    \item $\frac{1}{36} = \frac{44135}{1588860}$
    \item $\frac{1}{1261} = \frac{1260}{1588860}$
    \item $\frac{1}{1588860} = \frac{1}{1588860}$
\end{itemize}
La somme des numérateurs est :
$$ \frac{1}{36} + \frac{1}{1261} + \frac{1}{1588860} = \frac{44135 + 1260 + 1}{1588860} = \frac{45396}{1588860} $$
Le PGCD du numérateur et du dénominateur est $\text{PGCD}(45396, 1588860) = 45396$.
La fraction irréductible est :
$$ \frac{45396}{1588860} = \frac{45396 \div 45396}{1588860 \div 45396} = \frac{1}{35} $$
Cette fraction est bien égale à $\frac{5}{175}$.

\subsection{Démonstration pour $n = 176$}
Soit $n = 176$. Nous cherchons $x, y, z \in \mathbb{Z}^{+}$ tels que $\frac{5}{176} = \frac{1}{x} + \frac{1}{y} + \frac{1}{z}$.
Posons $x = 36$, $y = 1585$, $z = 2510640$.
Les conditions $x > 0$, $y > 0$ et $z > 0$ sont satisfaites car $x, y, z \geq 1$.
Le PPCM des dénominateurs est $\text{PPCM}(36, 1585, 2510640) = 2510640$.
En réduisant au même dénominateur :
\begin{itemize}
    \item $\frac{1}{36} = \frac{69740}{2510640}$
    \item $\frac{1}{1585} = \frac{1584}{2510640}$
    \item $\frac{1}{2510640} = \frac{1}{2510640}$
\end{itemize}
La somme des numérateurs est :
$$ \frac{1}{36} + \frac{1}{1585} + \frac{1}{2510640} = \frac{69740 + 1584 + 1}{2510640} = \frac{71325}{2510640} $$
Le PGCD du numérateur et du dénominateur est $\text{PGCD}(71325, 2510640) = 14265$.
La fraction irréductible est :
$$ \frac{71325}{2510640} = \frac{71325 \div 14265}{2510640 \div 14265} = \frac{5}{176} $$
Cette fraction est bien égale à $\frac{5}{176}$.

\subsection{Démonstration pour $n = 177$}
Soit $n = 177$. Nous cherchons $x, y, z \in \mathbb{Z}^{+}$ tels que $\frac{5}{177} = \frac{1}{x} + \frac{1}{y} + \frac{1}{z}$.
Posons $x = 36$, $y = 2125$, $z = 4513500$.
Les conditions $x > 0$, $y > 0$ et $z > 0$ sont satisfaites car $x, y, z \geq 1$.
Le PPCM des dénominateurs est $\text{PPCM}(36, 2125, 4513500) = 4513500$.
En réduisant au même dénominateur :
\begin{itemize}
    \item $\frac{1}{36} = \frac{125375}{4513500}$
    \item $\frac{1}{2125} = \frac{2124}{4513500}$
    \item $\frac{1}{4513500} = \frac{1}{4513500}$
\end{itemize}
La somme des numérateurs est :
$$ \frac{1}{36} + \frac{1}{2125} + \frac{1}{4513500} = \frac{125375 + 2124 + 1}{4513500} = \frac{127500}{4513500} $$
Le PGCD du numérateur et du dénominateur est $\text{PGCD}(127500, 4513500) = 25500$.
La fraction irréductible est :
$$ \frac{127500}{4513500} = \frac{127500 \div 25500}{4513500 \div 25500} = \frac{5}{177} $$
Cette fraction est bien égale à $\frac{5}{177}$.

\subsection{Démonstration pour $n = 178$}
Soit $n = 178$. Nous cherchons $x, y, z \in \mathbb{Z}^{+}$ tels que $\frac{5}{178} = \frac{1}{x} + \frac{1}{y} + \frac{1}{z}$.
Posons $x = 36$, $y = 3205$, $z = 10268820$.
Les conditions $x > 0$, $y > 0$ et $z > 0$ sont satisfaites car $x, y, z \geq 1$.
Le PPCM des dénominateurs est $\text{PPCM}(36, 3205, 10268820) = 10268820$.
En réduisant au même dénominateur :
\begin{itemize}
    \item $\frac{1}{36} = \frac{285245}{10268820}$
    \item $\frac{1}{3205} = \frac{3204}{10268820}$
    \item $\frac{1}{10268820} = \frac{1}{10268820}$
\end{itemize}
La somme des numérateurs est :
$$ \frac{1}{36} + \frac{1}{3205} + \frac{1}{10268820} = \frac{285245 + 3204 + 1}{10268820} = \frac{288450}{10268820} $$
Le PGCD du numérateur et du dénominateur est $\text{PGCD}(288450, 10268820) = 57690$.
La fraction irréductible est :
$$ \frac{288450}{10268820} = \frac{288450 \div 57690}{10268820 \div 57690} = \frac{5}{178} $$
Cette fraction est bien égale à $\frac{5}{178}$.

\subsection{Démonstration pour $n = 179$}
Soit $n = 179$. Nous cherchons $x, y, z \in \mathbb{Z}^{+}$ tels que $\frac{5}{179} = \frac{1}{x} + \frac{1}{y} + \frac{1}{z}$.
Posons $x = 36$, $y = 6445$, $z = 41531580$.
Les conditions $x > 0$, $y > 0$ et $z > 0$ sont satisfaites car $x, y, z \geq 1$.
Le PPCM des dénominateurs est $\text{PPCM}(36, 6445, 41531580) = 41531580$.
En réduisant au même dénominateur :
\begin{itemize}
    \item $\frac{1}{36} = \frac{1153655}{41531580}$
    \item $\frac{1}{6445} = \frac{6444}{41531580}$
    \item $\frac{1}{41531580} = \frac{1}{41531580}$
\end{itemize}
La somme des numérateurs est :
$$ \frac{1}{36} + \frac{1}{6445} + \frac{1}{41531580} = \frac{1153655 + 6444 + 1}{41531580} = \frac{1160100}{41531580} $$
Le PGCD du numérateur et du dénominateur est $\text{PGCD}(1160100, 41531580) = 232020$.
La fraction irréductible est :
$$ \frac{1160100}{41531580} = \frac{1160100 \div 232020}{41531580 \div 232020} = \frac{5}{179} $$
Cette fraction est bien égale à $\frac{5}{179}$.

\subsection{Démonstration pour $n = 180$}
Soit $n = 180$. Nous cherchons $x, y, z \in \mathbb{Z}^{+}$ tels que $\frac{5}{180} = \frac{1}{x} + \frac{1}{y} + \frac{1}{z}$.
Posons $x = 37$, $y = 1333$, $z = 1775556$.
Les conditions $x > 0$, $y > 0$ et $z > 0$ sont satisfaites car $x, y, z \geq 1$.
Le PPCM des dénominateurs est $\text{PPCM}(37, 1333, 1775556) = 1775556$.
En réduisant au même dénominateur :
\begin{itemize}
    \item $\frac{1}{37} = \frac{47988}{1775556}$
    \item $\frac{1}{1333} = \frac{1332}{1775556}$
    \item $\frac{1}{1775556} = \frac{1}{1775556}$
\end{itemize}
La somme des numérateurs est :
$$ \frac{1}{37} + \frac{1}{1333} + \frac{1}{1775556} = \frac{47988 + 1332 + 1}{1775556} = \frac{49321}{1775556} $$
Le PGCD du numérateur et du dénominateur est $\text{PGCD}(49321, 1775556) = 49321$.
La fraction irréductible est :
$$ \frac{49321}{1775556} = \frac{49321 \div 49321}{1775556 \div 49321} = \frac{1}{36} $$
Cette fraction est bien égale à $\frac{5}{180}$.

\subsection{Démonstration pour $n = 181$}
Soit $n = 181$. Nous cherchons $x, y, z \in \mathbb{Z}^{+}$ tels que $\frac{5}{181} = \frac{1}{x} + \frac{1}{y} + \frac{1}{z}$.
Posons $x = 39$, $y = 507$, $z = 91767$.
Les conditions $x > 0$, $y > 0$ et $z > 0$ sont satisfaites car $x, y, z \geq 1$.
Le PPCM des dénominateurs est $\text{PPCM}(39, 507, 91767) = 91767$.
En réduisant au même dénominateur :
\begin{itemize}
    \item $\frac{1}{39} = \frac{2353}{91767}$
    \item $\frac{1}{507} = \frac{181}{91767}$
    \item $\frac{1}{91767} = \frac{1}{91767}$
\end{itemize}
La somme des numérateurs est :
$$ \frac{1}{39} + \frac{1}{507} + \frac{1}{91767} = \frac{2353 + 181 + 1}{91767} = \frac{2535}{91767} $$
Le PGCD du numérateur et du dénominateur est $\text{PGCD}(2535, 91767) = 507$.
La fraction irréductible est :
$$ \frac{2535}{91767} = \frac{2535 \div 507}{91767 \div 507} = \frac{5}{181} $$
Cette fraction est bien égale à $\frac{5}{181}$.

\subsection{Démonstration pour $n = 182$}
Soit $n = 182$. Nous cherchons $x, y, z \in \mathbb{Z}^{+}$ tels que $\frac{5}{182} = \frac{1}{x} + \frac{1}{y} + \frac{1}{z}$.
Posons $x = 37$, $y = 2245$, $z = 15117830$.
Les conditions $x > 0$, $y > 0$ et $z > 0$ sont satisfaites car $x, y, z \geq 1$.
Le PPCM des dénominateurs est $\text{PPCM}(37, 2245, 15117830) = 15117830$.
En réduisant au même dénominateur :
\begin{itemize}
    \item $\frac{1}{37} = \frac{408590}{15117830}$
    \item $\frac{1}{2245} = \frac{6734}{15117830}$
    \item $\frac{1}{15117830} = \frac{1}{15117830}$
\end{itemize}
La somme des numérateurs est :
$$ \frac{1}{37} + \frac{1}{2245} + \frac{1}{15117830} = \frac{408590 + 6734 + 1}{15117830} = \frac{415325}{15117830} $$
Le PGCD du numérateur et du dénominateur est $\text{PGCD}(415325, 15117830) = 83065$.
La fraction irréductible est :
$$ \frac{415325}{15117830} = \frac{415325 \div 83065}{15117830 \div 83065} = \frac{5}{182} $$
Cette fraction est bien égale à $\frac{5}{182}$.

\subsection{Démonstration pour $n = 183$}
Soit $n = 183$. Nous cherchons $x, y, z \in \mathbb{Z}^{+}$ tels que $\frac{5}{183} = \frac{1}{x} + \frac{1}{y} + \frac{1}{z}$.
Posons $x = 37$, $y = 3386$, $z = 22926606$.
Les conditions $x > 0$, $y > 0$ et $z > 0$ sont satisfaites car $x, y, z \geq 1$.
Le PPCM des dénominateurs est $\text{PPCM}(37, 3386, 22926606) = 22926606$.
En réduisant au même dénominateur :
\begin{itemize}
    \item $\frac{1}{37} = \frac{619638}{22926606}$
    \item $\frac{1}{3386} = \frac{6771}{22926606}$
    \item $\frac{1}{22926606} = \frac{1}{22926606}$
\end{itemize}
La somme des numérateurs est :
$$ \frac{1}{37} + \frac{1}{3386} + \frac{1}{22926606} = \frac{619638 + 6771 + 1}{22926606} = \frac{626410}{22926606} $$
Le PGCD du numérateur et du dénominateur est $\text{PGCD}(626410, 22926606) = 125282$.
La fraction irréductible est :
$$ \frac{626410}{22926606} = \frac{626410 \div 125282}{22926606 \div 125282} = \frac{5}{183} $$
Cette fraction est bien égale à $\frac{5}{183}$.

\subsection{Démonstration pour $n = 184$}
Soit $n = 184$. Nous cherchons $x, y, z \in \mathbb{Z}^{+}$ tels que $\frac{5}{184} = \frac{1}{x} + \frac{1}{y} + \frac{1}{z}$.
Posons $x = 37$, $y = 6809$, $z = 46355672$.
Les conditions $x > 0$, $y > 0$ et $z > 0$ sont satisfaites car $x, y, z \geq 1$.
Le PPCM des dénominateurs est $\text{PPCM}(37, 6809, 46355672) = 46355672$.
En réduisant au même dénominateur :
\begin{itemize}
    \item $\frac{1}{37} = \frac{1252856}{46355672}$
    \item $\frac{1}{6809} = \frac{6808}{46355672}$
    \item $\frac{1}{46355672} = \frac{1}{46355672}$
\end{itemize}
La somme des numérateurs est :
$$ \frac{1}{37} + \frac{1}{6809} + \frac{1}{46355672} = \frac{1252856 + 6808 + 1}{46355672} = \frac{1259665}{46355672} $$
Le PGCD du numérateur et du dénominateur est $\text{PGCD}(1259665, 46355672) = 251933$.
La fraction irréductible est :
$$ \frac{1259665}{46355672} = \frac{1259665 \div 251933}{46355672 \div 251933} = \frac{5}{184} $$
Cette fraction est bien égale à $\frac{5}{184}$.

\subsection{Démonstration pour $n = 185$}
Soit $n = 185$. Nous cherchons $x, y, z \in \mathbb{Z}^{+}$ tels que $\frac{5}{185} = \frac{1}{x} + \frac{1}{y} + \frac{1}{z}$.
Posons $x = 38$, $y = 1407$, $z = 1978242$.
Les conditions $x > 0$, $y > 0$ et $z > 0$ sont satisfaites car $x, y, z \geq 1$.
Le PPCM des dénominateurs est $\text{PPCM}(38, 1407, 1978242) = 1978242$.
En réduisant au même dénominateur :
\begin{itemize}
    \item $\frac{1}{38} = \frac{52059}{1978242}$
    \item $\frac{1}{1407} = \frac{1406}{1978242}$
    \item $\frac{1}{1978242} = \frac{1}{1978242}$
\end{itemize}
La somme des numérateurs est :
$$ \frac{1}{38} + \frac{1}{1407} + \frac{1}{1978242} = \frac{52059 + 1406 + 1}{1978242} = \frac{53466}{1978242} $$
Le PGCD du numérateur et du dénominateur est $\text{PGCD}(53466, 1978242) = 53466$.
La fraction irréductible est :
$$ \frac{53466}{1978242} = \frac{53466 \div 53466}{1978242 \div 53466} = \frac{1}{37} $$
Cette fraction est bien égale à $\frac{5}{185}$.

\subsection{Démonstration pour $n = 186$}
Soit $n = 186$. Nous cherchons $x, y, z \in \mathbb{Z}^{+}$ tels que $\frac{5}{186} = \frac{1}{x} + \frac{1}{y} + \frac{1}{z}$.
Posons $x = 38$, $y = 1768$, $z = 3124056$.
Les conditions $x > 0$, $y > 0$ et $z > 0$ sont satisfaites car $x, y, z \geq 1$.
Le PPCM des dénominateurs est $\text{PPCM}(38, 1768, 3124056) = 3124056$.
En réduisant au même dénominateur :
\begin{itemize}
    \item $\frac{1}{38} = \frac{82212}{3124056}$
    \item $\frac{1}{1768} = \frac{1767}{3124056}$
    \item $\frac{1}{3124056} = \frac{1}{3124056}$
\end{itemize}
La somme des numérateurs est :
$$ \frac{1}{38} + \frac{1}{1768} + \frac{1}{3124056} = \frac{82212 + 1767 + 1}{3124056} = \frac{83980}{3124056} $$
Le PGCD du numérateur et du dénominateur est $\text{PGCD}(83980, 3124056) = 16796$.
La fraction irréductible est :
$$ \frac{83980}{3124056} = \frac{83980 \div 16796}{3124056 \div 16796} = \frac{5}{186} $$
Cette fraction est bien égale à $\frac{5}{186}$.

\subsection{Démonstration pour $n = 187$}
Soit $n = 187$. Nous cherchons $x, y, z \in \mathbb{Z}^{+}$ tels que $\frac{5}{187} = \frac{1}{x} + \frac{1}{y} + \frac{1}{z}$.
Posons $x = 38$, $y = 2369$, $z = 16834114$.
Les conditions $x > 0$, $y > 0$ et $z > 0$ sont satisfaites car $x, y, z \geq 1$.
Le PPCM des dénominateurs est $\text{PPCM}(38, 2369, 16834114) = 16834114$.
En réduisant au même dénominateur :
\begin{itemize}
    \item $\frac{1}{38} = \frac{443003}{16834114}$
    \item $\frac{1}{2369} = \frac{7106}{16834114}$
    \item $\frac{1}{16834114} = \frac{1}{16834114}$
\end{itemize}
La somme des numérateurs est :
$$ \frac{1}{38} + \frac{1}{2369} + \frac{1}{16834114} = \frac{443003 + 7106 + 1}{16834114} = \frac{450110}{16834114} $$
Le PGCD du numérateur et du dénominateur est $\text{PGCD}(450110, 16834114) = 90022$.
La fraction irréductible est :
$$ \frac{450110}{16834114} = \frac{450110 \div 90022}{16834114 \div 90022} = \frac{5}{187} $$
Cette fraction est bien égale à $\frac{5}{187}$.

\subsection{Démonstration pour $n = 188$}
Soit $n = 188$. Nous cherchons $x, y, z \in \mathbb{Z}^{+}$ tels que $\frac{5}{188} = \frac{1}{x} + \frac{1}{y} + \frac{1}{z}$.
Posons $x = 38$, $y = 3573$, $z = 12762756$.
Les conditions $x > 0$, $y > 0$ et $z > 0$ sont satisfaites car $x, y, z \geq 1$.
Le PPCM des dénominateurs est $\text{PPCM}(38, 3573, 12762756) = 12762756$.
En réduisant au même dénominateur :
\begin{itemize}
    \item $\frac{1}{38} = \frac{335862}{12762756}$
    \item $\frac{1}{3573} = \frac{3572}{12762756}$
    \item $\frac{1}{12762756} = \frac{1}{12762756}$
\end{itemize}
La somme des numérateurs est :
$$ \frac{1}{38} + \frac{1}{3573} + \frac{1}{12762756} = \frac{335862 + 3572 + 1}{12762756} = \frac{339435}{12762756} $$
Le PGCD du numérateur et du dénominateur est $\text{PGCD}(339435, 12762756) = 67887$.
La fraction irréductible est :
$$ \frac{339435}{12762756} = \frac{339435 \div 67887}{12762756 \div 67887} = \frac{5}{188} $$
Cette fraction est bien égale à $\frac{5}{188}$.

\subsection{Démonstration pour $n = 189$}
Soit $n = 189$. Nous cherchons $x, y, z \in \mathbb{Z}^{+}$ tels que $\frac{5}{189} = \frac{1}{x} + \frac{1}{y} + \frac{1}{z}$.
Posons $x = 38$, $y = 7183$, $z = 51588306$.
Les conditions $x > 0$, $y > 0$ et $z > 0$ sont satisfaites car $x, y, z \geq 1$.
Le PPCM des dénominateurs est $\text{PPCM}(38, 7183, 51588306) = 51588306$.
En réduisant au même dénominateur :
\begin{itemize}
    \item $\frac{1}{38} = \frac{1357587}{51588306}$
    \item $\frac{1}{7183} = \frac{7182}{51588306}$
    \item $\frac{1}{51588306} = \frac{1}{51588306}$
\end{itemize}
La somme des numérateurs est :
$$ \frac{1}{38} + \frac{1}{7183} + \frac{1}{51588306} = \frac{1357587 + 7182 + 1}{51588306} = \frac{1364770}{51588306} $$
Le PGCD du numérateur et du dénominateur est $\text{PGCD}(1364770, 51588306) = 272954$.
La fraction irréductible est :
$$ \frac{1364770}{51588306} = \frac{1364770 \div 272954}{51588306 \div 272954} = \frac{5}{189} $$
Cette fraction est bien égale à $\frac{5}{189}$.

\subsection{Démonstration pour $n = 190$}
Soit $n = 190$. Nous cherchons $x, y, z \in \mathbb{Z}^{+}$ tels que $\frac{5}{190} = \frac{1}{x} + \frac{1}{y} + \frac{1}{z}$.
Posons $x = 39$, $y = 1483$, $z = 2197806$.
Les conditions $x > 0$, $y > 0$ et $z > 0$ sont satisfaites car $x, y, z \geq 1$.
Le PPCM des dénominateurs est $\text{PPCM}(39, 1483, 2197806) = 2197806$.
En réduisant au même dénominateur :
\begin{itemize}
    \item $\frac{1}{39} = \frac{56354}{2197806}$
    \item $\frac{1}{1483} = \frac{1482}{2197806}$
    \item $\frac{1}{2197806} = \frac{1}{2197806}$
\end{itemize}
La somme des numérateurs est :
$$ \frac{1}{39} + \frac{1}{1483} + \frac{1}{2197806} = \frac{56354 + 1482 + 1}{2197806} = \frac{57837}{2197806} $$
Le PGCD du numérateur et du dénominateur est $\text{PGCD}(57837, 2197806) = 57837$.
La fraction irréductible est :
$$ \frac{57837}{2197806} = \frac{57837 \div 57837}{2197806 \div 57837} = \frac{1}{38} $$
Cette fraction est bien égale à $\frac{5}{190}$.

\subsection{Démonstration pour $n = 191$}
Soit $n = 191$. Nous cherchons $x, y, z \in \mathbb{Z}^{+}$ tels que $\frac{5}{191} = \frac{1}{x} + \frac{1}{y} + \frac{1}{z}$.
Posons $x = 39$, $y = 1863$, $z = 4625829$.
Les conditions $x > 0$, $y > 0$ et $z > 0$ sont satisfaites car $x, y, z \geq 1$.
Le PPCM des dénominateurs est $\text{PPCM}(39, 1863, 4625829) = 4625829$.
En réduisant au même dénominateur :
\begin{itemize}
    \item $\frac{1}{39} = \frac{118611}{4625829}$
    \item $\frac{1}{1863} = \frac{2483}{4625829}$
    \item $\frac{1}{4625829} = \frac{1}{4625829}$
\end{itemize}
La somme des numérateurs est :
$$ \frac{1}{39} + \frac{1}{1863} + \frac{1}{4625829} = \frac{118611 + 2483 + 1}{4625829} = \frac{121095}{4625829} $$
Le PGCD du numérateur et du dénominateur est $\text{PGCD}(121095, 4625829) = 24219$.
La fraction irréductible est :
$$ \frac{121095}{4625829} = \frac{121095 \div 24219}{4625829 \div 24219} = \frac{5}{191} $$
Cette fraction est bien égale à $\frac{5}{191}$.

\subsection{Démonstration pour $n = 192$}
Soit $n = 192$. Nous cherchons $x, y, z \in \mathbb{Z}^{+}$ tels que $\frac{5}{192} = \frac{1}{x} + \frac{1}{y} + \frac{1}{z}$.
Posons $x = 39$, $y = 2497$, $z = 6232512$.
Les conditions $x > 0$, $y > 0$ et $z > 0$ sont satisfaites car $x, y, z \geq 1$.
Le PPCM des dénominateurs est $\text{PPCM}(39, 2497, 6232512) = 6232512$.
En réduisant au même dénominateur :
\begin{itemize}
    \item $\frac{1}{39} = \frac{159808}{6232512}$
    \item $\frac{1}{2497} = \frac{2496}{6232512}$
    \item $\frac{1}{6232512} = \frac{1}{6232512}$
\end{itemize}
La somme des numérateurs est :
$$ \frac{1}{39} + \frac{1}{2497} + \frac{1}{6232512} = \frac{159808 + 2496 + 1}{6232512} = \frac{162305}{6232512} $$
Le PGCD du numérateur et du dénominateur est $\text{PGCD}(162305, 6232512) = 32461$.
La fraction irréductible est :
$$ \frac{162305}{6232512} = \frac{162305 \div 32461}{6232512 \div 32461} = \frac{5}{192} $$
Cette fraction est bien égale à $\frac{5}{192}$.

\subsection{Démonstration pour $n = 193$}
Soit $n = 193$. Nous cherchons $x, y, z \in \mathbb{Z}^{+}$ tels que $\frac{5}{193} = \frac{1}{x} + \frac{1}{y} + \frac{1}{z}$.
Posons $x = 39$, $y = 3764$, $z = 28331628$.
Les conditions $x > 0$, $y > 0$ et $z > 0$ sont satisfaites car $x, y, z \geq 1$.
Le PPCM des dénominateurs est $\text{PPCM}(39, 3764, 28331628) = 28331628$.
En réduisant au même dénominateur :
\begin{itemize}
    \item $\frac{1}{39} = \frac{726452}{28331628}$
    \item $\frac{1}{3764} = \frac{7527}{28331628}$
    \item $\frac{1}{28331628} = \frac{1}{28331628}$
\end{itemize}
La somme des numérateurs est :
$$ \frac{1}{39} + \frac{1}{3764} + \frac{1}{28331628} = \frac{726452 + 7527 + 1}{28331628} = \frac{733980}{28331628} $$
Le PGCD du numérateur et du dénominateur est $\text{PGCD}(733980, 28331628) = 146796$.
La fraction irréductible est :
$$ \frac{733980}{28331628} = \frac{733980 \div 146796}{28331628 \div 146796} = \frac{5}{193} $$
Cette fraction est bien égale à $\frac{5}{193}$.

\subsection{Démonstration pour $n = 194$}
Soit $n = 194$. Nous cherchons $x, y, z \in \mathbb{Z}^{+}$ tels que $\frac{5}{194} = \frac{1}{x} + \frac{1}{y} + \frac{1}{z}$.
Posons $x = 39$, $y = 7567$, $z = 57251922$.
Les conditions $x > 0$, $y > 0$ et $z > 0$ sont satisfaites car $x, y, z \geq 1$.
Le PPCM des dénominateurs est $\text{PPCM}(39, 7567, 57251922) = 57251922$.
En réduisant au même dénominateur :
\begin{itemize}
    \item $\frac{1}{39} = \frac{1467998}{57251922}$
    \item $\frac{1}{7567} = \frac{7566}{57251922}$
    \item $\frac{1}{57251922} = \frac{1}{57251922}$
\end{itemize}
La somme des numérateurs est :
$$ \frac{1}{39} + \frac{1}{7567} + \frac{1}{57251922} = \frac{1467998 + 7566 + 1}{57251922} = \frac{1475565}{57251922} $$
Le PGCD du numérateur et du dénominateur est $\text{PGCD}(1475565, 57251922) = 295113$.
La fraction irréductible est :
$$ \frac{1475565}{57251922} = \frac{1475565 \div 295113}{57251922 \div 295113} = \frac{5}{194} $$
Cette fraction est bien égale à $\frac{5}{194}$.

\subsection{Démonstration pour $n = 195$}
Soit $n = 195$. Nous cherchons $x, y, z \in \mathbb{Z}^{+}$ tels que $\frac{5}{195} = \frac{1}{x} + \frac{1}{y} + \frac{1}{z}$.
Posons $x = 40$, $y = 1561$, $z = 2435160$.
Les conditions $x > 0$, $y > 0$ et $z > 0$ sont satisfaites car $x, y, z \geq 1$.
Le PPCM des dénominateurs est $\text{PPCM}(40, 1561, 2435160) = 2435160$.
En réduisant au même dénominateur :
\begin{itemize}
    \item $\frac{1}{40} = \frac{60879}{2435160}$
    \item $\frac{1}{1561} = \frac{1560}{2435160}$
    \item $\frac{1}{2435160} = \frac{1}{2435160}$
\end{itemize}
La somme des numérateurs est :
$$ \frac{1}{40} + \frac{1}{1561} + \frac{1}{2435160} = \frac{60879 + 1560 + 1}{2435160} = \frac{62440}{2435160} $$
Le PGCD du numérateur et du dénominateur est $\text{PGCD}(62440, 2435160) = 62440$.
La fraction irréductible est :
$$ \frac{62440}{2435160} = \frac{62440 \div 62440}{2435160 \div 62440} = \frac{1}{39} $$
Cette fraction est bien égale à $\frac{5}{195}$.

\subsection{Démonstration pour $n = 196$}
Soit $n = 196$. Nous cherchons $x, y, z \in \mathbb{Z}^{+}$ tels que $\frac{5}{196} = \frac{1}{x} + \frac{1}{y} + \frac{1}{z}$.
Posons $x = 40$, $y = 1961$, $z = 3843560$.
Les conditions $x > 0$, $y > 0$ et $z > 0$ sont satisfaites car $x, y, z \geq 1$.
Le PPCM des dénominateurs est $\text{PPCM}(40, 1961, 3843560) = 3843560$.
En réduisant au même dénominateur :
\begin{itemize}
    \item $\frac{1}{40} = \frac{96089}{3843560}$
    \item $\frac{1}{1961} = \frac{1960}{3843560}$
    \item $\frac{1}{3843560} = \frac{1}{3843560}$
\end{itemize}
La somme des numérateurs est :
$$ \frac{1}{40} + \frac{1}{1961} + \frac{1}{3843560} = \frac{96089 + 1960 + 1}{3843560} = \frac{98050}{3843560} $$
Le PGCD du numérateur et du dénominateur est $\text{PGCD}(98050, 3843560) = 19610$.
La fraction irréductible est :
$$ \frac{98050}{3843560} = \frac{98050 \div 19610}{3843560 \div 19610} = \frac{5}{196} $$
Cette fraction est bien égale à $\frac{5}{196}$.

\subsection{Démonstration pour $n = 197$}
Soit $n = 197$. Nous cherchons $x, y, z \in \mathbb{Z}^{+}$ tels que $\frac{5}{197} = \frac{1}{x} + \frac{1}{y} + \frac{1}{z}$.
Posons $x = 40$, $y = 2627$, $z = 20700760$.
Les conditions $x > 0$, $y > 0$ et $z > 0$ sont satisfaites car $x, y, z \geq 1$.
Le PPCM des dénominateurs est $\text{PPCM}(40, 2627, 20700760) = 20700760$.
En réduisant au même dénominateur :
\begin{itemize}
    \item $\frac{1}{40} = \frac{517519}{20700760}$
    \item $\frac{1}{2627} = \frac{7880}{20700760}$
    \item $\frac{1}{20700760} = \frac{1}{20700760}$
\end{itemize}
La somme des numérateurs est :
$$ \frac{1}{40} + \frac{1}{2627} + \frac{1}{20700760} = \frac{517519 + 7880 + 1}{20700760} = \frac{525400}{20700760} $$
Le PGCD du numérateur et du dénominateur est $\text{PGCD}(525400, 20700760) = 105080$.
La fraction irréductible est :
$$ \frac{525400}{20700760} = \frac{525400 \div 105080}{20700760 \div 105080} = \frac{5}{197} $$
Cette fraction est bien égale à $\frac{5}{197}$.

\subsection{Démonstration pour $n = 198$}
Soit $n = 198$. Nous cherchons $x, y, z \in \mathbb{Z}^{+}$ tels que $\frac{5}{198} = \frac{1}{x} + \frac{1}{y} + \frac{1}{z}$.
Posons $x = 40$, $y = 3961$, $z = 15685560$.
Les conditions $x > 0$, $y > 0$ et $z > 0$ sont satisfaites car $x, y, z \geq 1$.
Le PPCM des dénominateurs est $\text{PPCM}(40, 3961, 15685560) = 15685560$.
En réduisant au même dénominateur :
\begin{itemize}
    \item $\frac{1}{40} = \frac{392139}{15685560}$
    \item $\frac{1}{3961} = \frac{3960}{15685560}$
    \item $\frac{1}{15685560} = \frac{1}{15685560}$
\end{itemize}
La somme des numérateurs est :
$$ \frac{1}{40} + \frac{1}{3961} + \frac{1}{15685560} = \frac{392139 + 3960 + 1}{15685560} = \frac{396100}{15685560} $$
Le PGCD du numérateur et du dénominateur est $\text{PGCD}(396100, 15685560) = 79220$.
La fraction irréductible est :
$$ \frac{396100}{15685560} = \frac{396100 \div 79220}{15685560 \div 79220} = \frac{5}{198} $$
Cette fraction est bien égale à $\frac{5}{198}$.

\subsection{Démonstration pour $n = 199$}
Soit $n = 199$. Nous cherchons $x, y, z \in \mathbb{Z}^{+}$ tels que $\frac{5}{199} = \frac{1}{x} + \frac{1}{y} + \frac{1}{z}$.
Posons $x = 40$, $y = 7961$, $z = 63369560$.
Les conditions $x > 0$, $y > 0$ et $z > 0$ sont satisfaites car $x, y, z \geq 1$.
Le PPCM des dénominateurs est $\text{PPCM}(40, 7961, 63369560) = 63369560$.
En réduisant au même dénominateur :
\begin{itemize}
    \item $\frac{1}{40} = \frac{1584239}{63369560}$
    \item $\frac{1}{7961} = \frac{7960}{63369560}$
    \item $\frac{1}{63369560} = \frac{1}{63369560}$
\end{itemize}
La somme des numérateurs est :
$$ \frac{1}{40} + \frac{1}{7961} + \frac{1}{63369560} = \frac{1584239 + 7960 + 1}{63369560} = \frac{1592200}{63369560} $$
Le PGCD du numérateur et du dénominateur est $\text{PGCD}(1592200, 63369560) = 318440$.
La fraction irréductible est :
$$ \frac{1592200}{63369560} = \frac{1592200 \div 318440}{63369560 \div 318440} = \frac{5}{199} $$
Cette fraction est bien égale à $\frac{5}{199}$.

\subsection{Démonstration pour $n = 200$}
Soit $n = 200$. Nous cherchons $x, y, z \in \mathbb{Z}^{+}$ tels que $\frac{5}{200} = \frac{1}{x} + \frac{1}{y} + \frac{1}{z}$.
Posons $x = 41$, $y = 1641$, $z = 2691240$.
Les conditions $x > 0$, $y > 0$ et $z > 0$ sont satisfaites car $x, y, z \geq 1$.
Le PPCM des dénominateurs est $\text{PPCM}(41, 1641, 2691240) = 2691240$.
En réduisant au même dénominateur :
\begin{itemize}
    \item $\frac{1}{41} = \frac{65640}{2691240}$
    \item $\frac{1}{1641} = \frac{1640}{2691240}$
    \item $\frac{1}{2691240} = \frac{1}{2691240}$
\end{itemize}
La somme des numérateurs est :
$$ \frac{1}{41} + \frac{1}{1641} + \frac{1}{2691240} = \frac{65640 + 1640 + 1}{2691240} = \frac{67281}{2691240} $$
Le PGCD du numérateur et du dénominateur est $\text{PGCD}(67281, 2691240) = 67281$.
La fraction irréductible est :
$$ \frac{67281}{2691240} = \frac{67281 \div 67281}{2691240 \div 67281} = \frac{1}{40} $$
Cette fraction est bien égale à $\frac{5}{200}$.

\subsection{Démonstration pour $n = 201$}
Soit $n = 201$. Nous cherchons $x, y, z \in \mathbb{Z}^{+}$ tels que $\frac{5}{201} = \frac{1}{x} + \frac{1}{y} + \frac{1}{z}$.
Posons $x = 41$, $y = 2061$, $z = 5661567$.
Les conditions $x > 0$, $y > 0$ et $z > 0$ sont satisfaites car $x, y, z \geq 1$.
Le PPCM des dénominateurs est $\text{PPCM}(41, 2061, 5661567) = 5661567$.
En réduisant au même dénominateur :
\begin{itemize}
    \item $\frac{1}{41} = \frac{138087}{5661567}$
    \item $\frac{1}{2061} = \frac{2747}{5661567}$
    \item $\frac{1}{5661567} = \frac{1}{5661567}$
\end{itemize}
La somme des numérateurs est :
$$ \frac{1}{41} + \frac{1}{2061} + \frac{1}{5661567} = \frac{138087 + 2747 + 1}{5661567} = \frac{140835}{5661567} $$
Le PGCD du numérateur et du dénominateur est $\text{PGCD}(140835, 5661567) = 28167$.
La fraction irréductible est :
$$ \frac{140835}{5661567} = \frac{140835 \div 28167}{5661567 \div 28167} = \frac{5}{201} $$
Cette fraction est bien égale à $\frac{5}{201}$.

\subsection{Démonstration pour $n = 202$}
Soit $n = 202$. Nous cherchons $x, y, z \in \mathbb{Z}^{+}$ tels que $\frac{5}{202} = \frac{1}{x} + \frac{1}{y} + \frac{1}{z}$.
Posons $x = 41$, $y = 2761$, $z = 22866602$.
Les conditions $x > 0$, $y > 0$ et $z > 0$ sont satisfaites car $x, y, z \geq 1$.
Le PPCM des dénominateurs est $\text{PPCM}(41, 2761, 22866602) = 22866602$.
En réduisant au même dénominateur :
\begin{itemize}
    \item $\frac{1}{41} = \frac{557722}{22866602}$
    \item $\frac{1}{2761} = \frac{8282}{22866602}$
    \item $\frac{1}{22866602} = \frac{1}{22866602}$
\end{itemize}
La somme des numérateurs est :
$$ \frac{1}{41} + \frac{1}{2761} + \frac{1}{22866602} = \frac{557722 + 8282 + 1}{22866602} = \frac{566005}{22866602} $$
Le PGCD du numérateur et du dénominateur est $\text{PGCD}(566005, 22866602) = 113201$.
La fraction irréductible est :
$$ \frac{566005}{22866602} = \frac{566005 \div 113201}{22866602 \div 113201} = \frac{5}{202} $$
Cette fraction est bien égale à $\frac{5}{202}$.

\subsection{Démonstration pour $n = 203$}
Soit $n = 203$. Nous cherchons $x, y, z \in \mathbb{Z}^{+}$ tels que $\frac{5}{203} = \frac{1}{x} + \frac{1}{y} + \frac{1}{z}$.
Posons $x = 41$, $y = 4162$, $z = 34640326$.
Les conditions $x > 0$, $y > 0$ et $z > 0$ sont satisfaites car $x, y, z \geq 1$.
Le PPCM des dénominateurs est $\text{PPCM}(41, 4162, 34640326) = 34640326$.
En réduisant au même dénominateur :
\begin{itemize}
    \item $\frac{1}{41} = \frac{844886}{34640326}$
    \item $\frac{1}{4162} = \frac{8323}{34640326}$
    \item $\frac{1}{34640326} = \frac{1}{34640326}$
\end{itemize}
La somme des numérateurs est :
$$ \frac{1}{41} + \frac{1}{4162} + \frac{1}{34640326} = \frac{844886 + 8323 + 1}{34640326} = \frac{853210}{34640326} $$
Le PGCD du numérateur et du dénominateur est $\text{PGCD}(853210, 34640326) = 170642$.
La fraction irréductible est :
$$ \frac{853210}{34640326} = \frac{853210 \div 170642}{34640326 \div 170642} = \frac{5}{203} $$
Cette fraction est bien égale à $\frac{5}{203}$.

\subsection{Démonstration pour $n = 204$}
Soit $n = 204$. Nous cherchons $x, y, z \in \mathbb{Z}^{+}$ tels que $\frac{5}{204} = \frac{1}{x} + \frac{1}{y} + \frac{1}{z}$.
Posons $x = 41$, $y = 8365$, $z = 69964860$.
Les conditions $x > 0$, $y > 0$ et $z > 0$ sont satisfaites car $x, y, z \geq 1$.
Le PPCM des dénominateurs est $\text{PPCM}(41, 8365, 69964860) = 69964860$.
En réduisant au même dénominateur :
\begin{itemize}
    \item $\frac{1}{41} = \frac{1706460}{69964860}$
    \item $\frac{1}{8365} = \frac{8364}{69964860}$
    \item $\frac{1}{69964860} = \frac{1}{69964860}$
\end{itemize}
La somme des numérateurs est :
$$ \frac{1}{41} + \frac{1}{8365} + \frac{1}{69964860} = \frac{1706460 + 8364 + 1}{69964860} = \frac{1714825}{69964860} $$
Le PGCD du numérateur et du dénominateur est $\text{PGCD}(1714825, 69964860) = 342965$.
La fraction irréductible est :
$$ \frac{1714825}{69964860} = \frac{1714825 \div 342965}{69964860 \div 342965} = \frac{5}{204} $$
Cette fraction est bien égale à $\frac{5}{204}$.

\subsection{Démonstration pour $n = 205$}
Soit $n = 205$. Nous cherchons $x, y, z \in \mathbb{Z}^{+}$ tels que $\frac{5}{205} = \frac{1}{x} + \frac{1}{y} + \frac{1}{z}$.
Posons $x = 42$, $y = 1723$, $z = 2967006$.
Les conditions $x > 0$, $y > 0$ et $z > 0$ sont satisfaites car $x, y, z \geq 1$.
Le PPCM des dénominateurs est $\text{PPCM}(42, 1723, 2967006) = 2967006$.
En réduisant au même dénominateur :
\begin{itemize}
    \item $\frac{1}{42} = \frac{70643}{2967006}$
    \item $\frac{1}{1723} = \frac{1722}{2967006}$
    \item $\frac{1}{2967006} = \frac{1}{2967006}$
\end{itemize}
La somme des numérateurs est :
$$ \frac{1}{42} + \frac{1}{1723} + \frac{1}{2967006} = \frac{70643 + 1722 + 1}{2967006} = \frac{72366}{2967006} $$
Le PGCD du numérateur et du dénominateur est $\text{PGCD}(72366, 2967006) = 72366$.
La fraction irréductible est :
$$ \frac{72366}{2967006} = \frac{72366 \div 72366}{2967006 \div 72366} = \frac{1}{41} $$
Cette fraction est bien égale à $\frac{5}{205}$.

\subsection{Démonstration pour $n = 206$}
Soit $n = 206$. Nous cherchons $x, y, z \in \mathbb{Z}^{+}$ tels que $\frac{5}{206} = \frac{1}{x} + \frac{1}{y} + \frac{1}{z}$.
Posons $x = 42$, $y = 2164$, $z = 4680732$.
Les conditions $x > 0$, $y > 0$ et $z > 0$ sont satisfaites car $x, y, z \geq 1$.
Le PPCM des dénominateurs est $\text{PPCM}(42, 2164, 4680732) = 4680732$.
En réduisant au même dénominateur :
\begin{itemize}
    \item $\frac{1}{42} = \frac{111446}{4680732}$
    \item $\frac{1}{2164} = \frac{2163}{4680732}$
    \item $\frac{1}{4680732} = \frac{1}{4680732}$
\end{itemize}
La somme des numérateurs est :
$$ \frac{1}{42} + \frac{1}{2164} + \frac{1}{4680732} = \frac{111446 + 2163 + 1}{4680732} = \frac{113610}{4680732} $$
Le PGCD du numérateur et du dénominateur est $\text{PGCD}(113610, 4680732) = 22722$.
La fraction irréductible est :
$$ \frac{113610}{4680732} = \frac{113610 \div 22722}{4680732 \div 22722} = \frac{5}{206} $$
Cette fraction est bien égale à $\frac{5}{206}$.

\subsection{Démonstration pour $n = 207$}
Soit $n = 207$. Nous cherchons $x, y, z \in \mathbb{Z}^{+}$ tels que $\frac{5}{207} = \frac{1}{x} + \frac{1}{y} + \frac{1}{z}$.
Posons $x = 42$, $y = 2899$, $z = 8401302$.
Les conditions $x > 0$, $y > 0$ et $z > 0$ sont satisfaites car $x, y, z \geq 1$.
Le PPCM des dénominateurs est $\text{PPCM}(42, 2899, 8401302) = 8401302$.
En réduisant au même dénominateur :
\begin{itemize}
    \item $\frac{1}{42} = \frac{200031}{8401302}$
    \item $\frac{1}{2899} = \frac{2898}{8401302}$
    \item $\frac{1}{8401302} = \frac{1}{8401302}$
\end{itemize}
La somme des numérateurs est :
$$ \frac{1}{42} + \frac{1}{2899} + \frac{1}{8401302} = \frac{200031 + 2898 + 1}{8401302} = \frac{202930}{8401302} $$
Le PGCD du numérateur et du dénominateur est $\text{PGCD}(202930, 8401302) = 40586$.
La fraction irréductible est :
$$ \frac{202930}{8401302} = \frac{202930 \div 40586}{8401302 \div 40586} = \frac{5}{207} $$
Cette fraction est bien égale à $\frac{5}{207}$.

\subsection{Démonstration pour $n = 208$}
Soit $n = 208$. Nous cherchons $x, y, z \in \mathbb{Z}^{+}$ tels que $\frac{5}{208} = \frac{1}{x} + \frac{1}{y} + \frac{1}{z}$.
Posons $x = 42$, $y = 4369$, $z = 19083792$.
Les conditions $x > 0$, $y > 0$ et $z > 0$ sont satisfaites car $x, y, z \geq 1$.
Le PPCM des dénominateurs est $\text{PPCM}(42, 4369, 19083792) = 19083792$.
En réduisant au même dénominateur :
\begin{itemize}
    \item $\frac{1}{42} = \frac{454376}{19083792}$
    \item $\frac{1}{4369} = \frac{4368}{19083792}$
    \item $\frac{1}{19083792} = \frac{1}{19083792}$
\end{itemize}
La somme des numérateurs est :
$$ \frac{1}{42} + \frac{1}{4369} + \frac{1}{19083792} = \frac{454376 + 4368 + 1}{19083792} = \frac{458745}{19083792} $$
Le PGCD du numérateur et du dénominateur est $\text{PGCD}(458745, 19083792) = 91749$.
La fraction irréductible est :
$$ \frac{458745}{19083792} = \frac{458745 \div 91749}{19083792 \div 91749} = \frac{5}{208} $$
Cette fraction est bien égale à $\frac{5}{208}$.

\subsection{Démonstration pour $n = 209$}
Soit $n = 209$. Nous cherchons $x, y, z \in \mathbb{Z}^{+}$ tels que $\frac{5}{209} = \frac{1}{x} + \frac{1}{y} + \frac{1}{z}$.
Posons $x = 42$, $y = 8779$, $z = 77062062$.
Les conditions $x > 0$, $y > 0$ et $z > 0$ sont satisfaites car $x, y, z \geq 1$.
Le PPCM des dénominateurs est $\text{PPCM}(42, 8779, 77062062) = 77062062$.
En réduisant au même dénominateur :
\begin{itemize}
    \item $\frac{1}{42} = \frac{1834811}{77062062}$
    \item $\frac{1}{8779} = \frac{8778}{77062062}$
    \item $\frac{1}{77062062} = \frac{1}{77062062}$
\end{itemize}
La somme des numérateurs est :
$$ \frac{1}{42} + \frac{1}{8779} + \frac{1}{77062062} = \frac{1834811 + 8778 + 1}{77062062} = \frac{1843590}{77062062} $$
Le PGCD du numérateur et du dénominateur est $\text{PGCD}(1843590, 77062062) = 368718$.
La fraction irréductible est :
$$ \frac{1843590}{77062062} = \frac{1843590 \div 368718}{77062062 \div 368718} = \frac{5}{209} $$
Cette fraction est bien égale à $\frac{5}{209}$.

\subsection{Démonstration pour $n = 210$}
Soit $n = 210$. Nous cherchons $x, y, z \in \mathbb{Z}^{+}$ tels que $\frac{5}{210} = \frac{1}{x} + \frac{1}{y} + \frac{1}{z}$.
Posons $x = 43$, $y = 1807$, $z = 3263442$.
Les conditions $x > 0$, $y > 0$ et $z > 0$ sont satisfaites car $x, y, z \geq 1$.
Le PPCM des dénominateurs est $\text{PPCM}(43, 1807, 3263442) = 3263442$.
En réduisant au même dénominateur :
\begin{itemize}
    \item $\frac{1}{43} = \frac{75894}{3263442}$
    \item $\frac{1}{1807} = \frac{1806}{3263442}$
    \item $\frac{1}{3263442} = \frac{1}{3263442}$
\end{itemize}
La somme des numérateurs est :
$$ \frac{1}{43} + \frac{1}{1807} + \frac{1}{3263442} = \frac{75894 + 1806 + 1}{3263442} = \frac{77701}{3263442} $$
Le PGCD du numérateur et du dénominateur est $\text{PGCD}(77701, 3263442) = 77701$.
La fraction irréductible est :
$$ \frac{77701}{3263442} = \frac{77701 \div 77701}{3263442 \div 77701} = \frac{1}{42} $$
Cette fraction est bien égale à $\frac{5}{210}$.

\subsection{Démonstration pour $n = 211$}
Soit $n = 211$. Nous cherchons $x, y, z \in \mathbb{Z}^{+}$ tels que $\frac{5}{211} = \frac{1}{x} + \frac{1}{y} + \frac{1}{z}$.
Posons $x = 43$, $y = 2279$, $z = 480869$.
Les conditions $x > 0$, $y > 0$ et $z > 0$ sont satisfaites car $x, y, z \geq 1$.
Le PPCM des dénominateurs est $\text{PPCM}(43, 2279, 480869) = 480869$.
En réduisant au même dénominateur :
\begin{itemize}
    \item $\frac{1}{43} = \frac{11183}{480869}$
    \item $\frac{1}{2279} = \frac{211}{480869}$
    \item $\frac{1}{480869} = \frac{1}{480869}$
\end{itemize}
La somme des numérateurs est :
$$ \frac{1}{43} + \frac{1}{2279} + \frac{1}{480869} = \frac{11183 + 211 + 1}{480869} = \frac{11395}{480869} $$
Le PGCD du numérateur et du dénominateur est $\text{PGCD}(11395, 480869) = 2279$.
La fraction irréductible est :
$$ \frac{11395}{480869} = \frac{11395 \div 2279}{480869 \div 2279} = \frac{5}{211} $$
Cette fraction est bien égale à $\frac{5}{211}$.

\subsection{Démonstration pour $n = 212$}
Soit $n = 212$. Nous cherchons $x, y, z \in \mathbb{Z}^{+}$ tels que $\frac{5}{212} = \frac{1}{x} + \frac{1}{y} + \frac{1}{z}$.
Posons $x = 43$, $y = 3039$, $z = 27703524$.
Les conditions $x > 0$, $y > 0$ et $z > 0$ sont satisfaites car $x, y, z \geq 1$.
Le PPCM des dénominateurs est $\text{PPCM}(43, 3039, 27703524) = 27703524$.
En réduisant au même dénominateur :
\begin{itemize}
    \item $\frac{1}{43} = \frac{644268}{27703524}$
    \item $\frac{1}{3039} = \frac{9116}{27703524}$
    \item $\frac{1}{27703524} = \frac{1}{27703524}$
\end{itemize}
La somme des numérateurs est :
$$ \frac{1}{43} + \frac{1}{3039} + \frac{1}{27703524} = \frac{644268 + 9116 + 1}{27703524} = \frac{653385}{27703524} $$
Le PGCD du numérateur et du dénominateur est $\text{PGCD}(653385, 27703524) = 130677$.
La fraction irréductible est :
$$ \frac{653385}{27703524} = \frac{653385 \div 130677}{27703524 \div 130677} = \frac{5}{212} $$
Cette fraction est bien égale à $\frac{5}{212}$.

\subsection{Démonstration pour $n = 213$}
Soit $n = 213$. Nous cherchons $x, y, z \in \mathbb{Z}^{+}$ tels que $\frac{5}{213} = \frac{1}{x} + \frac{1}{y} + \frac{1}{z}$.
Posons $x = 43$, $y = 4580$, $z = 41948220$.
Les conditions $x > 0$, $y > 0$ et $z > 0$ sont satisfaites car $x, y, z \geq 1$.
Le PPCM des dénominateurs est $\text{PPCM}(43, 4580, 41948220) = 41948220$.
En réduisant au même dénominateur :
\begin{itemize}
    \item $\frac{1}{43} = \frac{975540}{41948220}$
    \item $\frac{1}{4580} = \frac{9159}{41948220}$
    \item $\frac{1}{41948220} = \frac{1}{41948220}$
\end{itemize}
La somme des numérateurs est :
$$ \frac{1}{43} + \frac{1}{4580} + \frac{1}{41948220} = \frac{975540 + 9159 + 1}{41948220} = \frac{984700}{41948220} $$
Le PGCD du numérateur et du dénominateur est $\text{PGCD}(984700, 41948220) = 196940$.
La fraction irréductible est :
$$ \frac{984700}{41948220} = \frac{984700 \div 196940}{41948220 \div 196940} = \frac{5}{213} $$
Cette fraction est bien égale à $\frac{5}{213}$.

\subsection{Démonstration pour $n = 214$}
Soit $n = 214$. Nous cherchons $x, y, z \in \mathbb{Z}^{+}$ tels que $\frac{5}{214} = \frac{1}{x} + \frac{1}{y} + \frac{1}{z}$.
Posons $x = 43$, $y = 9203$, $z = 84686006$.
Les conditions $x > 0$, $y > 0$ et $z > 0$ sont satisfaites car $x, y, z \geq 1$.
Le PPCM des dénominateurs est $\text{PPCM}(43, 9203, 84686006) = 84686006$.
En réduisant au même dénominateur :
\begin{itemize}
    \item $\frac{1}{43} = \frac{1969442}{84686006}$
    \item $\frac{1}{9203} = \frac{9202}{84686006}$
    \item $\frac{1}{84686006} = \frac{1}{84686006}$
\end{itemize}
La somme des numérateurs est :
$$ \frac{1}{43} + \frac{1}{9203} + \frac{1}{84686006} = \frac{1969442 + 9202 + 1}{84686006} = \frac{1978645}{84686006} $$
Le PGCD du numérateur et du dénominateur est $\text{PGCD}(1978645, 84686006) = 395729$.
La fraction irréductible est :
$$ \frac{1978645}{84686006} = \frac{1978645 \div 395729}{84686006 \div 395729} = \frac{5}{214} $$
Cette fraction est bien égale à $\frac{5}{214}$.

\subsection{Démonstration pour $n = 215$}
Soit $n = 215$. Nous cherchons $x, y, z \in \mathbb{Z}^{+}$ tels que $\frac{5}{215} = \frac{1}{x} + \frac{1}{y} + \frac{1}{z}$.
Posons $x = 44$, $y = 1893$, $z = 3581556$.
Les conditions $x > 0$, $y > 0$ et $z > 0$ sont satisfaites car $x, y, z \geq 1$.
Le PPCM des dénominateurs est $\text{PPCM}(44, 1893, 3581556) = 3581556$.
En réduisant au même dénominateur :
\begin{itemize}
    \item $\frac{1}{44} = \frac{81399}{3581556}$
    \item $\frac{1}{1893} = \frac{1892}{3581556}$
    \item $\frac{1}{3581556} = \frac{1}{3581556}$
\end{itemize}
La somme des numérateurs est :
$$ \frac{1}{44} + \frac{1}{1893} + \frac{1}{3581556} = \frac{81399 + 1892 + 1}{3581556} = \frac{83292}{3581556} $$
Le PGCD du numérateur et du dénominateur est $\text{PGCD}(83292, 3581556) = 83292$.
La fraction irréductible est :
$$ \frac{83292}{3581556} = \frac{83292 \div 83292}{3581556 \div 83292} = \frac{1}{43} $$
Cette fraction est bien égale à $\frac{5}{215}$.

\subsection{Démonstration pour $n = 216$}
Soit $n = 216$. Nous cherchons $x, y, z \in \mathbb{Z}^{+}$ tels que $\frac{5}{216} = \frac{1}{x} + \frac{1}{y} + \frac{1}{z}$.
Posons $x = 44$, $y = 2377$, $z = 5647752$.
Les conditions $x > 0$, $y > 0$ et $z > 0$ sont satisfaites car $x, y, z \geq 1$.
Le PPCM des dénominateurs est $\text{PPCM}(44, 2377, 5647752) = 5647752$.
En réduisant au même dénominateur :
\begin{itemize}
    \item $\frac{1}{44} = \frac{128358}{5647752}$
    \item $\frac{1}{2377} = \frac{2376}{5647752}$
    \item $\frac{1}{5647752} = \frac{1}{5647752}$
\end{itemize}
La somme des numérateurs est :
$$ \frac{1}{44} + \frac{1}{2377} + \frac{1}{5647752} = \frac{128358 + 2376 + 1}{5647752} = \frac{130735}{5647752} $$
Le PGCD du numérateur et du dénominateur est $\text{PGCD}(130735, 5647752) = 26147$.
La fraction irréductible est :
$$ \frac{130735}{5647752} = \frac{130735 \div 26147}{5647752 \div 26147} = \frac{5}{216} $$
Cette fraction est bien égale à $\frac{5}{216}$.

\subsection{Démonstration pour $n = 217$}
Soit $n = 217$. Nous cherchons $x, y, z \in \mathbb{Z}^{+}$ tels que $\frac{5}{217} = \frac{1}{x} + \frac{1}{y} + \frac{1}{z}$.
Posons $x = 44$, $y = 3183$, $z = 30391284$.
Les conditions $x > 0$, $y > 0$ et $z > 0$ sont satisfaites car $x, y, z \geq 1$.
Le PPCM des dénominateurs est $\text{PPCM}(44, 3183, 30391284) = 30391284$.
En réduisant au même dénominateur :
\begin{itemize}
    \item $\frac{1}{44} = \frac{690711}{30391284}$
    \item $\frac{1}{3183} = \frac{9548}{30391284}$
    \item $\frac{1}{30391284} = \frac{1}{30391284}$
\end{itemize}
La somme des numérateurs est :
$$ \frac{1}{44} + \frac{1}{3183} + \frac{1}{30391284} = \frac{690711 + 9548 + 1}{30391284} = \frac{700260}{30391284} $$
Le PGCD du numérateur et du dénominateur est $\text{PGCD}(700260, 30391284) = 140052$.
La fraction irréductible est :
$$ \frac{700260}{30391284} = \frac{700260 \div 140052}{30391284 \div 140052} = \frac{5}{217} $$
Cette fraction est bien égale à $\frac{5}{217}$.

\subsection{Démonstration pour $n = 218$}
Soit $n = 218$. Nous cherchons $x, y, z \in \mathbb{Z}^{+}$ tels que $\frac{5}{218} = \frac{1}{x} + \frac{1}{y} + \frac{1}{z}$.
Posons $x = 44$, $y = 4797$, $z = 23006412$.
Les conditions $x > 0$, $y > 0$ et $z > 0$ sont satisfaites car $x, y, z \geq 1$.
Le PPCM des dénominateurs est $\text{PPCM}(44, 4797, 23006412) = 23006412$.
En réduisant au même dénominateur :
\begin{itemize}
    \item $\frac{1}{44} = \frac{522873}{23006412}$
    \item $\frac{1}{4797} = \frac{4796}{23006412}$
    \item $\frac{1}{23006412} = \frac{1}{23006412}$
\end{itemize}
La somme des numérateurs est :
$$ \frac{1}{44} + \frac{1}{4797} + \frac{1}{23006412} = \frac{522873 + 4796 + 1}{23006412} = \frac{527670}{23006412} $$
Le PGCD du numérateur et du dénominateur est $\text{PGCD}(527670, 23006412) = 105534$.
La fraction irréductible est :
$$ \frac{527670}{23006412} = \frac{527670 \div 105534}{23006412 \div 105534} = \frac{5}{218} $$
Cette fraction est bien égale à $\frac{5}{218}$.

\subsection{Démonstration pour $n = 219$}
Soit $n = 219$. Nous cherchons $x, y, z \in \mathbb{Z}^{+}$ tels que $\frac{5}{219} = \frac{1}{x} + \frac{1}{y} + \frac{1}{z}$.
Posons $x = 44$, $y = 9637$, $z = 92862132$.
Les conditions $x > 0$, $y > 0$ et $z > 0$ sont satisfaites car $x, y, z \geq 1$.
Le PPCM des dénominateurs est $\text{PPCM}(44, 9637, 92862132) = 92862132$.
En réduisant au même dénominateur :
\begin{itemize}
    \item $\frac{1}{44} = \frac{2110503}{92862132}$
    \item $\frac{1}{9637} = \frac{9636}{92862132}$
    \item $\frac{1}{92862132} = \frac{1}{92862132}$
\end{itemize}
La somme des numérateurs est :
$$ \frac{1}{44} + \frac{1}{9637} + \frac{1}{92862132} = \frac{2110503 + 9636 + 1}{92862132} = \frac{2120140}{92862132} $$
Le PGCD du numérateur et du dénominateur est $\text{PGCD}(2120140, 92862132) = 424028$.
La fraction irréductible est :
$$ \frac{2120140}{92862132} = \frac{2120140 \div 424028}{92862132 \div 424028} = \frac{5}{219} $$
Cette fraction est bien égale à $\frac{5}{219}$.

\subsection{Démonstration pour $n = 220$}
Soit $n = 220$. Nous cherchons $x, y, z \in \mathbb{Z}^{+}$ tels que $\frac{5}{220} = \frac{1}{x} + \frac{1}{y} + \frac{1}{z}$.
Posons $x = 45$, $y = 1981$, $z = 3922380$.
Les conditions $x > 0$, $y > 0$ et $z > 0$ sont satisfaites car $x, y, z \geq 1$.
Le PPCM des dénominateurs est $\text{PPCM}(45, 1981, 3922380) = 3922380$.
En réduisant au même dénominateur :
\begin{itemize}
    \item $\frac{1}{45} = \frac{87164}{3922380}$
    \item $\frac{1}{1981} = \frac{1980}{3922380}$
    \item $\frac{1}{3922380} = \frac{1}{3922380}$
\end{itemize}
La somme des numérateurs est :
$$ \frac{1}{45} + \frac{1}{1981} + \frac{1}{3922380} = \frac{87164 + 1980 + 1}{3922380} = \frac{89145}{3922380} $$
Le PGCD du numérateur et du dénominateur est $\text{PGCD}(89145, 3922380) = 89145$.
La fraction irréductible est :
$$ \frac{89145}{3922380} = \frac{89145 \div 89145}{3922380 \div 89145} = \frac{1}{44} $$
Cette fraction est bien égale à $\frac{5}{220}$.

\subsection{Démonstration pour $n = 221$}
Soit $n = 221$. Nous cherchons $x, y, z \in \mathbb{Z}^{+}$ tels que $\frac{5}{221} = \frac{1}{x} + \frac{1}{y} + \frac{1}{z}$.
Posons $x = 45$, $y = 2487$, $z = 8244405$.
Les conditions $x > 0$, $y > 0$ et $z > 0$ sont satisfaites car $x, y, z \geq 1$.
Le PPCM des dénominateurs est $\text{PPCM}(45, 2487, 8244405) = 8244405$.
En réduisant au même dénominateur :
\begin{itemize}
    \item $\frac{1}{45} = \frac{183209}{8244405}$
    \item $\frac{1}{2487} = \frac{3315}{8244405}$
    \item $\frac{1}{8244405} = \frac{1}{8244405}$
\end{itemize}
La somme des numérateurs est :
$$ \frac{1}{45} + \frac{1}{2487} + \frac{1}{8244405} = \frac{183209 + 3315 + 1}{8244405} = \frac{186525}{8244405} $$
Le PGCD du numérateur et du dénominateur est $\text{PGCD}(186525, 8244405) = 37305$.
La fraction irréductible est :
$$ \frac{186525}{8244405} = \frac{186525 \div 37305}{8244405 \div 37305} = \frac{5}{221} $$
Cette fraction est bien égale à $\frac{5}{221}$.

\subsection{Démonstration pour $n = 222$}
Soit $n = 222$. Nous cherchons $x, y, z \in \mathbb{Z}^{+}$ tels que $\frac{5}{222} = \frac{1}{x} + \frac{1}{y} + \frac{1}{z}$.
Posons $x = 45$, $y = 3331$, $z = 11092230$.
Les conditions $x > 0$, $y > 0$ et $z > 0$ sont satisfaites car $x, y, z \geq 1$.
Le PPCM des dénominateurs est $\text{PPCM}(45, 3331, 11092230) = 11092230$.
En réduisant au même dénominateur :
\begin{itemize}
    \item $\frac{1}{45} = \frac{246494}{11092230}$
    \item $\frac{1}{3331} = \frac{3330}{11092230}$
    \item $\frac{1}{11092230} = \frac{1}{11092230}$
\end{itemize}
La somme des numérateurs est :
$$ \frac{1}{45} + \frac{1}{3331} + \frac{1}{11092230} = \frac{246494 + 3330 + 1}{11092230} = \frac{249825}{11092230} $$
Le PGCD du numérateur et du dénominateur est $\text{PGCD}(249825, 11092230) = 49965$.
La fraction irréductible est :
$$ \frac{249825}{11092230} = \frac{249825 \div 49965}{11092230 \div 49965} = \frac{5}{222} $$
Cette fraction est bien égale à $\frac{5}{222}$.

\subsection{Démonstration pour $n = 223$}
Soit $n = 223$. Nous cherchons $x, y, z \in \mathbb{Z}^{+}$ tels que $\frac{5}{223} = \frac{1}{x} + \frac{1}{y} + \frac{1}{z}$.
Posons $x = 45$, $y = 5018$, $z = 50355630$.
Les conditions $x > 0$, $y > 0$ et $z > 0$ sont satisfaites car $x, y, z \geq 1$.
Le PPCM des dénominateurs est $\text{PPCM}(45, 5018, 50355630) = 50355630$.
En réduisant au même dénominateur :
\begin{itemize}
    \item $\frac{1}{45} = \frac{1119014}{50355630}$
    \item $\frac{1}{5018} = \frac{10035}{50355630}$
    \item $\frac{1}{50355630} = \frac{1}{50355630}$
\end{itemize}
La somme des numérateurs est :
$$ \frac{1}{45} + \frac{1}{5018} + \frac{1}{50355630} = \frac{1119014 + 10035 + 1}{50355630} = \frac{1129050}{50355630} $$
Le PGCD du numérateur et du dénominateur est $\text{PGCD}(1129050, 50355630) = 225810$.
La fraction irréductible est :
$$ \frac{1129050}{50355630} = \frac{1129050 \div 225810}{50355630 \div 225810} = \frac{5}{223} $$
Cette fraction est bien égale à $\frac{5}{223}$.

\subsection{Démonstration pour $n = 224$}
Soit $n = 224$. Nous cherchons $x, y, z \in \mathbb{Z}^{+}$ tels que $\frac{5}{224} = \frac{1}{x} + \frac{1}{y} + \frac{1}{z}$.
Posons $x = 45$, $y = 10081$, $z = 101616480$.
Les conditions $x > 0$, $y > 0$ et $z > 0$ sont satisfaites car $x, y, z \geq 1$.
Le PPCM des dénominateurs est $\text{PPCM}(45, 10081, 101616480) = 101616480$.
En réduisant au même dénominateur :
\begin{itemize}
    \item $\frac{1}{45} = \frac{2258144}{101616480}$
    \item $\frac{1}{10081} = \frac{10080}{101616480}$
    \item $\frac{1}{101616480} = \frac{1}{101616480}$
\end{itemize}
La somme des numérateurs est :
$$ \frac{1}{45} + \frac{1}{10081} + \frac{1}{101616480} = \frac{2258144 + 10080 + 1}{101616480} = \frac{2268225}{101616480} $$
Le PGCD du numérateur et du dénominateur est $\text{PGCD}(2268225, 101616480) = 453645$.
La fraction irréductible est :
$$ \frac{2268225}{101616480} = \frac{2268225 \div 453645}{101616480 \div 453645} = \frac{5}{224} $$
Cette fraction est bien égale à $\frac{5}{224}$.

\subsection{Démonstration pour $n = 225$}
Soit $n = 225$. Nous cherchons $x, y, z \in \mathbb{Z}^{+}$ tels que $\frac{5}{225} = \frac{1}{x} + \frac{1}{y} + \frac{1}{z}$.
Posons $x = 46$, $y = 2071$, $z = 4286970$.
Les conditions $x > 0$, $y > 0$ et $z > 0$ sont satisfaites car $x, y, z \geq 1$.
Le PPCM des dénominateurs est $\text{PPCM}(46, 2071, 4286970) = 4286970$.
En réduisant au même dénominateur :
\begin{itemize}
    \item $\frac{1}{46} = \frac{93195}{4286970}$
    \item $\frac{1}{2071} = \frac{2070}{4286970}$
    \item $\frac{1}{4286970} = \frac{1}{4286970}$
\end{itemize}
La somme des numérateurs est :
$$ \frac{1}{46} + \frac{1}{2071} + \frac{1}{4286970} = \frac{93195 + 2070 + 1}{4286970} = \frac{95266}{4286970} $$
Le PGCD du numérateur et du dénominateur est $\text{PGCD}(95266, 4286970) = 95266$.
La fraction irréductible est :
$$ \frac{95266}{4286970} = \frac{95266 \div 95266}{4286970 \div 95266} = \frac{1}{45} $$
Cette fraction est bien égale à $\frac{5}{225}$.

\subsection{Démonstration pour $n = 226$}
Soit $n = 226$. Nous cherchons $x, y, z \in \mathbb{Z}^{+}$ tels que $\frac{5}{226} = \frac{1}{x} + \frac{1}{y} + \frac{1}{z}$.
Posons $x = 46$, $y = 2600$, $z = 6757400$.
Les conditions $x > 0$, $y > 0$ et $z > 0$ sont satisfaites car $x, y, z \geq 1$.
Le PPCM des dénominateurs est $\text{PPCM}(46, 2600, 6757400) = 6757400$.
En réduisant au même dénominateur :
\begin{itemize}
    \item $\frac{1}{46} = \frac{146900}{6757400}$
    \item $\frac{1}{2600} = \frac{2599}{6757400}$
    \item $\frac{1}{6757400} = \frac{1}{6757400}$
\end{itemize}
La somme des numérateurs est :
$$ \frac{1}{46} + \frac{1}{2600} + \frac{1}{6757400} = \frac{146900 + 2599 + 1}{6757400} = \frac{149500}{6757400} $$
Le PGCD du numérateur et du dénominateur est $\text{PGCD}(149500, 6757400) = 29900$.
La fraction irréductible est :
$$ \frac{149500}{6757400} = \frac{149500 \div 29900}{6757400 \div 29900} = \frac{5}{226} $$
Cette fraction est bien égale à $\frac{5}{226}$.

\subsection{Démonstration pour $n = 227$}
Soit $n = 227$. Nous cherchons $x, y, z \in \mathbb{Z}^{+}$ tels que $\frac{5}{227} = \frac{1}{x} + \frac{1}{y} + \frac{1}{z}$.
Posons $x = 46$, $y = 3481$, $z = 36348602$.
Les conditions $x > 0$, $y > 0$ et $z > 0$ sont satisfaites car $x, y, z \geq 1$.
Le PPCM des dénominateurs est $\text{PPCM}(46, 3481, 36348602) = 36348602$.
En réduisant au même dénominateur :
\begin{itemize}
    \item $\frac{1}{46} = \frac{790187}{36348602}$
    \item $\frac{1}{3481} = \frac{10442}{36348602}$
    \item $\frac{1}{36348602} = \frac{1}{36348602}$
\end{itemize}
La somme des numérateurs est :
$$ \frac{1}{46} + \frac{1}{3481} + \frac{1}{36348602} = \frac{790187 + 10442 + 1}{36348602} = \frac{800630}{36348602} $$
Le PGCD du numérateur et du dénominateur est $\text{PGCD}(800630, 36348602) = 160126$.
La fraction irréductible est :
$$ \frac{800630}{36348602} = \frac{800630 \div 160126}{36348602 \div 160126} = \frac{5}{227} $$
Cette fraction est bien égale à $\frac{5}{227}$.

\subsection{Démonstration pour $n = 228$}
Soit $n = 228$. Nous cherchons $x, y, z \in \mathbb{Z}^{+}$ tels que $\frac{5}{228} = \frac{1}{x} + \frac{1}{y} + \frac{1}{z}$.
Posons $x = 46$, $y = 5245$, $z = 27504780$.
Les conditions $x > 0$, $y > 0$ et $z > 0$ sont satisfaites car $x, y, z \geq 1$.
Le PPCM des dénominateurs est $\text{PPCM}(46, 5245, 27504780) = 27504780$.
En réduisant au même dénominateur :
\begin{itemize}
    \item $\frac{1}{46} = \frac{597930}{27504780}$
    \item $\frac{1}{5245} = \frac{5244}{27504780}$
    \item $\frac{1}{27504780} = \frac{1}{27504780}$
\end{itemize}
La somme des numérateurs est :
$$ \frac{1}{46} + \frac{1}{5245} + \frac{1}{27504780} = \frac{597930 + 5244 + 1}{27504780} = \frac{603175}{27504780} $$
Le PGCD du numérateur et du dénominateur est $\text{PGCD}(603175, 27504780) = 120635$.
La fraction irréductible est :
$$ \frac{603175}{27504780} = \frac{603175 \div 120635}{27504780 \div 120635} = \frac{5}{228} $$
Cette fraction est bien égale à $\frac{5}{228}$.

\subsection{Démonstration pour $n = 229$}
Soit $n = 229$. Nous cherchons $x, y, z \in \mathbb{Z}^{+}$ tels que $\frac{5}{229} = \frac{1}{x} + \frac{1}{y} + \frac{1}{z}$.
Posons $x = 46$, $y = 10535$, $z = 110975690$.
Les conditions $x > 0$, $y > 0$ et $z > 0$ sont satisfaites car $x, y, z \geq 1$.
Le PPCM des dénominateurs est $\text{PPCM}(46, 10535, 110975690) = 110975690$.
En réduisant au même dénominateur :
\begin{itemize}
    \item $\frac{1}{46} = \frac{2412515}{110975690}$
    \item $\frac{1}{10535} = \frac{10534}{110975690}$
    \item $\frac{1}{110975690} = \frac{1}{110975690}$
\end{itemize}
La somme des numérateurs est :
$$ \frac{1}{46} + \frac{1}{10535} + \frac{1}{110975690} = \frac{2412515 + 10534 + 1}{110975690} = \frac{2423050}{110975690} $$
Le PGCD du numérateur et du dénominateur est $\text{PGCD}(2423050, 110975690) = 484610$.
La fraction irréductible est :
$$ \frac{2423050}{110975690} = \frac{2423050 \div 484610}{110975690 \div 484610} = \frac{5}{229} $$
Cette fraction est bien égale à $\frac{5}{229}$.

\subsection{Démonstration pour $n = 230$}
Soit $n = 230$. Nous cherchons $x, y, z \in \mathbb{Z}^{+}$ tels que $\frac{5}{230} = \frac{1}{x} + \frac{1}{y} + \frac{1}{z}$.
Posons $x = 47$, $y = 2163$, $z = 4676406$.
Les conditions $x > 0$, $y > 0$ et $z > 0$ sont satisfaites car $x, y, z \geq 1$.
Le PPCM des dénominateurs est $\text{PPCM}(47, 2163, 4676406) = 4676406$.
En réduisant au même dénominateur :
\begin{itemize}
    \item $\frac{1}{47} = \frac{99498}{4676406}$
    \item $\frac{1}{2163} = \frac{2162}{4676406}$
    \item $\frac{1}{4676406} = \frac{1}{4676406}$
\end{itemize}
La somme des numérateurs est :
$$ \frac{1}{47} + \frac{1}{2163} + \frac{1}{4676406} = \frac{99498 + 2162 + 1}{4676406} = \frac{101661}{4676406} $$
Le PGCD du numérateur et du dénominateur est $\text{PGCD}(101661, 4676406) = 101661$.
La fraction irréductible est :
$$ \frac{101661}{4676406} = \frac{101661 \div 101661}{4676406 \div 101661} = \frac{1}{46} $$
Cette fraction est bien égale à $\frac{5}{230}$.

\subsection{Démonstration pour $n = 231$}
Soit $n = 231$. Nous cherchons $x, y, z \in \mathbb{Z}^{+}$ tels que $\frac{5}{231} = \frac{1}{x} + \frac{1}{y} + \frac{1}{z}$.
Posons $x = 47$, $y = 2715$, $z = 9825585$.
Les conditions $x > 0$, $y > 0$ et $z > 0$ sont satisfaites car $x, y, z \geq 1$.
Le PPCM des dénominateurs est $\text{PPCM}(47, 2715, 9825585) = 9825585$.
En réduisant au même dénominateur :
\begin{itemize}
    \item $\frac{1}{47} = \frac{209055}{9825585}$
    \item $\frac{1}{2715} = \frac{3619}{9825585}$
    \item $\frac{1}{9825585} = \frac{1}{9825585}$
\end{itemize}
La somme des numérateurs est :
$$ \frac{1}{47} + \frac{1}{2715} + \frac{1}{9825585} = \frac{209055 + 3619 + 1}{9825585} = \frac{212675}{9825585} $$
Le PGCD du numérateur et du dénominateur est $\text{PGCD}(212675, 9825585) = 42535$.
La fraction irréductible est :
$$ \frac{212675}{9825585} = \frac{212675 \div 42535}{9825585 \div 42535} = \frac{5}{231} $$
Cette fraction est bien égale à $\frac{5}{231}$.

\subsection{Démonstration pour $n = 232$}
Soit $n = 232$. Nous cherchons $x, y, z \in \mathbb{Z}^{+}$ tels que $\frac{5}{232} = \frac{1}{x} + \frac{1}{y} + \frac{1}{z}$.
Posons $x = 47$, $y = 3635$, $z = 39636040$.
Les conditions $x > 0$, $y > 0$ et $z > 0$ sont satisfaites car $x, y, z \geq 1$.
Le PPCM des dénominateurs est $\text{PPCM}(47, 3635, 39636040) = 39636040$.
En réduisant au même dénominateur :
\begin{itemize}
    \item $\frac{1}{47} = \frac{843320}{39636040}$
    \item $\frac{1}{3635} = \frac{10904}{39636040}$
    \item $\frac{1}{39636040} = \frac{1}{39636040}$
\end{itemize}
La somme des numérateurs est :
$$ \frac{1}{47} + \frac{1}{3635} + \frac{1}{39636040} = \frac{843320 + 10904 + 1}{39636040} = \frac{854225}{39636040} $$
Le PGCD du numérateur et du dénominateur est $\text{PGCD}(854225, 39636040) = 170845$.
La fraction irréductible est :
$$ \frac{854225}{39636040} = \frac{854225 \div 170845}{39636040 \div 170845} = \frac{5}{232} $$
Cette fraction est bien égale à $\frac{5}{232}$.

\subsection{Démonstration pour $n = 233$}
Soit $n = 233$. Nous cherchons $x, y, z \in \mathbb{Z}^{+}$ tels que $\frac{5}{233} = \frac{1}{x} + \frac{1}{y} + \frac{1}{z}$.
Posons $x = 47$, $y = 5476$, $z = 59967676$.
Les conditions $x > 0$, $y > 0$ et $z > 0$ sont satisfaites car $x, y, z \geq 1$.
Le PPCM des dénominateurs est $\text{PPCM}(47, 5476, 59967676) = 59967676$.
En réduisant au même dénominateur :
\begin{itemize}
    \item $\frac{1}{47} = \frac{1275908}{59967676}$
    \item $\frac{1}{5476} = \frac{10951}{59967676}$
    \item $\frac{1}{59967676} = \frac{1}{59967676}$
\end{itemize}
La somme des numérateurs est :
$$ \frac{1}{47} + \frac{1}{5476} + \frac{1}{59967676} = \frac{1275908 + 10951 + 1}{59967676} = \frac{1286860}{59967676} $$
Le PGCD du numérateur et du dénominateur est $\text{PGCD}(1286860, 59967676) = 257372$.
La fraction irréductible est :
$$ \frac{1286860}{59967676} = \frac{1286860 \div 257372}{59967676 \div 257372} = \frac{5}{233} $$
Cette fraction est bien égale à $\frac{5}{233}$.

\subsection{Démonstration pour $n = 234$}
Soit $n = 234$. Nous cherchons $x, y, z \in \mathbb{Z}^{+}$ tels que $\frac{5}{234} = \frac{1}{x} + \frac{1}{y} + \frac{1}{z}$.
Posons $x = 47$, $y = 10999$, $z = 120967002$.
Les conditions $x > 0$, $y > 0$ et $z > 0$ sont satisfaites car $x, y, z \geq 1$.
Le PPCM des dénominateurs est $\text{PPCM}(47, 10999, 120967002) = 120967002$.
En réduisant au même dénominateur :
\begin{itemize}
    \item $\frac{1}{47} = \frac{2573766}{120967002}$
    \item $\frac{1}{10999} = \frac{10998}{120967002}$
    \item $\frac{1}{120967002} = \frac{1}{120967002}$
\end{itemize}
La somme des numérateurs est :
$$ \frac{1}{47} + \frac{1}{10999} + \frac{1}{120967002} = \frac{2573766 + 10998 + 1}{120967002} = \frac{2584765}{120967002} $$
Le PGCD du numérateur et du dénominateur est $\text{PGCD}(2584765, 120967002) = 516953$.
La fraction irréductible est :
$$ \frac{2584765}{120967002} = \frac{2584765 \div 516953}{120967002 \div 516953} = \frac{5}{234} $$
Cette fraction est bien égale à $\frac{5}{234}$.

\subsection{Démonstration pour $n = 235$}
Soit $n = 235$. Nous cherchons $x, y, z \in \mathbb{Z}^{+}$ tels que $\frac{5}{235} = \frac{1}{x} + \frac{1}{y} + \frac{1}{z}$.
Posons $x = 48$, $y = 2257$, $z = 5091792$.
Les conditions $x > 0$, $y > 0$ et $z > 0$ sont satisfaites car $x, y, z \geq 1$.
Le PPCM des dénominateurs est $\text{PPCM}(48, 2257, 5091792) = 5091792$.
En réduisant au même dénominateur :
\begin{itemize}
    \item $\frac{1}{48} = \frac{106079}{5091792}$
    \item $\frac{1}{2257} = \frac{2256}{5091792}$
    \item $\frac{1}{5091792} = \frac{1}{5091792}$
\end{itemize}
La somme des numérateurs est :
$$ \frac{1}{48} + \frac{1}{2257} + \frac{1}{5091792} = \frac{106079 + 2256 + 1}{5091792} = \frac{108336}{5091792} $$
Le PGCD du numérateur et du dénominateur est $\text{PGCD}(108336, 5091792) = 108336$.
La fraction irréductible est :
$$ \frac{108336}{5091792} = \frac{108336 \div 108336}{5091792 \div 108336} = \frac{1}{47} $$
Cette fraction est bien égale à $\frac{5}{235}$.

\subsection{Démonstration pour $n = 236$}
Soit $n = 236$. Nous cherchons $x, y, z \in \mathbb{Z}^{+}$ tels que $\frac{5}{236} = \frac{1}{x} + \frac{1}{y} + \frac{1}{z}$.
Posons $x = 48$, $y = 2833$, $z = 8023056$.
Les conditions $x > 0$, $y > 0$ et $z > 0$ sont satisfaites car $x, y, z \geq 1$.
Le PPCM des dénominateurs est $\text{PPCM}(48, 2833, 8023056) = 8023056$.
En réduisant au même dénominateur :
\begin{itemize}
    \item $\frac{1}{48} = \frac{167147}{8023056}$
    \item $\frac{1}{2833} = \frac{2832}{8023056}$
    \item $\frac{1}{8023056} = \frac{1}{8023056}$
\end{itemize}
La somme des numérateurs est :
$$ \frac{1}{48} + \frac{1}{2833} + \frac{1}{8023056} = \frac{167147 + 2832 + 1}{8023056} = \frac{169980}{8023056} $$
Le PGCD du numérateur et du dénominateur est $\text{PGCD}(169980, 8023056) = 33996$.
La fraction irréductible est :
$$ \frac{169980}{8023056} = \frac{169980 \div 33996}{8023056 \div 33996} = \frac{5}{236} $$
Cette fraction est bien égale à $\frac{5}{236}$.

\subsection{Démonstration pour $n = 237$}
Soit $n = 237$. Nous cherchons $x, y, z \in \mathbb{Z}^{+}$ tels que $\frac{5}{237} = \frac{1}{x} + \frac{1}{y} + \frac{1}{z}$.
Posons $x = 48$, $y = 3793$, $z = 14383056$.
Les conditions $x > 0$, $y > 0$ et $z > 0$ sont satisfaites car $x, y, z \geq 1$.
Le PPCM des dénominateurs est $\text{PPCM}(48, 3793, 14383056) = 14383056$.
En réduisant au même dénominateur :
\begin{itemize}
    \item $\frac{1}{48} = \frac{299647}{14383056}$
    \item $\frac{1}{3793} = \frac{3792}{14383056}$
    \item $\frac{1}{14383056} = \frac{1}{14383056}$
\end{itemize}
La somme des numérateurs est :
$$ \frac{1}{48} + \frac{1}{3793} + \frac{1}{14383056} = \frac{299647 + 3792 + 1}{14383056} = \frac{303440}{14383056} $$
Le PGCD du numérateur et du dénominateur est $\text{PGCD}(303440, 14383056) = 60688$.
La fraction irréductible est :
$$ \frac{303440}{14383056} = \frac{303440 \div 60688}{14383056 \div 60688} = \frac{5}{237} $$
Cette fraction est bien égale à $\frac{5}{237}$.

\subsection{Démonstration pour $n = 238$}
Soit $n = 238$. Nous cherchons $x, y, z \in \mathbb{Z}^{+}$ tels que $\frac{5}{238} = \frac{1}{x} + \frac{1}{y} + \frac{1}{z}$.
Posons $x = 48$, $y = 5713$, $z = 32632656$.
Les conditions $x > 0$, $y > 0$ et $z > 0$ sont satisfaites car $x, y, z \geq 1$.
Le PPCM des dénominateurs est $\text{PPCM}(48, 5713, 32632656) = 32632656$.
En réduisant au même dénominateur :
\begin{itemize}
    \item $\frac{1}{48} = \frac{679847}{32632656}$
    \item $\frac{1}{5713} = \frac{5712}{32632656}$
    \item $\frac{1}{32632656} = \frac{1}{32632656}$
\end{itemize}
La somme des numérateurs est :
$$ \frac{1}{48} + \frac{1}{5713} + \frac{1}{32632656} = \frac{679847 + 5712 + 1}{32632656} = \frac{685560}{32632656} $$
Le PGCD du numérateur et du dénominateur est $\text{PGCD}(685560, 32632656) = 137112$.
La fraction irréductible est :
$$ \frac{685560}{32632656} = \frac{685560 \div 137112}{32632656 \div 137112} = \frac{5}{238} $$
Cette fraction est bien égale à $\frac{5}{238}$.

\subsection{Démonstration pour $n = 239$}
Soit $n = 239$. Nous cherchons $x, y, z \in \mathbb{Z}^{+}$ tels que $\frac{5}{239} = \frac{1}{x} + \frac{1}{y} + \frac{1}{z}$.
Posons $x = 48$, $y = 11473$, $z = 131618256$.
Les conditions $x > 0$, $y > 0$ et $z > 0$ sont satisfaites car $x, y, z \geq 1$.
Le PPCM des dénominateurs est $\text{PPCM}(48, 11473, 131618256) = 131618256$.
En réduisant au même dénominateur :
\begin{itemize}
    \item $\frac{1}{48} = \frac{2742047}{131618256}$
    \item $\frac{1}{11473} = \frac{11472}{131618256}$
    \item $\frac{1}{131618256} = \frac{1}{131618256}$
\end{itemize}
La somme des numérateurs est :
$$ \frac{1}{48} + \frac{1}{11473} + \frac{1}{131618256} = \frac{2742047 + 11472 + 1}{131618256} = \frac{2753520}{131618256} $$
Le PGCD du numérateur et du dénominateur est $\text{PGCD}(2753520, 131618256) = 550704$.
La fraction irréductible est :
$$ \frac{2753520}{131618256} = \frac{2753520 \div 550704}{131618256 \div 550704} = \frac{5}{239} $$
Cette fraction est bien égale à $\frac{5}{239}$.

\subsection{Démonstration pour $n = 240$}
Soit $n = 240$. Nous cherchons $x, y, z \in \mathbb{Z}^{+}$ tels que $\frac{5}{240} = \frac{1}{x} + \frac{1}{y} + \frac{1}{z}$.
Posons $x = 49$, $y = 2353$, $z = 5534256$.
Les conditions $x > 0$, $y > 0$ et $z > 0$ sont satisfaites car $x, y, z \geq 1$.
Le PPCM des dénominateurs est $\text{PPCM}(49, 2353, 5534256) = 5534256$.
En réduisant au même dénominateur :
\begin{itemize}
    \item $\frac{1}{49} = \frac{112944}{5534256}$
    \item $\frac{1}{2353} = \frac{2352}{5534256}$
    \item $\frac{1}{5534256} = \frac{1}{5534256}$
\end{itemize}
La somme des numérateurs est :
$$ \frac{1}{49} + \frac{1}{2353} + \frac{1}{5534256} = \frac{112944 + 2352 + 1}{5534256} = \frac{115297}{5534256} $$
Le PGCD du numérateur et du dénominateur est $\text{PGCD}(115297, 5534256) = 115297$.
La fraction irréductible est :
$$ \frac{115297}{5534256} = \frac{115297 \div 115297}{5534256 \div 115297} = \frac{1}{48} $$
Cette fraction est bien égale à $\frac{5}{240}$.

\subsection{Démonstration pour $n = 241$}
Soit $n = 241$. Nous cherchons $x, y, z \in \mathbb{Z}^{+}$ tels que $\frac{5}{241} = \frac{1}{x} + \frac{1}{y} + \frac{1}{z}$.
Posons $x = 49$, $y = 2954$, $z = 4983398$.
Les conditions $x > 0$, $y > 0$ et $z > 0$ sont satisfaites car $x, y, z \geq 1$.
Le PPCM des dénominateurs est $\text{PPCM}(49, 2954, 4983398) = 4983398$.
En réduisant au même dénominateur :
\begin{itemize}
    \item $\frac{1}{49} = \frac{101702}{4983398}$
    \item $\frac{1}{2954} = \frac{1687}{4983398}$
    \item $\frac{1}{4983398} = \frac{1}{4983398}$
\end{itemize}
La somme des numérateurs est :
$$ \frac{1}{49} + \frac{1}{2954} + \frac{1}{4983398} = \frac{101702 + 1687 + 1}{4983398} = \frac{103390}{4983398} $$
Le PGCD du numérateur et du dénominateur est $\text{PGCD}(103390, 4983398) = 20678$.
La fraction irréductible est :
$$ \frac{103390}{4983398} = \frac{103390 \div 20678}{4983398 \div 20678} = \frac{5}{241} $$
Cette fraction est bien égale à $\frac{5}{241}$.

\subsection{Démonstration pour $n = 242$}
Soit $n = 242$. Nous cherchons $x, y, z \in \mathbb{Z}^{+}$ tels que $\frac{5}{242} = \frac{1}{x} + \frac{1}{y} + \frac{1}{z}$.
Posons $x = 49$, $y = 3953$, $z = 46874674$.
Les conditions $x > 0$, $y > 0$ et $z > 0$ sont satisfaites car $x, y, z \geq 1$.
Le PPCM des dénominateurs est $\text{PPCM}(49, 3953, 46874674) = 46874674$.
En réduisant au même dénominateur :
\begin{itemize}
    \item $\frac{1}{49} = \frac{956626}{46874674}$
    \item $\frac{1}{3953} = \frac{11858}{46874674}$
    \item $\frac{1}{46874674} = \frac{1}{46874674}$
\end{itemize}
La somme des numérateurs est :
$$ \frac{1}{49} + \frac{1}{3953} + \frac{1}{46874674} = \frac{956626 + 11858 + 1}{46874674} = \frac{968485}{46874674} $$
Le PGCD du numérateur et du dénominateur est $\text{PGCD}(968485, 46874674) = 193697$.
La fraction irréductible est :
$$ \frac{968485}{46874674} = \frac{968485 \div 193697}{46874674 \div 193697} = \frac{5}{242} $$
Cette fraction est bien égale à $\frac{5}{242}$.

\subsection{Démonstration pour $n = 243$}
Soit $n = 243$. Nous cherchons $x, y, z \in \mathbb{Z}^{+}$ tels que $\frac{5}{243} = \frac{1}{x} + \frac{1}{y} + \frac{1}{z}$.
Posons $x = 49$, $y = 5954$, $z = 70894278$.
Les conditions $x > 0$, $y > 0$ et $z > 0$ sont satisfaites car $x, y, z \geq 1$.
Le PPCM des dénominateurs est $\text{PPCM}(49, 5954, 70894278) = 70894278$.
En réduisant au même dénominateur :
\begin{itemize}
    \item $\frac{1}{49} = \frac{1446822}{70894278}$
    \item $\frac{1}{5954} = \frac{11907}{70894278}$
    \item $\frac{1}{70894278} = \frac{1}{70894278}$
\end{itemize}
La somme des numérateurs est :
$$ \frac{1}{49} + \frac{1}{5954} + \frac{1}{70894278} = \frac{1446822 + 11907 + 1}{70894278} = \frac{1458730}{70894278} $$
Le PGCD du numérateur et du dénominateur est $\text{PGCD}(1458730, 70894278) = 291746$.
La fraction irréductible est :
$$ \frac{1458730}{70894278} = \frac{1458730 \div 291746}{70894278 \div 291746} = \frac{5}{243} $$
Cette fraction est bien égale à $\frac{5}{243}$.

\subsection{Démonstration pour $n = 244$}
Soit $n = 244$. Nous cherchons $x, y, z \in \mathbb{Z}^{+}$ tels que $\frac{5}{244} = \frac{1}{x} + \frac{1}{y} + \frac{1}{z}$.
Posons $x = 49$, $y = 11957$, $z = 142957892$.
Les conditions $x > 0$, $y > 0$ et $z > 0$ sont satisfaites car $x, y, z \geq 1$.
Le PPCM des dénominateurs est $\text{PPCM}(49, 11957, 142957892) = 142957892$.
En réduisant au même dénominateur :
\begin{itemize}
    \item $\frac{1}{49} = \frac{2917508}{142957892}$
    \item $\frac{1}{11957} = \frac{11956}{142957892}$
    \item $\frac{1}{142957892} = \frac{1}{142957892}$
\end{itemize}
La somme des numérateurs est :
$$ \frac{1}{49} + \frac{1}{11957} + \frac{1}{142957892} = \frac{2917508 + 11956 + 1}{142957892} = \frac{2929465}{142957892} $$
Le PGCD du numérateur et du dénominateur est $\text{PGCD}(2929465, 142957892) = 585893$.
La fraction irréductible est :
$$ \frac{2929465}{142957892} = \frac{2929465 \div 585893}{142957892 \div 585893} = \frac{5}{244} $$
Cette fraction est bien égale à $\frac{5}{244}$.

\subsection{Démonstration pour $n = 245$}
Soit $n = 245$. Nous cherchons $x, y, z \in \mathbb{Z}^{+}$ tels que $\frac{5}{245} = \frac{1}{x} + \frac{1}{y} + \frac{1}{z}$.
Posons $x = 50$, $y = 2451$, $z = 6004950$.
Les conditions $x > 0$, $y > 0$ et $z > 0$ sont satisfaites car $x, y, z \geq 1$.
Le PPCM des dénominateurs est $\text{PPCM}(50, 2451, 6004950) = 6004950$.
En réduisant au même dénominateur :
\begin{itemize}
    \item $\frac{1}{50} = \frac{120099}{6004950}$
    \item $\frac{1}{2451} = \frac{2450}{6004950}$
    \item $\frac{1}{6004950} = \frac{1}{6004950}$
\end{itemize}
La somme des numérateurs est :
$$ \frac{1}{50} + \frac{1}{2451} + \frac{1}{6004950} = \frac{120099 + 2450 + 1}{6004950} = \frac{122550}{6004950} $$
Le PGCD du numérateur et du dénominateur est $\text{PGCD}(122550, 6004950) = 122550$.
La fraction irréductible est :
$$ \frac{122550}{6004950} = \frac{122550 \div 122550}{6004950 \div 122550} = \frac{1}{49} $$
Cette fraction est bien égale à $\frac{5}{245}$.

\subsection{Démonstration pour $n = 246$}
Soit $n = 246$. Nous cherchons $x, y, z \in \mathbb{Z}^{+}$ tels que $\frac{5}{246} = \frac{1}{x} + \frac{1}{y} + \frac{1}{z}$.
Posons $x = 50$, $y = 3076$, $z = 9458700$.
Les conditions $x > 0$, $y > 0$ et $z > 0$ sont satisfaites car $x, y, z \geq 1$.
Le PPCM des dénominateurs est $\text{PPCM}(50, 3076, 9458700) = 9458700$.
En réduisant au même dénominateur :
\begin{itemize}
    \item $\frac{1}{50} = \frac{189174}{9458700}$
    \item $\frac{1}{3076} = \frac{3075}{9458700}$
    \item $\frac{1}{9458700} = \frac{1}{9458700}$
\end{itemize}
La somme des numérateurs est :
$$ \frac{1}{50} + \frac{1}{3076} + \frac{1}{9458700} = \frac{189174 + 3075 + 1}{9458700} = \frac{192250}{9458700} $$
Le PGCD du numérateur et du dénominateur est $\text{PGCD}(192250, 9458700) = 38450$.
La fraction irréductible est :
$$ \frac{192250}{9458700} = \frac{192250 \div 38450}{9458700 \div 38450} = \frac{5}{246} $$
Cette fraction est bien égale à $\frac{5}{246}$.

\subsection{Démonstration pour $n = 247$}
Soit $n = 247$. Nous cherchons $x, y, z \in \mathbb{Z}^{+}$ tels que $\frac{5}{247} = \frac{1}{x} + \frac{1}{y} + \frac{1}{z}$.
Posons $x = 50$, $y = 4117$, $z = 50844950$.
Les conditions $x > 0$, $y > 0$ et $z > 0$ sont satisfaites car $x, y, z \geq 1$.
Le PPCM des dénominateurs est $\text{PPCM}(50, 4117, 50844950) = 50844950$.
En réduisant au même dénominateur :
\begin{itemize}
    \item $\frac{1}{50} = \frac{1016899}{50844950}$
    \item $\frac{1}{4117} = \frac{12350}{50844950}$
    \item $\frac{1}{50844950} = \frac{1}{50844950}$
\end{itemize}
La somme des numérateurs est :
$$ \frac{1}{50} + \frac{1}{4117} + \frac{1}{50844950} = \frac{1016899 + 12350 + 1}{50844950} = \frac{1029250}{50844950} $$
Le PGCD du numérateur et du dénominateur est $\text{PGCD}(1029250, 50844950) = 205850$.
La fraction irréductible est :
$$ \frac{1029250}{50844950} = \frac{1029250 \div 205850}{50844950 \div 205850} = \frac{5}{247} $$
Cette fraction est bien égale à $\frac{5}{247}$.

\subsection{Démonstration pour $n = 248$}
Soit $n = 248$. Nous cherchons $x, y, z \in \mathbb{Z}^{+}$ tels que $\frac{5}{248} = \frac{1}{x} + \frac{1}{y} + \frac{1}{z}$.
Posons $x = 50$, $y = 6201$, $z = 38446200$.
Les conditions $x > 0$, $y > 0$ et $z > 0$ sont satisfaites car $x, y, z \geq 1$.
Le PPCM des dénominateurs est $\text{PPCM}(50, 6201, 38446200) = 38446200$.
En réduisant au même dénominateur :
\begin{itemize}
    \item $\frac{1}{50} = \frac{768924}{38446200}$
    \item $\frac{1}{6201} = \frac{6200}{38446200}$
    \item $\frac{1}{38446200} = \frac{1}{38446200}$
\end{itemize}
La somme des numérateurs est :
$$ \frac{1}{50} + \frac{1}{6201} + \frac{1}{38446200} = \frac{768924 + 6200 + 1}{38446200} = \frac{775125}{38446200} $$
Le PGCD du numérateur et du dénominateur est $\text{PGCD}(775125, 38446200) = 155025$.
La fraction irréductible est :
$$ \frac{775125}{38446200} = \frac{775125 \div 155025}{38446200 \div 155025} = \frac{5}{248} $$
Cette fraction est bien égale à $\frac{5}{248}$.

\subsection{Démonstration pour $n = 249$}
Soit $n = 249$. Nous cherchons $x, y, z \in \mathbb{Z}^{+}$ tels que $\frac{5}{249} = \frac{1}{x} + \frac{1}{y} + \frac{1}{z}$.
Posons $x = 50$, $y = 12451$, $z = 155014950$.
Les conditions $x > 0$, $y > 0$ et $z > 0$ sont satisfaites car $x, y, z \geq 1$.
Le PPCM des dénominateurs est $\text{PPCM}(50, 12451, 155014950) = 155014950$.
En réduisant au même dénominateur :
\begin{itemize}
    \item $\frac{1}{50} = \frac{3100299}{155014950}$
    \item $\frac{1}{12451} = \frac{12450}{155014950}$
    \item $\frac{1}{155014950} = \frac{1}{155014950}$
\end{itemize}
La somme des numérateurs est :
$$ \frac{1}{50} + \frac{1}{12451} + \frac{1}{155014950} = \frac{3100299 + 12450 + 1}{155014950} = \frac{3112750}{155014950} $$
Le PGCD du numérateur et du dénominateur est $\text{PGCD}(3112750, 155014950) = 622550$.
La fraction irréductible est :
$$ \frac{3112750}{155014950} = \frac{3112750 \div 622550}{155014950 \div 622550} = \frac{5}{249} $$
Cette fraction est bien égale à $\frac{5}{249}$.

\subsection{Démonstration pour $n = 250$}
Soit $n = 250$. Nous cherchons $x, y, z \in \mathbb{Z}^{+}$ tels que $\frac{5}{250} = \frac{1}{x} + \frac{1}{y} + \frac{1}{z}$.
Posons $x = 51$, $y = 2551$, $z = 6505050$.
Les conditions $x > 0$, $y > 0$ et $z > 0$ sont satisfaites car $x, y, z \geq 1$.
Le PPCM des dénominateurs est $\text{PPCM}(51, 2551, 6505050) = 6505050$.
En réduisant au même dénominateur :
\begin{itemize}
    \item $\frac{1}{51} = \frac{127550}{6505050}$
    \item $\frac{1}{2551} = \frac{2550}{6505050}$
    \item $\frac{1}{6505050} = \frac{1}{6505050}$
\end{itemize}
La somme des numérateurs est :
$$ \frac{1}{51} + \frac{1}{2551} + \frac{1}{6505050} = \frac{127550 + 2550 + 1}{6505050} = \frac{130101}{6505050} $$
Le PGCD du numérateur et du dénominateur est $\text{PGCD}(130101, 6505050) = 130101$.
La fraction irréductible est :
$$ \frac{130101}{6505050} = \frac{130101 \div 130101}{6505050 \div 130101} = \frac{1}{50} $$
Cette fraction est bien égale à $\frac{5}{250}$.

\subsection{Démonstration pour $n = 251$}
Soit $n = 251$. Nous cherchons $x, y, z \in \mathbb{Z}^{+}$ tels que $\frac{5}{251} = \frac{1}{x} + \frac{1}{y} + \frac{1}{z}$.
Posons $x = 51$, $y = 3201$, $z = 13658667$.
Les conditions $x > 0$, $y > 0$ et $z > 0$ sont satisfaites car $x, y, z \geq 1$.
Le PPCM des dénominateurs est $\text{PPCM}(51, 3201, 13658667) = 13658667$.
En réduisant au même dénominateur :
\begin{itemize}
    \item $\frac{1}{51} = \frac{267817}{13658667}$
    \item $\frac{1}{3201} = \frac{4267}{13658667}$
    \item $\frac{1}{13658667} = \frac{1}{13658667}$
\end{itemize}
La somme des numérateurs est :
$$ \frac{1}{51} + \frac{1}{3201} + \frac{1}{13658667} = \frac{267817 + 4267 + 1}{13658667} = \frac{272085}{13658667} $$
Le PGCD du numérateur et du dénominateur est $\text{PGCD}(272085, 13658667) = 54417$.
La fraction irréductible est :
$$ \frac{272085}{13658667} = \frac{272085 \div 54417}{13658667 \div 54417} = \frac{5}{251} $$
Cette fraction est bien égale à $\frac{5}{251}$.

\subsection{Démonstration pour $n = 252$}
Soit $n = 252$. Nous cherchons $x, y, z \in \mathbb{Z}^{+}$ tels que $\frac{5}{252} = \frac{1}{x} + \frac{1}{y} + \frac{1}{z}$.
Posons $x = 51$, $y = 4285$, $z = 18356940$.
Les conditions $x > 0$, $y > 0$ et $z > 0$ sont satisfaites car $x, y, z \geq 1$.
Le PPCM des dénominateurs est $\text{PPCM}(51, 4285, 18356940) = 18356940$.
En réduisant au même dénominateur :
\begin{itemize}
    \item $\frac{1}{51} = \frac{359940}{18356940}$
    \item $\frac{1}{4285} = \frac{4284}{18356940}$
    \item $\frac{1}{18356940} = \frac{1}{18356940}$
\end{itemize}
La somme des numérateurs est :
$$ \frac{1}{51} + \frac{1}{4285} + \frac{1}{18356940} = \frac{359940 + 4284 + 1}{18356940} = \frac{364225}{18356940} $$
Le PGCD du numérateur et du dénominateur est $\text{PGCD}(364225, 18356940) = 72845$.
La fraction irréductible est :
$$ \frac{364225}{18356940} = \frac{364225 \div 72845}{18356940 \div 72845} = \frac{5}{252} $$
Cette fraction est bien égale à $\frac{5}{252}$.

\subsection{Démonstration pour $n = 253$}
Soit $n = 253$. Nous cherchons $x, y, z \in \mathbb{Z}^{+}$ tels que $\frac{5}{253} = \frac{1}{x} + \frac{1}{y} + \frac{1}{z}$.
Posons $x = 51$, $y = 6452$, $z = 83250156$.
Les conditions $x > 0$, $y > 0$ et $z > 0$ sont satisfaites car $x, y, z \geq 1$.
Le PPCM des dénominateurs est $\text{PPCM}(51, 6452, 83250156) = 83250156$.
En réduisant au même dénominateur :
\begin{itemize}
    \item $\frac{1}{51} = \frac{1632356}{83250156}$
    \item $\frac{1}{6452} = \frac{12903}{83250156}$
    \item $\frac{1}{83250156} = \frac{1}{83250156}$
\end{itemize}
La somme des numérateurs est :
$$ \frac{1}{51} + \frac{1}{6452} + \frac{1}{83250156} = \frac{1632356 + 12903 + 1}{83250156} = \frac{1645260}{83250156} $$
Le PGCD du numérateur et du dénominateur est $\text{PGCD}(1645260, 83250156) = 329052$.
La fraction irréductible est :
$$ \frac{1645260}{83250156} = \frac{1645260 \div 329052}{83250156 \div 329052} = \frac{5}{253} $$
Cette fraction est bien égale à $\frac{5}{253}$.

\subsection{Démonstration pour $n = 254$}
Soit $n = 254$. Nous cherchons $x, y, z \in \mathbb{Z}^{+}$ tels que $\frac{5}{254} = \frac{1}{x} + \frac{1}{y} + \frac{1}{z}$.
Posons $x = 51$, $y = 12955$, $z = 167819070$.
Les conditions $x > 0$, $y > 0$ et $z > 0$ sont satisfaites car $x, y, z \geq 1$.
Le PPCM des dénominateurs est $\text{PPCM}(51, 12955, 167819070) = 167819070$.
En réduisant au même dénominateur :
\begin{itemize}
    \item $\frac{1}{51} = \frac{3290570}{167819070}$
    \item $\frac{1}{12955} = \frac{12954}{167819070}$
    \item $\frac{1}{167819070} = \frac{1}{167819070}$
\end{itemize}
La somme des numérateurs est :
$$ \frac{1}{51} + \frac{1}{12955} + \frac{1}{167819070} = \frac{3290570 + 12954 + 1}{167819070} = \frac{3303525}{167819070} $$
Le PGCD du numérateur et du dénominateur est $\text{PGCD}(3303525, 167819070) = 660705$.
La fraction irréductible est :
$$ \frac{3303525}{167819070} = \frac{3303525 \div 660705}{167819070 \div 660705} = \frac{5}{254} $$
Cette fraction est bien égale à $\frac{5}{254}$.

\subsection{Démonstration pour $n = 255$}
Soit $n = 255$. Nous cherchons $x, y, z \in \mathbb{Z}^{+}$ tels que $\frac{5}{255} = \frac{1}{x} + \frac{1}{y} + \frac{1}{z}$.
Posons $x = 52$, $y = 2653$, $z = 7035756$.
Les conditions $x > 0$, $y > 0$ et $z > 0$ sont satisfaites car $x, y, z \geq 1$.
Le PPCM des dénominateurs est $\text{PPCM}(52, 2653, 7035756) = 7035756$.
En réduisant au même dénominateur :
\begin{itemize}
    \item $\frac{1}{52} = \frac{135303}{7035756}$
    \item $\frac{1}{2653} = \frac{2652}{7035756}$
    \item $\frac{1}{7035756} = \frac{1}{7035756}$
\end{itemize}
La somme des numérateurs est :
$$ \frac{1}{52} + \frac{1}{2653} + \frac{1}{7035756} = \frac{135303 + 2652 + 1}{7035756} = \frac{137956}{7035756} $$
Le PGCD du numérateur et du dénominateur est $\text{PGCD}(137956, 7035756) = 137956$.
La fraction irréductible est :
$$ \frac{137956}{7035756} = \frac{137956 \div 137956}{7035756 \div 137956} = \frac{1}{51} $$
Cette fraction est bien égale à $\frac{5}{255}$.

\subsection{Démonstration pour $n = 256$}
Soit $n = 256$. Nous cherchons $x, y, z \in \mathbb{Z}^{+}$ tels que $\frac{5}{256} = \frac{1}{x} + \frac{1}{y} + \frac{1}{z}$.
Posons $x = 52$, $y = 3329$, $z = 11078912$.
Les conditions $x > 0$, $y > 0$ et $z > 0$ sont satisfaites car $x, y, z \geq 1$.
Le PPCM des dénominateurs est $\text{PPCM}(52, 3329, 11078912) = 11078912$.
En réduisant au même dénominateur :
\begin{itemize}
    \item $\frac{1}{52} = \frac{213056}{11078912}$
    \item $\frac{1}{3329} = \frac{3328}{11078912}$
    \item $\frac{1}{11078912} = \frac{1}{11078912}$
\end{itemize}
La somme des numérateurs est :
$$ \frac{1}{52} + \frac{1}{3329} + \frac{1}{11078912} = \frac{213056 + 3328 + 1}{11078912} = \frac{216385}{11078912} $$
Le PGCD du numérateur et du dénominateur est $\text{PGCD}(216385, 11078912) = 43277$.
La fraction irréductible est :
$$ \frac{216385}{11078912} = \frac{216385 \div 43277}{11078912 \div 43277} = \frac{5}{256} $$
Cette fraction est bien égale à $\frac{5}{256}$.

\subsection{Démonstration pour $n = 257$}
Soit $n = 257$. Nous cherchons $x, y, z \in \mathbb{Z}^{+}$ tels que $\frac{5}{257} = \frac{1}{x} + \frac{1}{y} + \frac{1}{z}$.
Posons $x = 52$, $y = 4455$, $z = 59536620$.
Les conditions $x > 0$, $y > 0$ et $z > 0$ sont satisfaites car $x, y, z \geq 1$.
Le PPCM des dénominateurs est $\text{PPCM}(52, 4455, 59536620) = 59536620$.
En réduisant au même dénominateur :
\begin{itemize}
    \item $\frac{1}{52} = \frac{1144935}{59536620}$
    \item $\frac{1}{4455} = \frac{13364}{59536620}$
    \item $\frac{1}{59536620} = \frac{1}{59536620}$
\end{itemize}
La somme des numérateurs est :
$$ \frac{1}{52} + \frac{1}{4455} + \frac{1}{59536620} = \frac{1144935 + 13364 + 1}{59536620} = \frac{1158300}{59536620} $$
Le PGCD du numérateur et du dénominateur est $\text{PGCD}(1158300, 59536620) = 231660$.
La fraction irréductible est :
$$ \frac{1158300}{59536620} = \frac{1158300 \div 231660}{59536620 \div 231660} = \frac{5}{257} $$
Cette fraction est bien égale à $\frac{5}{257}$.

\subsection{Démonstration pour $n = 258$}
Soit $n = 258$. Nous cherchons $x, y, z \in \mathbb{Z}^{+}$ tels que $\frac{5}{258} = \frac{1}{x} + \frac{1}{y} + \frac{1}{z}$.
Posons $x = 52$, $y = 6709$, $z = 45003972$.
Les conditions $x > 0$, $y > 0$ et $z > 0$ sont satisfaites car $x, y, z \geq 1$.
Le PPCM des dénominateurs est $\text{PPCM}(52, 6709, 45003972) = 45003972$.
En réduisant au même dénominateur :
\begin{itemize}
    \item $\frac{1}{52} = \frac{865461}{45003972}$
    \item $\frac{1}{6709} = \frac{6708}{45003972}$
    \item $\frac{1}{45003972} = \frac{1}{45003972}$
\end{itemize}
La somme des numérateurs est :
$$ \frac{1}{52} + \frac{1}{6709} + \frac{1}{45003972} = \frac{865461 + 6708 + 1}{45003972} = \frac{872170}{45003972} $$
Le PGCD du numérateur et du dénominateur est $\text{PGCD}(872170, 45003972) = 174434$.
La fraction irréductible est :
$$ \frac{872170}{45003972} = \frac{872170 \div 174434}{45003972 \div 174434} = \frac{5}{258} $$
Cette fraction est bien égale à $\frac{5}{258}$.

\subsection{Démonstration pour $n = 259$}
Soit $n = 259$. Nous cherchons $x, y, z \in \mathbb{Z}^{+}$ tels que $\frac{5}{259} = \frac{1}{x} + \frac{1}{y} + \frac{1}{z}$.
Posons $x = 52$, $y = 13469$, $z = 181400492$.
Les conditions $x > 0$, $y > 0$ et $z > 0$ sont satisfaites car $x, y, z \geq 1$.
Le PPCM des dénominateurs est $\text{PPCM}(52, 13469, 181400492) = 181400492$.
En réduisant au même dénominateur :
\begin{itemize}
    \item $\frac{1}{52} = \frac{3488471}{181400492}$
    \item $\frac{1}{13469} = \frac{13468}{181400492}$
    \item $\frac{1}{181400492} = \frac{1}{181400492}$
\end{itemize}
La somme des numérateurs est :
$$ \frac{1}{52} + \frac{1}{13469} + \frac{1}{181400492} = \frac{3488471 + 13468 + 1}{181400492} = \frac{3501940}{181400492} $$
Le PGCD du numérateur et du dénominateur est $\text{PGCD}(3501940, 181400492) = 700388$.
La fraction irréductible est :
$$ \frac{3501940}{181400492} = \frac{3501940 \div 700388}{181400492 \div 700388} = \frac{5}{259} $$
Cette fraction est bien égale à $\frac{5}{259}$.

\subsection{Démonstration pour $n = 260$}
Soit $n = 260$. Nous cherchons $x, y, z \in \mathbb{Z}^{+}$ tels que $\frac{5}{260} = \frac{1}{x} + \frac{1}{y} + \frac{1}{z}$.
Posons $x = 53$, $y = 2757$, $z = 7598292$.
Les conditions $x > 0$, $y > 0$ et $z > 0$ sont satisfaites car $x, y, z \geq 1$.
Le PPCM des dénominateurs est $\text{PPCM}(53, 2757, 7598292) = 7598292$.
En réduisant au même dénominateur :
\begin{itemize}
    \item $\frac{1}{53} = \frac{143364}{7598292}$
    \item $\frac{1}{2757} = \frac{2756}{7598292}$
    \item $\frac{1}{7598292} = \frac{1}{7598292}$
\end{itemize}
La somme des numérateurs est :
$$ \frac{1}{53} + \frac{1}{2757} + \frac{1}{7598292} = \frac{143364 + 2756 + 1}{7598292} = \frac{146121}{7598292} $$
Le PGCD du numérateur et du dénominateur est $\text{PGCD}(146121, 7598292) = 146121$.
La fraction irréductible est :
$$ \frac{146121}{7598292} = \frac{146121 \div 146121}{7598292 \div 146121} = \frac{1}{52} $$
Cette fraction est bien égale à $\frac{5}{260}$.

\subsection{Démonstration pour $n = 261$}
Soit $n = 261$. Nous cherchons $x, y, z \in \mathbb{Z}^{+}$ tels que $\frac{5}{261} = \frac{1}{x} + \frac{1}{y} + \frac{1}{z}$.
Posons $x = 53$, $y = 3459$, $z = 15949449$.
Les conditions $x > 0$, $y > 0$ et $z > 0$ sont satisfaites car $x, y, z \geq 1$.
Le PPCM des dénominateurs est $\text{PPCM}(53, 3459, 15949449) = 15949449$.
En réduisant au même dénominateur :
\begin{itemize}
    \item $\frac{1}{53} = \frac{300933}{15949449}$
    \item $\frac{1}{3459} = \frac{4611}{15949449}$
    \item $\frac{1}{15949449} = \frac{1}{15949449}$
\end{itemize}
La somme des numérateurs est :
$$ \frac{1}{53} + \frac{1}{3459} + \frac{1}{15949449} = \frac{300933 + 4611 + 1}{15949449} = \frac{305545}{15949449} $$
Le PGCD du numérateur et du dénominateur est $\text{PGCD}(305545, 15949449) = 61109$.
La fraction irréductible est :
$$ \frac{305545}{15949449} = \frac{305545 \div 61109}{15949449 \div 61109} = \frac{5}{261} $$
Cette fraction est bien égale à $\frac{5}{261}$.

\subsection{Démonstration pour $n = 262$}
Soit $n = 262$. Nous cherchons $x, y, z \in \mathbb{Z}^{+}$ tels que $\frac{5}{262} = \frac{1}{x} + \frac{1}{y} + \frac{1}{z}$.
Posons $x = 53$, $y = 4629$, $z = 64278294$.
Les conditions $x > 0$, $y > 0$ et $z > 0$ sont satisfaites car $x, y, z \geq 1$.
Le PPCM des dénominateurs est $\text{PPCM}(53, 4629, 64278294) = 64278294$.
En réduisant au même dénominateur :
\begin{itemize}
    \item $\frac{1}{53} = \frac{1212798}{64278294}$
    \item $\frac{1}{4629} = \frac{13886}{64278294}$
    \item $\frac{1}{64278294} = \frac{1}{64278294}$
\end{itemize}
La somme des numérateurs est :
$$ \frac{1}{53} + \frac{1}{4629} + \frac{1}{64278294} = \frac{1212798 + 13886 + 1}{64278294} = \frac{1226685}{64278294} $$
Le PGCD du numérateur et du dénominateur est $\text{PGCD}(1226685, 64278294) = 245337$.
La fraction irréductible est :
$$ \frac{1226685}{64278294} = \frac{1226685 \div 245337}{64278294 \div 245337} = \frac{5}{262} $$
Cette fraction est bien égale à $\frac{5}{262}$.

\subsection{Démonstration pour $n = 263$}
Soit $n = 263$. Nous cherchons $x, y, z \in \mathbb{Z}^{+}$ tels que $\frac{5}{263} = \frac{1}{x} + \frac{1}{y} + \frac{1}{z}$.
Posons $x = 53$, $y = 6970$, $z = 97154830$.
Les conditions $x > 0$, $y > 0$ et $z > 0$ sont satisfaites car $x, y, z \geq 1$.
Le PPCM des dénominateurs est $\text{PPCM}(53, 6970, 97154830) = 97154830$.
En réduisant au même dénominateur :
\begin{itemize}
    \item $\frac{1}{53} = \frac{1833110}{97154830}$
    \item $\frac{1}{6970} = \frac{13939}{97154830}$
    \item $\frac{1}{97154830} = \frac{1}{97154830}$
\end{itemize}
La somme des numérateurs est :
$$ \frac{1}{53} + \frac{1}{6970} + \frac{1}{97154830} = \frac{1833110 + 13939 + 1}{97154830} = \frac{1847050}{97154830} $$
Le PGCD du numérateur et du dénominateur est $\text{PGCD}(1847050, 97154830) = 369410$.
La fraction irréductible est :
$$ \frac{1847050}{97154830} = \frac{1847050 \div 369410}{97154830 \div 369410} = \frac{5}{263} $$
Cette fraction est bien égale à $\frac{5}{263}$.

\subsection{Démonstration pour $n = 264$}
Soit $n = 264$. Nous cherchons $x, y, z \in \mathbb{Z}^{+}$ tels que $\frac{5}{264} = \frac{1}{x} + \frac{1}{y} + \frac{1}{z}$.
Posons $x = 53$, $y = 13993$, $z = 195790056$.
Les conditions $x > 0$, $y > 0$ et $z > 0$ sont satisfaites car $x, y, z \geq 1$.
Le PPCM des dénominateurs est $\text{PPCM}(53, 13993, 195790056) = 195790056$.
En réduisant au même dénominateur :
\begin{itemize}
    \item $\frac{1}{53} = \frac{3694152}{195790056}$
    \item $\frac{1}{13993} = \frac{13992}{195790056}$
    \item $\frac{1}{195790056} = \frac{1}{195790056}$
\end{itemize}
La somme des numérateurs est :
$$ \frac{1}{53} + \frac{1}{13993} + \frac{1}{195790056} = \frac{3694152 + 13992 + 1}{195790056} = \frac{3708145}{195790056} $$
Le PGCD du numérateur et du dénominateur est $\text{PGCD}(3708145, 195790056) = 741629$.
La fraction irréductible est :
$$ \frac{3708145}{195790056} = \frac{3708145 \div 741629}{195790056 \div 741629} = \frac{5}{264} $$
Cette fraction est bien égale à $\frac{5}{264}$.

\subsection{Démonstration pour $n = 265$}
Soit $n = 265$. Nous cherchons $x, y, z \in \mathbb{Z}^{+}$ tels que $\frac{5}{265} = \frac{1}{x} + \frac{1}{y} + \frac{1}{z}$.
Posons $x = 54$, $y = 2863$, $z = 8193906$.
Les conditions $x > 0$, $y > 0$ et $z > 0$ sont satisfaites car $x, y, z \geq 1$.
Le PPCM des dénominateurs est $\text{PPCM}(54, 2863, 8193906) = 8193906$.
En réduisant au même dénominateur :
\begin{itemize}
    \item $\frac{1}{54} = \frac{151739}{8193906}$
    \item $\frac{1}{2863} = \frac{2862}{8193906}$
    \item $\frac{1}{8193906} = \frac{1}{8193906}$
\end{itemize}
La somme des numérateurs est :
$$ \frac{1}{54} + \frac{1}{2863} + \frac{1}{8193906} = \frac{151739 + 2862 + 1}{8193906} = \frac{154602}{8193906} $$
Le PGCD du numérateur et du dénominateur est $\text{PGCD}(154602, 8193906) = 154602$.
La fraction irréductible est :
$$ \frac{154602}{8193906} = \frac{154602 \div 154602}{8193906 \div 154602} = \frac{1}{53} $$
Cette fraction est bien égale à $\frac{5}{265}$.

\subsection{Démonstration pour $n = 266$}
Soit $n = 266$. Nous cherchons $x, y, z \in \mathbb{Z}^{+}$ tels que $\frac{5}{266} = \frac{1}{x} + \frac{1}{y} + \frac{1}{z}$.
Posons $x = 54$, $y = 3592$, $z = 12898872$.
Les conditions $x > 0$, $y > 0$ et $z > 0$ sont satisfaites car $x, y, z \geq 1$.
Le PPCM des dénominateurs est $\text{PPCM}(54, 3592, 12898872) = 12898872$.
En réduisant au même dénominateur :
\begin{itemize}
    \item $\frac{1}{54} = \frac{238868}{12898872}$
    \item $\frac{1}{3592} = \frac{3591}{12898872}$
    \item $\frac{1}{12898872} = \frac{1}{12898872}$
\end{itemize}
La somme des numérateurs est :
$$ \frac{1}{54} + \frac{1}{3592} + \frac{1}{12898872} = \frac{238868 + 3591 + 1}{12898872} = \frac{242460}{12898872} $$
Le PGCD du numérateur et du dénominateur est $\text{PGCD}(242460, 12898872) = 48492$.
La fraction irréductible est :
$$ \frac{242460}{12898872} = \frac{242460 \div 48492}{12898872 \div 48492} = \frac{5}{266} $$
Cette fraction est bien égale à $\frac{5}{266}$.

\subsection{Démonstration pour $n = 267$}
Soit $n = 267$. Nous cherchons $x, y, z \in \mathbb{Z}^{+}$ tels que $\frac{5}{267} = \frac{1}{x} + \frac{1}{y} + \frac{1}{z}$.
Posons $x = 54$, $y = 4807$, $z = 23102442$.
Les conditions $x > 0$, $y > 0$ et $z > 0$ sont satisfaites car $x, y, z \geq 1$.
Le PPCM des dénominateurs est $\text{PPCM}(54, 4807, 23102442) = 23102442$.
En réduisant au même dénominateur :
\begin{itemize}
    \item $\frac{1}{54} = \frac{427823}{23102442}$
    \item $\frac{1}{4807} = \frac{4806}{23102442}$
    \item $\frac{1}{23102442} = \frac{1}{23102442}$
\end{itemize}
La somme des numérateurs est :
$$ \frac{1}{54} + \frac{1}{4807} + \frac{1}{23102442} = \frac{427823 + 4806 + 1}{23102442} = \frac{432630}{23102442} $$
Le PGCD du numérateur et du dénominateur est $\text{PGCD}(432630, 23102442) = 86526$.
La fraction irréductible est :
$$ \frac{432630}{23102442} = \frac{432630 \div 86526}{23102442 \div 86526} = \frac{5}{267} $$
Cette fraction est bien égale à $\frac{5}{267}$.

\subsection{Démonstration pour $n = 268$}
Soit $n = 268$. Nous cherchons $x, y, z \in \mathbb{Z}^{+}$ tels que $\frac{5}{268} = \frac{1}{x} + \frac{1}{y} + \frac{1}{z}$.
Posons $x = 54$, $y = 7237$, $z = 52366932$.
Les conditions $x > 0$, $y > 0$ et $z > 0$ sont satisfaites car $x, y, z \geq 1$.
Le PPCM des dénominateurs est $\text{PPCM}(54, 7237, 52366932) = 52366932$.
En réduisant au même dénominateur :
\begin{itemize}
    \item $\frac{1}{54} = \frac{969758}{52366932}$
    \item $\frac{1}{7237} = \frac{7236}{52366932}$
    \item $\frac{1}{52366932} = \frac{1}{52366932}$
\end{itemize}
La somme des numérateurs est :
$$ \frac{1}{54} + \frac{1}{7237} + \frac{1}{52366932} = \frac{969758 + 7236 + 1}{52366932} = \frac{976995}{52366932} $$
Le PGCD du numérateur et du dénominateur est $\text{PGCD}(976995, 52366932) = 195399$.
La fraction irréductible est :
$$ \frac{976995}{52366932} = \frac{976995 \div 195399}{52366932 \div 195399} = \frac{5}{268} $$
Cette fraction est bien égale à $\frac{5}{268}$.

\subsection{Démonstration pour $n = 269$}
Soit $n = 269$. Nous cherchons $x, y, z \in \mathbb{Z}^{+}$ tels que $\frac{5}{269} = \frac{1}{x} + \frac{1}{y} + \frac{1}{z}$.
Posons $x = 54$, $y = 14527$, $z = 211019202$.
Les conditions $x > 0$, $y > 0$ et $z > 0$ sont satisfaites car $x, y, z \geq 1$.
Le PPCM des dénominateurs est $\text{PPCM}(54, 14527, 211019202) = 211019202$.
En réduisant au même dénominateur :
\begin{itemize}
    \item $\frac{1}{54} = \frac{3907763}{211019202}$
    \item $\frac{1}{14527} = \frac{14526}{211019202}$
    \item $\frac{1}{211019202} = \frac{1}{211019202}$
\end{itemize}
La somme des numérateurs est :
$$ \frac{1}{54} + \frac{1}{14527} + \frac{1}{211019202} = \frac{3907763 + 14526 + 1}{211019202} = \frac{3922290}{211019202} $$
Le PGCD du numérateur et du dénominateur est $\text{PGCD}(3922290, 211019202) = 784458$.
La fraction irréductible est :
$$ \frac{3922290}{211019202} = \frac{3922290 \div 784458}{211019202 \div 784458} = \frac{5}{269} $$
Cette fraction est bien égale à $\frac{5}{269}$.

\subsection{Démonstration pour $n = 270$}
Soit $n = 270$. Nous cherchons $x, y, z \in \mathbb{Z}^{+}$ tels que $\frac{5}{270} = \frac{1}{x} + \frac{1}{y} + \frac{1}{z}$.
Posons $x = 55$, $y = 2971$, $z = 8823870$.
Les conditions $x > 0$, $y > 0$ et $z > 0$ sont satisfaites car $x, y, z \geq 1$.
Le PPCM des dénominateurs est $\text{PPCM}(55, 2971, 8823870) = 8823870$.
En réduisant au même dénominateur :
\begin{itemize}
    \item $\frac{1}{55} = \frac{160434}{8823870}$
    \item $\frac{1}{2971} = \frac{2970}{8823870}$
    \item $\frac{1}{8823870} = \frac{1}{8823870}$
\end{itemize}
La somme des numérateurs est :
$$ \frac{1}{55} + \frac{1}{2971} + \frac{1}{8823870} = \frac{160434 + 2970 + 1}{8823870} = \frac{163405}{8823870} $$
Le PGCD du numérateur et du dénominateur est $\text{PGCD}(163405, 8823870) = 163405$.
La fraction irréductible est :
$$ \frac{163405}{8823870} = \frac{163405 \div 163405}{8823870 \div 163405} = \frac{1}{54} $$
Cette fraction est bien égale à $\frac{5}{270}$.

\subsection{Démonstration pour $n = 271$}
Soit $n = 271$. Nous cherchons $x, y, z \in \mathbb{Z}^{+}$ tels que $\frac{5}{271} = \frac{1}{x} + \frac{1}{y} + \frac{1}{z}$.
Posons $x = 55$, $y = 3729$, $z = 5052795$.
Les conditions $x > 0$, $y > 0$ et $z > 0$ sont satisfaites car $x, y, z \geq 1$.
Le PPCM des dénominateurs est $\text{PPCM}(55, 3729, 5052795) = 5052795$.
En réduisant au même dénominateur :
\begin{itemize}
    \item $\frac{1}{55} = \frac{91869}{5052795}$
    \item $\frac{1}{3729} = \frac{1355}{5052795}$
    \item $\frac{1}{5052795} = \frac{1}{5052795}$
\end{itemize}
La somme des numérateurs est :
$$ \frac{1}{55} + \frac{1}{3729} + \frac{1}{5052795} = \frac{91869 + 1355 + 1}{5052795} = \frac{93225}{5052795} $$
Le PGCD du numérateur et du dénominateur est $\text{PGCD}(93225, 5052795) = 18645$.
La fraction irréductible est :
$$ \frac{93225}{5052795} = \frac{93225 \div 18645}{5052795 \div 18645} = \frac{5}{271} $$
Cette fraction est bien égale à $\frac{5}{271}$.

\subsection{Démonstration pour $n = 272$}
Soit $n = 272$. Nous cherchons $x, y, z \in \mathbb{Z}^{+}$ tels que $\frac{5}{272} = \frac{1}{x} + \frac{1}{y} + \frac{1}{z}$.
Posons $x = 55$, $y = 4987$, $z = 74605520$.
Les conditions $x > 0$, $y > 0$ et $z > 0$ sont satisfaites car $x, y, z \geq 1$.
Le PPCM des dénominateurs est $\text{PPCM}(55, 4987, 74605520) = 74605520$.
En réduisant au même dénominateur :
\begin{itemize}
    \item $\frac{1}{55} = \frac{1356464}{74605520}$
    \item $\frac{1}{4987} = \frac{14960}{74605520}$
    \item $\frac{1}{74605520} = \frac{1}{74605520}$
\end{itemize}
La somme des numérateurs est :
$$ \frac{1}{55} + \frac{1}{4987} + \frac{1}{74605520} = \frac{1356464 + 14960 + 1}{74605520} = \frac{1371425}{74605520} $$
Le PGCD du numérateur et du dénominateur est $\text{PGCD}(1371425, 74605520) = 274285$.
La fraction irréductible est :
$$ \frac{1371425}{74605520} = \frac{1371425 \div 274285}{74605520 \div 274285} = \frac{5}{272} $$
Cette fraction est bien égale à $\frac{5}{272}$.

\subsection{Démonstration pour $n = 273$}
Soit $n = 273$. Nous cherchons $x, y, z \in \mathbb{Z}^{+}$ tels que $\frac{5}{273} = \frac{1}{x} + \frac{1}{y} + \frac{1}{z}$.
Posons $x = 55$, $y = 7508$, $z = 112732620$.
Les conditions $x > 0$, $y > 0$ et $z > 0$ sont satisfaites car $x, y, z \geq 1$.
Le PPCM des dénominateurs est $\text{PPCM}(55, 7508, 112732620) = 112732620$.
En réduisant au même dénominateur :
\begin{itemize}
    \item $\frac{1}{55} = \frac{2049684}{112732620}$
    \item $\frac{1}{7508} = \frac{15015}{112732620}$
    \item $\frac{1}{112732620} = \frac{1}{112732620}$
\end{itemize}
La somme des numérateurs est :
$$ \frac{1}{55} + \frac{1}{7508} + \frac{1}{112732620} = \frac{2049684 + 15015 + 1}{112732620} = \frac{2064700}{112732620} $$
Le PGCD du numérateur et du dénominateur est $\text{PGCD}(2064700, 112732620) = 412940$.
La fraction irréductible est :
$$ \frac{2064700}{112732620} = \frac{2064700 \div 412940}{112732620 \div 412940} = \frac{5}{273} $$
Cette fraction est bien égale à $\frac{5}{273}$.

\subsection{Démonstration pour $n = 274$}
Soit $n = 274$. Nous cherchons $x, y, z \in \mathbb{Z}^{+}$ tels que $\frac{5}{274} = \frac{1}{x} + \frac{1}{y} + \frac{1}{z}$.
Posons $x = 55$, $y = 15071$, $z = 227119970$.
Les conditions $x > 0$, $y > 0$ et $z > 0$ sont satisfaites car $x, y, z \geq 1$.
Le PPCM des dénominateurs est $\text{PPCM}(55, 15071, 227119970) = 227119970$.
En réduisant au même dénominateur :
\begin{itemize}
    \item $\frac{1}{55} = \frac{4129454}{227119970}$
    \item $\frac{1}{15071} = \frac{15070}{227119970}$
    \item $\frac{1}{227119970} = \frac{1}{227119970}$
\end{itemize}
La somme des numérateurs est :
$$ \frac{1}{55} + \frac{1}{15071} + \frac{1}{227119970} = \frac{4129454 + 15070 + 1}{227119970} = \frac{4144525}{227119970} $$
Le PGCD du numérateur et du dénominateur est $\text{PGCD}(4144525, 227119970) = 828905$.
La fraction irréductible est :
$$ \frac{4144525}{227119970} = \frac{4144525 \div 828905}{227119970 \div 828905} = \frac{5}{274} $$
Cette fraction est bien égale à $\frac{5}{274}$.

\subsection{Démonstration pour $n = 275$}
Soit $n = 275$. Nous cherchons $x, y, z \in \mathbb{Z}^{+}$ tels que $\frac{5}{275} = \frac{1}{x} + \frac{1}{y} + \frac{1}{z}$.
Posons $x = 56$, $y = 3081$, $z = 9489480$.
Les conditions $x > 0$, $y > 0$ et $z > 0$ sont satisfaites car $x, y, z \geq 1$.
Le PPCM des dénominateurs est $\text{PPCM}(56, 3081, 9489480) = 9489480$.
En réduisant au même dénominateur :
\begin{itemize}
    \item $\frac{1}{56} = \frac{169455}{9489480}$
    \item $\frac{1}{3081} = \frac{3080}{9489480}$
    \item $\frac{1}{9489480} = \frac{1}{9489480}$
\end{itemize}
La somme des numérateurs est :
$$ \frac{1}{56} + \frac{1}{3081} + \frac{1}{9489480} = \frac{169455 + 3080 + 1}{9489480} = \frac{172536}{9489480} $$
Le PGCD du numérateur et du dénominateur est $\text{PGCD}(172536, 9489480) = 172536$.
La fraction irréductible est :
$$ \frac{172536}{9489480} = \frac{172536 \div 172536}{9489480 \div 172536} = \frac{1}{55} $$
Cette fraction est bien égale à $\frac{5}{275}$.

\subsection{Démonstration pour $n = 276$}
Soit $n = 276$. Nous cherchons $x, y, z \in \mathbb{Z}^{+}$ tels que $\frac{5}{276} = \frac{1}{x} + \frac{1}{y} + \frac{1}{z}$.
Posons $x = 56$, $y = 3865$, $z = 14934360$.
Les conditions $x > 0$, $y > 0$ et $z > 0$ sont satisfaites car $x, y, z \geq 1$.
Le PPCM des dénominateurs est $\text{PPCM}(56, 3865, 14934360) = 14934360$.
En réduisant au même dénominateur :
\begin{itemize}
    \item $\frac{1}{56} = \frac{266685}{14934360}$
    \item $\frac{1}{3865} = \frac{3864}{14934360}$
    \item $\frac{1}{14934360} = \frac{1}{14934360}$
\end{itemize}
La somme des numérateurs est :
$$ \frac{1}{56} + \frac{1}{3865} + \frac{1}{14934360} = \frac{266685 + 3864 + 1}{14934360} = \frac{270550}{14934360} $$
Le PGCD du numérateur et du dénominateur est $\text{PGCD}(270550, 14934360) = 54110$.
La fraction irréductible est :
$$ \frac{270550}{14934360} = \frac{270550 \div 54110}{14934360 \div 54110} = \frac{5}{276} $$
Cette fraction est bien égale à $\frac{5}{276}$.

\subsection{Démonstration pour $n = 277$}
Soit $n = 277$. Nous cherchons $x, y, z \in \mathbb{Z}^{+}$ tels que $\frac{5}{277} = \frac{1}{x} + \frac{1}{y} + \frac{1}{z}$.
Posons $x = 56$, $y = 5171$, $z = 80212552$.
Les conditions $x > 0$, $y > 0$ et $z > 0$ sont satisfaites car $x, y, z \geq 1$.
Le PPCM des dénominateurs est $\text{PPCM}(56, 5171, 80212552) = 80212552$.
En réduisant au même dénominateur :
\begin{itemize}
    \item $\frac{1}{56} = \frac{1432367}{80212552}$
    \item $\frac{1}{5171} = \frac{15512}{80212552}$
    \item $\frac{1}{80212552} = \frac{1}{80212552}$
\end{itemize}
La somme des numérateurs est :
$$ \frac{1}{56} + \frac{1}{5171} + \frac{1}{80212552} = \frac{1432367 + 15512 + 1}{80212552} = \frac{1447880}{80212552} $$
Le PGCD du numérateur et du dénominateur est $\text{PGCD}(1447880, 80212552) = 289576$.
La fraction irréductible est :
$$ \frac{1447880}{80212552} = \frac{1447880 \div 289576}{80212552 \div 289576} = \frac{5}{277} $$
Cette fraction est bien égale à $\frac{5}{277}$.

\subsection{Démonstration pour $n = 278$}
Soit $n = 278$. Nous cherchons $x, y, z \in \mathbb{Z}^{+}$ tels que $\frac{5}{278} = \frac{1}{x} + \frac{1}{y} + \frac{1}{z}$.
Posons $x = 56$, $y = 7785$, $z = 60598440$.
Les conditions $x > 0$, $y > 0$ et $z > 0$ sont satisfaites car $x, y, z \geq 1$.
Le PPCM des dénominateurs est $\text{PPCM}(56, 7785, 60598440) = 60598440$.
En réduisant au même dénominateur :
\begin{itemize}
    \item $\frac{1}{56} = \frac{1082115}{60598440}$
    \item $\frac{1}{7785} = \frac{7784}{60598440}$
    \item $\frac{1}{60598440} = \frac{1}{60598440}$
\end{itemize}
La somme des numérateurs est :
$$ \frac{1}{56} + \frac{1}{7785} + \frac{1}{60598440} = \frac{1082115 + 7784 + 1}{60598440} = \frac{1089900}{60598440} $$
Le PGCD du numérateur et du dénominateur est $\text{PGCD}(1089900, 60598440) = 217980$.
La fraction irréductible est :
$$ \frac{1089900}{60598440} = \frac{1089900 \div 217980}{60598440 \div 217980} = \frac{5}{278} $$
Cette fraction est bien égale à $\frac{5}{278}$.

\subsection{Démonstration pour $n = 279$}
Soit $n = 279$. Nous cherchons $x, y, z \in \mathbb{Z}^{+}$ tels que $\frac{5}{279} = \frac{1}{x} + \frac{1}{y} + \frac{1}{z}$.
Posons $x = 56$, $y = 15625$, $z = 244125000$.
Les conditions $x > 0$, $y > 0$ et $z > 0$ sont satisfaites car $x, y, z \geq 1$.
Le PPCM des dénominateurs est $\text{PPCM}(56, 15625, 244125000) = 244125000$.
En réduisant au même dénominateur :
\begin{itemize}
    \item $\frac{1}{56} = \frac{4359375}{244125000}$
    \item $\frac{1}{15625} = \frac{15624}{244125000}$
    \item $\frac{1}{244125000} = \frac{1}{244125000}$
\end{itemize}
La somme des numérateurs est :
$$ \frac{1}{56} + \frac{1}{15625} + \frac{1}{244125000} = \frac{4359375 + 15624 + 1}{244125000} = \frac{4375000}{244125000} $$
Le PGCD du numérateur et du dénominateur est $\text{PGCD}(4375000, 244125000) = 875000$.
La fraction irréductible est :
$$ \frac{4375000}{244125000} = \frac{4375000 \div 875000}{244125000 \div 875000} = \frac{5}{279} $$
Cette fraction est bien égale à $\frac{5}{279}$.

\subsection{Démonstration pour $n = 280$}
Soit $n = 280$. Nous cherchons $x, y, z \in \mathbb{Z}^{+}$ tels que $\frac{5}{280} = \frac{1}{x} + \frac{1}{y} + \frac{1}{z}$.
Posons $x = 57$, $y = 3193$, $z = 10192056$.
Les conditions $x > 0$, $y > 0$ et $z > 0$ sont satisfaites car $x, y, z \geq 1$.
Le PPCM des dénominateurs est $\text{PPCM}(57, 3193, 10192056) = 10192056$.
En réduisant au même dénominateur :
\begin{itemize}
    \item $\frac{1}{57} = \frac{178808}{10192056}$
    \item $\frac{1}{3193} = \frac{3192}{10192056}$
    \item $\frac{1}{10192056} = \frac{1}{10192056}$
\end{itemize}
La somme des numérateurs est :
$$ \frac{1}{57} + \frac{1}{3193} + \frac{1}{10192056} = \frac{178808 + 3192 + 1}{10192056} = \frac{182001}{10192056} $$
Le PGCD du numérateur et du dénominateur est $\text{PGCD}(182001, 10192056) = 182001$.
La fraction irréductible est :
$$ \frac{182001}{10192056} = \frac{182001 \div 182001}{10192056 \div 182001} = \frac{1}{56} $$
Cette fraction est bien égale à $\frac{5}{280}$.

\subsection{Démonstration pour $n = 281$}
Soit $n = 281$. Nous cherchons $x, y, z \in \mathbb{Z}^{+}$ tels que $\frac{5}{281} = \frac{1}{x} + \frac{1}{y} + \frac{1}{z}$.
Posons $x = 57$, $y = 4005$, $z = 21382695$.
Les conditions $x > 0$, $y > 0$ et $z > 0$ sont satisfaites car $x, y, z \geq 1$.
Le PPCM des dénominateurs est $\text{PPCM}(57, 4005, 21382695) = 21382695$.
En réduisant au même dénominateur :
\begin{itemize}
    \item $\frac{1}{57} = \frac{375135}{21382695}$
    \item $\frac{1}{4005} = \frac{5339}{21382695}$
    \item $\frac{1}{21382695} = \frac{1}{21382695}$
\end{itemize}
La somme des numérateurs est :
$$ \frac{1}{57} + \frac{1}{4005} + \frac{1}{21382695} = \frac{375135 + 5339 + 1}{21382695} = \frac{380475}{21382695} $$
Le PGCD du numérateur et du dénominateur est $\text{PGCD}(380475, 21382695) = 76095$.
La fraction irréductible est :
$$ \frac{380475}{21382695} = \frac{380475 \div 76095}{21382695 \div 76095} = \frac{5}{281} $$
Cette fraction est bien égale à $\frac{5}{281}$.

\subsection{Démonstration pour $n = 282$}
Soit $n = 282$. Nous cherchons $x, y, z \in \mathbb{Z}^{+}$ tels que $\frac{5}{282} = \frac{1}{x} + \frac{1}{y} + \frac{1}{z}$.
Posons $x = 57$, $y = 5359$, $z = 28713522$.
Les conditions $x > 0$, $y > 0$ et $z > 0$ sont satisfaites car $x, y, z \geq 1$.
Le PPCM des dénominateurs est $\text{PPCM}(57, 5359, 28713522) = 28713522$.
En réduisant au même dénominateur :
\begin{itemize}
    \item $\frac{1}{57} = \frac{503746}{28713522}$
    \item $\frac{1}{5359} = \frac{5358}{28713522}$
    \item $\frac{1}{28713522} = \frac{1}{28713522}$
\end{itemize}
La somme des numérateurs est :
$$ \frac{1}{57} + \frac{1}{5359} + \frac{1}{28713522} = \frac{503746 + 5358 + 1}{28713522} = \frac{509105}{28713522} $$
Le PGCD du numérateur et du dénominateur est $\text{PGCD}(509105, 28713522) = 101821$.
La fraction irréductible est :
$$ \frac{509105}{28713522} = \frac{509105 \div 101821}{28713522 \div 101821} = \frac{5}{282} $$
Cette fraction est bien égale à $\frac{5}{282}$.

\subsection{Démonstration pour $n = 283$}
Soit $n = 283$. Nous cherchons $x, y, z \in \mathbb{Z}^{+}$ tels que $\frac{5}{283} = \frac{1}{x} + \frac{1}{y} + \frac{1}{z}$.
Posons $x = 57$, $y = 8066$, $z = 130112646$.
Les conditions $x > 0$, $y > 0$ et $z > 0$ sont satisfaites car $x, y, z \geq 1$.
Le PPCM des dénominateurs est $\text{PPCM}(57, 8066, 130112646) = 130112646$.
En réduisant au même dénominateur :
\begin{itemize}
    \item $\frac{1}{57} = \frac{2282678}{130112646}$
    \item $\frac{1}{8066} = \frac{16131}{130112646}$
    \item $\frac{1}{130112646} = \frac{1}{130112646}$
\end{itemize}
La somme des numérateurs est :
$$ \frac{1}{57} + \frac{1}{8066} + \frac{1}{130112646} = \frac{2282678 + 16131 + 1}{130112646} = \frac{2298810}{130112646} $$
Le PGCD du numérateur et du dénominateur est $\text{PGCD}(2298810, 130112646) = 459762$.
La fraction irréductible est :
$$ \frac{2298810}{130112646} = \frac{2298810 \div 459762}{130112646 \div 459762} = \frac{5}{283} $$
Cette fraction est bien égale à $\frac{5}{283}$.

\subsection{Démonstration pour $n = 284$}
Soit $n = 284$. Nous cherchons $x, y, z \in \mathbb{Z}^{+}$ tels que $\frac{5}{284} = \frac{1}{x} + \frac{1}{y} + \frac{1}{z}$.
Posons $x = 57$, $y = 16189$, $z = 262067532$.
Les conditions $x > 0$, $y > 0$ et $z > 0$ sont satisfaites car $x, y, z \geq 1$.
Le PPCM des dénominateurs est $\text{PPCM}(57, 16189, 262067532) = 262067532$.
En réduisant au même dénominateur :
\begin{itemize}
    \item $\frac{1}{57} = \frac{4597676}{262067532}$
    \item $\frac{1}{16189} = \frac{16188}{262067532}$
    \item $\frac{1}{262067532} = \frac{1}{262067532}$
\end{itemize}
La somme des numérateurs est :
$$ \frac{1}{57} + \frac{1}{16189} + \frac{1}{262067532} = \frac{4597676 + 16188 + 1}{262067532} = \frac{4613865}{262067532} $$
Le PGCD du numérateur et du dénominateur est $\text{PGCD}(4613865, 262067532) = 922773$.
La fraction irréductible est :
$$ \frac{4613865}{262067532} = \frac{4613865 \div 922773}{262067532 \div 922773} = \frac{5}{284} $$
Cette fraction est bien égale à $\frac{5}{284}$.

\subsection{Démonstration pour $n = 285$}
Soit $n = 285$. Nous cherchons $x, y, z \in \mathbb{Z}^{+}$ tels que $\frac{5}{285} = \frac{1}{x} + \frac{1}{y} + \frac{1}{z}$.
Posons $x = 58$, $y = 3307$, $z = 10932942$.
Les conditions $x > 0$, $y > 0$ et $z > 0$ sont satisfaites car $x, y, z \geq 1$.
Le PPCM des dénominateurs est $\text{PPCM}(58, 3307, 10932942) = 10932942$.
En réduisant au même dénominateur :
\begin{itemize}
    \item $\frac{1}{58} = \frac{188499}{10932942}$
    \item $\frac{1}{3307} = \frac{3306}{10932942}$
    \item $\frac{1}{10932942} = \frac{1}{10932942}$
\end{itemize}
La somme des numérateurs est :
$$ \frac{1}{58} + \frac{1}{3307} + \frac{1}{10932942} = \frac{188499 + 3306 + 1}{10932942} = \frac{191806}{10932942} $$
Le PGCD du numérateur et du dénominateur est $\text{PGCD}(191806, 10932942) = 191806$.
La fraction irréductible est :
$$ \frac{191806}{10932942} = \frac{191806 \div 191806}{10932942 \div 191806} = \frac{1}{57} $$
Cette fraction est bien égale à $\frac{5}{285}$.

\subsection{Démonstration pour $n = 286$}
Soit $n = 286$. Nous cherchons $x, y, z \in \mathbb{Z}^{+}$ tels que $\frac{5}{286} = \frac{1}{x} + \frac{1}{y} + \frac{1}{z}$.
Posons $x = 58$, $y = 4148$, $z = 17201756$.
Les conditions $x > 0$, $y > 0$ et $z > 0$ sont satisfaites car $x, y, z \geq 1$.
Le PPCM des dénominateurs est $\text{PPCM}(58, 4148, 17201756) = 17201756$.
En réduisant au même dénominateur :
\begin{itemize}
    \item $\frac{1}{58} = \frac{296582}{17201756}$
    \item $\frac{1}{4148} = \frac{4147}{17201756}$
    \item $\frac{1}{17201756} = \frac{1}{17201756}$
\end{itemize}
La somme des numérateurs est :
$$ \frac{1}{58} + \frac{1}{4148} + \frac{1}{17201756} = \frac{296582 + 4147 + 1}{17201756} = \frac{300730}{17201756} $$
Le PGCD du numérateur et du dénominateur est $\text{PGCD}(300730, 17201756) = 60146$.
La fraction irréductible est :
$$ \frac{300730}{17201756} = \frac{300730 \div 60146}{17201756 \div 60146} = \frac{5}{286} $$
Cette fraction est bien égale à $\frac{5}{286}$.

\subsection{Démonstration pour $n = 287$}
Soit $n = 287$. Nous cherchons $x, y, z \in \mathbb{Z}^{+}$ tels que $\frac{5}{287} = \frac{1}{x} + \frac{1}{y} + \frac{1}{z}$.
Posons $x = 58$, $y = 5549$, $z = 92368654$.
Les conditions $x > 0$, $y > 0$ et $z > 0$ sont satisfaites car $x, y, z \geq 1$.
Le PPCM des dénominateurs est $\text{PPCM}(58, 5549, 92368654) = 92368654$.
En réduisant au même dénominateur :
\begin{itemize}
    \item $\frac{1}{58} = \frac{1592563}{92368654}$
    \item $\frac{1}{5549} = \frac{16646}{92368654}$
    \item $\frac{1}{92368654} = \frac{1}{92368654}$
\end{itemize}
La somme des numérateurs est :
$$ \frac{1}{58} + \frac{1}{5549} + \frac{1}{92368654} = \frac{1592563 + 16646 + 1}{92368654} = \frac{1609210}{92368654} $$
Le PGCD du numérateur et du dénominateur est $\text{PGCD}(1609210, 92368654) = 321842$.
La fraction irréductible est :
$$ \frac{1609210}{92368654} = \frac{1609210 \div 321842}{92368654 \div 321842} = \frac{5}{287} $$
Cette fraction est bien égale à $\frac{5}{287}$.

\subsection{Démonstration pour $n = 288$}
Soit $n = 288$. Nous cherchons $x, y, z \in \mathbb{Z}^{+}$ tels que $\frac{5}{288} = \frac{1}{x} + \frac{1}{y} + \frac{1}{z}$.
Posons $x = 58$, $y = 8353$, $z = 69764256$.
Les conditions $x > 0$, $y > 0$ et $z > 0$ sont satisfaites car $x, y, z \geq 1$.
Le PPCM des dénominateurs est $\text{PPCM}(58, 8353, 69764256) = 69764256$.
En réduisant au même dénominateur :
\begin{itemize}
    \item $\frac{1}{58} = \frac{1202832}{69764256}$
    \item $\frac{1}{8353} = \frac{8352}{69764256}$
    \item $\frac{1}{69764256} = \frac{1}{69764256}$
\end{itemize}
La somme des numérateurs est :
$$ \frac{1}{58} + \frac{1}{8353} + \frac{1}{69764256} = \frac{1202832 + 8352 + 1}{69764256} = \frac{1211185}{69764256} $$
Le PGCD du numérateur et du dénominateur est $\text{PGCD}(1211185, 69764256) = 242237$.
La fraction irréductible est :
$$ \frac{1211185}{69764256} = \frac{1211185 \div 242237}{69764256 \div 242237} = \frac{5}{288} $$
Cette fraction est bien égale à $\frac{5}{288}$.

\subsection{Démonstration pour $n = 289$}
Soit $n = 289$. Nous cherchons $x, y, z \in \mathbb{Z}^{+}$ tels que $\frac{5}{289} = \frac{1}{x} + \frac{1}{y} + \frac{1}{z}$.
Posons $x = 58$, $y = 16763$, $z = 280981406$.
Les conditions $x > 0$, $y > 0$ et $z > 0$ sont satisfaites car $x, y, z \geq 1$.
Le PPCM des dénominateurs est $\text{PPCM}(58, 16763, 280981406) = 280981406$.
En réduisant au même dénominateur :
\begin{itemize}
    \item $\frac{1}{58} = \frac{4844507}{280981406}$
    \item $\frac{1}{16763} = \frac{16762}{280981406}$
    \item $\frac{1}{280981406} = \frac{1}{280981406}$
\end{itemize}
La somme des numérateurs est :
$$ \frac{1}{58} + \frac{1}{16763} + \frac{1}{280981406} = \frac{4844507 + 16762 + 1}{280981406} = \frac{4861270}{280981406} $$
Le PGCD du numérateur et du dénominateur est $\text{PGCD}(4861270, 280981406) = 972254$.
La fraction irréductible est :
$$ \frac{4861270}{280981406} = \frac{4861270 \div 972254}{280981406 \div 972254} = \frac{5}{289} $$
Cette fraction est bien égale à $\frac{5}{289}$.

\subsection{Démonstration pour $n = 290$}
Soit $n = 290$. Nous cherchons $x, y, z \in \mathbb{Z}^{+}$ tels que $\frac{5}{290} = \frac{1}{x} + \frac{1}{y} + \frac{1}{z}$.
Posons $x = 59$, $y = 3423$, $z = 11713506$.
Les conditions $x > 0$, $y > 0$ et $z > 0$ sont satisfaites car $x, y, z \geq 1$.
Le PPCM des dénominateurs est $\text{PPCM}(59, 3423, 11713506) = 11713506$.
En réduisant au même dénominateur :
\begin{itemize}
    \item $\frac{1}{59} = \frac{198534}{11713506}$
    \item $\frac{1}{3423} = \frac{3422}{11713506}$
    \item $\frac{1}{11713506} = \frac{1}{11713506}$
\end{itemize}
La somme des numérateurs est :
$$ \frac{1}{59} + \frac{1}{3423} + \frac{1}{11713506} = \frac{198534 + 3422 + 1}{11713506} = \frac{201957}{11713506} $$
Le PGCD du numérateur et du dénominateur est $\text{PGCD}(201957, 11713506) = 201957$.
La fraction irréductible est :
$$ \frac{201957}{11713506} = \frac{201957 \div 201957}{11713506 \div 201957} = \frac{1}{58} $$
Cette fraction est bien égale à $\frac{5}{290}$.

\subsection{Démonstration pour $n = 291$}
Soit $n = 291$. Nous cherchons $x, y, z \in \mathbb{Z}^{+}$ tels que $\frac{5}{291} = \frac{1}{x} + \frac{1}{y} + \frac{1}{z}$.
Posons $x = 59$, $y = 4293$, $z = 24568839$.
Les conditions $x > 0$, $y > 0$ et $z > 0$ sont satisfaites car $x, y, z \geq 1$.
Le PPCM des dénominateurs est $\text{PPCM}(59, 4293, 24568839) = 24568839$.
En réduisant au même dénominateur :
\begin{itemize}
    \item $\frac{1}{59} = \frac{416421}{24568839}$
    \item $\frac{1}{4293} = \frac{5723}{24568839}$
    \item $\frac{1}{24568839} = \frac{1}{24568839}$
\end{itemize}
La somme des numérateurs est :
$$ \frac{1}{59} + \frac{1}{4293} + \frac{1}{24568839} = \frac{416421 + 5723 + 1}{24568839} = \frac{422145}{24568839} $$
Le PGCD du numérateur et du dénominateur est $\text{PGCD}(422145, 24568839) = 84429$.
La fraction irréductible est :
$$ \frac{422145}{24568839} = \frac{422145 \div 84429}{24568839 \div 84429} = \frac{5}{291} $$
Cette fraction est bien égale à $\frac{5}{291}$.

\subsection{Démonstration pour $n = 292$}
Soit $n = 292$. Nous cherchons $x, y, z \in \mathbb{Z}^{+}$ tels que $\frac{5}{292} = \frac{1}{x} + \frac{1}{y} + \frac{1}{z}$.
Posons $x = 59$, $y = 5743$, $z = 98940404$.
Les conditions $x > 0$, $y > 0$ et $z > 0$ sont satisfaites car $x, y, z \geq 1$.
Le PPCM des dénominateurs est $\text{PPCM}(59, 5743, 98940404) = 98940404$.
En réduisant au même dénominateur :
\begin{itemize}
    \item $\frac{1}{59} = \frac{1676956}{98940404}$
    \item $\frac{1}{5743} = \frac{17228}{98940404}$
    \item $\frac{1}{98940404} = \frac{1}{98940404}$
\end{itemize}
La somme des numérateurs est :
$$ \frac{1}{59} + \frac{1}{5743} + \frac{1}{98940404} = \frac{1676956 + 17228 + 1}{98940404} = \frac{1694185}{98940404} $$
Le PGCD du numérateur et du dénominateur est $\text{PGCD}(1694185, 98940404) = 338837$.
La fraction irréductible est :
$$ \frac{1694185}{98940404} = \frac{1694185 \div 338837}{98940404 \div 338837} = \frac{5}{292} $$
Cette fraction est bien égale à $\frac{5}{292}$.

\subsection{Démonstration pour $n = 293$}
Soit $n = 293$. Nous cherchons $x, y, z \in \mathbb{Z}^{+}$ tels que $\frac{5}{293} = \frac{1}{x} + \frac{1}{y} + \frac{1}{z}$.
Posons $x = 59$, $y = 8644$, $z = 149428828$.
Les conditions $x > 0$, $y > 0$ et $z > 0$ sont satisfaites car $x, y, z \geq 1$.
Le PPCM des dénominateurs est $\text{PPCM}(59, 8644, 149428828) = 149428828$.
En réduisant au même dénominateur :
\begin{itemize}
    \item $\frac{1}{59} = \frac{2532692}{149428828}$
    \item $\frac{1}{8644} = \frac{17287}{149428828}$
    \item $\frac{1}{149428828} = \frac{1}{149428828}$
\end{itemize}
La somme des numérateurs est :
$$ \frac{1}{59} + \frac{1}{8644} + \frac{1}{149428828} = \frac{2532692 + 17287 + 1}{149428828} = \frac{2549980}{149428828} $$
Le PGCD du numérateur et du dénominateur est $\text{PGCD}(2549980, 149428828) = 509996$.
La fraction irréductible est :
$$ \frac{2549980}{149428828} = \frac{2549980 \div 509996}{149428828 \div 509996} = \frac{5}{293} $$
Cette fraction est bien égale à $\frac{5}{293}$.

\subsection{Démonstration pour $n = 294$}
Soit $n = 294$. Nous cherchons $x, y, z \in \mathbb{Z}^{+}$ tels que $\frac{5}{294} = \frac{1}{x} + \frac{1}{y} + \frac{1}{z}$.
Posons $x = 59$, $y = 17347$, $z = 300901062$.
Les conditions $x > 0$, $y > 0$ et $z > 0$ sont satisfaites car $x, y, z \geq 1$.
Le PPCM des dénominateurs est $\text{PPCM}(59, 17347, 300901062) = 300901062$.
En réduisant au même dénominateur :
\begin{itemize}
    \item $\frac{1}{59} = \frac{5100018}{300901062}$
    \item $\frac{1}{17347} = \frac{17346}{300901062}$
    \item $\frac{1}{300901062} = \frac{1}{300901062}$
\end{itemize}
La somme des numérateurs est :
$$ \frac{1}{59} + \frac{1}{17347} + \frac{1}{300901062} = \frac{5100018 + 17346 + 1}{300901062} = \frac{5117365}{300901062} $$
Le PGCD du numérateur et du dénominateur est $\text{PGCD}(5117365, 300901062) = 1023473$.
La fraction irréductible est :
$$ \frac{5117365}{300901062} = \frac{5117365 \div 1023473}{300901062 \div 1023473} = \frac{5}{294} $$
Cette fraction est bien égale à $\frac{5}{294}$.

\subsection{Démonstration pour $n = 295$}
Soit $n = 295$. Nous cherchons $x, y, z \in \mathbb{Z}^{+}$ tels que $\frac{5}{295} = \frac{1}{x} + \frac{1}{y} + \frac{1}{z}$.
Posons $x = 60$, $y = 3541$, $z = 12535140$.
Les conditions $x > 0$, $y > 0$ et $z > 0$ sont satisfaites car $x, y, z \geq 1$.
Le PPCM des dénominateurs est $\text{PPCM}(60, 3541, 12535140) = 12535140$.
En réduisant au même dénominateur :
\begin{itemize}
    \item $\frac{1}{60} = \frac{208919}{12535140}$
    \item $\frac{1}{3541} = \frac{3540}{12535140}$
    \item $\frac{1}{12535140} = \frac{1}{12535140}$
\end{itemize}
La somme des numérateurs est :
$$ \frac{1}{60} + \frac{1}{3541} + \frac{1}{12535140} = \frac{208919 + 3540 + 1}{12535140} = \frac{212460}{12535140} $$
Le PGCD du numérateur et du dénominateur est $\text{PGCD}(212460, 12535140) = 212460$.
La fraction irréductible est :
$$ \frac{212460}{12535140} = \frac{212460 \div 212460}{12535140 \div 212460} = \frac{1}{59} $$
Cette fraction est bien égale à $\frac{5}{295}$.

\subsection{Démonstration pour $n = 296$}
Soit $n = 296$. Nous cherchons $x, y, z \in \mathbb{Z}^{+}$ tels que $\frac{5}{296} = \frac{1}{x} + \frac{1}{y} + \frac{1}{z}$.
Posons $x = 60$, $y = 4441$, $z = 19718040$.
Les conditions $x > 0$, $y > 0$ et $z > 0$ sont satisfaites car $x, y, z \geq 1$.
Le PPCM des dénominateurs est $\text{PPCM}(60, 4441, 19718040) = 19718040$.
En réduisant au même dénominateur :
\begin{itemize}
    \item $\frac{1}{60} = \frac{328634}{19718040}$
    \item $\frac{1}{4441} = \frac{4440}{19718040}$
    \item $\frac{1}{19718040} = \frac{1}{19718040}$
\end{itemize}
La somme des numérateurs est :
$$ \frac{1}{60} + \frac{1}{4441} + \frac{1}{19718040} = \frac{328634 + 4440 + 1}{19718040} = \frac{333075}{19718040} $$
Le PGCD du numérateur et du dénominateur est $\text{PGCD}(333075, 19718040) = 66615$.
La fraction irréductible est :
$$ \frac{333075}{19718040} = \frac{333075 \div 66615}{19718040 \div 66615} = \frac{5}{296} $$
Cette fraction est bien égale à $\frac{5}{296}$.

\subsection{Démonstration pour $n = 297$}
Soit $n = 297$. Nous cherchons $x, y, z \in \mathbb{Z}^{+}$ tels que $\frac{5}{297} = \frac{1}{x} + \frac{1}{y} + \frac{1}{z}$.
Posons $x = 60$, $y = 5941$, $z = 35289540$.
Les conditions $x > 0$, $y > 0$ et $z > 0$ sont satisfaites car $x, y, z \geq 1$.
Le PPCM des dénominateurs est $\text{PPCM}(60, 5941, 35289540) = 35289540$.
En réduisant au même dénominateur :
\begin{itemize}
    \item $\frac{1}{60} = \frac{588159}{35289540}$
    \item $\frac{1}{5941} = \frac{5940}{35289540}$
    \item $\frac{1}{35289540} = \frac{1}{35289540}$
\end{itemize}
La somme des numérateurs est :
$$ \frac{1}{60} + \frac{1}{5941} + \frac{1}{35289540} = \frac{588159 + 5940 + 1}{35289540} = \frac{594100}{35289540} $$
Le PGCD du numérateur et du dénominateur est $\text{PGCD}(594100, 35289540) = 118820$.
La fraction irréductible est :
$$ \frac{594100}{35289540} = \frac{594100 \div 118820}{35289540 \div 118820} = \frac{5}{297} $$
Cette fraction est bien égale à $\frac{5}{297}$.

\subsection{Démonstration pour $n = 298$}
Soit $n = 298$. Nous cherchons $x, y, z \in \mathbb{Z}^{+}$ tels que $\frac{5}{298} = \frac{1}{x} + \frac{1}{y} + \frac{1}{z}$.
Posons $x = 60$, $y = 8941$, $z = 79932540$.
Les conditions $x > 0$, $y > 0$ et $z > 0$ sont satisfaites car $x, y, z \geq 1$.
Le PPCM des dénominateurs est $\text{PPCM}(60, 8941, 79932540) = 79932540$.
En réduisant au même dénominateur :
\begin{itemize}
    \item $\frac{1}{60} = \frac{1332209}{79932540}$
    \item $\frac{1}{8941} = \frac{8940}{79932540}$
    \item $\frac{1}{79932540} = \frac{1}{79932540}$
\end{itemize}
La somme des numérateurs est :
$$ \frac{1}{60} + \frac{1}{8941} + \frac{1}{79932540} = \frac{1332209 + 8940 + 1}{79932540} = \frac{1341150}{79932540} $$
Le PGCD du numérateur et du dénominateur est $\text{PGCD}(1341150, 79932540) = 268230$.
La fraction irréductible est :
$$ \frac{1341150}{79932540} = \frac{1341150 \div 268230}{79932540 \div 268230} = \frac{5}{298} $$
Cette fraction est bien égale à $\frac{5}{298}$.

\subsection{Démonstration pour $n = 299$}
Soit $n = 299$. Nous cherchons $x, y, z \in \mathbb{Z}^{+}$ tels que $\frac{5}{299} = \frac{1}{x} + \frac{1}{y} + \frac{1}{z}$.
Posons $x = 60$, $y = 17941$, $z = 321861540$.
Les conditions $x > 0$, $y > 0$ et $z > 0$ sont satisfaites car $x, y, z \geq 1$.
Le PPCM des dénominateurs est $\text{PPCM}(60, 17941, 321861540) = 321861540$.
En réduisant au même dénominateur :
\begin{itemize}
    \item $\frac{1}{60} = \frac{5364359}{321861540}$
    \item $\frac{1}{17941} = \frac{17940}{321861540}$
    \item $\frac{1}{321861540} = \frac{1}{321861540}$
\end{itemize}
La somme des numérateurs est :
$$ \frac{1}{60} + \frac{1}{17941} + \frac{1}{321861540} = \frac{5364359 + 17940 + 1}{321861540} = \frac{5382300}{321861540} $$
Le PGCD du numérateur et du dénominateur est $\text{PGCD}(5382300, 321861540) = 1076460$.
La fraction irréductible est :
$$ \frac{5382300}{321861540} = \frac{5382300 \div 1076460}{321861540 \div 1076460} = \frac{5}{299} $$
Cette fraction est bien égale à $\frac{5}{299}$.

\subsection{Démonstration pour $n = 300$}
Soit $n = 300$. Nous cherchons $x, y, z \in \mathbb{Z}^{+}$ tels que $\frac{5}{300} = \frac{1}{x} + \frac{1}{y} + \frac{1}{z}$.
Posons $x = 61$, $y = 3661$, $z = 13399260$.
Les conditions $x > 0$, $y > 0$ et $z > 0$ sont satisfaites car $x, y, z \geq 1$.
Le PPCM des dénominateurs est $\text{PPCM}(61, 3661, 13399260) = 13399260$.
En réduisant au même dénominateur :
\begin{itemize}
    \item $\frac{1}{61} = \frac{219660}{13399260}$
    \item $\frac{1}{3661} = \frac{3660}{13399260}$
    \item $\frac{1}{13399260} = \frac{1}{13399260}$
\end{itemize}
La somme des numérateurs est :
$$ \frac{1}{61} + \frac{1}{3661} + \frac{1}{13399260} = \frac{219660 + 3660 + 1}{13399260} = \frac{223321}{13399260} $$
Le PGCD du numérateur et du dénominateur est $\text{PGCD}(223321, 13399260) = 223321$.
La fraction irréductible est :
$$ \frac{223321}{13399260} = \frac{223321 \div 223321}{13399260 \div 223321} = \frac{1}{60} $$
Cette fraction est bien égale à $\frac{5}{300}$.

\subsection{Démonstration pour $n = 301$}
Soit $n = 301$. Nous cherchons $x, y, z \in \mathbb{Z}^{+}$ tels que $\frac{5}{301} = \frac{1}{x} + \frac{1}{y} + \frac{1}{z}$.
Posons $x = 61$, $y = 4592$, $z = 12044816$.
Les conditions $x > 0$, $y > 0$ et $z > 0$ sont satisfaites car $x, y, z \geq 1$.
Le PPCM des dénominateurs est $\text{PPCM}(61, 4592, 12044816) = 12044816$.
En réduisant au même dénominateur :
\begin{itemize}
    \item $\frac{1}{61} = \frac{197456}{12044816}$
    \item $\frac{1}{4592} = \frac{2623}{12044816}$
    \item $\frac{1}{12044816} = \frac{1}{12044816}$
\end{itemize}
La somme des numérateurs est :
$$ \frac{1}{61} + \frac{1}{4592} + \frac{1}{12044816} = \frac{197456 + 2623 + 1}{12044816} = \frac{200080}{12044816} $$
Le PGCD du numérateur et du dénominateur est $\text{PGCD}(200080, 12044816) = 40016$.
La fraction irréductible est :
$$ \frac{200080}{12044816} = \frac{200080 \div 40016}{12044816 \div 40016} = \frac{5}{301} $$
Cette fraction est bien égale à $\frac{5}{301}$.

\subsection{Démonstration pour $n = 302$}
Soit $n = 302$. Nous cherchons $x, y, z \in \mathbb{Z}^{+}$ tels que $\frac{5}{302} = \frac{1}{x} + \frac{1}{y} + \frac{1}{z}$.
Posons $x = 61$, $y = 6141$, $z = 113129502$.
Les conditions $x > 0$, $y > 0$ et $z > 0$ sont satisfaites car $x, y, z \geq 1$.
Le PPCM des dénominateurs est $\text{PPCM}(61, 6141, 113129502) = 113129502$.
En réduisant au même dénominateur :
\begin{itemize}
    \item $\frac{1}{61} = \frac{1854582}{113129502}$
    \item $\frac{1}{6141} = \frac{18422}{113129502}$
    \item $\frac{1}{113129502} = \frac{1}{113129502}$
\end{itemize}
La somme des numérateurs est :
$$ \frac{1}{61} + \frac{1}{6141} + \frac{1}{113129502} = \frac{1854582 + 18422 + 1}{113129502} = \frac{1873005}{113129502} $$
Le PGCD du numérateur et du dénominateur est $\text{PGCD}(1873005, 113129502) = 374601$.
La fraction irréductible est :
$$ \frac{1873005}{113129502} = \frac{1873005 \div 374601}{113129502 \div 374601} = \frac{5}{302} $$
Cette fraction est bien égale à $\frac{5}{302}$.

\subsection{Démonstration pour $n = 303$}
Soit $n = 303$. Nous cherchons $x, y, z \in \mathbb{Z}^{+}$ tels que $\frac{5}{303} = \frac{1}{x} + \frac{1}{y} + \frac{1}{z}$.
Posons $x = 61$, $y = 9242$, $z = 170819886$.
Les conditions $x > 0$, $y > 0$ et $z > 0$ sont satisfaites car $x, y, z \geq 1$.
Le PPCM des dénominateurs est $\text{PPCM}(61, 9242, 170819886) = 170819886$.
En réduisant au même dénominateur :
\begin{itemize}
    \item $\frac{1}{61} = \frac{2800326}{170819886}$
    \item $\frac{1}{9242} = \frac{18483}{170819886}$
    \item $\frac{1}{170819886} = \frac{1}{170819886}$
\end{itemize}
La somme des numérateurs est :
$$ \frac{1}{61} + \frac{1}{9242} + \frac{1}{170819886} = \frac{2800326 + 18483 + 1}{170819886} = \frac{2818810}{170819886} $$
Le PGCD du numérateur et du dénominateur est $\text{PGCD}(2818810, 170819886) = 563762$.
La fraction irréductible est :
$$ \frac{2818810}{170819886} = \frac{2818810 \div 563762}{170819886 \div 563762} = \frac{5}{303} $$
Cette fraction est bien égale à $\frac{5}{303}$.

\subsection{Démonstration pour $n = 304$}
Soit $n = 304$. Nous cherchons $x, y, z \in \mathbb{Z}^{+}$ tels que $\frac{5}{304} = \frac{1}{x} + \frac{1}{y} + \frac{1}{z}$.
Posons $x = 61$, $y = 18545$, $z = 343898480$.
Les conditions $x > 0$, $y > 0$ et $z > 0$ sont satisfaites car $x, y, z \geq 1$.
Le PPCM des dénominateurs est $\text{PPCM}(61, 18545, 343898480) = 343898480$.
En réduisant au même dénominateur :
\begin{itemize}
    \item $\frac{1}{61} = \frac{5637680}{343898480}$
    \item $\frac{1}{18545} = \frac{18544}{343898480}$
    \item $\frac{1}{343898480} = \frac{1}{343898480}$
\end{itemize}
La somme des numérateurs est :
$$ \frac{1}{61} + \frac{1}{18545} + \frac{1}{343898480} = \frac{5637680 + 18544 + 1}{343898480} = \frac{5656225}{343898480} $$
Le PGCD du numérateur et du dénominateur est $\text{PGCD}(5656225, 343898480) = 1131245$.
La fraction irréductible est :
$$ \frac{5656225}{343898480} = \frac{5656225 \div 1131245}{343898480 \div 1131245} = \frac{5}{304} $$
Cette fraction est bien égale à $\frac{5}{304}$.

\subsection{Démonstration pour $n = 305$}
Soit $n = 305$. Nous cherchons $x, y, z \in \mathbb{Z}^{+}$ tels que $\frac{5}{305} = \frac{1}{x} + \frac{1}{y} + \frac{1}{z}$.
Posons $x = 62$, $y = 3783$, $z = 14307306$.
Les conditions $x > 0$, $y > 0$ et $z > 0$ sont satisfaites car $x, y, z \geq 1$.
Le PPCM des dénominateurs est $\text{PPCM}(62, 3783, 14307306) = 14307306$.
En réduisant au même dénominateur :
\begin{itemize}
    \item $\frac{1}{62} = \frac{230763}{14307306}$
    \item $\frac{1}{3783} = \frac{3782}{14307306}$
    \item $\frac{1}{14307306} = \frac{1}{14307306}$
\end{itemize}
La somme des numérateurs est :
$$ \frac{1}{62} + \frac{1}{3783} + \frac{1}{14307306} = \frac{230763 + 3782 + 1}{14307306} = \frac{234546}{14307306} $$
Le PGCD du numérateur et du dénominateur est $\text{PGCD}(234546, 14307306) = 234546$.
La fraction irréductible est :
$$ \frac{234546}{14307306} = \frac{234546 \div 234546}{14307306 \div 234546} = \frac{1}{61} $$
Cette fraction est bien égale à $\frac{5}{305}$.

\subsection{Démonstration pour $n = 306$}
Soit $n = 306$. Nous cherchons $x, y, z \in \mathbb{Z}^{+}$ tels que $\frac{5}{306} = \frac{1}{x} + \frac{1}{y} + \frac{1}{z}$.
Posons $x = 62$, $y = 4744$, $z = 22500792$.
Les conditions $x > 0$, $y > 0$ et $z > 0$ sont satisfaites car $x, y, z \geq 1$.
Le PPCM des dénominateurs est $\text{PPCM}(62, 4744, 22500792) = 22500792$.
En réduisant au même dénominateur :
\begin{itemize}
    \item $\frac{1}{62} = \frac{362916}{22500792}$
    \item $\frac{1}{4744} = \frac{4743}{22500792}$
    \item $\frac{1}{22500792} = \frac{1}{22500792}$
\end{itemize}
La somme des numérateurs est :
$$ \frac{1}{62} + \frac{1}{4744} + \frac{1}{22500792} = \frac{362916 + 4743 + 1}{22500792} = \frac{367660}{22500792} $$
Le PGCD du numérateur et du dénominateur est $\text{PGCD}(367660, 22500792) = 73532$.
La fraction irréductible est :
$$ \frac{367660}{22500792} = \frac{367660 \div 73532}{22500792 \div 73532} = \frac{5}{306} $$
Cette fraction est bien égale à $\frac{5}{306}$.

\subsection{Démonstration pour $n = 307$}
Soit $n = 307$. Nous cherchons $x, y, z \in \mathbb{Z}^{+}$ tels que $\frac{5}{307} = \frac{1}{x} + \frac{1}{y} + \frac{1}{z}$.
Posons $x = 62$, $y = 6345$, $z = 120770730$.
Les conditions $x > 0$, $y > 0$ et $z > 0$ sont satisfaites car $x, y, z \geq 1$.
Le PPCM des dénominateurs est $\text{PPCM}(62, 6345, 120770730) = 120770730$.
En réduisant au même dénominateur :
\begin{itemize}
    \item $\frac{1}{62} = \frac{1947915}{120770730}$
    \item $\frac{1}{6345} = \frac{19034}{120770730}$
    \item $\frac{1}{120770730} = \frac{1}{120770730}$
\end{itemize}
La somme des numérateurs est :
$$ \frac{1}{62} + \frac{1}{6345} + \frac{1}{120770730} = \frac{1947915 + 19034 + 1}{120770730} = \frac{1966950}{120770730} $$
Le PGCD du numérateur et du dénominateur est $\text{PGCD}(1966950, 120770730) = 393390$.
La fraction irréductible est :
$$ \frac{1966950}{120770730} = \frac{1966950 \div 393390}{120770730 \div 393390} = \frac{5}{307} $$
Cette fraction est bien égale à $\frac{5}{307}$.

\subsection{Démonstration pour $n = 308$}
Soit $n = 308$. Nous cherchons $x, y, z \in \mathbb{Z}^{+}$ tels que $\frac{5}{308} = \frac{1}{x} + \frac{1}{y} + \frac{1}{z}$.
Posons $x = 62$, $y = 9549$, $z = 91173852$.
Les conditions $x > 0$, $y > 0$ et $z > 0$ sont satisfaites car $x, y, z \geq 1$.
Le PPCM des dénominateurs est $\text{PPCM}(62, 9549, 91173852) = 91173852$.
En réduisant au même dénominateur :
\begin{itemize}
    \item $\frac{1}{62} = \frac{1470546}{91173852}$
    \item $\frac{1}{9549} = \frac{9548}{91173852}$
    \item $\frac{1}{91173852} = \frac{1}{91173852}$
\end{itemize}
La somme des numérateurs est :
$$ \frac{1}{62} + \frac{1}{9549} + \frac{1}{91173852} = \frac{1470546 + 9548 + 1}{91173852} = \frac{1480095}{91173852} $$
Le PGCD du numérateur et du dénominateur est $\text{PGCD}(1480095, 91173852) = 296019$.
La fraction irréductible est :
$$ \frac{1480095}{91173852} = \frac{1480095 \div 296019}{91173852 \div 296019} = \frac{5}{308} $$
Cette fraction est bien égale à $\frac{5}{308}$.

\subsection{Démonstration pour $n = 309$}
Soit $n = 309$. Nous cherchons $x, y, z \in \mathbb{Z}^{+}$ tels que $\frac{5}{309} = \frac{1}{x} + \frac{1}{y} + \frac{1}{z}$.
Posons $x = 62$, $y = 19159$, $z = 367048122$.
Les conditions $x > 0$, $y > 0$ et $z > 0$ sont satisfaites car $x, y, z \geq 1$.
Le PPCM des dénominateurs est $\text{PPCM}(62, 19159, 367048122) = 367048122$.
En réduisant au même dénominateur :
\begin{itemize}
    \item $\frac{1}{62} = \frac{5920131}{367048122}$
    \item $\frac{1}{19159} = \frac{19158}{367048122}$
    \item $\frac{1}{367048122} = \frac{1}{367048122}$
\end{itemize}
La somme des numérateurs est :
$$ \frac{1}{62} + \frac{1}{19159} + \frac{1}{367048122} = \frac{5920131 + 19158 + 1}{367048122} = \frac{5939290}{367048122} $$
Le PGCD du numérateur et du dénominateur est $\text{PGCD}(5939290, 367048122) = 1187858$.
La fraction irréductible est :
$$ \frac{5939290}{367048122} = \frac{5939290 \div 1187858}{367048122 \div 1187858} = \frac{5}{309} $$
Cette fraction est bien égale à $\frac{5}{309}$.

\subsection{Démonstration pour $n = 310$}
Soit $n = 310$. Nous cherchons $x, y, z \in \mathbb{Z}^{+}$ tels que $\frac{5}{310} = \frac{1}{x} + \frac{1}{y} + \frac{1}{z}$.
Posons $x = 63$, $y = 3907$, $z = 15260742$.
Les conditions $x > 0$, $y > 0$ et $z > 0$ sont satisfaites car $x, y, z \geq 1$.
Le PPCM des dénominateurs est $\text{PPCM}(63, 3907, 15260742) = 15260742$.
En réduisant au même dénominateur :
\begin{itemize}
    \item $\frac{1}{63} = \frac{242234}{15260742}$
    \item $\frac{1}{3907} = \frac{3906}{15260742}$
    \item $\frac{1}{15260742} = \frac{1}{15260742}$
\end{itemize}
La somme des numérateurs est :
$$ \frac{1}{63} + \frac{1}{3907} + \frac{1}{15260742} = \frac{242234 + 3906 + 1}{15260742} = \frac{246141}{15260742} $$
Le PGCD du numérateur et du dénominateur est $\text{PGCD}(246141, 15260742) = 246141$.
La fraction irréductible est :
$$ \frac{246141}{15260742} = \frac{246141 \div 246141}{15260742 \div 246141} = \frac{1}{62} $$
Cette fraction est bien égale à $\frac{5}{310}$.

\subsection{Démonstration pour $n = 311$}
Soit $n = 311$. Nous cherchons $x, y, z \in \mathbb{Z}^{+}$ tels que $\frac{5}{311} = \frac{1}{x} + \frac{1}{y} + \frac{1}{z}$.
Posons $x = 63$, $y = 4899$, $z = 31995369$.
Les conditions $x > 0$, $y > 0$ et $z > 0$ sont satisfaites car $x, y, z \geq 1$.
Le PPCM des dénominateurs est $\text{PPCM}(63, 4899, 31995369) = 31995369$.
En réduisant au même dénominateur :
\begin{itemize}
    \item $\frac{1}{63} = \frac{507863}{31995369}$
    \item $\frac{1}{4899} = \frac{6531}{31995369}$
    \item $\frac{1}{31995369} = \frac{1}{31995369}$
\end{itemize}
La somme des numérateurs est :
$$ \frac{1}{63} + \frac{1}{4899} + \frac{1}{31995369} = \frac{507863 + 6531 + 1}{31995369} = \frac{514395}{31995369} $$
Le PGCD du numérateur et du dénominateur est $\text{PGCD}(514395, 31995369) = 102879$.
La fraction irréductible est :
$$ \frac{514395}{31995369} = \frac{514395 \div 102879}{31995369 \div 102879} = \frac{5}{311} $$
Cette fraction est bien égale à $\frac{5}{311}$.

\subsection{Démonstration pour $n = 312$}
Soit $n = 312$. Nous cherchons $x, y, z \in \mathbb{Z}^{+}$ tels que $\frac{5}{312} = \frac{1}{x} + \frac{1}{y} + \frac{1}{z}$.
Posons $x = 63$, $y = 6553$, $z = 42935256$.
Les conditions $x > 0$, $y > 0$ et $z > 0$ sont satisfaites car $x, y, z \geq 1$.
Le PPCM des dénominateurs est $\text{PPCM}(63, 6553, 42935256) = 42935256$.
En réduisant au même dénominateur :
\begin{itemize}
    \item $\frac{1}{63} = \frac{681512}{42935256}$
    \item $\frac{1}{6553} = \frac{6552}{42935256}$
    \item $\frac{1}{42935256} = \frac{1}{42935256}$
\end{itemize}
La somme des numérateurs est :
$$ \frac{1}{63} + \frac{1}{6553} + \frac{1}{42935256} = \frac{681512 + 6552 + 1}{42935256} = \frac{688065}{42935256} $$
Le PGCD du numérateur et du dénominateur est $\text{PGCD}(688065, 42935256) = 137613$.
La fraction irréductible est :
$$ \frac{688065}{42935256} = \frac{688065 \div 137613}{42935256 \div 137613} = \frac{5}{312} $$
Cette fraction est bien égale à $\frac{5}{312}$.

\subsection{Démonstration pour $n = 313$}
Soit $n = 313$. Nous cherchons $x, y, z \in \mathbb{Z}^{+}$ tels que $\frac{5}{313} = \frac{1}{x} + \frac{1}{y} + \frac{1}{z}$.
Posons $x = 63$, $y = 9860$, $z = 194429340$.
Les conditions $x > 0$, $y > 0$ et $z > 0$ sont satisfaites car $x, y, z \geq 1$.
Le PPCM des dénominateurs est $\text{PPCM}(63, 9860, 194429340) = 194429340$.
En réduisant au même dénominateur :
\begin{itemize}
    \item $\frac{1}{63} = \frac{3086180}{194429340}$
    \item $\frac{1}{9860} = \frac{19719}{194429340}$
    \item $\frac{1}{194429340} = \frac{1}{194429340}$
\end{itemize}
La somme des numérateurs est :
$$ \frac{1}{63} + \frac{1}{9860} + \frac{1}{194429340} = \frac{3086180 + 19719 + 1}{194429340} = \frac{3105900}{194429340} $$
Le PGCD du numérateur et du dénominateur est $\text{PGCD}(3105900, 194429340) = 621180$.
La fraction irréductible est :
$$ \frac{3105900}{194429340} = \frac{3105900 \div 621180}{194429340 \div 621180} = \frac{5}{313} $$
Cette fraction est bien égale à $\frac{5}{313}$.

\subsection{Démonstration pour $n = 314$}
Soit $n = 314$. Nous cherchons $x, y, z \in \mathbb{Z}^{+}$ tels que $\frac{5}{314} = \frac{1}{x} + \frac{1}{y} + \frac{1}{z}$.
Posons $x = 63$, $y = 19783$, $z = 391347306$.
Les conditions $x > 0$, $y > 0$ et $z > 0$ sont satisfaites car $x, y, z \geq 1$.
Le PPCM des dénominateurs est $\text{PPCM}(63, 19783, 391347306) = 391347306$.
En réduisant au même dénominateur :
\begin{itemize}
    \item $\frac{1}{63} = \frac{6211862}{391347306}$
    \item $\frac{1}{19783} = \frac{19782}{391347306}$
    \item $\frac{1}{391347306} = \frac{1}{391347306}$
\end{itemize}
La somme des numérateurs est :
$$ \frac{1}{63} + \frac{1}{19783} + \frac{1}{391347306} = \frac{6211862 + 19782 + 1}{391347306} = \frac{6231645}{391347306} $$
Le PGCD du numérateur et du dénominateur est $\text{PGCD}(6231645, 391347306) = 1246329$.
La fraction irréductible est :
$$ \frac{6231645}{391347306} = \frac{6231645 \div 1246329}{391347306 \div 1246329} = \frac{5}{314} $$
Cette fraction est bien égale à $\frac{5}{314}$.

\subsection{Démonstration pour $n = 315$}
Soit $n = 315$. Nous cherchons $x, y, z \in \mathbb{Z}^{+}$ tels que $\frac{5}{315} = \frac{1}{x} + \frac{1}{y} + \frac{1}{z}$.
Posons $x = 64$, $y = 4033$, $z = 16261056$.
Les conditions $x > 0$, $y > 0$ et $z > 0$ sont satisfaites car $x, y, z \geq 1$.
Le PPCM des dénominateurs est $\text{PPCM}(64, 4033, 16261056) = 16261056$.
En réduisant au même dénominateur :
\begin{itemize}
    \item $\frac{1}{64} = \frac{254079}{16261056}$
    \item $\frac{1}{4033} = \frac{4032}{16261056}$
    \item $\frac{1}{16261056} = \frac{1}{16261056}$
\end{itemize}
La somme des numérateurs est :
$$ \frac{1}{64} + \frac{1}{4033} + \frac{1}{16261056} = \frac{254079 + 4032 + 1}{16261056} = \frac{258112}{16261056} $$
Le PGCD du numérateur et du dénominateur est $\text{PGCD}(258112, 16261056) = 258112$.
La fraction irréductible est :
$$ \frac{258112}{16261056} = \frac{258112 \div 258112}{16261056 \div 258112} = \frac{1}{63} $$
Cette fraction est bien égale à $\frac{5}{315}$.

\subsection{Démonstration pour $n = 316$}
Soit $n = 316$. Nous cherchons $x, y, z \in \mathbb{Z}^{+}$ tels que $\frac{5}{316} = \frac{1}{x} + \frac{1}{y} + \frac{1}{z}$.
Posons $x = 64$, $y = 5057$, $z = 25568192$.
Les conditions $x > 0$, $y > 0$ et $z > 0$ sont satisfaites car $x, y, z \geq 1$.
Le PPCM des dénominateurs est $\text{PPCM}(64, 5057, 25568192) = 25568192$.
En réduisant au même dénominateur :
\begin{itemize}
    \item $\frac{1}{64} = \frac{399503}{25568192}$
    \item $\frac{1}{5057} = \frac{5056}{25568192}$
    \item $\frac{1}{25568192} = \frac{1}{25568192}$
\end{itemize}
La somme des numérateurs est :
$$ \frac{1}{64} + \frac{1}{5057} + \frac{1}{25568192} = \frac{399503 + 5056 + 1}{25568192} = \frac{404560}{25568192} $$
Le PGCD du numérateur et du dénominateur est $\text{PGCD}(404560, 25568192) = 80912$.
La fraction irréductible est :
$$ \frac{404560}{25568192} = \frac{404560 \div 80912}{25568192 \div 80912} = \frac{5}{316} $$
Cette fraction est bien égale à $\frac{5}{316}$.

\subsection{Démonstration pour $n = 317$}
Soit $n = 317$. Nous cherchons $x, y, z \in \mathbb{Z}^{+}$ tels que $\frac{5}{317} = \frac{1}{x} + \frac{1}{y} + \frac{1}{z}$.
Posons $x = 64$, $y = 6763$, $z = 137207744$.
Les conditions $x > 0$, $y > 0$ et $z > 0$ sont satisfaites car $x, y, z \geq 1$.
Le PPCM des dénominateurs est $\text{PPCM}(64, 6763, 137207744) = 137207744$.
En réduisant au même dénominateur :
\begin{itemize}
    \item $\frac{1}{64} = \frac{2143871}{137207744}$
    \item $\frac{1}{6763} = \frac{20288}{137207744}$
    \item $\frac{1}{137207744} = \frac{1}{137207744}$
\end{itemize}
La somme des numérateurs est :
$$ \frac{1}{64} + \frac{1}{6763} + \frac{1}{137207744} = \frac{2143871 + 20288 + 1}{137207744} = \frac{2164160}{137207744} $$
Le PGCD du numérateur et du dénominateur est $\text{PGCD}(2164160, 137207744) = 432832$.
La fraction irréductible est :
$$ \frac{2164160}{137207744} = \frac{2164160 \div 432832}{137207744 \div 432832} = \frac{5}{317} $$
Cette fraction est bien égale à $\frac{5}{317}$.

\subsection{Démonstration pour $n = 318$}
Soit $n = 318$. Nous cherchons $x, y, z \in \mathbb{Z}^{+}$ tels que $\frac{5}{318} = \frac{1}{x} + \frac{1}{y} + \frac{1}{z}$.
Posons $x = 64$, $y = 10177$, $z = 103561152$.
Les conditions $x > 0$, $y > 0$ et $z > 0$ sont satisfaites car $x, y, z \geq 1$.
Le PPCM des dénominateurs est $\text{PPCM}(64, 10177, 103561152) = 103561152$.
En réduisant au même dénominateur :
\begin{itemize}
    \item $\frac{1}{64} = \frac{1618143}{103561152}$
    \item $\frac{1}{10177} = \frac{10176}{103561152}$
    \item $\frac{1}{103561152} = \frac{1}{103561152}$
\end{itemize}
La somme des numérateurs est :
$$ \frac{1}{64} + \frac{1}{10177} + \frac{1}{103561152} = \frac{1618143 + 10176 + 1}{103561152} = \frac{1628320}{103561152} $$
Le PGCD du numérateur et du dénominateur est $\text{PGCD}(1628320, 103561152) = 325664$.
La fraction irréductible est :
$$ \frac{1628320}{103561152} = \frac{1628320 \div 325664}{103561152 \div 325664} = \frac{5}{318} $$
Cette fraction est bien égale à $\frac{5}{318}$.

\subsection{Démonstration pour $n = 319$}
Soit $n = 319$. Nous cherchons $x, y, z \in \mathbb{Z}^{+}$ tels que $\frac{5}{319} = \frac{1}{x} + \frac{1}{y} + \frac{1}{z}$.
Posons $x = 64$, $y = 20417$, $z = 416833472$.
Les conditions $x > 0$, $y > 0$ et $z > 0$ sont satisfaites car $x, y, z \geq 1$.
Le PPCM des dénominateurs est $\text{PPCM}(64, 20417, 416833472) = 416833472$.
En réduisant au même dénominateur :
\begin{itemize}
    \item $\frac{1}{64} = \frac{6513023}{416833472}$
    \item $\frac{1}{20417} = \frac{20416}{416833472}$
    \item $\frac{1}{416833472} = \frac{1}{416833472}$
\end{itemize}
La somme des numérateurs est :
$$ \frac{1}{64} + \frac{1}{20417} + \frac{1}{416833472} = \frac{6513023 + 20416 + 1}{416833472} = \frac{6533440}{416833472} $$
Le PGCD du numérateur et du dénominateur est $\text{PGCD}(6533440, 416833472) = 1306688$.
La fraction irréductible est :
$$ \frac{6533440}{416833472} = \frac{6533440 \div 1306688}{416833472 \div 1306688} = \frac{5}{319} $$
Cette fraction est bien égale à $\frac{5}{319}$.

\subsection{Démonstration pour $n = 320$}
Soit $n = 320$. Nous cherchons $x, y, z \in \mathbb{Z}^{+}$ tels que $\frac{5}{320} = \frac{1}{x} + \frac{1}{y} + \frac{1}{z}$.
Posons $x = 65$, $y = 4161$, $z = 17309760$.
Les conditions $x > 0$, $y > 0$ et $z > 0$ sont satisfaites car $x, y, z \geq 1$.
Le PPCM des dénominateurs est $\text{PPCM}(65, 4161, 17309760) = 17309760$.
En réduisant au même dénominateur :
\begin{itemize}
    \item $\frac{1}{65} = \frac{266304}{17309760}$
    \item $\frac{1}{4161} = \frac{4160}{17309760}$
    \item $\frac{1}{17309760} = \frac{1}{17309760}$
\end{itemize}
La somme des numérateurs est :
$$ \frac{1}{65} + \frac{1}{4161} + \frac{1}{17309760} = \frac{266304 + 4160 + 1}{17309760} = \frac{270465}{17309760} $$
Le PGCD du numérateur et du dénominateur est $\text{PGCD}(270465, 17309760) = 270465$.
La fraction irréductible est :
$$ \frac{270465}{17309760} = \frac{270465 \div 270465}{17309760 \div 270465} = \frac{1}{64} $$
Cette fraction est bien égale à $\frac{5}{320}$.

\subsection{Démonstration pour $n = 321$}
Soit $n = 321$. Nous cherchons $x, y, z \in \mathbb{Z}^{+}$ tels que $\frac{5}{321} = \frac{1}{x} + \frac{1}{y} + \frac{1}{z}$.
Posons $x = 65$, $y = 5217$, $z = 36284235$.
Les conditions $x > 0$, $y > 0$ et $z > 0$ sont satisfaites car $x, y, z \geq 1$.
Le PPCM des dénominateurs est $\text{PPCM}(65, 5217, 36284235) = 36284235$.
En réduisant au même dénominateur :
\begin{itemize}
    \item $\frac{1}{65} = \frac{558219}{36284235}$
    \item $\frac{1}{5217} = \frac{6955}{36284235}$
    \item $\frac{1}{36284235} = \frac{1}{36284235}$
\end{itemize}
La somme des numérateurs est :
$$ \frac{1}{65} + \frac{1}{5217} + \frac{1}{36284235} = \frac{558219 + 6955 + 1}{36284235} = \frac{565175}{36284235} $$
Le PGCD du numérateur et du dénominateur est $\text{PGCD}(565175, 36284235) = 113035$.
La fraction irréductible est :
$$ \frac{565175}{36284235} = \frac{565175 \div 113035}{36284235 \div 113035} = \frac{5}{321} $$
Cette fraction est bien égale à $\frac{5}{321}$.

\subsection{Démonstration pour $n = 322$}
Soit $n = 322$. Nous cherchons $x, y, z \in \mathbb{Z}^{+}$ tels que $\frac{5}{322} = \frac{1}{x} + \frac{1}{y} + \frac{1}{z}$.
Posons $x = 65$, $y = 6977$, $z = 146028610$.
Les conditions $x > 0$, $y > 0$ et $z > 0$ sont satisfaites car $x, y, z \geq 1$.
Le PPCM des dénominateurs est $\text{PPCM}(65, 6977, 146028610) = 146028610$.
En réduisant au même dénominateur :
\begin{itemize}
    \item $\frac{1}{65} = \frac{2246594}{146028610}$
    \item $\frac{1}{6977} = \frac{20930}{146028610}$
    \item $\frac{1}{146028610} = \frac{1}{146028610}$
\end{itemize}
La somme des numérateurs est :
$$ \frac{1}{65} + \frac{1}{6977} + \frac{1}{146028610} = \frac{2246594 + 20930 + 1}{146028610} = \frac{2267525}{146028610} $$
Le PGCD du numérateur et du dénominateur est $\text{PGCD}(2267525, 146028610) = 453505$.
La fraction irréductible est :
$$ \frac{2267525}{146028610} = \frac{2267525 \div 453505}{146028610 \div 453505} = \frac{5}{322} $$
Cette fraction est bien égale à $\frac{5}{322}$.

\subsection{Démonstration pour $n = 323$}
Soit $n = 323$. Nous cherchons $x, y, z \in \mathbb{Z}^{+}$ tels que $\frac{5}{323} = \frac{1}{x} + \frac{1}{y} + \frac{1}{z}$.
Posons $x = 65$, $y = 10498$, $z = 220405510$.
Les conditions $x > 0$, $y > 0$ et $z > 0$ sont satisfaites car $x, y, z \geq 1$.
Le PPCM des dénominateurs est $\text{PPCM}(65, 10498, 220405510) = 220405510$.
En réduisant au même dénominateur :
\begin{itemize}
    \item $\frac{1}{65} = \frac{3390854}{220405510}$
    \item $\frac{1}{10498} = \frac{20995}{220405510}$
    \item $\frac{1}{220405510} = \frac{1}{220405510}$
\end{itemize}
La somme des numérateurs est :
$$ \frac{1}{65} + \frac{1}{10498} + \frac{1}{220405510} = \frac{3390854 + 20995 + 1}{220405510} = \frac{3411850}{220405510} $$
Le PGCD du numérateur et du dénominateur est $\text{PGCD}(3411850, 220405510) = 682370$.
La fraction irréductible est :
$$ \frac{3411850}{220405510} = \frac{3411850 \div 682370}{220405510 \div 682370} = \frac{5}{323} $$
Cette fraction est bien égale à $\frac{5}{323}$.

\subsection{Démonstration pour $n = 324$}
Soit $n = 324$. Nous cherchons $x, y, z \in \mathbb{Z}^{+}$ tels que $\frac{5}{324} = \frac{1}{x} + \frac{1}{y} + \frac{1}{z}$.
Posons $x = 65$, $y = 21061$, $z = 443544660$.
Les conditions $x > 0$, $y > 0$ et $z > 0$ sont satisfaites car $x, y, z \geq 1$.
Le PPCM des dénominateurs est $\text{PPCM}(65, 21061, 443544660) = 443544660$.
En réduisant au même dénominateur :
\begin{itemize}
    \item $\frac{1}{65} = \frac{6823764}{443544660}$
    \item $\frac{1}{21061} = \frac{21060}{443544660}$
    \item $\frac{1}{443544660} = \frac{1}{443544660}$
\end{itemize}
La somme des numérateurs est :
$$ \frac{1}{65} + \frac{1}{21061} + \frac{1}{443544660} = \frac{6823764 + 21060 + 1}{443544660} = \frac{6844825}{443544660} $$
Le PGCD du numérateur et du dénominateur est $\text{PGCD}(6844825, 443544660) = 1368965$.
La fraction irréductible est :
$$ \frac{6844825}{443544660} = \frac{6844825 \div 1368965}{443544660 \div 1368965} = \frac{5}{324} $$
Cette fraction est bien égale à $\frac{5}{324}$.

\subsection{Démonstration pour $n = 325$}
Soit $n = 325$. Nous cherchons $x, y, z \in \mathbb{Z}^{+}$ tels que $\frac{5}{325} = \frac{1}{x} + \frac{1}{y} + \frac{1}{z}$.
Posons $x = 66$, $y = 4291$, $z = 18408390$.
Les conditions $x > 0$, $y > 0$ et $z > 0$ sont satisfaites car $x, y, z \geq 1$.
Le PPCM des dénominateurs est $\text{PPCM}(66, 4291, 18408390) = 18408390$.
En réduisant au même dénominateur :
\begin{itemize}
    \item $\frac{1}{66} = \frac{278915}{18408390}$
    \item $\frac{1}{4291} = \frac{4290}{18408390}$
    \item $\frac{1}{18408390} = \frac{1}{18408390}$
\end{itemize}
La somme des numérateurs est :
$$ \frac{1}{66} + \frac{1}{4291} + \frac{1}{18408390} = \frac{278915 + 4290 + 1}{18408390} = \frac{283206}{18408390} $$
Le PGCD du numérateur et du dénominateur est $\text{PGCD}(283206, 18408390) = 283206$.
La fraction irréductible est :
$$ \frac{283206}{18408390} = \frac{283206 \div 283206}{18408390 \div 283206} = \frac{1}{65} $$
Cette fraction est bien égale à $\frac{5}{325}$.

\subsection{Démonstration pour $n = 326$}
Soit $n = 326$. Nous cherchons $x, y, z \in \mathbb{Z}^{+}$ tels que $\frac{5}{326} = \frac{1}{x} + \frac{1}{y} + \frac{1}{z}$.
Posons $x = 66$, $y = 5380$, $z = 28939020$.
Les conditions $x > 0$, $y > 0$ et $z > 0$ sont satisfaites car $x, y, z \geq 1$.
Le PPCM des dénominateurs est $\text{PPCM}(66, 5380, 28939020) = 28939020$.
En réduisant au même dénominateur :
\begin{itemize}
    \item $\frac{1}{66} = \frac{438470}{28939020}$
    \item $\frac{1}{5380} = \frac{5379}{28939020}$
    \item $\frac{1}{28939020} = \frac{1}{28939020}$
\end{itemize}
La somme des numérateurs est :
$$ \frac{1}{66} + \frac{1}{5380} + \frac{1}{28939020} = \frac{438470 + 5379 + 1}{28939020} = \frac{443850}{28939020} $$
Le PGCD du numérateur et du dénominateur est $\text{PGCD}(443850, 28939020) = 88770$.
La fraction irréductible est :
$$ \frac{443850}{28939020} = \frac{443850 \div 88770}{28939020 \div 88770} = \frac{5}{326} $$
Cette fraction est bien égale à $\frac{5}{326}$.

\subsection{Démonstration pour $n = 327$}
Soit $n = 327$. Nous cherchons $x, y, z \in \mathbb{Z}^{+}$ tels que $\frac{5}{327} = \frac{1}{x} + \frac{1}{y} + \frac{1}{z}$.
Posons $x = 66$, $y = 7195$, $z = 51760830$.
Les conditions $x > 0$, $y > 0$ et $z > 0$ sont satisfaites car $x, y, z \geq 1$.
Le PPCM des dénominateurs est $\text{PPCM}(66, 7195, 51760830) = 51760830$.
En réduisant au même dénominateur :
\begin{itemize}
    \item $\frac{1}{66} = \frac{784255}{51760830}$
    \item $\frac{1}{7195} = \frac{7194}{51760830}$
    \item $\frac{1}{51760830} = \frac{1}{51760830}$
\end{itemize}
La somme des numérateurs est :
$$ \frac{1}{66} + \frac{1}{7195} + \frac{1}{51760830} = \frac{784255 + 7194 + 1}{51760830} = \frac{791450}{51760830} $$
Le PGCD du numérateur et du dénominateur est $\text{PGCD}(791450, 51760830) = 158290$.
La fraction irréductible est :
$$ \frac{791450}{51760830} = \frac{791450 \div 158290}{51760830 \div 158290} = \frac{5}{327} $$
Cette fraction est bien égale à $\frac{5}{327}$.

\subsection{Démonstration pour $n = 328$}
Soit $n = 328$. Nous cherchons $x, y, z \in \mathbb{Z}^{+}$ tels que $\frac{5}{328} = \frac{1}{x} + \frac{1}{y} + \frac{1}{z}$.
Posons $x = 66$, $y = 10825$, $z = 117169800$.
Les conditions $x > 0$, $y > 0$ et $z > 0$ sont satisfaites car $x, y, z \geq 1$.
Le PPCM des dénominateurs est $\text{PPCM}(66, 10825, 117169800) = 117169800$.
En réduisant au même dénominateur :
\begin{itemize}
    \item $\frac{1}{66} = \frac{1775300}{117169800}$
    \item $\frac{1}{10825} = \frac{10824}{117169800}$
    \item $\frac{1}{117169800} = \frac{1}{117169800}$
\end{itemize}
La somme des numérateurs est :
$$ \frac{1}{66} + \frac{1}{10825} + \frac{1}{117169800} = \frac{1775300 + 10824 + 1}{117169800} = \frac{1786125}{117169800} $$
Le PGCD du numérateur et du dénominateur est $\text{PGCD}(1786125, 117169800) = 357225$.
La fraction irréductible est :
$$ \frac{1786125}{117169800} = \frac{1786125 \div 357225}{117169800 \div 357225} = \frac{5}{328} $$
Cette fraction est bien égale à $\frac{5}{328}$.

\subsection{Démonstration pour $n = 329$}
Soit $n = 329$. Nous cherchons $x, y, z \in \mathbb{Z}^{+}$ tels que $\frac{5}{329} = \frac{1}{x} + \frac{1}{y} + \frac{1}{z}$.
Posons $x = 66$, $y = 21715$, $z = 471519510$.
Les conditions $x > 0$, $y > 0$ et $z > 0$ sont satisfaites car $x, y, z \geq 1$.
Le PPCM des dénominateurs est $\text{PPCM}(66, 21715, 471519510) = 471519510$.
En réduisant au même dénominateur :
\begin{itemize}
    \item $\frac{1}{66} = \frac{7144235}{471519510}$
    \item $\frac{1}{21715} = \frac{21714}{471519510}$
    \item $\frac{1}{471519510} = \frac{1}{471519510}$
\end{itemize}
La somme des numérateurs est :
$$ \frac{1}{66} + \frac{1}{21715} + \frac{1}{471519510} = \frac{7144235 + 21714 + 1}{471519510} = \frac{7165950}{471519510} $$
Le PGCD du numérateur et du dénominateur est $\text{PGCD}(7165950, 471519510) = 1433190$.
La fraction irréductible est :
$$ \frac{7165950}{471519510} = \frac{7165950 \div 1433190}{471519510 \div 1433190} = \frac{5}{329} $$
Cette fraction est bien égale à $\frac{5}{329}$.

\subsection{Démonstration pour $n = 330$}
Soit $n = 330$. Nous cherchons $x, y, z \in \mathbb{Z}^{+}$ tels que $\frac{5}{330} = \frac{1}{x} + \frac{1}{y} + \frac{1}{z}$.
Posons $x = 67$, $y = 4423$, $z = 19558506$.
Les conditions $x > 0$, $y > 0$ et $z > 0$ sont satisfaites car $x, y, z \geq 1$.
Le PPCM des dénominateurs est $\text{PPCM}(67, 4423, 19558506) = 19558506$.
En réduisant au même dénominateur :
\begin{itemize}
    \item $\frac{1}{67} = \frac{291918}{19558506}$
    \item $\frac{1}{4423} = \frac{4422}{19558506}$
    \item $\frac{1}{19558506} = \frac{1}{19558506}$
\end{itemize}
La somme des numérateurs est :
$$ \frac{1}{67} + \frac{1}{4423} + \frac{1}{19558506} = \frac{291918 + 4422 + 1}{19558506} = \frac{296341}{19558506} $$
Le PGCD du numérateur et du dénominateur est $\text{PGCD}(296341, 19558506) = 296341$.
La fraction irréductible est :
$$ \frac{296341}{19558506} = \frac{296341 \div 296341}{19558506 \div 296341} = \frac{1}{66} $$
Cette fraction est bien égale à $\frac{5}{330}$.

\subsection{Démonstration pour $n = 331$}
Soit $n = 331$. Nous cherchons $x, y, z \in \mathbb{Z}^{+}$ tels que $\frac{5}{331} = \frac{1}{x} + \frac{1}{y} + \frac{1}{z}$.
Posons $x = 67$, $y = 5561$, $z = 1840691$.
Les conditions $x > 0$, $y > 0$ et $z > 0$ sont satisfaites car $x, y, z \geq 1$.
Le PPCM des dénominateurs est $\text{PPCM}(67, 5561, 1840691) = 1840691$.
En réduisant au même dénominateur :
\begin{itemize}
    \item $\frac{1}{67} = \frac{27473}{1840691}$
    \item $\frac{1}{5561} = \frac{331}{1840691}$
    \item $\frac{1}{1840691} = \frac{1}{1840691}$
\end{itemize}
La somme des numérateurs est :
$$ \frac{1}{67} + \frac{1}{5561} + \frac{1}{1840691} = \frac{27473 + 331 + 1}{1840691} = \frac{27805}{1840691} $$
Le PGCD du numérateur et du dénominateur est $\text{PGCD}(27805, 1840691) = 5561$.
La fraction irréductible est :
$$ \frac{27805}{1840691} = \frac{27805 \div 5561}{1840691 \div 5561} = \frac{5}{331} $$
Cette fraction est bien égale à $\frac{5}{331}$.

\subsection{Démonstration pour $n = 332$}
Soit $n = 332$. Nous cherchons $x, y, z \in \mathbb{Z}^{+}$ tels que $\frac{5}{332} = \frac{1}{x} + \frac{1}{y} + \frac{1}{z}$.
Posons $x = 67$, $y = 7415$, $z = 164939260$.
Les conditions $x > 0$, $y > 0$ et $z > 0$ sont satisfaites car $x, y, z \geq 1$.
Le PPCM des dénominateurs est $\text{PPCM}(67, 7415, 164939260) = 164939260$.
En réduisant au même dénominateur :
\begin{itemize}
    \item $\frac{1}{67} = \frac{2461780}{164939260}$
    \item $\frac{1}{7415} = \frac{22244}{164939260}$
    \item $\frac{1}{164939260} = \frac{1}{164939260}$
\end{itemize}
La somme des numérateurs est :
$$ \frac{1}{67} + \frac{1}{7415} + \frac{1}{164939260} = \frac{2461780 + 22244 + 1}{164939260} = \frac{2484025}{164939260} $$
Le PGCD du numérateur et du dénominateur est $\text{PGCD}(2484025, 164939260) = 496805$.
La fraction irréductible est :
$$ \frac{2484025}{164939260} = \frac{2484025 \div 496805}{164939260 \div 496805} = \frac{5}{332} $$
Cette fraction est bien égale à $\frac{5}{332}$.

\subsection{Démonstration pour $n = 333$}
Soit $n = 333$. Nous cherchons $x, y, z \in \mathbb{Z}^{+}$ tels que $\frac{5}{333} = \frac{1}{x} + \frac{1}{y} + \frac{1}{z}$.
Posons $x = 67$, $y = 11156$, $z = 248901516$.
Les conditions $x > 0$, $y > 0$ et $z > 0$ sont satisfaites car $x, y, z \geq 1$.
Le PPCM des dénominateurs est $\text{PPCM}(67, 11156, 248901516) = 248901516$.
En réduisant au même dénominateur :
\begin{itemize}
    \item $\frac{1}{67} = \frac{3714948}{248901516}$
    \item $\frac{1}{11156} = \frac{22311}{248901516}$
    \item $\frac{1}{248901516} = \frac{1}{248901516}$
\end{itemize}
La somme des numérateurs est :
$$ \frac{1}{67} + \frac{1}{11156} + \frac{1}{248901516} = \frac{3714948 + 22311 + 1}{248901516} = \frac{3737260}{248901516} $$
Le PGCD du numérateur et du dénominateur est $\text{PGCD}(3737260, 248901516) = 747452$.
La fraction irréductible est :
$$ \frac{3737260}{248901516} = \frac{3737260 \div 747452}{248901516 \div 747452} = \frac{5}{333} $$
Cette fraction est bien égale à $\frac{5}{333}$.

\subsection{Démonstration pour $n = 334$}
Soit $n = 334$. Nous cherchons $x, y, z \in \mathbb{Z}^{+}$ tels que $\frac{5}{334} = \frac{1}{x} + \frac{1}{y} + \frac{1}{z}$.
Posons $x = 67$, $y = 22379$, $z = 500797262$.
Les conditions $x > 0$, $y > 0$ et $z > 0$ sont satisfaites car $x, y, z \geq 1$.
Le PPCM des dénominateurs est $\text{PPCM}(67, 22379, 500797262) = 500797262$.
En réduisant au même dénominateur :
\begin{itemize}
    \item $\frac{1}{67} = \frac{7474586}{500797262}$
    \item $\frac{1}{22379} = \frac{22378}{500797262}$
    \item $\frac{1}{500797262} = \frac{1}{500797262}$
\end{itemize}
La somme des numérateurs est :
$$ \frac{1}{67} + \frac{1}{22379} + \frac{1}{500797262} = \frac{7474586 + 22378 + 1}{500797262} = \frac{7496965}{500797262} $$
Le PGCD du numérateur et du dénominateur est $\text{PGCD}(7496965, 500797262) = 1499393$.
La fraction irréductible est :
$$ \frac{7496965}{500797262} = \frac{7496965 \div 1499393}{500797262 \div 1499393} = \frac{5}{334} $$
Cette fraction est bien égale à $\frac{5}{334}$.

\subsection{Démonstration pour $n = 335$}
Soit $n = 335$. Nous cherchons $x, y, z \in \mathbb{Z}^{+}$ tels que $\frac{5}{335} = \frac{1}{x} + \frac{1}{y} + \frac{1}{z}$.
Posons $x = 68$, $y = 4557$, $z = 20761692$.
Les conditions $x > 0$, $y > 0$ et $z > 0$ sont satisfaites car $x, y, z \geq 1$.
Le PPCM des dénominateurs est $\text{PPCM}(68, 4557, 20761692) = 20761692$.
En réduisant au même dénominateur :
\begin{itemize}
    \item $\frac{1}{68} = \frac{305319}{20761692}$
    \item $\frac{1}{4557} = \frac{4556}{20761692}$
    \item $\frac{1}{20761692} = \frac{1}{20761692}$
\end{itemize}
La somme des numérateurs est :
$$ \frac{1}{68} + \frac{1}{4557} + \frac{1}{20761692} = \frac{305319 + 4556 + 1}{20761692} = \frac{309876}{20761692} $$
Le PGCD du numérateur et du dénominateur est $\text{PGCD}(309876, 20761692) = 309876$.
La fraction irréductible est :
$$ \frac{309876}{20761692} = \frac{309876 \div 309876}{20761692 \div 309876} = \frac{1}{67} $$
Cette fraction est bien égale à $\frac{5}{335}$.

\subsection{Démonstration pour $n = 336$}
Soit $n = 336$. Nous cherchons $x, y, z \in \mathbb{Z}^{+}$ tels que $\frac{5}{336} = \frac{1}{x} + \frac{1}{y} + \frac{1}{z}$.
Posons $x = 68$, $y = 5713$, $z = 32632656$.
Les conditions $x > 0$, $y > 0$ et $z > 0$ sont satisfaites car $x, y, z \geq 1$.
Le PPCM des dénominateurs est $\text{PPCM}(68, 5713, 32632656) = 32632656$.
En réduisant au même dénominateur :
\begin{itemize}
    \item $\frac{1}{68} = \frac{479892}{32632656}$
    \item $\frac{1}{5713} = \frac{5712}{32632656}$
    \item $\frac{1}{32632656} = \frac{1}{32632656}$
\end{itemize}
La somme des numérateurs est :
$$ \frac{1}{68} + \frac{1}{5713} + \frac{1}{32632656} = \frac{479892 + 5712 + 1}{32632656} = \frac{485605}{32632656} $$
Le PGCD du numérateur et du dénominateur est $\text{PGCD}(485605, 32632656) = 97121$.
La fraction irréductible est :
$$ \frac{485605}{32632656} = \frac{485605 \div 97121}{32632656 \div 97121} = \frac{5}{336} $$
Cette fraction est bien égale à $\frac{5}{336}$.

\subsection{Démonstration pour $n = 337$}
Soit $n = 337$. Nous cherchons $x, y, z \in \mathbb{Z}^{+}$ tels que $\frac{5}{337} = \frac{1}{x} + \frac{1}{y} + \frac{1}{z}$.
Posons $x = 68$, $y = 7639$, $z = 175055324$.
Les conditions $x > 0$, $y > 0$ et $z > 0$ sont satisfaites car $x, y, z \geq 1$.
Le PPCM des dénominateurs est $\text{PPCM}(68, 7639, 175055324) = 175055324$.
En réduisant au même dénominateur :
\begin{itemize}
    \item $\frac{1}{68} = \frac{2574343}{175055324}$
    \item $\frac{1}{7639} = \frac{22916}{175055324}$
    \item $\frac{1}{175055324} = \frac{1}{175055324}$
\end{itemize}
La somme des numérateurs est :
$$ \frac{1}{68} + \frac{1}{7639} + \frac{1}{175055324} = \frac{2574343 + 22916 + 1}{175055324} = \frac{2597260}{175055324} $$
Le PGCD du numérateur et du dénominateur est $\text{PGCD}(2597260, 175055324) = 519452$.
La fraction irréductible est :
$$ \frac{2597260}{175055324} = \frac{2597260 \div 519452}{175055324 \div 519452} = \frac{5}{337} $$
Cette fraction est bien égale à $\frac{5}{337}$.

\subsection{Démonstration pour $n = 338$}
Soit $n = 338$. Nous cherchons $x, y, z \in \mathbb{Z}^{+}$ tels que $\frac{5}{338} = \frac{1}{x} + \frac{1}{y} + \frac{1}{z}$.
Posons $x = 68$, $y = 11493$, $z = 132077556$.
Les conditions $x > 0$, $y > 0$ et $z > 0$ sont satisfaites car $x, y, z \geq 1$.
Le PPCM des dénominateurs est $\text{PPCM}(68, 11493, 132077556) = 132077556$.
En réduisant au même dénominateur :
\begin{itemize}
    \item $\frac{1}{68} = \frac{1942317}{132077556}$
    \item $\frac{1}{11493} = \frac{11492}{132077556}$
    \item $\frac{1}{132077556} = \frac{1}{132077556}$
\end{itemize}
La somme des numérateurs est :
$$ \frac{1}{68} + \frac{1}{11493} + \frac{1}{132077556} = \frac{1942317 + 11492 + 1}{132077556} = \frac{1953810}{132077556} $$
Le PGCD du numérateur et du dénominateur est $\text{PGCD}(1953810, 132077556) = 390762$.
La fraction irréductible est :
$$ \frac{1953810}{132077556} = \frac{1953810 \div 390762}{132077556 \div 390762} = \frac{5}{338} $$
Cette fraction est bien égale à $\frac{5}{338}$.

\subsection{Démonstration pour $n = 339$}
Soit $n = 339$. Nous cherchons $x, y, z \in \mathbb{Z}^{+}$ tels que $\frac{5}{339} = \frac{1}{x} + \frac{1}{y} + \frac{1}{z}$.
Posons $x = 68$, $y = 23053$, $z = 531417756$.
Les conditions $x > 0$, $y > 0$ et $z > 0$ sont satisfaites car $x, y, z \geq 1$.
Le PPCM des dénominateurs est $\text{PPCM}(68, 23053, 531417756) = 531417756$.
En réduisant au même dénominateur :
\begin{itemize}
    \item $\frac{1}{68} = \frac{7814967}{531417756}$
    \item $\frac{1}{23053} = \frac{23052}{531417756}$
    \item $\frac{1}{531417756} = \frac{1}{531417756}$
\end{itemize}
La somme des numérateurs est :
$$ \frac{1}{68} + \frac{1}{23053} + \frac{1}{531417756} = \frac{7814967 + 23052 + 1}{531417756} = \frac{7838020}{531417756} $$
Le PGCD du numérateur et du dénominateur est $\text{PGCD}(7838020, 531417756) = 1567604$.
La fraction irréductible est :
$$ \frac{7838020}{531417756} = \frac{7838020 \div 1567604}{531417756 \div 1567604} = \frac{5}{339} $$
Cette fraction est bien égale à $\frac{5}{339}$.

\subsection{Démonstration pour $n = 340$}
Soit $n = 340$. Nous cherchons $x, y, z \in \mathbb{Z}^{+}$ tels que $\frac{5}{340} = \frac{1}{x} + \frac{1}{y} + \frac{1}{z}$.
Posons $x = 69$, $y = 4693$, $z = 22019556$.
Les conditions $x > 0$, $y > 0$ et $z > 0$ sont satisfaites car $x, y, z \geq 1$.
Le PPCM des dénominateurs est $\text{PPCM}(69, 4693, 22019556) = 22019556$.
En réduisant au même dénominateur :
\begin{itemize}
    \item $\frac{1}{69} = \frac{319124}{22019556}$
    \item $\frac{1}{4693} = \frac{4692}{22019556}$
    \item $\frac{1}{22019556} = \frac{1}{22019556}$
\end{itemize}
La somme des numérateurs est :
$$ \frac{1}{69} + \frac{1}{4693} + \frac{1}{22019556} = \frac{319124 + 4692 + 1}{22019556} = \frac{323817}{22019556} $$
Le PGCD du numérateur et du dénominateur est $\text{PGCD}(323817, 22019556) = 323817$.
La fraction irréductible est :
$$ \frac{323817}{22019556} = \frac{323817 \div 323817}{22019556 \div 323817} = \frac{1}{68} $$
Cette fraction est bien égale à $\frac{5}{340}$.

\subsection{Démonstration pour $n = 341$}
Soit $n = 341$. Nous cherchons $x, y, z \in \mathbb{Z}^{+}$ tels que $\frac{5}{341} = \frac{1}{x} + \frac{1}{y} + \frac{1}{z}$.
Posons $x = 69$, $y = 5883$, $z = 46140369$.
Les conditions $x > 0$, $y > 0$ et $z > 0$ sont satisfaites car $x, y, z \geq 1$.
Le PPCM des dénominateurs est $\text{PPCM}(69, 5883, 46140369) = 46140369$.
En réduisant au même dénominateur :
\begin{itemize}
    \item $\frac{1}{69} = \frac{668701}{46140369}$
    \item $\frac{1}{5883} = \frac{7843}{46140369}$
    \item $\frac{1}{46140369} = \frac{1}{46140369}$
\end{itemize}
La somme des numérateurs est :
$$ \frac{1}{69} + \frac{1}{5883} + \frac{1}{46140369} = \frac{668701 + 7843 + 1}{46140369} = \frac{676545}{46140369} $$
Le PGCD du numérateur et du dénominateur est $\text{PGCD}(676545, 46140369) = 135309$.
La fraction irréductible est :
$$ \frac{676545}{46140369} = \frac{676545 \div 135309}{46140369 \div 135309} = \frac{5}{341} $$
Cette fraction est bien égale à $\frac{5}{341}$.

\subsection{Démonstration pour $n = 342$}
Soit $n = 342$. Nous cherchons $x, y, z \in \mathbb{Z}^{+}$ tels que $\frac{5}{342} = \frac{1}{x} + \frac{1}{y} + \frac{1}{z}$.
Posons $x = 69$, $y = 7867$, $z = 61881822$.
Les conditions $x > 0$, $y > 0$ et $z > 0$ sont satisfaites car $x, y, z \geq 1$.
Le PPCM des dénominateurs est $\text{PPCM}(69, 7867, 61881822) = 61881822$.
En réduisant au même dénominateur :
\begin{itemize}
    \item $\frac{1}{69} = \frac{896838}{61881822}$
    \item $\frac{1}{7867} = \frac{7866}{61881822}$
    \item $\frac{1}{61881822} = \frac{1}{61881822}$
\end{itemize}
La somme des numérateurs est :
$$ \frac{1}{69} + \frac{1}{7867} + \frac{1}{61881822} = \frac{896838 + 7866 + 1}{61881822} = \frac{904705}{61881822} $$
Le PGCD du numérateur et du dénominateur est $\text{PGCD}(904705, 61881822) = 180941$.
La fraction irréductible est :
$$ \frac{904705}{61881822} = \frac{904705 \div 180941}{61881822 \div 180941} = \frac{5}{342} $$
Cette fraction est bien égale à $\frac{5}{342}$.

\subsection{Démonstration pour $n = 343$}
Soit $n = 343$. Nous cherchons $x, y, z \in \mathbb{Z}^{+}$ tels que $\frac{5}{343} = \frac{1}{x} + \frac{1}{y} + \frac{1}{z}$.
Posons $x = 69$, $y = 11834$, $z = 280075278$.
Les conditions $x > 0$, $y > 0$ et $z > 0$ sont satisfaites car $x, y, z \geq 1$.
Le PPCM des dénominateurs est $\text{PPCM}(69, 11834, 280075278) = 280075278$.
En réduisant au même dénominateur :
\begin{itemize}
    \item $\frac{1}{69} = \frac{4059062}{280075278}$
    \item $\frac{1}{11834} = \frac{23667}{280075278}$
    \item $\frac{1}{280075278} = \frac{1}{280075278}$
\end{itemize}
La somme des numérateurs est :
$$ \frac{1}{69} + \frac{1}{11834} + \frac{1}{280075278} = \frac{4059062 + 23667 + 1}{280075278} = \frac{4082730}{280075278} $$
Le PGCD du numérateur et du dénominateur est $\text{PGCD}(4082730, 280075278) = 816546$.
La fraction irréductible est :
$$ \frac{4082730}{280075278} = \frac{4082730 \div 816546}{280075278 \div 816546} = \frac{5}{343} $$
Cette fraction est bien égale à $\frac{5}{343}$.

\subsection{Démonstration pour $n = 344$}
Soit $n = 344$. Nous cherchons $x, y, z \in \mathbb{Z}^{+}$ tels que $\frac{5}{344} = \frac{1}{x} + \frac{1}{y} + \frac{1}{z}$.
Posons $x = 69$, $y = 23737$, $z = 563421432$.
Les conditions $x > 0$, $y > 0$ et $z > 0$ sont satisfaites car $x, y, z \geq 1$.
Le PPCM des dénominateurs est $\text{PPCM}(69, 23737, 563421432) = 563421432$.
En réduisant au même dénominateur :
\begin{itemize}
    \item $\frac{1}{69} = \frac{8165528}{563421432}$
    \item $\frac{1}{23737} = \frac{23736}{563421432}$
    \item $\frac{1}{563421432} = \frac{1}{563421432}$
\end{itemize}
La somme des numérateurs est :
$$ \frac{1}{69} + \frac{1}{23737} + \frac{1}{563421432} = \frac{8165528 + 23736 + 1}{563421432} = \frac{8189265}{563421432} $$
Le PGCD du numérateur et du dénominateur est $\text{PGCD}(8189265, 563421432) = 1637853$.
La fraction irréductible est :
$$ \frac{8189265}{563421432} = \frac{8189265 \div 1637853}{563421432 \div 1637853} = \frac{5}{344} $$
Cette fraction est bien égale à $\frac{5}{344}$.

\subsection{Démonstration pour $n = 345$}
Soit $n = 345$. Nous cherchons $x, y, z \in \mathbb{Z}^{+}$ tels que $\frac{5}{345} = \frac{1}{x} + \frac{1}{y} + \frac{1}{z}$.
Posons $x = 70$, $y = 4831$, $z = 23333730$.
Les conditions $x > 0$, $y > 0$ et $z > 0$ sont satisfaites car $x, y, z \geq 1$.
Le PPCM des dénominateurs est $\text{PPCM}(70, 4831, 23333730) = 23333730$.
En réduisant au même dénominateur :
\begin{itemize}
    \item $\frac{1}{70} = \frac{333339}{23333730}$
    \item $\frac{1}{4831} = \frac{4830}{23333730}$
    \item $\frac{1}{23333730} = \frac{1}{23333730}$
\end{itemize}
La somme des numérateurs est :
$$ \frac{1}{70} + \frac{1}{4831} + \frac{1}{23333730} = \frac{333339 + 4830 + 1}{23333730} = \frac{338170}{23333730} $$
Le PGCD du numérateur et du dénominateur est $\text{PGCD}(338170, 23333730) = 338170$.
La fraction irréductible est :
$$ \frac{338170}{23333730} = \frac{338170 \div 338170}{23333730 \div 338170} = \frac{1}{69} $$
Cette fraction est bien égale à $\frac{5}{345}$.

\subsection{Démonstration pour $n = 346$}
Soit $n = 346$. Nous cherchons $x, y, z \in \mathbb{Z}^{+}$ tels que $\frac{5}{346} = \frac{1}{x} + \frac{1}{y} + \frac{1}{z}$.
Posons $x = 70$, $y = 6056$, $z = 36669080$.
Les conditions $x > 0$, $y > 0$ et $z > 0$ sont satisfaites car $x, y, z \geq 1$.
Le PPCM des dénominateurs est $\text{PPCM}(70, 6056, 36669080) = 36669080$.
En réduisant au même dénominateur :
\begin{itemize}
    \item $\frac{1}{70} = \frac{523844}{36669080}$
    \item $\frac{1}{6056} = \frac{6055}{36669080}$
    \item $\frac{1}{36669080} = \frac{1}{36669080}$
\end{itemize}
La somme des numérateurs est :
$$ \frac{1}{70} + \frac{1}{6056} + \frac{1}{36669080} = \frac{523844 + 6055 + 1}{36669080} = \frac{529900}{36669080} $$
Le PGCD du numérateur et du dénominateur est $\text{PGCD}(529900, 36669080) = 105980$.
La fraction irréductible est :
$$ \frac{529900}{36669080} = \frac{529900 \div 105980}{36669080 \div 105980} = \frac{5}{346} $$
Cette fraction est bien égale à $\frac{5}{346}$.

\subsection{Démonstration pour $n = 347$}
Soit $n = 347$. Nous cherchons $x, y, z \in \mathbb{Z}^{+}$ tels que $\frac{5}{347} = \frac{1}{x} + \frac{1}{y} + \frac{1}{z}$.
Posons $x = 70$, $y = 8097$, $z = 196676130$.
Les conditions $x > 0$, $y > 0$ et $z > 0$ sont satisfaites car $x, y, z \geq 1$.
Le PPCM des dénominateurs est $\text{PPCM}(70, 8097, 196676130) = 196676130$.
En réduisant au même dénominateur :
\begin{itemize}
    \item $\frac{1}{70} = \frac{2809659}{196676130}$
    \item $\frac{1}{8097} = \frac{24290}{196676130}$
    \item $\frac{1}{196676130} = \frac{1}{196676130}$
\end{itemize}
La somme des numérateurs est :
$$ \frac{1}{70} + \frac{1}{8097} + \frac{1}{196676130} = \frac{2809659 + 24290 + 1}{196676130} = \frac{2833950}{196676130} $$
Le PGCD du numérateur et du dénominateur est $\text{PGCD}(2833950, 196676130) = 566790$.
La fraction irréductible est :
$$ \frac{2833950}{196676130} = \frac{2833950 \div 566790}{196676130 \div 566790} = \frac{5}{347} $$
Cette fraction est bien égale à $\frac{5}{347}$.

\subsection{Démonstration pour $n = 348$}
Soit $n = 348$. Nous cherchons $x, y, z \in \mathbb{Z}^{+}$ tels que $\frac{5}{348} = \frac{1}{x} + \frac{1}{y} + \frac{1}{z}$.
Posons $x = 70$, $y = 12181$, $z = 148364580$.
Les conditions $x > 0$, $y > 0$ et $z > 0$ sont satisfaites car $x, y, z \geq 1$.
Le PPCM des dénominateurs est $\text{PPCM}(70, 12181, 148364580) = 148364580$.
En réduisant au même dénominateur :
\begin{itemize}
    \item $\frac{1}{70} = \frac{2119494}{148364580}$
    \item $\frac{1}{12181} = \frac{12180}{148364580}$
    \item $\frac{1}{148364580} = \frac{1}{148364580}$
\end{itemize}
La somme des numérateurs est :
$$ \frac{1}{70} + \frac{1}{12181} + \frac{1}{148364580} = \frac{2119494 + 12180 + 1}{148364580} = \frac{2131675}{148364580} $$
Le PGCD du numérateur et du dénominateur est $\text{PGCD}(2131675, 148364580) = 426335$.
La fraction irréductible est :
$$ \frac{2131675}{148364580} = \frac{2131675 \div 426335}{148364580 \div 426335} = \frac{5}{348} $$
Cette fraction est bien égale à $\frac{5}{348}$.

\subsection{Démonstration pour $n = 349$}
Soit $n = 349$. Nous cherchons $x, y, z \in \mathbb{Z}^{+}$ tels que $\frac{5}{349} = \frac{1}{x} + \frac{1}{y} + \frac{1}{z}$.
Posons $x = 70$, $y = 24431$, $z = 596849330$.
Les conditions $x > 0$, $y > 0$ et $z > 0$ sont satisfaites car $x, y, z \geq 1$.
Le PPCM des dénominateurs est $\text{PPCM}(70, 24431, 596849330) = 596849330$.
En réduisant au même dénominateur :
\begin{itemize}
    \item $\frac{1}{70} = \frac{8526419}{596849330}$
    \item $\frac{1}{24431} = \frac{24430}{596849330}$
    \item $\frac{1}{596849330} = \frac{1}{596849330}$
\end{itemize}
La somme des numérateurs est :
$$ \frac{1}{70} + \frac{1}{24431} + \frac{1}{596849330} = \frac{8526419 + 24430 + 1}{596849330} = \frac{8550850}{596849330} $$
Le PGCD du numérateur et du dénominateur est $\text{PGCD}(8550850, 596849330) = 1710170$.
La fraction irréductible est :
$$ \frac{8550850}{596849330} = \frac{8550850 \div 1710170}{596849330 \div 1710170} = \frac{5}{349} $$
Cette fraction est bien égale à $\frac{5}{349}$.

\subsection{Démonstration pour $n = 350$}
Soit $n = 350$. Nous cherchons $x, y, z \in \mathbb{Z}^{+}$ tels que $\frac{5}{350} = \frac{1}{x} + \frac{1}{y} + \frac{1}{z}$.
Posons $x = 71$, $y = 4971$, $z = 24705870$.
Les conditions $x > 0$, $y > 0$ et $z > 0$ sont satisfaites car $x, y, z \geq 1$.
Le PPCM des dénominateurs est $\text{PPCM}(71, 4971, 24705870) = 24705870$.
En réduisant au même dénominateur :
\begin{itemize}
    \item $\frac{1}{71} = \frac{347970}{24705870}$
    \item $\frac{1}{4971} = \frac{4970}{24705870}$
    \item $\frac{1}{24705870} = \frac{1}{24705870}$
\end{itemize}
La somme des numérateurs est :
$$ \frac{1}{71} + \frac{1}{4971} + \frac{1}{24705870} = \frac{347970 + 4970 + 1}{24705870} = \frac{352941}{24705870} $$
Le PGCD du numérateur et du dénominateur est $\text{PGCD}(352941, 24705870) = 352941$.
La fraction irréductible est :
$$ \frac{352941}{24705870} = \frac{352941 \div 352941}{24705870 \div 352941} = \frac{1}{70} $$
Cette fraction est bien égale à $\frac{5}{350}$.

\subsection{Démonstration pour $n = 351$}
Soit $n = 351$. Nous cherchons $x, y, z \in \mathbb{Z}^{+}$ tels que $\frac{5}{351} = \frac{1}{x} + \frac{1}{y} + \frac{1}{z}$.
Posons $x = 71$, $y = 6231$, $z = 51760917$.
Les conditions $x > 0$, $y > 0$ et $z > 0$ sont satisfaites car $x, y, z \geq 1$.
Le PPCM des dénominateurs est $\text{PPCM}(71, 6231, 51760917) = 51760917$.
En réduisant au même dénominateur :
\begin{itemize}
    \item $\frac{1}{71} = \frac{729027}{51760917}$
    \item $\frac{1}{6231} = \frac{8307}{51760917}$
    \item $\frac{1}{51760917} = \frac{1}{51760917}$
\end{itemize}
La somme des numérateurs est :
$$ \frac{1}{71} + \frac{1}{6231} + \frac{1}{51760917} = \frac{729027 + 8307 + 1}{51760917} = \frac{737335}{51760917} $$
Le PGCD du numérateur et du dénominateur est $\text{PGCD}(737335, 51760917) = 147467$.
La fraction irréductible est :
$$ \frac{737335}{51760917} = \frac{737335 \div 147467}{51760917 \div 147467} = \frac{5}{351} $$
Cette fraction est bien égale à $\frac{5}{351}$.

\subsection{Démonstration pour $n = 352$}
Soit $n = 352$. Nous cherchons $x, y, z \in \mathbb{Z}^{+}$ tels que $\frac{5}{352} = \frac{1}{x} + \frac{1}{y} + \frac{1}{z}$.
Posons $x = 71$, $y = 8331$, $z = 208208352$.
Les conditions $x > 0$, $y > 0$ et $z > 0$ sont satisfaites car $x, y, z \geq 1$.
Le PPCM des dénominateurs est $\text{PPCM}(71, 8331, 208208352) = 208208352$.
En réduisant au même dénominateur :
\begin{itemize}
    \item $\frac{1}{71} = \frac{2932512}{208208352}$
    \item $\frac{1}{8331} = \frac{24992}{208208352}$
    \item $\frac{1}{208208352} = \frac{1}{208208352}$
\end{itemize}
La somme des numérateurs est :
$$ \frac{1}{71} + \frac{1}{8331} + \frac{1}{208208352} = \frac{2932512 + 24992 + 1}{208208352} = \frac{2957505}{208208352} $$
Le PGCD du numérateur et du dénominateur est $\text{PGCD}(2957505, 208208352) = 591501$.
La fraction irréductible est :
$$ \frac{2957505}{208208352} = \frac{2957505 \div 591501}{208208352 \div 591501} = \frac{5}{352} $$
Cette fraction est bien égale à $\frac{5}{352}$.

\subsection{Démonstration pour $n = 353$}
Soit $n = 353$. Nous cherchons $x, y, z \in \mathbb{Z}^{+}$ tels que $\frac{5}{353} = \frac{1}{x} + \frac{1}{y} + \frac{1}{z}$.
Posons $x = 71$, $y = 12532$, $z = 314089516$.
Les conditions $x > 0$, $y > 0$ et $z > 0$ sont satisfaites car $x, y, z \geq 1$.
Le PPCM des dénominateurs est $\text{PPCM}(71, 12532, 314089516) = 314089516$.
En réduisant au même dénominateur :
\begin{itemize}
    \item $\frac{1}{71} = \frac{4423796}{314089516}$
    \item $\frac{1}{12532} = \frac{25063}{314089516}$
    \item $\frac{1}{314089516} = \frac{1}{314089516}$
\end{itemize}
La somme des numérateurs est :
$$ \frac{1}{71} + \frac{1}{12532} + \frac{1}{314089516} = \frac{4423796 + 25063 + 1}{314089516} = \frac{4448860}{314089516} $$
Le PGCD du numérateur et du dénominateur est $\text{PGCD}(4448860, 314089516) = 889772$.
La fraction irréductible est :
$$ \frac{4448860}{314089516} = \frac{4448860 \div 889772}{314089516 \div 889772} = \frac{5}{353} $$
Cette fraction est bien égale à $\frac{5}{353}$.

\subsection{Démonstration pour $n = 354$}
Soit $n = 354$. Nous cherchons $x, y, z \in \mathbb{Z}^{+}$ tels que $\frac{5}{354} = \frac{1}{x} + \frac{1}{y} + \frac{1}{z}$.
Posons $x = 71$, $y = 25135$, $z = 631743090$.
Les conditions $x > 0$, $y > 0$ et $z > 0$ sont satisfaites car $x, y, z \geq 1$.
Le PPCM des dénominateurs est $\text{PPCM}(71, 25135, 631743090) = 631743090$.
En réduisant au même dénominateur :
\begin{itemize}
    \item $\frac{1}{71} = \frac{8897790}{631743090}$
    \item $\frac{1}{25135} = \frac{25134}{631743090}$
    \item $\frac{1}{631743090} = \frac{1}{631743090}$
\end{itemize}
La somme des numérateurs est :
$$ \frac{1}{71} + \frac{1}{25135} + \frac{1}{631743090} = \frac{8897790 + 25134 + 1}{631743090} = \frac{8922925}{631743090} $$
Le PGCD du numérateur et du dénominateur est $\text{PGCD}(8922925, 631743090) = 1784585$.
La fraction irréductible est :
$$ \frac{8922925}{631743090} = \frac{8922925 \div 1784585}{631743090 \div 1784585} = \frac{5}{354} $$
Cette fraction est bien égale à $\frac{5}{354}$.

\subsection{Démonstration pour $n = 355$}
Soit $n = 355$. Nous cherchons $x, y, z \in \mathbb{Z}^{+}$ tels que $\frac{5}{355} = \frac{1}{x} + \frac{1}{y} + \frac{1}{z}$.
Posons $x = 72$, $y = 5113$, $z = 26137656$.
Les conditions $x > 0$, $y > 0$ et $z > 0$ sont satisfaites car $x, y, z \geq 1$.
Le PPCM des dénominateurs est $\text{PPCM}(72, 5113, 26137656) = 26137656$.
En réduisant au même dénominateur :
\begin{itemize}
    \item $\frac{1}{72} = \frac{363023}{26137656}$
    \item $\frac{1}{5113} = \frac{5112}{26137656}$
    \item $\frac{1}{26137656} = \frac{1}{26137656}$
\end{itemize}
La somme des numérateurs est :
$$ \frac{1}{72} + \frac{1}{5113} + \frac{1}{26137656} = \frac{363023 + 5112 + 1}{26137656} = \frac{368136}{26137656} $$
Le PGCD du numérateur et du dénominateur est $\text{PGCD}(368136, 26137656) = 368136$.
La fraction irréductible est :
$$ \frac{368136}{26137656} = \frac{368136 \div 368136}{26137656 \div 368136} = \frac{1}{71} $$
Cette fraction est bien égale à $\frac{5}{355}$.

\subsection{Démonstration pour $n = 356$}
Soit $n = 356$. Nous cherchons $x, y, z \in \mathbb{Z}^{+}$ tels que $\frac{5}{356} = \frac{1}{x} + \frac{1}{y} + \frac{1}{z}$.
Posons $x = 72$, $y = 6409$, $z = 41068872$.
Les conditions $x > 0$, $y > 0$ et $z > 0$ sont satisfaites car $x, y, z \geq 1$.
Le PPCM des dénominateurs est $\text{PPCM}(72, 6409, 41068872) = 41068872$.
En réduisant au même dénominateur :
\begin{itemize}
    \item $\frac{1}{72} = \frac{570401}{41068872}$
    \item $\frac{1}{6409} = \frac{6408}{41068872}$
    \item $\frac{1}{41068872} = \frac{1}{41068872}$
\end{itemize}
La somme des numérateurs est :
$$ \frac{1}{72} + \frac{1}{6409} + \frac{1}{41068872} = \frac{570401 + 6408 + 1}{41068872} = \frac{576810}{41068872} $$
Le PGCD du numérateur et du dénominateur est $\text{PGCD}(576810, 41068872) = 115362$.
La fraction irréductible est :
$$ \frac{576810}{41068872} = \frac{576810 \div 115362}{41068872 \div 115362} = \frac{5}{356} $$
Cette fraction est bien égale à $\frac{5}{356}$.

\subsection{Démonstration pour $n = 357$}
Soit $n = 357$. Nous cherchons $x, y, z \in \mathbb{Z}^{+}$ tels que $\frac{5}{357} = \frac{1}{x} + \frac{1}{y} + \frac{1}{z}$.
Posons $x = 72$, $y = 8569$, $z = 73419192$.
Les conditions $x > 0$, $y > 0$ et $z > 0$ sont satisfaites car $x, y, z \geq 1$.
Le PPCM des dénominateurs est $\text{PPCM}(72, 8569, 73419192) = 73419192$.
En réduisant au même dénominateur :
\begin{itemize}
    \item $\frac{1}{72} = \frac{1019711}{73419192}$
    \item $\frac{1}{8569} = \frac{8568}{73419192}$
    \item $\frac{1}{73419192} = \frac{1}{73419192}$
\end{itemize}
La somme des numérateurs est :
$$ \frac{1}{72} + \frac{1}{8569} + \frac{1}{73419192} = \frac{1019711 + 8568 + 1}{73419192} = \frac{1028280}{73419192} $$
Le PGCD du numérateur et du dénominateur est $\text{PGCD}(1028280, 73419192) = 205656$.
La fraction irréductible est :
$$ \frac{1028280}{73419192} = \frac{1028280 \div 205656}{73419192 \div 205656} = \frac{5}{357} $$
Cette fraction est bien égale à $\frac{5}{357}$.

\subsection{Démonstration pour $n = 358$}
Soit $n = 358$. Nous cherchons $x, y, z \in \mathbb{Z}^{+}$ tels que $\frac{5}{358} = \frac{1}{x} + \frac{1}{y} + \frac{1}{z}$.
Posons $x = 72$, $y = 12889$, $z = 166113432$.
Les conditions $x > 0$, $y > 0$ et $z > 0$ sont satisfaites car $x, y, z \geq 1$.
Le PPCM des dénominateurs est $\text{PPCM}(72, 12889, 166113432) = 166113432$.
En réduisant au même dénominateur :
\begin{itemize}
    \item $\frac{1}{72} = \frac{2307131}{166113432}$
    \item $\frac{1}{12889} = \frac{12888}{166113432}$
    \item $\frac{1}{166113432} = \frac{1}{166113432}$
\end{itemize}
La somme des numérateurs est :
$$ \frac{1}{72} + \frac{1}{12889} + \frac{1}{166113432} = \frac{2307131 + 12888 + 1}{166113432} = \frac{2320020}{166113432} $$
Le PGCD du numérateur et du dénominateur est $\text{PGCD}(2320020, 166113432) = 464004$.
La fraction irréductible est :
$$ \frac{2320020}{166113432} = \frac{2320020 \div 464004}{166113432 \div 464004} = \frac{5}{358} $$
Cette fraction est bien égale à $\frac{5}{358}$.

\subsection{Démonstration pour $n = 359$}
Soit $n = 359$. Nous cherchons $x, y, z \in \mathbb{Z}^{+}$ tels que $\frac{5}{359} = \frac{1}{x} + \frac{1}{y} + \frac{1}{z}$.
Posons $x = 72$, $y = 25849$, $z = 668144952$.
Les conditions $x > 0$, $y > 0$ et $z > 0$ sont satisfaites car $x, y, z \geq 1$.
Le PPCM des dénominateurs est $\text{PPCM}(72, 25849, 668144952) = 668144952$.
En réduisant au même dénominateur :
\begin{itemize}
    \item $\frac{1}{72} = \frac{9279791}{668144952}$
    \item $\frac{1}{25849} = \frac{25848}{668144952}$
    \item $\frac{1}{668144952} = \frac{1}{668144952}$
\end{itemize}
La somme des numérateurs est :
$$ \frac{1}{72} + \frac{1}{25849} + \frac{1}{668144952} = \frac{9279791 + 25848 + 1}{668144952} = \frac{9305640}{668144952} $$
Le PGCD du numérateur et du dénominateur est $\text{PGCD}(9305640, 668144952) = 1861128$.
La fraction irréductible est :
$$ \frac{9305640}{668144952} = \frac{9305640 \div 1861128}{668144952 \div 1861128} = \frac{5}{359} $$
Cette fraction est bien égale à $\frac{5}{359}$.

\subsection{Démonstration pour $n = 360$}
Soit $n = 360$. Nous cherchons $x, y, z \in \mathbb{Z}^{+}$ tels que $\frac{5}{360} = \frac{1}{x} + \frac{1}{y} + \frac{1}{z}$.
Posons $x = 73$, $y = 5257$, $z = 27630792$.
Les conditions $x > 0$, $y > 0$ et $z > 0$ sont satisfaites car $x, y, z \geq 1$.
Le PPCM des dénominateurs est $\text{PPCM}(73, 5257, 27630792) = 27630792$.
En réduisant au même dénominateur :
\begin{itemize}
    \item $\frac{1}{73} = \frac{378504}{27630792}$
    \item $\frac{1}{5257} = \frac{5256}{27630792}$
    \item $\frac{1}{27630792} = \frac{1}{27630792}$
\end{itemize}
La somme des numérateurs est :
$$ \frac{1}{73} + \frac{1}{5257} + \frac{1}{27630792} = \frac{378504 + 5256 + 1}{27630792} = \frac{383761}{27630792} $$
Le PGCD du numérateur et du dénominateur est $\text{PGCD}(383761, 27630792) = 383761$.
La fraction irréductible est :
$$ \frac{383761}{27630792} = \frac{383761 \div 383761}{27630792 \div 383761} = \frac{1}{72} $$
Cette fraction est bien égale à $\frac{5}{360}$.

\subsection{Démonstration pour $n = 361$}
Soit $n = 361$. Nous cherchons $x, y, z \in \mathbb{Z}^{+}$ tels que $\frac{5}{361} = \frac{1}{x} + \frac{1}{y} + \frac{1}{z}$.
Posons $x = 73$, $y = 6593$, $z = 9144491$.
Les conditions $x > 0$, $y > 0$ et $z > 0$ sont satisfaites car $x, y, z \geq 1$.
Le PPCM des dénominateurs est $\text{PPCM}(73, 6593, 9144491) = 9144491$.
En réduisant au même dénominateur :
\begin{itemize}
    \item $\frac{1}{73} = \frac{125267}{9144491}$
    \item $\frac{1}{6593} = \frac{1387}{9144491}$
    \item $\frac{1}{9144491} = \frac{1}{9144491}$
\end{itemize}
La somme des numérateurs est :
$$ \frac{1}{73} + \frac{1}{6593} + \frac{1}{9144491} = \frac{125267 + 1387 + 1}{9144491} = \frac{126655}{9144491} $$
Le PGCD du numérateur et du dénominateur est $\text{PGCD}(126655, 9144491) = 25331$.
La fraction irréductible est :
$$ \frac{126655}{9144491} = \frac{126655 \div 25331}{9144491 \div 25331} = \frac{5}{361} $$
Cette fraction est bien égale à $\frac{5}{361}$.

\subsection{Démonstration pour $n = 362$}
Soit $n = 362$. Nous cherchons $x, y, z \in \mathbb{Z}^{+}$ tels que $\frac{5}{362} = \frac{1}{x} + \frac{1}{y} + \frac{1}{z}$.
Posons $x = 73$, $y = 8809$, $z = 232786634$.
Les conditions $x > 0$, $y > 0$ et $z > 0$ sont satisfaites car $x, y, z \geq 1$.
Le PPCM des dénominateurs est $\text{PPCM}(73, 8809, 232786634) = 232786634$.
En réduisant au même dénominateur :
\begin{itemize}
    \item $\frac{1}{73} = \frac{3188858}{232786634}$
    \item $\frac{1}{8809} = \frac{26426}{232786634}$
    \item $\frac{1}{232786634} = \frac{1}{232786634}$
\end{itemize}
La somme des numérateurs est :
$$ \frac{1}{73} + \frac{1}{8809} + \frac{1}{232786634} = \frac{3188858 + 26426 + 1}{232786634} = \frac{3215285}{232786634} $$
Le PGCD du numérateur et du dénominateur est $\text{PGCD}(3215285, 232786634) = 643057$.
La fraction irréductible est :
$$ \frac{3215285}{232786634} = \frac{3215285 \div 643057}{232786634 \div 643057} = \frac{5}{362} $$
Cette fraction est bien égale à $\frac{5}{362}$.

\subsection{Démonstration pour $n = 363$}
Soit $n = 363$. Nous cherchons $x, y, z \in \mathbb{Z}^{+}$ tels que $\frac{5}{363} = \frac{1}{x} + \frac{1}{y} + \frac{1}{z}$.
Posons $x = 73$, $y = 13250$, $z = 351111750$.
Les conditions $x > 0$, $y > 0$ et $z > 0$ sont satisfaites car $x, y, z \geq 1$.
Le PPCM des dénominateurs est $\text{PPCM}(73, 13250, 351111750) = 351111750$.
En réduisant au même dénominateur :
\begin{itemize}
    \item $\frac{1}{73} = \frac{4809750}{351111750}$
    \item $\frac{1}{13250} = \frac{26499}{351111750}$
    \item $\frac{1}{351111750} = \frac{1}{351111750}$
\end{itemize}
La somme des numérateurs est :
$$ \frac{1}{73} + \frac{1}{13250} + \frac{1}{351111750} = \frac{4809750 + 26499 + 1}{351111750} = \frac{4836250}{351111750} $$
Le PGCD du numérateur et du dénominateur est $\text{PGCD}(4836250, 351111750) = 967250$.
La fraction irréductible est :
$$ \frac{4836250}{351111750} = \frac{4836250 \div 967250}{351111750 \div 967250} = \frac{5}{363} $$
Cette fraction est bien égale à $\frac{5}{363}$.

\subsection{Démonstration pour $n = 364$}
Soit $n = 364$. Nous cherchons $x, y, z \in \mathbb{Z}^{+}$ tels que $\frac{5}{364} = \frac{1}{x} + \frac{1}{y} + \frac{1}{z}$.
Posons $x = 73$, $y = 26573$, $z = 706097756$.
Les conditions $x > 0$, $y > 0$ et $z > 0$ sont satisfaites car $x, y, z \geq 1$.
Le PPCM des dénominateurs est $\text{PPCM}(73, 26573, 706097756) = 706097756$.
En réduisant au même dénominateur :
\begin{itemize}
    \item $\frac{1}{73} = \frac{9672572}{706097756}$
    \item $\frac{1}{26573} = \frac{26572}{706097756}$
    \item $\frac{1}{706097756} = \frac{1}{706097756}$
\end{itemize}
La somme des numérateurs est :
$$ \frac{1}{73} + \frac{1}{26573} + \frac{1}{706097756} = \frac{9672572 + 26572 + 1}{706097756} = \frac{9699145}{706097756} $$
Le PGCD du numérateur et du dénominateur est $\text{PGCD}(9699145, 706097756) = 1939829$.
La fraction irréductible est :
$$ \frac{9699145}{706097756} = \frac{9699145 \div 1939829}{706097756 \div 1939829} = \frac{5}{364} $$
Cette fraction est bien égale à $\frac{5}{364}$.

\subsection{Démonstration pour $n = 365$}
Soit $n = 365$. Nous cherchons $x, y, z \in \mathbb{Z}^{+}$ tels que $\frac{5}{365} = \frac{1}{x} + \frac{1}{y} + \frac{1}{z}$.
Posons $x = 74$, $y = 5403$, $z = 29187006$.
Les conditions $x > 0$, $y > 0$ et $z > 0$ sont satisfaites car $x, y, z \geq 1$.
Le PPCM des dénominateurs est $\text{PPCM}(74, 5403, 29187006) = 29187006$.
En réduisant au même dénominateur :
\begin{itemize}
    \item $\frac{1}{74} = \frac{394419}{29187006}$
    \item $\frac{1}{5403} = \frac{5402}{29187006}$
    \item $\frac{1}{29187006} = \frac{1}{29187006}$
\end{itemize}
La somme des numérateurs est :
$$ \frac{1}{74} + \frac{1}{5403} + \frac{1}{29187006} = \frac{394419 + 5402 + 1}{29187006} = \frac{399822}{29187006} $$
Le PGCD du numérateur et du dénominateur est $\text{PGCD}(399822, 29187006) = 399822$.
La fraction irréductible est :
$$ \frac{399822}{29187006} = \frac{399822 \div 399822}{29187006 \div 399822} = \frac{1}{73} $$
Cette fraction est bien égale à $\frac{5}{365}$.

\subsection{Démonstration pour $n = 366$}
Soit $n = 366$. Nous cherchons $x, y, z \in \mathbb{Z}^{+}$ tels que $\frac{5}{366} = \frac{1}{x} + \frac{1}{y} + \frac{1}{z}$.
Posons $x = 74$, $y = 6772$, $z = 45853212$.
Les conditions $x > 0$, $y > 0$ et $z > 0$ sont satisfaites car $x, y, z \geq 1$.
Le PPCM des dénominateurs est $\text{PPCM}(74, 6772, 45853212) = 45853212$.
En réduisant au même dénominateur :
\begin{itemize}
    \item $\frac{1}{74} = \frac{619638}{45853212}$
    \item $\frac{1}{6772} = \frac{6771}{45853212}$
    \item $\frac{1}{45853212} = \frac{1}{45853212}$
\end{itemize}
La somme des numérateurs est :
$$ \frac{1}{74} + \frac{1}{6772} + \frac{1}{45853212} = \frac{619638 + 6771 + 1}{45853212} = \frac{626410}{45853212} $$
Le PGCD du numérateur et du dénominateur est $\text{PGCD}(626410, 45853212) = 125282$.
La fraction irréductible est :
$$ \frac{626410}{45853212} = \frac{626410 \div 125282}{45853212 \div 125282} = \frac{5}{366} $$
Cette fraction est bien égale à $\frac{5}{366}$.

\subsection{Démonstration pour $n = 367$}
Soit $n = 367$. Nous cherchons $x, y, z \in \mathbb{Z}^{+}$ tels que $\frac{5}{367} = \frac{1}{x} + \frac{1}{y} + \frac{1}{z}$.
Posons $x = 74$, $y = 9053$, $z = 245861374$.
Les conditions $x > 0$, $y > 0$ et $z > 0$ sont satisfaites car $x, y, z \geq 1$.
Le PPCM des dénominateurs est $\text{PPCM}(74, 9053, 245861374) = 245861374$.
En réduisant au même dénominateur :
\begin{itemize}
    \item $\frac{1}{74} = \frac{3322451}{245861374}$
    \item $\frac{1}{9053} = \frac{27158}{245861374}$
    \item $\frac{1}{245861374} = \frac{1}{245861374}$
\end{itemize}
La somme des numérateurs est :
$$ \frac{1}{74} + \frac{1}{9053} + \frac{1}{245861374} = \frac{3322451 + 27158 + 1}{245861374} = \frac{3349610}{245861374} $$
Le PGCD du numérateur et du dénominateur est $\text{PGCD}(3349610, 245861374) = 669922$.
La fraction irréductible est :
$$ \frac{3349610}{245861374} = \frac{3349610 \div 669922}{245861374 \div 669922} = \frac{5}{367} $$
Cette fraction est bien égale à $\frac{5}{367}$.

\subsection{Démonstration pour $n = 368$}
Soit $n = 368$. Nous cherchons $x, y, z \in \mathbb{Z}^{+}$ tels que $\frac{5}{368} = \frac{1}{x} + \frac{1}{y} + \frac{1}{z}$.
Posons $x = 74$, $y = 13617$, $z = 185409072$.
Les conditions $x > 0$, $y > 0$ et $z > 0$ sont satisfaites car $x, y, z \geq 1$.
Le PPCM des dénominateurs est $\text{PPCM}(74, 13617, 185409072) = 185409072$.
En réduisant au même dénominateur :
\begin{itemize}
    \item $\frac{1}{74} = \frac{2505528}{185409072}$
    \item $\frac{1}{13617} = \frac{13616}{185409072}$
    \item $\frac{1}{185409072} = \frac{1}{185409072}$
\end{itemize}
La somme des numérateurs est :
$$ \frac{1}{74} + \frac{1}{13617} + \frac{1}{185409072} = \frac{2505528 + 13616 + 1}{185409072} = \frac{2519145}{185409072} $$
Le PGCD du numérateur et du dénominateur est $\text{PGCD}(2519145, 185409072) = 503829$.
La fraction irréductible est :
$$ \frac{2519145}{185409072} = \frac{2519145 \div 503829}{185409072 \div 503829} = \frac{5}{368} $$
Cette fraction est bien égale à $\frac{5}{368}$.

\subsection{Démonstration pour $n = 369$}
Soit $n = 369$. Nous cherchons $x, y, z \in \mathbb{Z}^{+}$ tels que $\frac{5}{369} = \frac{1}{x} + \frac{1}{y} + \frac{1}{z}$.
Posons $x = 74$, $y = 27307$, $z = 745644942$.
Les conditions $x > 0$, $y > 0$ et $z > 0$ sont satisfaites car $x, y, z \geq 1$.
Le PPCM des dénominateurs est $\text{PPCM}(74, 27307, 745644942) = 745644942$.
En réduisant au même dénominateur :
\begin{itemize}
    \item $\frac{1}{74} = \frac{10076283}{745644942}$
    \item $\frac{1}{27307} = \frac{27306}{745644942}$
    \item $\frac{1}{745644942} = \frac{1}{745644942}$
\end{itemize}
La somme des numérateurs est :
$$ \frac{1}{74} + \frac{1}{27307} + \frac{1}{745644942} = \frac{10076283 + 27306 + 1}{745644942} = \frac{10103590}{745644942} $$
Le PGCD du numérateur et du dénominateur est $\text{PGCD}(10103590, 745644942) = 2020718$.
La fraction irréductible est :
$$ \frac{10103590}{745644942} = \frac{10103590 \div 2020718}{745644942 \div 2020718} = \frac{5}{369} $$
Cette fraction est bien égale à $\frac{5}{369}$.

\subsection{Démonstration pour $n = 370$}
Soit $n = 370$. Nous cherchons $x, y, z \in \mathbb{Z}^{+}$ tels que $\frac{5}{370} = \frac{1}{x} + \frac{1}{y} + \frac{1}{z}$.
Posons $x = 75$, $y = 5551$, $z = 30808050$.
Les conditions $x > 0$, $y > 0$ et $z > 0$ sont satisfaites car $x, y, z \geq 1$.
Le PPCM des dénominateurs est $\text{PPCM}(75, 5551, 30808050) = 30808050$.
En réduisant au même dénominateur :
\begin{itemize}
    \item $\frac{1}{75} = \frac{410774}{30808050}$
    \item $\frac{1}{5551} = \frac{5550}{30808050}$
    \item $\frac{1}{30808050} = \frac{1}{30808050}$
\end{itemize}
La somme des numérateurs est :
$$ \frac{1}{75} + \frac{1}{5551} + \frac{1}{30808050} = \frac{410774 + 5550 + 1}{30808050} = \frac{416325}{30808050} $$
Le PGCD du numérateur et du dénominateur est $\text{PGCD}(416325, 30808050) = 416325$.
La fraction irréductible est :
$$ \frac{416325}{30808050} = \frac{416325 \div 416325}{30808050 \div 416325} = \frac{1}{74} $$
Cette fraction est bien égale à $\frac{5}{370}$.

\subsection{Démonstration pour $n = 371$}
Soit $n = 371$. Nous cherchons $x, y, z \in \mathbb{Z}^{+}$ tels que $\frac{5}{371} = \frac{1}{x} + \frac{1}{y} + \frac{1}{z}$.
Posons $x = 75$, $y = 6957$, $z = 64526175$.
Les conditions $x > 0$, $y > 0$ et $z > 0$ sont satisfaites car $x, y, z \geq 1$.
Le PPCM des dénominateurs est $\text{PPCM}(75, 6957, 64526175) = 64526175$.
En réduisant au même dénominateur :
\begin{itemize}
    \item $\frac{1}{75} = \frac{860349}{64526175}$
    \item $\frac{1}{6957} = \frac{9275}{64526175}$
    \item $\frac{1}{64526175} = \frac{1}{64526175}$
\end{itemize}
La somme des numérateurs est :
$$ \frac{1}{75} + \frac{1}{6957} + \frac{1}{64526175} = \frac{860349 + 9275 + 1}{64526175} = \frac{869625}{64526175} $$
Le PGCD du numérateur et du dénominateur est $\text{PGCD}(869625, 64526175) = 173925$.
La fraction irréductible est :
$$ \frac{869625}{64526175} = \frac{869625 \div 173925}{64526175 \div 173925} = \frac{5}{371} $$
Cette fraction est bien égale à $\frac{5}{371}$.

\subsection{Démonstration pour $n = 372$}
Soit $n = 372$. Nous cherchons $x, y, z \in \mathbb{Z}^{+}$ tels que $\frac{5}{372} = \frac{1}{x} + \frac{1}{y} + \frac{1}{z}$.
Posons $x = 75$, $y = 9301$, $z = 86499300$.
Les conditions $x > 0$, $y > 0$ et $z > 0$ sont satisfaites car $x, y, z \geq 1$.
Le PPCM des dénominateurs est $\text{PPCM}(75, 9301, 86499300) = 86499300$.
En réduisant au même dénominateur :
\begin{itemize}
    \item $\frac{1}{75} = \frac{1153324}{86499300}$
    \item $\frac{1}{9301} = \frac{9300}{86499300}$
    \item $\frac{1}{86499300} = \frac{1}{86499300}$
\end{itemize}
La somme des numérateurs est :
$$ \frac{1}{75} + \frac{1}{9301} + \frac{1}{86499300} = \frac{1153324 + 9300 + 1}{86499300} = \frac{1162625}{86499300} $$
Le PGCD du numérateur et du dénominateur est $\text{PGCD}(1162625, 86499300) = 232525$.
La fraction irréductible est :
$$ \frac{1162625}{86499300} = \frac{1162625 \div 232525}{86499300 \div 232525} = \frac{5}{372} $$
Cette fraction est bien égale à $\frac{5}{372}$.

\subsection{Démonstration pour $n = 373$}
Soit $n = 373$. Nous cherchons $x, y, z \in \mathbb{Z}^{+}$ tels que $\frac{5}{373} = \frac{1}{x} + \frac{1}{y} + \frac{1}{z}$.
Posons $x = 75$, $y = 13988$, $z = 391314300$.
Les conditions $x > 0$, $y > 0$ et $z > 0$ sont satisfaites car $x, y, z \geq 1$.
Le PPCM des dénominateurs est $\text{PPCM}(75, 13988, 391314300) = 391314300$.
En réduisant au même dénominateur :
\begin{itemize}
    \item $\frac{1}{75} = \frac{5217524}{391314300}$
    \item $\frac{1}{13988} = \frac{27975}{391314300}$
    \item $\frac{1}{391314300} = \frac{1}{391314300}$
\end{itemize}
La somme des numérateurs est :
$$ \frac{1}{75} + \frac{1}{13988} + \frac{1}{391314300} = \frac{5217524 + 27975 + 1}{391314300} = \frac{5245500}{391314300} $$
Le PGCD du numérateur et du dénominateur est $\text{PGCD}(5245500, 391314300) = 1049100$.
La fraction irréductible est :
$$ \frac{5245500}{391314300} = \frac{5245500 \div 1049100}{391314300 \div 1049100} = \frac{5}{373} $$
Cette fraction est bien égale à $\frac{5}{373}$.

\subsection{Démonstration pour $n = 374$}
Soit $n = 374$. Nous cherchons $x, y, z \in \mathbb{Z}^{+}$ tels que $\frac{5}{374} = \frac{1}{x} + \frac{1}{y} + \frac{1}{z}$.
Posons $x = 75$, $y = 28051$, $z = 786830550$.
Les conditions $x > 0$, $y > 0$ et $z > 0$ sont satisfaites car $x, y, z \geq 1$.
Le PPCM des dénominateurs est $\text{PPCM}(75, 28051, 786830550) = 786830550$.
En réduisant au même dénominateur :
\begin{itemize}
    \item $\frac{1}{75} = \frac{10491074}{786830550}$
    \item $\frac{1}{28051} = \frac{28050}{786830550}$
    \item $\frac{1}{786830550} = \frac{1}{786830550}$
\end{itemize}
La somme des numérateurs est :
$$ \frac{1}{75} + \frac{1}{28051} + \frac{1}{786830550} = \frac{10491074 + 28050 + 1}{786830550} = \frac{10519125}{786830550} $$
Le PGCD du numérateur et du dénominateur est $\text{PGCD}(10519125, 786830550) = 2103825$.
La fraction irréductible est :
$$ \frac{10519125}{786830550} = \frac{10519125 \div 2103825}{786830550 \div 2103825} = \frac{5}{374} $$
Cette fraction est bien égale à $\frac{5}{374}$.

\subsection{Démonstration pour $n = 375$}
Soit $n = 375$. Nous cherchons $x, y, z \in \mathbb{Z}^{+}$ tels que $\frac{5}{375} = \frac{1}{x} + \frac{1}{y} + \frac{1}{z}$.
Posons $x = 76$, $y = 5701$, $z = 32495700$.
Les conditions $x > 0$, $y > 0$ et $z > 0$ sont satisfaites car $x, y, z \geq 1$.
Le PPCM des dénominateurs est $\text{PPCM}(76, 5701, 32495700) = 32495700$.
En réduisant au même dénominateur :
\begin{itemize}
    \item $\frac{1}{76} = \frac{427575}{32495700}$
    \item $\frac{1}{5701} = \frac{5700}{32495700}$
    \item $\frac{1}{32495700} = \frac{1}{32495700}$
\end{itemize}
La somme des numérateurs est :
$$ \frac{1}{76} + \frac{1}{5701} + \frac{1}{32495700} = \frac{427575 + 5700 + 1}{32495700} = \frac{433276}{32495700} $$
Le PGCD du numérateur et du dénominateur est $\text{PGCD}(433276, 32495700) = 433276$.
La fraction irréductible est :
$$ \frac{433276}{32495700} = \frac{433276 \div 433276}{32495700 \div 433276} = \frac{1}{75} $$
Cette fraction est bien égale à $\frac{5}{375}$.

\subsection{Démonstration pour $n = 376$}
Soit $n = 376$. Nous cherchons $x, y, z \in \mathbb{Z}^{+}$ tels que $\frac{5}{376} = \frac{1}{x} + \frac{1}{y} + \frac{1}{z}$.
Posons $x = 76$, $y = 7145$, $z = 51043880$.
Les conditions $x > 0$, $y > 0$ et $z > 0$ sont satisfaites car $x, y, z \geq 1$.
Le PPCM des dénominateurs est $\text{PPCM}(76, 7145, 51043880) = 51043880$.
En réduisant au même dénominateur :
\begin{itemize}
    \item $\frac{1}{76} = \frac{671630}{51043880}$
    \item $\frac{1}{7145} = \frac{7144}{51043880}$
    \item $\frac{1}{51043880} = \frac{1}{51043880}$
\end{itemize}
La somme des numérateurs est :
$$ \frac{1}{76} + \frac{1}{7145} + \frac{1}{51043880} = \frac{671630 + 7144 + 1}{51043880} = \frac{678775}{51043880} $$
Le PGCD du numérateur et du dénominateur est $\text{PGCD}(678775, 51043880) = 135755$.
La fraction irréductible est :
$$ \frac{678775}{51043880} = \frac{678775 \div 135755}{51043880 \div 135755} = \frac{5}{376} $$
Cette fraction est bien égale à $\frac{5}{376}$.

\subsection{Démonstration pour $n = 377$}
Soit $n = 377$. Nous cherchons $x, y, z \in \mathbb{Z}^{+}$ tels que $\frac{5}{377} = \frac{1}{x} + \frac{1}{y} + \frac{1}{z}$.
Posons $x = 76$, $y = 9551$, $z = 273655252$.
Les conditions $x > 0$, $y > 0$ et $z > 0$ sont satisfaites car $x, y, z \geq 1$.
Le PPCM des dénominateurs est $\text{PPCM}(76, 9551, 273655252) = 273655252$.
En réduisant au même dénominateur :
\begin{itemize}
    \item $\frac{1}{76} = \frac{3600727}{273655252}$
    \item $\frac{1}{9551} = \frac{28652}{273655252}$
    \item $\frac{1}{273655252} = \frac{1}{273655252}$
\end{itemize}
La somme des numérateurs est :
$$ \frac{1}{76} + \frac{1}{9551} + \frac{1}{273655252} = \frac{3600727 + 28652 + 1}{273655252} = \frac{3629380}{273655252} $$
Le PGCD du numérateur et du dénominateur est $\text{PGCD}(3629380, 273655252) = 725876$.
La fraction irréductible est :
$$ \frac{3629380}{273655252} = \frac{3629380 \div 725876}{273655252 \div 725876} = \frac{5}{377} $$
Cette fraction est bien égale à $\frac{5}{377}$.

\subsection{Démonstration pour $n = 378$}
Soit $n = 378$. Nous cherchons $x, y, z \in \mathbb{Z}^{+}$ tels que $\frac{5}{378} = \frac{1}{x} + \frac{1}{y} + \frac{1}{z}$.
Posons $x = 76$, $y = 14365$, $z = 206338860$.
Les conditions $x > 0$, $y > 0$ et $z > 0$ sont satisfaites car $x, y, z \geq 1$.
Le PPCM des dénominateurs est $\text{PPCM}(76, 14365, 206338860) = 206338860$.
En réduisant au même dénominateur :
\begin{itemize}
    \item $\frac{1}{76} = \frac{2714985}{206338860}$
    \item $\frac{1}{14365} = \frac{14364}{206338860}$
    \item $\frac{1}{206338860} = \frac{1}{206338860}$
\end{itemize}
La somme des numérateurs est :
$$ \frac{1}{76} + \frac{1}{14365} + \frac{1}{206338860} = \frac{2714985 + 14364 + 1}{206338860} = \frac{2729350}{206338860} $$
Le PGCD du numérateur et du dénominateur est $\text{PGCD}(2729350, 206338860) = 545870$.
La fraction irréductible est :
$$ \frac{2729350}{206338860} = \frac{2729350 \div 545870}{206338860 \div 545870} = \frac{5}{378} $$
Cette fraction est bien égale à $\frac{5}{378}$.

\subsection{Démonstration pour $n = 379$}
Soit $n = 379$. Nous cherchons $x, y, z \in \mathbb{Z}^{+}$ tels que $\frac{5}{379} = \frac{1}{x} + \frac{1}{y} + \frac{1}{z}$.
Posons $x = 76$, $y = 28805$, $z = 829699220$.
Les conditions $x > 0$, $y > 0$ et $z > 0$ sont satisfaites car $x, y, z \geq 1$.
Le PPCM des dénominateurs est $\text{PPCM}(76, 28805, 829699220) = 829699220$.
En réduisant au même dénominateur :
\begin{itemize}
    \item $\frac{1}{76} = \frac{10917095}{829699220}$
    \item $\frac{1}{28805} = \frac{28804}{829699220}$
    \item $\frac{1}{829699220} = \frac{1}{829699220}$
\end{itemize}
La somme des numérateurs est :
$$ \frac{1}{76} + \frac{1}{28805} + \frac{1}{829699220} = \frac{10917095 + 28804 + 1}{829699220} = \frac{10945900}{829699220} $$
Le PGCD du numérateur et du dénominateur est $\text{PGCD}(10945900, 829699220) = 2189180$.
La fraction irréductible est :
$$ \frac{10945900}{829699220} = \frac{10945900 \div 2189180}{829699220 \div 2189180} = \frac{5}{379} $$
Cette fraction est bien égale à $\frac{5}{379}$.

\subsection{Démonstration pour $n = 380$}
Soit $n = 380$. Nous cherchons $x, y, z \in \mathbb{Z}^{+}$ tels que $\frac{5}{380} = \frac{1}{x} + \frac{1}{y} + \frac{1}{z}$.
Posons $x = 77$, $y = 5853$, $z = 34251756$.
Les conditions $x > 0$, $y > 0$ et $z > 0$ sont satisfaites car $x, y, z \geq 1$.
Le PPCM des dénominateurs est $\text{PPCM}(77, 5853, 34251756) = 34251756$.
En réduisant au même dénominateur :
\begin{itemize}
    \item $\frac{1}{77} = \frac{444828}{34251756}$
    \item $\frac{1}{5853} = \frac{5852}{34251756}$
    \item $\frac{1}{34251756} = \frac{1}{34251756}$
\end{itemize}
La somme des numérateurs est :
$$ \frac{1}{77} + \frac{1}{5853} + \frac{1}{34251756} = \frac{444828 + 5852 + 1}{34251756} = \frac{450681}{34251756} $$
Le PGCD du numérateur et du dénominateur est $\text{PGCD}(450681, 34251756) = 450681$.
La fraction irréductible est :
$$ \frac{450681}{34251756} = \frac{450681 \div 450681}{34251756 \div 450681} = \frac{1}{76} $$
Cette fraction est bien égale à $\frac{5}{380}$.

\subsection{Démonstration pour $n = 381$}
Soit $n = 381$. Nous cherchons $x, y, z \in \mathbb{Z}^{+}$ tels que $\frac{5}{381} = \frac{1}{x} + \frac{1}{y} + \frac{1}{z}$.
Posons $x = 77$, $y = 7335$, $z = 71728965$.
Les conditions $x > 0$, $y > 0$ et $z > 0$ sont satisfaites car $x, y, z \geq 1$.
Le PPCM des dénominateurs est $\text{PPCM}(77, 7335, 71728965) = 71728965$.
En réduisant au même dénominateur :
\begin{itemize}
    \item $\frac{1}{77} = \frac{931545}{71728965}$
    \item $\frac{1}{7335} = \frac{9779}{71728965}$
    \item $\frac{1}{71728965} = \frac{1}{71728965}$
\end{itemize}
La somme des numérateurs est :
$$ \frac{1}{77} + \frac{1}{7335} + \frac{1}{71728965} = \frac{931545 + 9779 + 1}{71728965} = \frac{941325}{71728965} $$
Le PGCD du numérateur et du dénominateur est $\text{PGCD}(941325, 71728965) = 188265$.
La fraction irréductible est :
$$ \frac{941325}{71728965} = \frac{941325 \div 188265}{71728965 \div 188265} = \frac{5}{381} $$
Cette fraction est bien égale à $\frac{5}{381}$.

\subsection{Démonstration pour $n = 382$}
Soit $n = 382$. Nous cherchons $x, y, z \in \mathbb{Z}^{+}$ tels que $\frac{5}{382} = \frac{1}{x} + \frac{1}{y} + \frac{1}{z}$.
Posons $x = 77$, $y = 9805$, $z = 288404270$.
Les conditions $x > 0$, $y > 0$ et $z > 0$ sont satisfaites car $x, y, z \geq 1$.
Le PPCM des dénominateurs est $\text{PPCM}(77, 9805, 288404270) = 288404270$.
En réduisant au même dénominateur :
\begin{itemize}
    \item $\frac{1}{77} = \frac{3745510}{288404270}$
    \item $\frac{1}{9805} = \frac{29414}{288404270}$
    \item $\frac{1}{288404270} = \frac{1}{288404270}$
\end{itemize}
La somme des numérateurs est :
$$ \frac{1}{77} + \frac{1}{9805} + \frac{1}{288404270} = \frac{3745510 + 29414 + 1}{288404270} = \frac{3774925}{288404270} $$
Le PGCD du numérateur et du dénominateur est $\text{PGCD}(3774925, 288404270) = 754985$.
La fraction irréductible est :
$$ \frac{3774925}{288404270} = \frac{3774925 \div 754985}{288404270 \div 754985} = \frac{5}{382} $$
Cette fraction est bien égale à $\frac{5}{382}$.

\subsection{Démonstration pour $n = 383$}
Soit $n = 383$. Nous cherchons $x, y, z \in \mathbb{Z}^{+}$ tels que $\frac{5}{383} = \frac{1}{x} + \frac{1}{y} + \frac{1}{z}$.
Posons $x = 77$, $y = 14746$, $z = 434874286$.
Les conditions $x > 0$, $y > 0$ et $z > 0$ sont satisfaites car $x, y, z \geq 1$.
Le PPCM des dénominateurs est $\text{PPCM}(77, 14746, 434874286) = 434874286$.
En réduisant au même dénominateur :
\begin{itemize}
    \item $\frac{1}{77} = \frac{5647718}{434874286}$
    \item $\frac{1}{14746} = \frac{29491}{434874286}$
    \item $\frac{1}{434874286} = \frac{1}{434874286}$
\end{itemize}
La somme des numérateurs est :
$$ \frac{1}{77} + \frac{1}{14746} + \frac{1}{434874286} = \frac{5647718 + 29491 + 1}{434874286} = \frac{5677210}{434874286} $$
Le PGCD du numérateur et du dénominateur est $\text{PGCD}(5677210, 434874286) = 1135442$.
La fraction irréductible est :
$$ \frac{5677210}{434874286} = \frac{5677210 \div 1135442}{434874286 \div 1135442} = \frac{5}{383} $$
Cette fraction est bien égale à $\frac{5}{383}$.

\subsection{Démonstration pour $n = 384$}
Soit $n = 384$. Nous cherchons $x, y, z \in \mathbb{Z}^{+}$ tels que $\frac{5}{384} = \frac{1}{x} + \frac{1}{y} + \frac{1}{z}$.
Posons $x = 77$, $y = 29569$, $z = 874296192$.
Les conditions $x > 0$, $y > 0$ et $z > 0$ sont satisfaites car $x, y, z \geq 1$.
Le PPCM des dénominateurs est $\text{PPCM}(77, 29569, 874296192) = 874296192$.
En réduisant au même dénominateur :
\begin{itemize}
    \item $\frac{1}{77} = \frac{11354496}{874296192}$
    \item $\frac{1}{29569} = \frac{29568}{874296192}$
    \item $\frac{1}{874296192} = \frac{1}{874296192}$
\end{itemize}
La somme des numérateurs est :
$$ \frac{1}{77} + \frac{1}{29569} + \frac{1}{874296192} = \frac{11354496 + 29568 + 1}{874296192} = \frac{11384065}{874296192} $$
Le PGCD du numérateur et du dénominateur est $\text{PGCD}(11384065, 874296192) = 2276813$.
La fraction irréductible est :
$$ \frac{11384065}{874296192} = \frac{11384065 \div 2276813}{874296192 \div 2276813} = \frac{5}{384} $$
Cette fraction est bien égale à $\frac{5}{384}$.

\subsection{Démonstration pour $n = 385$}
Soit $n = 385$. Nous cherchons $x, y, z \in \mathbb{Z}^{+}$ tels que $\frac{5}{385} = \frac{1}{x} + \frac{1}{y} + \frac{1}{z}$.
Posons $x = 78$, $y = 6007$, $z = 36078042$.
Les conditions $x > 0$, $y > 0$ et $z > 0$ sont satisfaites car $x, y, z \geq 1$.
Le PPCM des dénominateurs est $\text{PPCM}(78, 6007, 36078042) = 36078042$.
En réduisant au même dénominateur :
\begin{itemize}
    \item $\frac{1}{78} = \frac{462539}{36078042}$
    \item $\frac{1}{6007} = \frac{6006}{36078042}$
    \item $\frac{1}{36078042} = \frac{1}{36078042}$
\end{itemize}
La somme des numérateurs est :
$$ \frac{1}{78} + \frac{1}{6007} + \frac{1}{36078042} = \frac{462539 + 6006 + 1}{36078042} = \frac{468546}{36078042} $$
Le PGCD du numérateur et du dénominateur est $\text{PGCD}(468546, 36078042) = 468546$.
La fraction irréductible est :
$$ \frac{468546}{36078042} = \frac{468546 \div 468546}{36078042 \div 468546} = \frac{1}{77} $$
Cette fraction est bien égale à $\frac{5}{385}$.

\subsection{Démonstration pour $n = 386$}
Soit $n = 386$. Nous cherchons $x, y, z \in \mathbb{Z}^{+}$ tels que $\frac{5}{386} = \frac{1}{x} + \frac{1}{y} + \frac{1}{z}$.
Posons $x = 78$, $y = 7528$, $z = 56663256$.
Les conditions $x > 0$, $y > 0$ et $z > 0$ sont satisfaites car $x, y, z \geq 1$.
Le PPCM des dénominateurs est $\text{PPCM}(78, 7528, 56663256) = 56663256$.
En réduisant au même dénominateur :
\begin{itemize}
    \item $\frac{1}{78} = \frac{726452}{56663256}$
    \item $\frac{1}{7528} = \frac{7527}{56663256}$
    \item $\frac{1}{56663256} = \frac{1}{56663256}$
\end{itemize}
La somme des numérateurs est :
$$ \frac{1}{78} + \frac{1}{7528} + \frac{1}{56663256} = \frac{726452 + 7527 + 1}{56663256} = \frac{733980}{56663256} $$
Le PGCD du numérateur et du dénominateur est $\text{PGCD}(733980, 56663256) = 146796$.
La fraction irréductible est :
$$ \frac{733980}{56663256} = \frac{733980 \div 146796}{56663256 \div 146796} = \frac{5}{386} $$
Cette fraction est bien égale à $\frac{5}{386}$.

\subsection{Démonstration pour $n = 387$}
Soit $n = 387$. Nous cherchons $x, y, z \in \mathbb{Z}^{+}$ tels que $\frac{5}{387} = \frac{1}{x} + \frac{1}{y} + \frac{1}{z}$.
Posons $x = 78$, $y = 10063$, $z = 101253906$.
Les conditions $x > 0$, $y > 0$ et $z > 0$ sont satisfaites car $x, y, z \geq 1$.
Le PPCM des dénominateurs est $\text{PPCM}(78, 10063, 101253906) = 101253906$.
En réduisant au même dénominateur :
\begin{itemize}
    \item $\frac{1}{78} = \frac{1298127}{101253906}$
    \item $\frac{1}{10063} = \frac{10062}{101253906}$
    \item $\frac{1}{101253906} = \frac{1}{101253906}$
\end{itemize}
La somme des numérateurs est :
$$ \frac{1}{78} + \frac{1}{10063} + \frac{1}{101253906} = \frac{1298127 + 10062 + 1}{101253906} = \frac{1308190}{101253906} $$
Le PGCD du numérateur et du dénominateur est $\text{PGCD}(1308190, 101253906) = 261638$.
La fraction irréductible est :
$$ \frac{1308190}{101253906} = \frac{1308190 \div 261638}{101253906 \div 261638} = \frac{5}{387} $$
Cette fraction est bien égale à $\frac{5}{387}$.

\subsection{Démonstration pour $n = 388$}
Soit $n = 388$. Nous cherchons $x, y, z \in \mathbb{Z}^{+}$ tels que $\frac{5}{388} = \frac{1}{x} + \frac{1}{y} + \frac{1}{z}$.
Posons $x = 78$, $y = 15133$, $z = 228992556$.
Les conditions $x > 0$, $y > 0$ et $z > 0$ sont satisfaites car $x, y, z \geq 1$.
Le PPCM des dénominateurs est $\text{PPCM}(78, 15133, 228992556) = 228992556$.
En réduisant au même dénominateur :
\begin{itemize}
    \item $\frac{1}{78} = \frac{2935802}{228992556}$
    \item $\frac{1}{15133} = \frac{15132}{228992556}$
    \item $\frac{1}{228992556} = \frac{1}{228992556}$
\end{itemize}
La somme des numérateurs est :
$$ \frac{1}{78} + \frac{1}{15133} + \frac{1}{228992556} = \frac{2935802 + 15132 + 1}{228992556} = \frac{2950935}{228992556} $$
Le PGCD du numérateur et du dénominateur est $\text{PGCD}(2950935, 228992556) = 590187$.
La fraction irréductible est :
$$ \frac{2950935}{228992556} = \frac{2950935 \div 590187}{228992556 \div 590187} = \frac{5}{388} $$
Cette fraction est bien égale à $\frac{5}{388}$.

\subsection{Démonstration pour $n = 389$}
Soit $n = 389$. Nous cherchons $x, y, z \in \mathbb{Z}^{+}$ tels que $\frac{5}{389} = \frac{1}{x} + \frac{1}{y} + \frac{1}{z}$.
Posons $x = 78$, $y = 30343$, $z = 920667306$.
Les conditions $x > 0$, $y > 0$ et $z > 0$ sont satisfaites car $x, y, z \geq 1$.
Le PPCM des dénominateurs est $\text{PPCM}(78, 30343, 920667306) = 920667306$.
En réduisant au même dénominateur :
\begin{itemize}
    \item $\frac{1}{78} = \frac{11803427}{920667306}$
    \item $\frac{1}{30343} = \frac{30342}{920667306}$
    \item $\frac{1}{920667306} = \frac{1}{920667306}$
\end{itemize}
La somme des numérateurs est :
$$ \frac{1}{78} + \frac{1}{30343} + \frac{1}{920667306} = \frac{11803427 + 30342 + 1}{920667306} = \frac{11833770}{920667306} $$
Le PGCD du numérateur et du dénominateur est $\text{PGCD}(11833770, 920667306) = 2366754$.
La fraction irréductible est :
$$ \frac{11833770}{920667306} = \frac{11833770 \div 2366754}{920667306 \div 2366754} = \frac{5}{389} $$
Cette fraction est bien égale à $\frac{5}{389}$.

\subsection{Démonstration pour $n = 390$}
Soit $n = 390$. Nous cherchons $x, y, z \in \mathbb{Z}^{+}$ tels que $\frac{5}{390} = \frac{1}{x} + \frac{1}{y} + \frac{1}{z}$.
Posons $x = 79$, $y = 6163$, $z = 37976406$.
Les conditions $x > 0$, $y > 0$ et $z > 0$ sont satisfaites car $x, y, z \geq 1$.
Le PPCM des dénominateurs est $\text{PPCM}(79, 6163, 37976406) = 37976406$.
En réduisant au même dénominateur :
\begin{itemize}
    \item $\frac{1}{79} = \frac{480714}{37976406}$
    \item $\frac{1}{6163} = \frac{6162}{37976406}$
    \item $\frac{1}{37976406} = \frac{1}{37976406}$
\end{itemize}
La somme des numérateurs est :
$$ \frac{1}{79} + \frac{1}{6163} + \frac{1}{37976406} = \frac{480714 + 6162 + 1}{37976406} = \frac{486877}{37976406} $$
Le PGCD du numérateur et du dénominateur est $\text{PGCD}(486877, 37976406) = 486877$.
La fraction irréductible est :
$$ \frac{486877}{37976406} = \frac{486877 \div 486877}{37976406 \div 486877} = \frac{1}{78} $$
Cette fraction est bien égale à $\frac{5}{390}$.

\subsection{Démonstration pour $n = 391$}
Soit $n = 391$. Nous cherchons $x, y, z \in \mathbb{Z}^{+}$ tels que $\frac{5}{391} = \frac{1}{x} + \frac{1}{y} + \frac{1}{z}$.
Posons $x = 79$, $y = 7728$, $z = 10378704$.
Les conditions $x > 0$, $y > 0$ et $z > 0$ sont satisfaites car $x, y, z \geq 1$.
Le PPCM des dénominateurs est $\text{PPCM}(79, 7728, 10378704) = 10378704$.
En réduisant au même dénominateur :
\begin{itemize}
    \item $\frac{1}{79} = \frac{131376}{10378704}$
    \item $\frac{1}{7728} = \frac{1343}{10378704}$
    \item $\frac{1}{10378704} = \frac{1}{10378704}$
\end{itemize}
La somme des numérateurs est :
$$ \frac{1}{79} + \frac{1}{7728} + \frac{1}{10378704} = \frac{131376 + 1343 + 1}{10378704} = \frac{132720}{10378704} $$
Le PGCD du numérateur et du dénominateur est $\text{PGCD}(132720, 10378704) = 26544$.
La fraction irréductible est :
$$ \frac{132720}{10378704} = \frac{132720 \div 26544}{10378704 \div 26544} = \frac{5}{391} $$
Cette fraction est bien égale à $\frac{5}{391}$.

\subsection{Démonstration pour $n = 392$}
Soit $n = 392$. Nous cherchons $x, y, z \in \mathbb{Z}^{+}$ tels que $\frac{5}{392} = \frac{1}{x} + \frac{1}{y} + \frac{1}{z}$.
Posons $x = 79$, $y = 10323$, $z = 319682664$.
Les conditions $x > 0$, $y > 0$ et $z > 0$ sont satisfaites car $x, y, z \geq 1$.
Le PPCM des dénominateurs est $\text{PPCM}(79, 10323, 319682664) = 319682664$.
En réduisant au même dénominateur :
\begin{itemize}
    \item $\frac{1}{79} = \frac{4046616}{319682664}$
    \item $\frac{1}{10323} = \frac{30968}{319682664}$
    \item $\frac{1}{319682664} = \frac{1}{319682664}$
\end{itemize}
La somme des numérateurs est :
$$ \frac{1}{79} + \frac{1}{10323} + \frac{1}{319682664} = \frac{4046616 + 30968 + 1}{319682664} = \frac{4077585}{319682664} $$
Le PGCD du numérateur et du dénominateur est $\text{PGCD}(4077585, 319682664) = 815517$.
La fraction irréductible est :
$$ \frac{4077585}{319682664} = \frac{4077585 \div 815517}{319682664 \div 815517} = \frac{5}{392} $$
Cette fraction est bien égale à $\frac{5}{392}$.

\subsection{Démonstration pour $n = 393$}
Soit $n = 393$. Nous cherchons $x, y, z \in \mathbb{Z}^{+}$ tels que $\frac{5}{393} = \frac{1}{x} + \frac{1}{y} + \frac{1}{z}$.
Posons $x = 79$, $y = 15524$, $z = 481973628$.
Les conditions $x > 0$, $y > 0$ et $z > 0$ sont satisfaites car $x, y, z \geq 1$.
Le PPCM des dénominateurs est $\text{PPCM}(79, 15524, 481973628) = 481973628$.
En réduisant au même dénominateur :
\begin{itemize}
    \item $\frac{1}{79} = \frac{6100932}{481973628}$
    \item $\frac{1}{15524} = \frac{31047}{481973628}$
    \item $\frac{1}{481973628} = \frac{1}{481973628}$
\end{itemize}
La somme des numérateurs est :
$$ \frac{1}{79} + \frac{1}{15524} + \frac{1}{481973628} = \frac{6100932 + 31047 + 1}{481973628} = \frac{6131980}{481973628} $$
Le PGCD du numérateur et du dénominateur est $\text{PGCD}(6131980, 481973628) = 1226396$.
La fraction irréductible est :
$$ \frac{6131980}{481973628} = \frac{6131980 \div 1226396}{481973628 \div 1226396} = \frac{5}{393} $$
Cette fraction est bien égale à $\frac{5}{393}$.

\subsection{Démonstration pour $n = 394$}
Soit $n = 394$. Nous cherchons $x, y, z \in \mathbb{Z}^{+}$ tels que $\frac{5}{394} = \frac{1}{x} + \frac{1}{y} + \frac{1}{z}$.
Posons $x = 79$, $y = 31127$, $z = 968859002$.
Les conditions $x > 0$, $y > 0$ et $z > 0$ sont satisfaites car $x, y, z \geq 1$.
Le PPCM des dénominateurs est $\text{PPCM}(79, 31127, 968859002) = 968859002$.
En réduisant au même dénominateur :
\begin{itemize}
    \item $\frac{1}{79} = \frac{12264038}{968859002}$
    \item $\frac{1}{31127} = \frac{31126}{968859002}$
    \item $\frac{1}{968859002} = \frac{1}{968859002}$
\end{itemize}
La somme des numérateurs est :
$$ \frac{1}{79} + \frac{1}{31127} + \frac{1}{968859002} = \frac{12264038 + 31126 + 1}{968859002} = \frac{12295165}{968859002} $$
Le PGCD du numérateur et du dénominateur est $\text{PGCD}(12295165, 968859002) = 2459033$.
La fraction irréductible est :
$$ \frac{12295165}{968859002} = \frac{12295165 \div 2459033}{968859002 \div 2459033} = \frac{5}{394} $$
Cette fraction est bien égale à $\frac{5}{394}$.

\subsection{Démonstration pour $n = 395$}
Soit $n = 395$. Nous cherchons $x, y, z \in \mathbb{Z}^{+}$ tels que $\frac{5}{395} = \frac{1}{x} + \frac{1}{y} + \frac{1}{z}$.
Posons $x = 80$, $y = 6321$, $z = 39948720$.
Les conditions $x > 0$, $y > 0$ et $z > 0$ sont satisfaites car $x, y, z \geq 1$.
Le PPCM des dénominateurs est $\text{PPCM}(80, 6321, 39948720) = 39948720$.
En réduisant au même dénominateur :
\begin{itemize}
    \item $\frac{1}{80} = \frac{499359}{39948720}$
    \item $\frac{1}{6321} = \frac{6320}{39948720}$
    \item $\frac{1}{39948720} = \frac{1}{39948720}$
\end{itemize}
La somme des numérateurs est :
$$ \frac{1}{80} + \frac{1}{6321} + \frac{1}{39948720} = \frac{499359 + 6320 + 1}{39948720} = \frac{505680}{39948720} $$
Le PGCD du numérateur et du dénominateur est $\text{PGCD}(505680, 39948720) = 505680$.
La fraction irréductible est :
$$ \frac{505680}{39948720} = \frac{505680 \div 505680}{39948720 \div 505680} = \frac{1}{79} $$
Cette fraction est bien égale à $\frac{5}{395}$.

\subsection{Démonstration pour $n = 396$}
Soit $n = 396$. Nous cherchons $x, y, z \in \mathbb{Z}^{+}$ tels que $\frac{5}{396} = \frac{1}{x} + \frac{1}{y} + \frac{1}{z}$.
Posons $x = 80$, $y = 7921$, $z = 62734320$.
Les conditions $x > 0$, $y > 0$ et $z > 0$ sont satisfaites car $x, y, z \geq 1$.
Le PPCM des dénominateurs est $\text{PPCM}(80, 7921, 62734320) = 62734320$.
En réduisant au même dénominateur :
\begin{itemize}
    \item $\frac{1}{80} = \frac{784179}{62734320}$
    \item $\frac{1}{7921} = \frac{7920}{62734320}$
    \item $\frac{1}{62734320} = \frac{1}{62734320}$
\end{itemize}
La somme des numérateurs est :
$$ \frac{1}{80} + \frac{1}{7921} + \frac{1}{62734320} = \frac{784179 + 7920 + 1}{62734320} = \frac{792100}{62734320} $$
Le PGCD du numérateur et du dénominateur est $\text{PGCD}(792100, 62734320) = 158420$.
La fraction irréductible est :
$$ \frac{792100}{62734320} = \frac{792100 \div 158420}{62734320 \div 158420} = \frac{5}{396} $$
Cette fraction est bien égale à $\frac{5}{396}$.

\subsection{Démonstration pour $n = 397$}
Soit $n = 397$. Nous cherchons $x, y, z \in \mathbb{Z}^{+}$ tels que $\frac{5}{397} = \frac{1}{x} + \frac{1}{y} + \frac{1}{z}$.
Posons $x = 80$, $y = 10587$, $z = 336243120$.
Les conditions $x > 0$, $y > 0$ et $z > 0$ sont satisfaites car $x, y, z \geq 1$.
Le PPCM des dénominateurs est $\text{PPCM}(80, 10587, 336243120) = 336243120$.
En réduisant au même dénominateur :
\begin{itemize}
    \item $\frac{1}{80} = \frac{4203039}{336243120}$
    \item $\frac{1}{10587} = \frac{31760}{336243120}$
    \item $\frac{1}{336243120} = \frac{1}{336243120}$
\end{itemize}
La somme des numérateurs est :
$$ \frac{1}{80} + \frac{1}{10587} + \frac{1}{336243120} = \frac{4203039 + 31760 + 1}{336243120} = \frac{4234800}{336243120} $$
Le PGCD du numérateur et du dénominateur est $\text{PGCD}(4234800, 336243120) = 846960$.
La fraction irréductible est :
$$ \frac{4234800}{336243120} = \frac{4234800 \div 846960}{336243120 \div 846960} = \frac{5}{397} $$
Cette fraction est bien égale à $\frac{5}{397}$.

\subsection{Démonstration pour $n = 398$}
Soit $n = 398$. Nous cherchons $x, y, z \in \mathbb{Z}^{+}$ tels que $\frac{5}{398} = \frac{1}{x} + \frac{1}{y} + \frac{1}{z}$.
Posons $x = 80$, $y = 15921$, $z = 253462320$.
Les conditions $x > 0$, $y > 0$ et $z > 0$ sont satisfaites car $x, y, z \geq 1$.
Le PPCM des dénominateurs est $\text{PPCM}(80, 15921, 253462320) = 253462320$.
En réduisant au même dénominateur :
\begin{itemize}
    \item $\frac{1}{80} = \frac{3168279}{253462320}$
    \item $\frac{1}{15921} = \frac{15920}{253462320}$
    \item $\frac{1}{253462320} = \frac{1}{253462320}$
\end{itemize}
La somme des numérateurs est :
$$ \frac{1}{80} + \frac{1}{15921} + \frac{1}{253462320} = \frac{3168279 + 15920 + 1}{253462320} = \frac{3184200}{253462320} $$
Le PGCD du numérateur et du dénominateur est $\text{PGCD}(3184200, 253462320) = 636840$.
La fraction irréductible est :
$$ \frac{3184200}{253462320} = \frac{3184200 \div 636840}{253462320 \div 636840} = \frac{5}{398} $$
Cette fraction est bien égale à $\frac{5}{398}$.

\subsection{Démonstration pour $n = 399$}
Soit $n = 399$. Nous cherchons $x, y, z \in \mathbb{Z}^{+}$ tels que $\frac{5}{399} = \frac{1}{x} + \frac{1}{y} + \frac{1}{z}$.
Posons $x = 80$, $y = 31921$, $z = 1018918320$.
Les conditions $x > 0$, $y > 0$ et $z > 0$ sont satisfaites car $x, y, z \geq 1$.
Le PPCM des dénominateurs est $\text{PPCM}(80, 31921, 1018918320) = 1018918320$.
En réduisant au même dénominateur :
\begin{itemize}
    \item $\frac{1}{80} = \frac{12736479}{1018918320}$
    \item $\frac{1}{31921} = \frac{31920}{1018918320}$
    \item $\frac{1}{1018918320} = \frac{1}{1018918320}$
\end{itemize}
La somme des numérateurs est :
$$ \frac{1}{80} + \frac{1}{31921} + \frac{1}{1018918320} = \frac{12736479 + 31920 + 1}{1018918320} = \frac{12768400}{1018918320} $$
Le PGCD du numérateur et du dénominateur est $\text{PGCD}(12768400, 1018918320) = 2553680$.
La fraction irréductible est :
$$ \frac{12768400}{1018918320} = \frac{12768400 \div 2553680}{1018918320 \div 2553680} = \frac{5}{399} $$
Cette fraction est bien égale à $\frac{5}{399}$.

\subsection{Démonstration pour $n = 400$}
Soit $n = 400$. Nous cherchons $x, y, z \in \mathbb{Z}^{+}$ tels que $\frac{5}{400} = \frac{1}{x} + \frac{1}{y} + \frac{1}{z}$.
Posons $x = 81$, $y = 6481$, $z = 41996880$.
Les conditions $x > 0$, $y > 0$ et $z > 0$ sont satisfaites car $x, y, z \geq 1$.
Le PPCM des dénominateurs est $\text{PPCM}(81, 6481, 41996880) = 41996880$.
En réduisant au même dénominateur :
\begin{itemize}
    \item $\frac{1}{81} = \frac{518480}{41996880}$
    \item $\frac{1}{6481} = \frac{6480}{41996880}$
    \item $\frac{1}{41996880} = \frac{1}{41996880}$
\end{itemize}
La somme des numérateurs est :
$$ \frac{1}{81} + \frac{1}{6481} + \frac{1}{41996880} = \frac{518480 + 6480 + 1}{41996880} = \frac{524961}{41996880} $$
Le PGCD du numérateur et du dénominateur est $\text{PGCD}(524961, 41996880) = 524961$.
La fraction irréductible est :
$$ \frac{524961}{41996880} = \frac{524961 \div 524961}{41996880 \div 524961} = \frac{1}{80} $$
Cette fraction est bien égale à $\frac{5}{400}$.

\subsection{Démonstration pour $n = 401$}
Soit $n = 401$. Nous cherchons $x, y, z \in \mathbb{Z}^{+}$ tels que $\frac{5}{401} = \frac{1}{x} + \frac{1}{y} + \frac{1}{z}$.
Posons $x = 81$, $y = 8121$, $z = 87926067$.
Les conditions $x > 0$, $y > 0$ et $z > 0$ sont satisfaites car $x, y, z \geq 1$.
Le PPCM des dénominateurs est $\text{PPCM}(81, 8121, 87926067) = 87926067$.
En réduisant au même dénominateur :
\begin{itemize}
    \item $\frac{1}{81} = \frac{1085507}{87926067}$
    \item $\frac{1}{8121} = \frac{10827}{87926067}$
    \item $\frac{1}{87926067} = \frac{1}{87926067}$
\end{itemize}
La somme des numérateurs est :
$$ \frac{1}{81} + \frac{1}{8121} + \frac{1}{87926067} = \frac{1085507 + 10827 + 1}{87926067} = \frac{1096335}{87926067} $$
Le PGCD du numérateur et du dénominateur est $\text{PGCD}(1096335, 87926067) = 219267$.
La fraction irréductible est :
$$ \frac{1096335}{87926067} = \frac{1096335 \div 219267}{87926067 \div 219267} = \frac{5}{401} $$
Cette fraction est bien égale à $\frac{5}{401}$.

\subsection{Démonstration pour $n = 402$}
Soit $n = 402$. Nous cherchons $x, y, z \in \mathbb{Z}^{+}$ tels que $\frac{5}{402} = \frac{1}{x} + \frac{1}{y} + \frac{1}{z}$.
Posons $x = 81$, $y = 10855$, $z = 117820170$.
Les conditions $x > 0$, $y > 0$ et $z > 0$ sont satisfaites car $x, y, z \geq 1$.
Le PPCM des dénominateurs est $\text{PPCM}(81, 10855, 117820170) = 117820170$.
En réduisant au même dénominateur :
\begin{itemize}
    \item $\frac{1}{81} = \frac{1454570}{117820170}$
    \item $\frac{1}{10855} = \frac{10854}{117820170}$
    \item $\frac{1}{117820170} = \frac{1}{117820170}$
\end{itemize}
La somme des numérateurs est :
$$ \frac{1}{81} + \frac{1}{10855} + \frac{1}{117820170} = \frac{1454570 + 10854 + 1}{117820170} = \frac{1465425}{117820170} $$
Le PGCD du numérateur et du dénominateur est $\text{PGCD}(1465425, 117820170) = 293085$.
La fraction irréductible est :
$$ \frac{1465425}{117820170} = \frac{1465425 \div 293085}{117820170 \div 293085} = \frac{5}{402} $$
Cette fraction est bien égale à $\frac{5}{402}$.

\subsection{Démonstration pour $n = 403$}
Soit $n = 403$. Nous cherchons $x, y, z \in \mathbb{Z}^{+}$ tels que $\frac{5}{403} = \frac{1}{x} + \frac{1}{y} + \frac{1}{z}$.
Posons $x = 81$, $y = 16322$, $z = 532799046$.
Les conditions $x > 0$, $y > 0$ et $z > 0$ sont satisfaites car $x, y, z \geq 1$.
Le PPCM des dénominateurs est $\text{PPCM}(81, 16322, 532799046) = 532799046$.
En réduisant au même dénominateur :
\begin{itemize}
    \item $\frac{1}{81} = \frac{6577766}{532799046}$
    \item $\frac{1}{16322} = \frac{32643}{532799046}$
    \item $\frac{1}{532799046} = \frac{1}{532799046}$
\end{itemize}
La somme des numérateurs est :
$$ \frac{1}{81} + \frac{1}{16322} + \frac{1}{532799046} = \frac{6577766 + 32643 + 1}{532799046} = \frac{6610410}{532799046} $$
Le PGCD du numérateur et du dénominateur est $\text{PGCD}(6610410, 532799046) = 1322082$.
La fraction irréductible est :
$$ \frac{6610410}{532799046} = \frac{6610410 \div 1322082}{532799046 \div 1322082} = \frac{5}{403} $$
Cette fraction est bien égale à $\frac{5}{403}$.

\subsection{Démonstration pour $n = 404$}
Soit $n = 404$. Nous cherchons $x, y, z \in \mathbb{Z}^{+}$ tels que $\frac{5}{404} = \frac{1}{x} + \frac{1}{y} + \frac{1}{z}$.
Posons $x = 81$, $y = 32725$, $z = 1070892900$.
Les conditions $x > 0$, $y > 0$ et $z > 0$ sont satisfaites car $x, y, z \geq 1$.
Le PPCM des dénominateurs est $\text{PPCM}(81, 32725, 1070892900) = 1070892900$.
En réduisant au même dénominateur :
\begin{itemize}
    \item $\frac{1}{81} = \frac{13220900}{1070892900}$
    \item $\frac{1}{32725} = \frac{32724}{1070892900}$
    \item $\frac{1}{1070892900} = \frac{1}{1070892900}$
\end{itemize}
La somme des numérateurs est :
$$ \frac{1}{81} + \frac{1}{32725} + \frac{1}{1070892900} = \frac{13220900 + 32724 + 1}{1070892900} = \frac{13253625}{1070892900} $$
Le PGCD du numérateur et du dénominateur est $\text{PGCD}(13253625, 1070892900) = 2650725$.
La fraction irréductible est :
$$ \frac{13253625}{1070892900} = \frac{13253625 \div 2650725}{1070892900 \div 2650725} = \frac{5}{404} $$
Cette fraction est bien égale à $\frac{5}{404}$.

\subsection{Démonstration pour $n = 405$}
Soit $n = 405$. Nous cherchons $x, y, z \in \mathbb{Z}^{+}$ tels que $\frac{5}{405} = \frac{1}{x} + \frac{1}{y} + \frac{1}{z}$.
Posons $x = 82$, $y = 6643$, $z = 44122806$.
Les conditions $x > 0$, $y > 0$ et $z > 0$ sont satisfaites car $x, y, z \geq 1$.
Le PPCM des dénominateurs est $\text{PPCM}(82, 6643, 44122806) = 44122806$.
En réduisant au même dénominateur :
\begin{itemize}
    \item $\frac{1}{82} = \frac{538083}{44122806}$
    \item $\frac{1}{6643} = \frac{6642}{44122806}$
    \item $\frac{1}{44122806} = \frac{1}{44122806}$
\end{itemize}
La somme des numérateurs est :
$$ \frac{1}{82} + \frac{1}{6643} + \frac{1}{44122806} = \frac{538083 + 6642 + 1}{44122806} = \frac{544726}{44122806} $$
Le PGCD du numérateur et du dénominateur est $\text{PGCD}(544726, 44122806) = 544726$.
La fraction irréductible est :
$$ \frac{544726}{44122806} = \frac{544726 \div 544726}{44122806 \div 544726} = \frac{1}{81} $$
Cette fraction est bien égale à $\frac{5}{405}$.

\subsection{Démonstration pour $n = 406$}
Soit $n = 406$. Nous cherchons $x, y, z \in \mathbb{Z}^{+}$ tels que $\frac{5}{406} = \frac{1}{x} + \frac{1}{y} + \frac{1}{z}$.
Posons $x = 82$, $y = 8324$, $z = 69280652$.
Les conditions $x > 0$, $y > 0$ et $z > 0$ sont satisfaites car $x, y, z \geq 1$.
Le PPCM des dénominateurs est $\text{PPCM}(82, 8324, 69280652) = 69280652$.
En réduisant au même dénominateur :
\begin{itemize}
    \item $\frac{1}{82} = \frac{844886}{69280652}$
    \item $\frac{1}{8324} = \frac{8323}{69280652}$
    \item $\frac{1}{69280652} = \frac{1}{69280652}$
\end{itemize}
La somme des numérateurs est :
$$ \frac{1}{82} + \frac{1}{8324} + \frac{1}{69280652} = \frac{844886 + 8323 + 1}{69280652} = \frac{853210}{69280652} $$
Le PGCD du numérateur et du dénominateur est $\text{PGCD}(853210, 69280652) = 170642$.
La fraction irréductible est :
$$ \frac{853210}{69280652} = \frac{853210 \div 170642}{69280652 \div 170642} = \frac{5}{406} $$
Cette fraction est bien égale à $\frac{5}{406}$.

\subsection{Démonstration pour $n = 407$}
Soit $n = 407$. Nous cherchons $x, y, z \in \mathbb{Z}^{+}$ tels que $\frac{5}{407} = \frac{1}{x} + \frac{1}{y} + \frac{1}{z}$.
Posons $x = 82$, $y = 11125$, $z = 371285750$.
Les conditions $x > 0$, $y > 0$ et $z > 0$ sont satisfaites car $x, y, z \geq 1$.
Le PPCM des dénominateurs est $\text{PPCM}(82, 11125, 371285750) = 371285750$.
En réduisant au même dénominateur :
\begin{itemize}
    \item $\frac{1}{82} = \frac{4527875}{371285750}$
    \item $\frac{1}{11125} = \frac{33374}{371285750}$
    \item $\frac{1}{371285750} = \frac{1}{371285750}$
\end{itemize}
La somme des numérateurs est :
$$ \frac{1}{82} + \frac{1}{11125} + \frac{1}{371285750} = \frac{4527875 + 33374 + 1}{371285750} = \frac{4561250}{371285750} $$
Le PGCD du numérateur et du dénominateur est $\text{PGCD}(4561250, 371285750) = 912250$.
La fraction irréductible est :
$$ \frac{4561250}{371285750} = \frac{4561250 \div 912250}{371285750 \div 912250} = \frac{5}{407} $$
Cette fraction est bien égale à $\frac{5}{407}$.

\subsection{Démonstration pour $n = 408$}
Soit $n = 408$. Nous cherchons $x, y, z \in \mathbb{Z}^{+}$ tels que $\frac{5}{408} = \frac{1}{x} + \frac{1}{y} + \frac{1}{z}$.
Posons $x = 82$, $y = 16729$, $z = 279842712$.
Les conditions $x > 0$, $y > 0$ et $z > 0$ sont satisfaites car $x, y, z \geq 1$.
Le PPCM des dénominateurs est $\text{PPCM}(82, 16729, 279842712) = 279842712$.
En réduisant au même dénominateur :
\begin{itemize}
    \item $\frac{1}{82} = \frac{3412716}{279842712}$
    \item $\frac{1}{16729} = \frac{16728}{279842712}$
    \item $\frac{1}{279842712} = \frac{1}{279842712}$
\end{itemize}
La somme des numérateurs est :
$$ \frac{1}{82} + \frac{1}{16729} + \frac{1}{279842712} = \frac{3412716 + 16728 + 1}{279842712} = \frac{3429445}{279842712} $$
Le PGCD du numérateur et du dénominateur est $\text{PGCD}(3429445, 279842712) = 685889$.
La fraction irréductible est :
$$ \frac{3429445}{279842712} = \frac{3429445 \div 685889}{279842712 \div 685889} = \frac{5}{408} $$
Cette fraction est bien égale à $\frac{5}{408}$.

\subsection{Démonstration pour $n = 409$}
Soit $n = 409$. Nous cherchons $x, y, z \in \mathbb{Z}^{+}$ tels que $\frac{5}{409} = \frac{1}{x} + \frac{1}{y} + \frac{1}{z}$.
Posons $x = 82$, $y = 33539$, $z = 1124830982$.
Les conditions $x > 0$, $y > 0$ et $z > 0$ sont satisfaites car $x, y, z \geq 1$.
Le PPCM des dénominateurs est $\text{PPCM}(82, 33539, 1124830982) = 1124830982$.
En réduisant au même dénominateur :
\begin{itemize}
    \item $\frac{1}{82} = \frac{13717451}{1124830982}$
    \item $\frac{1}{33539} = \frac{33538}{1124830982}$
    \item $\frac{1}{1124830982} = \frac{1}{1124830982}$
\end{itemize}
La somme des numérateurs est :
$$ \frac{1}{82} + \frac{1}{33539} + \frac{1}{1124830982} = \frac{13717451 + 33538 + 1}{1124830982} = \frac{13750990}{1124830982} $$
Le PGCD du numérateur et du dénominateur est $\text{PGCD}(13750990, 1124830982) = 2750198$.
La fraction irréductible est :
$$ \frac{13750990}{1124830982} = \frac{13750990 \div 2750198}{1124830982 \div 2750198} = \frac{5}{409} $$
Cette fraction est bien égale à $\frac{5}{409}$.

\subsection{Démonstration pour $n = 410$}
Soit $n = 410$. Nous cherchons $x, y, z \in \mathbb{Z}^{+}$ tels que $\frac{5}{410} = \frac{1}{x} + \frac{1}{y} + \frac{1}{z}$.
Posons $x = 83$, $y = 6807$, $z = 46328442$.
Les conditions $x > 0$, $y > 0$ et $z > 0$ sont satisfaites car $x, y, z \geq 1$.
Le PPCM des dénominateurs est $\text{PPCM}(83, 6807, 46328442) = 46328442$.
En réduisant au même dénominateur :
\begin{itemize}
    \item $\frac{1}{83} = \frac{558174}{46328442}$
    \item $\frac{1}{6807} = \frac{6806}{46328442}$
    \item $\frac{1}{46328442} = \frac{1}{46328442}$
\end{itemize}
La somme des numérateurs est :
$$ \frac{1}{83} + \frac{1}{6807} + \frac{1}{46328442} = \frac{558174 + 6806 + 1}{46328442} = \frac{564981}{46328442} $$
Le PGCD du numérateur et du dénominateur est $\text{PGCD}(564981, 46328442) = 564981$.
La fraction irréductible est :
$$ \frac{564981}{46328442} = \frac{564981 \div 564981}{46328442 \div 564981} = \frac{1}{82} $$
Cette fraction est bien égale à $\frac{5}{410}$.

\subsection{Démonstration pour $n = 411$}
Soit $n = 411$. Nous cherchons $x, y, z \in \mathbb{Z}^{+}$ tels que $\frac{5}{411} = \frac{1}{x} + \frac{1}{y} + \frac{1}{z}$.
Posons $x = 83$, $y = 8529$, $z = 96983259$.
Les conditions $x > 0$, $y > 0$ et $z > 0$ sont satisfaites car $x, y, z \geq 1$.
Le PPCM des dénominateurs est $\text{PPCM}(83, 8529, 96983259) = 96983259$.
En réduisant au même dénominateur :
\begin{itemize}
    \item $\frac{1}{83} = \frac{1168473}{96983259}$
    \item $\frac{1}{8529} = \frac{11371}{96983259}$
    \item $\frac{1}{96983259} = \frac{1}{96983259}$
\end{itemize}
La somme des numérateurs est :
$$ \frac{1}{83} + \frac{1}{8529} + \frac{1}{96983259} = \frac{1168473 + 11371 + 1}{96983259} = \frac{1179845}{96983259} $$
Le PGCD du numérateur et du dénominateur est $\text{PGCD}(1179845, 96983259) = 235969$.
La fraction irréductible est :
$$ \frac{1179845}{96983259} = \frac{1179845 \div 235969}{96983259 \div 235969} = \frac{5}{411} $$
Cette fraction est bien égale à $\frac{5}{411}$.

\subsection{Démonstration pour $n = 412$}
Soit $n = 412$. Nous cherchons $x, y, z \in \mathbb{Z}^{+}$ tels que $\frac{5}{412} = \frac{1}{x} + \frac{1}{y} + \frac{1}{z}$.
Posons $x = 83$, $y = 11399$, $z = 389800204$.
Les conditions $x > 0$, $y > 0$ et $z > 0$ sont satisfaites car $x, y, z \geq 1$.
Le PPCM des dénominateurs est $\text{PPCM}(83, 11399, 389800204) = 389800204$.
En réduisant au même dénominateur :
\begin{itemize}
    \item $\frac{1}{83} = \frac{4696388}{389800204}$
    \item $\frac{1}{11399} = \frac{34196}{389800204}$
    \item $\frac{1}{389800204} = \frac{1}{389800204}$
\end{itemize}
La somme des numérateurs est :
$$ \frac{1}{83} + \frac{1}{11399} + \frac{1}{389800204} = \frac{4696388 + 34196 + 1}{389800204} = \frac{4730585}{389800204} $$
Le PGCD du numérateur et du dénominateur est $\text{PGCD}(4730585, 389800204) = 946117$.
La fraction irréductible est :
$$ \frac{4730585}{389800204} = \frac{4730585 \div 946117}{389800204 \div 946117} = \frac{5}{412} $$
Cette fraction est bien égale à $\frac{5}{412}$.

\subsection{Démonstration pour $n = 413$}
Soit $n = 413$. Nous cherchons $x, y, z \in \mathbb{Z}^{+}$ tels que $\frac{5}{413} = \frac{1}{x} + \frac{1}{y} + \frac{1}{z}$.
Posons $x = 83$, $y = 17140$, $z = 587542060$.
Les conditions $x > 0$, $y > 0$ et $z > 0$ sont satisfaites car $x, y, z \geq 1$.
Le PPCM des dénominateurs est $\text{PPCM}(83, 17140, 587542060) = 587542060$.
En réduisant au même dénominateur :
\begin{itemize}
    \item $\frac{1}{83} = \frac{7078820}{587542060}$
    \item $\frac{1}{17140} = \frac{34279}{587542060}$
    \item $\frac{1}{587542060} = \frac{1}{587542060}$
\end{itemize}
La somme des numérateurs est :
$$ \frac{1}{83} + \frac{1}{17140} + \frac{1}{587542060} = \frac{7078820 + 34279 + 1}{587542060} = \frac{7113100}{587542060} $$
Le PGCD du numérateur et du dénominateur est $\text{PGCD}(7113100, 587542060) = 1422620$.
La fraction irréductible est :
$$ \frac{7113100}{587542060} = \frac{7113100 \div 1422620}{587542060 \div 1422620} = \frac{5}{413} $$
Cette fraction est bien égale à $\frac{5}{413}$.

\subsection{Démonstration pour $n = 414$}
Soit $n = 414$. Nous cherchons $x, y, z \in \mathbb{Z}^{+}$ tels que $\frac{5}{414} = \frac{1}{x} + \frac{1}{y} + \frac{1}{z}$.
Posons $x = 83$, $y = 34363$, $z = 1180781406$.
Les conditions $x > 0$, $y > 0$ et $z > 0$ sont satisfaites car $x, y, z \geq 1$.
Le PPCM des dénominateurs est $\text{PPCM}(83, 34363, 1180781406) = 1180781406$.
En réduisant au même dénominateur :
\begin{itemize}
    \item $\frac{1}{83} = \frac{14226282}{1180781406}$
    \item $\frac{1}{34363} = \frac{34362}{1180781406}$
    \item $\frac{1}{1180781406} = \frac{1}{1180781406}$
\end{itemize}
La somme des numérateurs est :
$$ \frac{1}{83} + \frac{1}{34363} + \frac{1}{1180781406} = \frac{14226282 + 34362 + 1}{1180781406} = \frac{14260645}{1180781406} $$
Le PGCD du numérateur et du dénominateur est $\text{PGCD}(14260645, 1180781406) = 2852129$.
La fraction irréductible est :
$$ \frac{14260645}{1180781406} = \frac{14260645 \div 2852129}{1180781406 \div 2852129} = \frac{5}{414} $$
Cette fraction est bien égale à $\frac{5}{414}$.

\subsection{Démonstration pour $n = 415$}
Soit $n = 415$. Nous cherchons $x, y, z \in \mathbb{Z}^{+}$ tels que $\frac{5}{415} = \frac{1}{x} + \frac{1}{y} + \frac{1}{z}$.
Posons $x = 84$, $y = 6973$, $z = 48615756$.
Les conditions $x > 0$, $y > 0$ et $z > 0$ sont satisfaites car $x, y, z \geq 1$.
Le PPCM des dénominateurs est $\text{PPCM}(84, 6973, 48615756) = 48615756$.
En réduisant au même dénominateur :
\begin{itemize}
    \item $\frac{1}{84} = \frac{578759}{48615756}$
    \item $\frac{1}{6973} = \frac{6972}{48615756}$
    \item $\frac{1}{48615756} = \frac{1}{48615756}$
\end{itemize}
La somme des numérateurs est :
$$ \frac{1}{84} + \frac{1}{6973} + \frac{1}{48615756} = \frac{578759 + 6972 + 1}{48615756} = \frac{585732}{48615756} $$
Le PGCD du numérateur et du dénominateur est $\text{PGCD}(585732, 48615756) = 585732$.
La fraction irréductible est :
$$ \frac{585732}{48615756} = \frac{585732 \div 585732}{48615756 \div 585732} = \frac{1}{83} $$
Cette fraction est bien égale à $\frac{5}{415}$.

\subsection{Démonstration pour $n = 416$}
Soit $n = 416$. Nous cherchons $x, y, z \in \mathbb{Z}^{+}$ tels que $\frac{5}{416} = \frac{1}{x} + \frac{1}{y} + \frac{1}{z}$.
Posons $x = 84$, $y = 8737$, $z = 76326432$.
Les conditions $x > 0$, $y > 0$ et $z > 0$ sont satisfaites car $x, y, z \geq 1$.
Le PPCM des dénominateurs est $\text{PPCM}(84, 8737, 76326432) = 76326432$.
En réduisant au même dénominateur :
\begin{itemize}
    \item $\frac{1}{84} = \frac{908648}{76326432}$
    \item $\frac{1}{8737} = \frac{8736}{76326432}$
    \item $\frac{1}{76326432} = \frac{1}{76326432}$
\end{itemize}
La somme des numérateurs est :
$$ \frac{1}{84} + \frac{1}{8737} + \frac{1}{76326432} = \frac{908648 + 8736 + 1}{76326432} = \frac{917385}{76326432} $$
Le PGCD du numérateur et du dénominateur est $\text{PGCD}(917385, 76326432) = 183477$.
La fraction irréductible est :
$$ \frac{917385}{76326432} = \frac{917385 \div 183477}{76326432 \div 183477} = \frac{5}{416} $$
Cette fraction est bien égale à $\frac{5}{416}$.

\subsection{Démonstration pour $n = 417$}
Soit $n = 417$. Nous cherchons $x, y, z \in \mathbb{Z}^{+}$ tels que $\frac{5}{417} = \frac{1}{x} + \frac{1}{y} + \frac{1}{z}$.
Posons $x = 84$, $y = 11677$, $z = 136340652$.
Les conditions $x > 0$, $y > 0$ et $z > 0$ sont satisfaites car $x, y, z \geq 1$.
Le PPCM des dénominateurs est $\text{PPCM}(84, 11677, 136340652) = 136340652$.
En réduisant au même dénominateur :
\begin{itemize}
    \item $\frac{1}{84} = \frac{1623103}{136340652}$
    \item $\frac{1}{11677} = \frac{11676}{136340652}$
    \item $\frac{1}{136340652} = \frac{1}{136340652}$
\end{itemize}
La somme des numérateurs est :
$$ \frac{1}{84} + \frac{1}{11677} + \frac{1}{136340652} = \frac{1623103 + 11676 + 1}{136340652} = \frac{1634780}{136340652} $$
Le PGCD du numérateur et du dénominateur est $\text{PGCD}(1634780, 136340652) = 326956$.
La fraction irréductible est :
$$ \frac{1634780}{136340652} = \frac{1634780 \div 326956}{136340652 \div 326956} = \frac{5}{417} $$
Cette fraction est bien égale à $\frac{5}{417}$.

\subsection{Démonstration pour $n = 418$}
Soit $n = 418$. Nous cherchons $x, y, z \in \mathbb{Z}^{+}$ tels que $\frac{5}{418} = \frac{1}{x} + \frac{1}{y} + \frac{1}{z}$.
Posons $x = 84$, $y = 17557$, $z = 308230692$.
Les conditions $x > 0$, $y > 0$ et $z > 0$ sont satisfaites car $x, y, z \geq 1$.
Le PPCM des dénominateurs est $\text{PPCM}(84, 17557, 308230692) = 308230692$.
En réduisant au même dénominateur :
\begin{itemize}
    \item $\frac{1}{84} = \frac{3669413}{308230692}$
    \item $\frac{1}{17557} = \frac{17556}{308230692}$
    \item $\frac{1}{308230692} = \frac{1}{308230692}$
\end{itemize}
La somme des numérateurs est :
$$ \frac{1}{84} + \frac{1}{17557} + \frac{1}{308230692} = \frac{3669413 + 17556 + 1}{308230692} = \frac{3686970}{308230692} $$
Le PGCD du numérateur et du dénominateur est $\text{PGCD}(3686970, 308230692) = 737394$.
La fraction irréductible est :
$$ \frac{3686970}{308230692} = \frac{3686970 \div 737394}{308230692 \div 737394} = \frac{5}{418} $$
Cette fraction est bien égale à $\frac{5}{418}$.

\subsection{Démonstration pour $n = 419$}
Soit $n = 419$. Nous cherchons $x, y, z \in \mathbb{Z}^{+}$ tels que $\frac{5}{419} = \frac{1}{x} + \frac{1}{y} + \frac{1}{z}$.
Posons $x = 84$, $y = 35197$, $z = 1238793612$.
Les conditions $x > 0$, $y > 0$ et $z > 0$ sont satisfaites car $x, y, z \geq 1$.
Le PPCM des dénominateurs est $\text{PPCM}(84, 35197, 1238793612) = 1238793612$.
En réduisant au même dénominateur :
\begin{itemize}
    \item $\frac{1}{84} = \frac{14747543}{1238793612}$
    \item $\frac{1}{35197} = \frac{35196}{1238793612}$
    \item $\frac{1}{1238793612} = \frac{1}{1238793612}$
\end{itemize}
La somme des numérateurs est :
$$ \frac{1}{84} + \frac{1}{35197} + \frac{1}{1238793612} = \frac{14747543 + 35196 + 1}{1238793612} = \frac{14782740}{1238793612} $$
Le PGCD du numérateur et du dénominateur est $\text{PGCD}(14782740, 1238793612) = 2956548$.
La fraction irréductible est :
$$ \frac{14782740}{1238793612} = \frac{14782740 \div 2956548}{1238793612 \div 2956548} = \frac{5}{419} $$
Cette fraction est bien égale à $\frac{5}{419}$.

\subsection{Démonstration pour $n = 420$}
Soit $n = 420$. Nous cherchons $x, y, z \in \mathbb{Z}^{+}$ tels que $\frac{5}{420} = \frac{1}{x} + \frac{1}{y} + \frac{1}{z}$.
Posons $x = 85$, $y = 7141$, $z = 50986740$.
Les conditions $x > 0$, $y > 0$ et $z > 0$ sont satisfaites car $x, y, z \geq 1$.
Le PPCM des dénominateurs est $\text{PPCM}(85, 7141, 50986740) = 50986740$.
En réduisant au même dénominateur :
\begin{itemize}
    \item $\frac{1}{85} = \frac{599844}{50986740}$
    \item $\frac{1}{7141} = \frac{7140}{50986740}$
    \item $\frac{1}{50986740} = \frac{1}{50986740}$
\end{itemize}
La somme des numérateurs est :
$$ \frac{1}{85} + \frac{1}{7141} + \frac{1}{50986740} = \frac{599844 + 7140 + 1}{50986740} = \frac{606985}{50986740} $$
Le PGCD du numérateur et du dénominateur est $\text{PGCD}(606985, 50986740) = 606985$.
La fraction irréductible est :
$$ \frac{606985}{50986740} = \frac{606985 \div 606985}{50986740 \div 606985} = \frac{1}{84} $$
Cette fraction est bien égale à $\frac{5}{420}$.

\subsection{Démonstration pour $n = 421$}
Soit $n = 421$. Nous cherchons $x, y, z \in \mathbb{Z}^{+}$ tels que $\frac{5}{421} = \frac{1}{x} + \frac{1}{y} + \frac{1}{z}$.
Posons $x = 86$, $y = 4023$, $z = 145656738$.
Les conditions $x > 0$, $y > 0$ et $z > 0$ sont satisfaites car $x, y, z \geq 1$.
Le PPCM des dénominateurs est $\text{PPCM}(86, 4023, 145656738) = 145656738$.
En réduisant au même dénominateur :
\begin{itemize}
    \item $\frac{1}{86} = \frac{1693683}{145656738}$
    \item $\frac{1}{4023} = \frac{36206}{145656738}$
    \item $\frac{1}{145656738} = \frac{1}{145656738}$
\end{itemize}
La somme des numérateurs est :
$$ \frac{1}{86} + \frac{1}{4023} + \frac{1}{145656738} = \frac{1693683 + 36206 + 1}{145656738} = \frac{1729890}{145656738} $$
Le PGCD du numérateur et du dénominateur est $\text{PGCD}(1729890, 145656738) = 345978$.
La fraction irréductible est :
$$ \frac{1729890}{145656738} = \frac{1729890 \div 345978}{145656738 \div 345978} = \frac{5}{421} $$
Cette fraction est bien égale à $\frac{5}{421}$.

\subsection{Démonstration pour $n = 422$}
Soit $n = 422$. Nous cherchons $x, y, z \in \mathbb{Z}^{+}$ tels que $\frac{5}{422} = \frac{1}{x} + \frac{1}{y} + \frac{1}{z}$.
Posons $x = 85$, $y = 11957$, $z = 428897590$.
Les conditions $x > 0$, $y > 0$ et $z > 0$ sont satisfaites car $x, y, z \geq 1$.
Le PPCM des dénominateurs est $\text{PPCM}(85, 11957, 428897590) = 428897590$.
En réduisant au même dénominateur :
\begin{itemize}
    \item $\frac{1}{85} = \frac{5045854}{428897590}$
    \item $\frac{1}{11957} = \frac{35870}{428897590}$
    \item $\frac{1}{428897590} = \frac{1}{428897590}$
\end{itemize}
La somme des numérateurs est :
$$ \frac{1}{85} + \frac{1}{11957} + \frac{1}{428897590} = \frac{5045854 + 35870 + 1}{428897590} = \frac{5081725}{428897590} $$
Le PGCD du numérateur et du dénominateur est $\text{PGCD}(5081725, 428897590) = 1016345$.
La fraction irréductible est :
$$ \frac{5081725}{428897590} = \frac{5081725 \div 1016345}{428897590 \div 1016345} = \frac{5}{422} $$
Cette fraction est bien égale à $\frac{5}{422}$.

\subsection{Démonstration pour $n = 423$}
Soit $n = 423$. Nous cherchons $x, y, z \in \mathbb{Z}^{+}$ tels que $\frac{5}{423} = \frac{1}{x} + \frac{1}{y} + \frac{1}{z}$.
Posons $x = 85$, $y = 17978$, $z = 646398990$.
Les conditions $x > 0$, $y > 0$ et $z > 0$ sont satisfaites car $x, y, z \geq 1$.
Le PPCM des dénominateurs est $\text{PPCM}(85, 17978, 646398990) = 646398990$.
En réduisant au même dénominateur :
\begin{itemize}
    \item $\frac{1}{85} = \frac{7604694}{646398990}$
    \item $\frac{1}{17978} = \frac{35955}{646398990}$
    \item $\frac{1}{646398990} = \frac{1}{646398990}$
\end{itemize}
La somme des numérateurs est :
$$ \frac{1}{85} + \frac{1}{17978} + \frac{1}{646398990} = \frac{7604694 + 35955 + 1}{646398990} = \frac{7640650}{646398990} $$
Le PGCD du numérateur et du dénominateur est $\text{PGCD}(7640650, 646398990) = 1528130$.
La fraction irréductible est :
$$ \frac{7640650}{646398990} = \frac{7640650 \div 1528130}{646398990 \div 1528130} = \frac{5}{423} $$
Cette fraction est bien égale à $\frac{5}{423}$.

\subsection{Démonstration pour $n = 424$}
Soit $n = 424$. Nous cherchons $x, y, z \in \mathbb{Z}^{+}$ tels que $\frac{5}{424} = \frac{1}{x} + \frac{1}{y} + \frac{1}{z}$.
Posons $x = 85$, $y = 36041$, $z = 1298917640$.
Les conditions $x > 0$, $y > 0$ et $z > 0$ sont satisfaites car $x, y, z \geq 1$.
Le PPCM des dénominateurs est $\text{PPCM}(85, 36041, 1298917640) = 1298917640$.
En réduisant au même dénominateur :
\begin{itemize}
    \item $\frac{1}{85} = \frac{15281384}{1298917640}$
    \item $\frac{1}{36041} = \frac{36040}{1298917640}$
    \item $\frac{1}{1298917640} = \frac{1}{1298917640}$
\end{itemize}
La somme des numérateurs est :
$$ \frac{1}{85} + \frac{1}{36041} + \frac{1}{1298917640} = \frac{15281384 + 36040 + 1}{1298917640} = \frac{15317425}{1298917640} $$
Le PGCD du numérateur et du dénominateur est $\text{PGCD}(15317425, 1298917640) = 3063485$.
La fraction irréductible est :
$$ \frac{15317425}{1298917640} = \frac{15317425 \div 3063485}{1298917640 \div 3063485} = \frac{5}{424} $$
Cette fraction est bien égale à $\frac{5}{424}$.

\subsection{Démonstration pour $n = 425$}
Soit $n = 425$. Nous cherchons $x, y, z \in \mathbb{Z}^{+}$ tels que $\frac{5}{425} = \frac{1}{x} + \frac{1}{y} + \frac{1}{z}$.
Posons $x = 86$, $y = 7311$, $z = 53443410$.
Les conditions $x > 0$, $y > 0$ et $z > 0$ sont satisfaites car $x, y, z \geq 1$.
Le PPCM des dénominateurs est $\text{PPCM}(86, 7311, 53443410) = 53443410$.
En réduisant au même dénominateur :
\begin{itemize}
    \item $\frac{1}{86} = \frac{621435}{53443410}$
    \item $\frac{1}{7311} = \frac{7310}{53443410}$
    \item $\frac{1}{53443410} = \frac{1}{53443410}$
\end{itemize}
La somme des numérateurs est :
$$ \frac{1}{86} + \frac{1}{7311} + \frac{1}{53443410} = \frac{621435 + 7310 + 1}{53443410} = \frac{628746}{53443410} $$
Le PGCD du numérateur et du dénominateur est $\text{PGCD}(628746, 53443410) = 628746$.
La fraction irréductible est :
$$ \frac{628746}{53443410} = \frac{628746 \div 628746}{53443410 \div 628746} = \frac{1}{85} $$
Cette fraction est bien égale à $\frac{5}{425}$.

\subsection{Démonstration pour $n = 426$}
Soit $n = 426$. Nous cherchons $x, y, z \in \mathbb{Z}^{+}$ tels que $\frac{5}{426} = \frac{1}{x} + \frac{1}{y} + \frac{1}{z}$.
Posons $x = 86$, $y = 9160$, $z = 83896440$.
Les conditions $x > 0$, $y > 0$ et $z > 0$ sont satisfaites car $x, y, z \geq 1$.
Le PPCM des dénominateurs est $\text{PPCM}(86, 9160, 83896440) = 83896440$.
En réduisant au même dénominateur :
\begin{itemize}
    \item $\frac{1}{86} = \frac{975540}{83896440}$
    \item $\frac{1}{9160} = \frac{9159}{83896440}$
    \item $\frac{1}{83896440} = \frac{1}{83896440}$
\end{itemize}
La somme des numérateurs est :
$$ \frac{1}{86} + \frac{1}{9160} + \frac{1}{83896440} = \frac{975540 + 9159 + 1}{83896440} = \frac{984700}{83896440} $$
Le PGCD du numérateur et du dénominateur est $\text{PGCD}(984700, 83896440) = 196940$.
La fraction irréductible est :
$$ \frac{984700}{83896440} = \frac{984700 \div 196940}{83896440 \div 196940} = \frac{5}{426} $$
Cette fraction est bien égale à $\frac{5}{426}$.

\subsection{Démonstration pour $n = 427$}
Soit $n = 427$. Nous cherchons $x, y, z \in \mathbb{Z}^{+}$ tels que $\frac{5}{427} = \frac{1}{x} + \frac{1}{y} + \frac{1}{z}$.
Posons $x = 86$, $y = 12241$, $z = 449514002$.
Les conditions $x > 0$, $y > 0$ et $z > 0$ sont satisfaites car $x, y, z \geq 1$.
Le PPCM des dénominateurs est $\text{PPCM}(86, 12241, 449514002) = 449514002$.
En réduisant au même dénominateur :
\begin{itemize}
    \item $\frac{1}{86} = \frac{5226907}{449514002}$
    \item $\frac{1}{12241} = \frac{36722}{449514002}$
    \item $\frac{1}{449514002} = \frac{1}{449514002}$
\end{itemize}
La somme des numérateurs est :
$$ \frac{1}{86} + \frac{1}{12241} + \frac{1}{449514002} = \frac{5226907 + 36722 + 1}{449514002} = \frac{5263630}{449514002} $$
Le PGCD du numérateur et du dénominateur est $\text{PGCD}(5263630, 449514002) = 1052726$.
La fraction irréductible est :
$$ \frac{5263630}{449514002} = \frac{5263630 \div 1052726}{449514002 \div 1052726} = \frac{5}{427} $$
Cette fraction est bien égale à $\frac{5}{427}$.

\subsection{Démonstration pour $n = 428$}
Soit $n = 428$. Nous cherchons $x, y, z \in \mathbb{Z}^{+}$ tels que $\frac{5}{428} = \frac{1}{x} + \frac{1}{y} + \frac{1}{z}$.
Posons $x = 86$, $y = 18405$, $z = 338725620$.
Les conditions $x > 0$, $y > 0$ et $z > 0$ sont satisfaites car $x, y, z \geq 1$.
Le PPCM des dénominateurs est $\text{PPCM}(86, 18405, 338725620) = 338725620$.
En réduisant au même dénominateur :
\begin{itemize}
    \item $\frac{1}{86} = \frac{3938670}{338725620}$
    \item $\frac{1}{18405} = \frac{18404}{338725620}$
    \item $\frac{1}{338725620} = \frac{1}{338725620}$
\end{itemize}
La somme des numérateurs est :
$$ \frac{1}{86} + \frac{1}{18405} + \frac{1}{338725620} = \frac{3938670 + 18404 + 1}{338725620} = \frac{3957075}{338725620} $$
Le PGCD du numérateur et du dénominateur est $\text{PGCD}(3957075, 338725620) = 791415$.
La fraction irréductible est :
$$ \frac{3957075}{338725620} = \frac{3957075 \div 791415}{338725620 \div 791415} = \frac{5}{428} $$
Cette fraction est bien égale à $\frac{5}{428}$.

\subsection{Démonstration pour $n = 429$}
Soit $n = 429$. Nous cherchons $x, y, z \in \mathbb{Z}^{+}$ tels que $\frac{5}{429} = \frac{1}{x} + \frac{1}{y} + \frac{1}{z}$.
Posons $x = 86$, $y = 36895$, $z = 1361204130$.
Les conditions $x > 0$, $y > 0$ et $z > 0$ sont satisfaites car $x, y, z \geq 1$.
Le PPCM des dénominateurs est $\text{PPCM}(86, 36895, 1361204130) = 1361204130$.
En réduisant au même dénominateur :
\begin{itemize}
    \item $\frac{1}{86} = \frac{15827955}{1361204130}$
    \item $\frac{1}{36895} = \frac{36894}{1361204130}$
    \item $\frac{1}{1361204130} = \frac{1}{1361204130}$
\end{itemize}
La somme des numérateurs est :
$$ \frac{1}{86} + \frac{1}{36895} + \frac{1}{1361204130} = \frac{15827955 + 36894 + 1}{1361204130} = \frac{15864850}{1361204130} $$
Le PGCD du numérateur et du dénominateur est $\text{PGCD}(15864850, 1361204130) = 3172970$.
La fraction irréductible est :
$$ \frac{15864850}{1361204130} = \frac{15864850 \div 3172970}{1361204130 \div 3172970} = \frac{5}{429} $$
Cette fraction est bien égale à $\frac{5}{429}$.

\subsection{Démonstration pour $n = 430$}
Soit $n = 430$. Nous cherchons $x, y, z \in \mathbb{Z}^{+}$ tels que $\frac{5}{430} = \frac{1}{x} + \frac{1}{y} + \frac{1}{z}$.
Posons $x = 87$, $y = 7483$, $z = 55987806$.
Les conditions $x > 0$, $y > 0$ et $z > 0$ sont satisfaites car $x, y, z \geq 1$.
Le PPCM des dénominateurs est $\text{PPCM}(87, 7483, 55987806) = 55987806$.
En réduisant au même dénominateur :
\begin{itemize}
    \item $\frac{1}{87} = \frac{643538}{55987806}$
    \item $\frac{1}{7483} = \frac{7482}{55987806}$
    \item $\frac{1}{55987806} = \frac{1}{55987806}$
\end{itemize}
La somme des numérateurs est :
$$ \frac{1}{87} + \frac{1}{7483} + \frac{1}{55987806} = \frac{643538 + 7482 + 1}{55987806} = \frac{651021}{55987806} $$
Le PGCD du numérateur et du dénominateur est $\text{PGCD}(651021, 55987806) = 651021$.
La fraction irréductible est :
$$ \frac{651021}{55987806} = \frac{651021 \div 651021}{55987806 \div 651021} = \frac{1}{86} $$
Cette fraction est bien égale à $\frac{5}{430}$.

\subsection{Démonstration pour $n = 431$}
Soit $n = 431$. Nous cherchons $x, y, z \in \mathbb{Z}^{+}$ tels que $\frac{5}{431} = \frac{1}{x} + \frac{1}{y} + \frac{1}{z}$.
Posons $x = 87$, $y = 9375$, $z = 117178125$.
Les conditions $x > 0$, $y > 0$ et $z > 0$ sont satisfaites car $x, y, z \geq 1$.
Le PPCM des dénominateurs est $\text{PPCM}(87, 9375, 117178125) = 117178125$.
En réduisant au même dénominateur :
\begin{itemize}
    \item $\frac{1}{87} = \frac{1346875}{117178125}$
    \item $\frac{1}{9375} = \frac{12499}{117178125}$
    \item $\frac{1}{117178125} = \frac{1}{117178125}$
\end{itemize}
La somme des numérateurs est :
$$ \frac{1}{87} + \frac{1}{9375} + \frac{1}{117178125} = \frac{1346875 + 12499 + 1}{117178125} = \frac{1359375}{117178125} $$
Le PGCD du numérateur et du dénominateur est $\text{PGCD}(1359375, 117178125) = 271875$.
La fraction irréductible est :
$$ \frac{1359375}{117178125} = \frac{1359375 \div 271875}{117178125 \div 271875} = \frac{5}{431} $$
Cette fraction est bien égale à $\frac{5}{431}$.

\subsection{Démonstration pour $n = 432$}
Soit $n = 432$. Nous cherchons $x, y, z \in \mathbb{Z}^{+}$ tels que $\frac{5}{432} = \frac{1}{x} + \frac{1}{y} + \frac{1}{z}$.
Posons $x = 87$, $y = 12529$, $z = 156963312$.
Les conditions $x > 0$, $y > 0$ et $z > 0$ sont satisfaites car $x, y, z \geq 1$.
Le PPCM des dénominateurs est $\text{PPCM}(87, 12529, 156963312) = 156963312$.
En réduisant au même dénominateur :
\begin{itemize}
    \item $\frac{1}{87} = \frac{1804176}{156963312}$
    \item $\frac{1}{12529} = \frac{12528}{156963312}$
    \item $\frac{1}{156963312} = \frac{1}{156963312}$
\end{itemize}
La somme des numérateurs est :
$$ \frac{1}{87} + \frac{1}{12529} + \frac{1}{156963312} = \frac{1804176 + 12528 + 1}{156963312} = \frac{1816705}{156963312} $$
Le PGCD du numérateur et du dénominateur est $\text{PGCD}(1816705, 156963312) = 363341$.
La fraction irréductible est :
$$ \frac{1816705}{156963312} = \frac{1816705 \div 363341}{156963312 \div 363341} = \frac{5}{432} $$
Cette fraction est bien égale à $\frac{5}{432}$.

\subsection{Démonstration pour $n = 433$}
Soit $n = 433$. Nous cherchons $x, y, z \in \mathbb{Z}^{+}$ tels que $\frac{5}{433} = \frac{1}{x} + \frac{1}{y} + \frac{1}{z}$.
Posons $x = 87$, $y = 18836$, $z = 709570956$.
Les conditions $x > 0$, $y > 0$ et $z > 0$ sont satisfaites car $x, y, z \geq 1$.
Le PPCM des dénominateurs est $\text{PPCM}(87, 18836, 709570956) = 709570956$.
En réduisant au même dénominateur :
\begin{itemize}
    \item $\frac{1}{87} = \frac{8155988}{709570956}$
    \item $\frac{1}{18836} = \frac{37671}{709570956}$
    \item $\frac{1}{709570956} = \frac{1}{709570956}$
\end{itemize}
La somme des numérateurs est :
$$ \frac{1}{87} + \frac{1}{18836} + \frac{1}{709570956} = \frac{8155988 + 37671 + 1}{709570956} = \frac{8193660}{709570956} $$
Le PGCD du numérateur et du dénominateur est $\text{PGCD}(8193660, 709570956) = 1638732$.
La fraction irréductible est :
$$ \frac{8193660}{709570956} = \frac{8193660 \div 1638732}{709570956 \div 1638732} = \frac{5}{433} $$
Cette fraction est bien égale à $\frac{5}{433}$.

\subsection{Démonstration pour $n = 434$}
Soit $n = 434$. Nous cherchons $x, y, z \in \mathbb{Z}^{+}$ tels que $\frac{5}{434} = \frac{1}{x} + \frac{1}{y} + \frac{1}{z}$.
Posons $x = 87$, $y = 37759$, $z = 1425704322$.
Les conditions $x > 0$, $y > 0$ et $z > 0$ sont satisfaites car $x, y, z \geq 1$.
Le PPCM des dénominateurs est $\text{PPCM}(87, 37759, 1425704322) = 1425704322$.
En réduisant au même dénominateur :
\begin{itemize}
    \item $\frac{1}{87} = \frac{16387406}{1425704322}$
    \item $\frac{1}{37759} = \frac{37758}{1425704322}$
    \item $\frac{1}{1425704322} = \frac{1}{1425704322}$
\end{itemize}
La somme des numérateurs est :
$$ \frac{1}{87} + \frac{1}{37759} + \frac{1}{1425704322} = \frac{16387406 + 37758 + 1}{1425704322} = \frac{16425165}{1425704322} $$
Le PGCD du numérateur et du dénominateur est $\text{PGCD}(16425165, 1425704322) = 3285033$.
La fraction irréductible est :
$$ \frac{16425165}{1425704322} = \frac{16425165 \div 3285033}{1425704322 \div 3285033} = \frac{5}{434} $$
Cette fraction est bien égale à $\frac{5}{434}$.

\subsection{Démonstration pour $n = 435$}
Soit $n = 435$. Nous cherchons $x, y, z \in \mathbb{Z}^{+}$ tels que $\frac{5}{435} = \frac{1}{x} + \frac{1}{y} + \frac{1}{z}$.
Posons $x = 88$, $y = 7657$, $z = 58621992$.
Les conditions $x > 0$, $y > 0$ et $z > 0$ sont satisfaites car $x, y, z \geq 1$.
Le PPCM des dénominateurs est $\text{PPCM}(88, 7657, 58621992) = 58621992$.
En réduisant au même dénominateur :
\begin{itemize}
    \item $\frac{1}{88} = \frac{666159}{58621992}$
    \item $\frac{1}{7657} = \frac{7656}{58621992}$
    \item $\frac{1}{58621992} = \frac{1}{58621992}$
\end{itemize}
La somme des numérateurs est :
$$ \frac{1}{88} + \frac{1}{7657} + \frac{1}{58621992} = \frac{666159 + 7656 + 1}{58621992} = \frac{673816}{58621992} $$
Le PGCD du numérateur et du dénominateur est $\text{PGCD}(673816, 58621992) = 673816$.
La fraction irréductible est :
$$ \frac{673816}{58621992} = \frac{673816 \div 673816}{58621992 \div 673816} = \frac{1}{87} $$
Cette fraction est bien égale à $\frac{5}{435}$.

\subsection{Démonstration pour $n = 436$}
Soit $n = 436$. Nous cherchons $x, y, z \in \mathbb{Z}^{+}$ tels que $\frac{5}{436} = \frac{1}{x} + \frac{1}{y} + \frac{1}{z}$.
Posons $x = 88$, $y = 9593$, $z = 92016056$.
Les conditions $x > 0$, $y > 0$ et $z > 0$ sont satisfaites car $x, y, z \geq 1$.
Le PPCM des dénominateurs est $\text{PPCM}(88, 9593, 92016056) = 92016056$.
En réduisant au même dénominateur :
\begin{itemize}
    \item $\frac{1}{88} = \frac{1045637}{92016056}$
    \item $\frac{1}{9593} = \frac{9592}{92016056}$
    \item $\frac{1}{92016056} = \frac{1}{92016056}$
\end{itemize}
La somme des numérateurs est :
$$ \frac{1}{88} + \frac{1}{9593} + \frac{1}{92016056} = \frac{1045637 + 9592 + 1}{92016056} = \frac{1055230}{92016056} $$
Le PGCD du numérateur et du dénominateur est $\text{PGCD}(1055230, 92016056) = 211046$.
La fraction irréductible est :
$$ \frac{1055230}{92016056} = \frac{1055230 \div 211046}{92016056 \div 211046} = \frac{5}{436} $$
Cette fraction est bien égale à $\frac{5}{436}$.

\subsection{Démonstration pour $n = 437$}
Soit $n = 437$. Nous cherchons $x, y, z \in \mathbb{Z}^{+}$ tels que $\frac{5}{437} = \frac{1}{x} + \frac{1}{y} + \frac{1}{z}$.
Posons $x = 88$, $y = 12819$, $z = 492967464$.
Les conditions $x > 0$, $y > 0$ et $z > 0$ sont satisfaites car $x, y, z \geq 1$.
Le PPCM des dénominateurs est $\text{PPCM}(88, 12819, 492967464) = 492967464$.
En réduisant au même dénominateur :
\begin{itemize}
    \item $\frac{1}{88} = \frac{5601903}{492967464}$
    \item $\frac{1}{12819} = \frac{38456}{492967464}$
    \item $\frac{1}{492967464} = \frac{1}{492967464}$
\end{itemize}
La somme des numérateurs est :
$$ \frac{1}{88} + \frac{1}{12819} + \frac{1}{492967464} = \frac{5601903 + 38456 + 1}{492967464} = \frac{5640360}{492967464} $$
Le PGCD du numérateur et du dénominateur est $\text{PGCD}(5640360, 492967464) = 1128072$.
La fraction irréductible est :
$$ \frac{5640360}{492967464} = \frac{5640360 \div 1128072}{492967464 \div 1128072} = \frac{5}{437} $$
Cette fraction est bien égale à $\frac{5}{437}$.

\subsection{Démonstration pour $n = 438$}
Soit $n = 438$. Nous cherchons $x, y, z \in \mathbb{Z}^{+}$ tels que $\frac{5}{438} = \frac{1}{x} + \frac{1}{y} + \frac{1}{z}$.
Posons $x = 88$, $y = 19273$, $z = 371429256$.
Les conditions $x > 0$, $y > 0$ et $z > 0$ sont satisfaites car $x, y, z \geq 1$.
Le PPCM des dénominateurs est $\text{PPCM}(88, 19273, 371429256) = 371429256$.
En réduisant au même dénominateur :
\begin{itemize}
    \item $\frac{1}{88} = \frac{4220787}{371429256}$
    \item $\frac{1}{19273} = \frac{19272}{371429256}$
    \item $\frac{1}{371429256} = \frac{1}{371429256}$
\end{itemize}
La somme des numérateurs est :
$$ \frac{1}{88} + \frac{1}{19273} + \frac{1}{371429256} = \frac{4220787 + 19272 + 1}{371429256} = \frac{4240060}{371429256} $$
Le PGCD du numérateur et du dénominateur est $\text{PGCD}(4240060, 371429256) = 848012$.
La fraction irréductible est :
$$ \frac{4240060}{371429256} = \frac{4240060 \div 848012}{371429256 \div 848012} = \frac{5}{438} $$
Cette fraction est bien égale à $\frac{5}{438}$.

\subsection{Démonstration pour $n = 439$}
Soit $n = 439$. Nous cherchons $x, y, z \in \mathbb{Z}^{+}$ tels que $\frac{5}{439} = \frac{1}{x} + \frac{1}{y} + \frac{1}{z}$.
Posons $x = 88$, $y = 38633$, $z = 1492470056$.
Les conditions $x > 0$, $y > 0$ et $z > 0$ sont satisfaites car $x, y, z \geq 1$.
Le PPCM des dénominateurs est $\text{PPCM}(88, 38633, 1492470056) = 1492470056$.
En réduisant au même dénominateur :
\begin{itemize}
    \item $\frac{1}{88} = \frac{16959887}{1492470056}$
    \item $\frac{1}{38633} = \frac{38632}{1492470056}$
    \item $\frac{1}{1492470056} = \frac{1}{1492470056}$
\end{itemize}
La somme des numérateurs est :
$$ \frac{1}{88} + \frac{1}{38633} + \frac{1}{1492470056} = \frac{16959887 + 38632 + 1}{1492470056} = \frac{16998520}{1492470056} $$
Le PGCD du numérateur et du dénominateur est $\text{PGCD}(16998520, 1492470056) = 3399704$.
La fraction irréductible est :
$$ \frac{16998520}{1492470056} = \frac{16998520 \div 3399704}{1492470056 \div 3399704} = \frac{5}{439} $$
Cette fraction est bien égale à $\frac{5}{439}$.

\subsection{Démonstration pour $n = 440$}
Soit $n = 440$. Nous cherchons $x, y, z \in \mathbb{Z}^{+}$ tels que $\frac{5}{440} = \frac{1}{x} + \frac{1}{y} + \frac{1}{z}$.
Posons $x = 89$, $y = 7833$, $z = 61348056$.
Les conditions $x > 0$, $y > 0$ et $z > 0$ sont satisfaites car $x, y, z \geq 1$.
Le PPCM des dénominateurs est $\text{PPCM}(89, 7833, 61348056) = 61348056$.
En réduisant au même dénominateur :
\begin{itemize}
    \item $\frac{1}{89} = \frac{689304}{61348056}$
    \item $\frac{1}{7833} = \frac{7832}{61348056}$
    \item $\frac{1}{61348056} = \frac{1}{61348056}$
\end{itemize}
La somme des numérateurs est :
$$ \frac{1}{89} + \frac{1}{7833} + \frac{1}{61348056} = \frac{689304 + 7832 + 1}{61348056} = \frac{697137}{61348056} $$
Le PGCD du numérateur et du dénominateur est $\text{PGCD}(697137, 61348056) = 697137$.
La fraction irréductible est :
$$ \frac{697137}{61348056} = \frac{697137 \div 697137}{61348056 \div 697137} = \frac{1}{88} $$
Cette fraction est bien égale à $\frac{5}{440}$.

\subsection{Démonstration pour $n = 441$}
Soit $n = 441$. Nous cherchons $x, y, z \in \mathbb{Z}^{+}$ tels que $\frac{5}{441} = \frac{1}{x} + \frac{1}{y} + \frac{1}{z}$.
Posons $x = 89$, $y = 9813$, $z = 128383479$.
Les conditions $x > 0$, $y > 0$ et $z > 0$ sont satisfaites car $x, y, z \geq 1$.
Le PPCM des dénominateurs est $\text{PPCM}(89, 9813, 128383479) = 128383479$.
En réduisant au même dénominateur :
\begin{itemize}
    \item $\frac{1}{89} = \frac{1442511}{128383479}$
    \item $\frac{1}{9813} = \frac{13083}{128383479}$
    \item $\frac{1}{128383479} = \frac{1}{128383479}$
\end{itemize}
La somme des numérateurs est :
$$ \frac{1}{89} + \frac{1}{9813} + \frac{1}{128383479} = \frac{1442511 + 13083 + 1}{128383479} = \frac{1455595}{128383479} $$
Le PGCD du numérateur et du dénominateur est $\text{PGCD}(1455595, 128383479) = 291119$.
La fraction irréductible est :
$$ \frac{1455595}{128383479} = \frac{1455595 \div 291119}{128383479 \div 291119} = \frac{5}{441} $$
Cette fraction est bien égale à $\frac{5}{441}$.

\subsection{Démonstration pour $n = 442$}
Soit $n = 442$. Nous cherchons $x, y, z \in \mathbb{Z}^{+}$ tels que $\frac{5}{442} = \frac{1}{x} + \frac{1}{y} + \frac{1}{z}$.
Posons $x = 89$, $y = 13113$, $z = 515839194$.
Les conditions $x > 0$, $y > 0$ et $z > 0$ sont satisfaites car $x, y, z \geq 1$.
Le PPCM des dénominateurs est $\text{PPCM}(89, 13113, 515839194) = 515839194$.
En réduisant au même dénominateur :
\begin{itemize}
    \item $\frac{1}{89} = \frac{5795946}{515839194}$
    \item $\frac{1}{13113} = \frac{39338}{515839194}$
    \item $\frac{1}{515839194} = \frac{1}{515839194}$
\end{itemize}
La somme des numérateurs est :
$$ \frac{1}{89} + \frac{1}{13113} + \frac{1}{515839194} = \frac{5795946 + 39338 + 1}{515839194} = \frac{5835285}{515839194} $$
Le PGCD du numérateur et du dénominateur est $\text{PGCD}(5835285, 515839194) = 1167057$.
La fraction irréductible est :
$$ \frac{5835285}{515839194} = \frac{5835285 \div 1167057}{515839194 \div 1167057} = \frac{5}{442} $$
Cette fraction est bien égale à $\frac{5}{442}$.

\subsection{Démonstration pour $n = 443$}
Soit $n = 443$. Nous cherchons $x, y, z \in \mathbb{Z}^{+}$ tels que $\frac{5}{443} = \frac{1}{x} + \frac{1}{y} + \frac{1}{z}$.
Posons $x = 89$, $y = 19714$, $z = 777263878$.
Les conditions $x > 0$, $y > 0$ et $z > 0$ sont satisfaites car $x, y, z \geq 1$.
Le PPCM des dénominateurs est $\text{PPCM}(89, 19714, 777263878) = 777263878$.
En réduisant au même dénominateur :
\begin{itemize}
    \item $\frac{1}{89} = \frac{8733302}{777263878}$
    \item $\frac{1}{19714} = \frac{39427}{777263878}$
    \item $\frac{1}{777263878} = \frac{1}{777263878}$
\end{itemize}
La somme des numérateurs est :
$$ \frac{1}{89} + \frac{1}{19714} + \frac{1}{777263878} = \frac{8733302 + 39427 + 1}{777263878} = \frac{8772730}{777263878} $$
Le PGCD du numérateur et du dénominateur est $\text{PGCD}(8772730, 777263878) = 1754546$.
La fraction irréductible est :
$$ \frac{8772730}{777263878} = \frac{8772730 \div 1754546}{777263878 \div 1754546} = \frac{5}{443} $$
Cette fraction est bien égale à $\frac{5}{443}$.

\subsection{Démonstration pour $n = 444$}
Soit $n = 444$. Nous cherchons $x, y, z \in \mathbb{Z}^{+}$ tels que $\frac{5}{444} = \frac{1}{x} + \frac{1}{y} + \frac{1}{z}$.
Posons $x = 89$, $y = 39517$, $z = 1561553772$.
Les conditions $x > 0$, $y > 0$ et $z > 0$ sont satisfaites car $x, y, z \geq 1$.
Le PPCM des dénominateurs est $\text{PPCM}(89, 39517, 1561553772) = 1561553772$.
En réduisant au même dénominateur :
\begin{itemize}
    \item $\frac{1}{89} = \frac{17545548}{1561553772}$
    \item $\frac{1}{39517} = \frac{39516}{1561553772}$
    \item $\frac{1}{1561553772} = \frac{1}{1561553772}$
\end{itemize}
La somme des numérateurs est :
$$ \frac{1}{89} + \frac{1}{39517} + \frac{1}{1561553772} = \frac{17545548 + 39516 + 1}{1561553772} = \frac{17585065}{1561553772} $$
Le PGCD du numérateur et du dénominateur est $\text{PGCD}(17585065, 1561553772) = 3517013$.
La fraction irréductible est :
$$ \frac{17585065}{1561553772} = \frac{17585065 \div 3517013}{1561553772 \div 3517013} = \frac{5}{444} $$
Cette fraction est bien égale à $\frac{5}{444}$.

\subsection{Démonstration pour $n = 445$}
Soit $n = 445$. Nous cherchons $x, y, z \in \mathbb{Z}^{+}$ tels que $\frac{5}{445} = \frac{1}{x} + \frac{1}{y} + \frac{1}{z}$.
Posons $x = 90$, $y = 8011$, $z = 64168110$.
Les conditions $x > 0$, $y > 0$ et $z > 0$ sont satisfaites car $x, y, z \geq 1$.
Le PPCM des dénominateurs est $\text{PPCM}(90, 8011, 64168110) = 64168110$.
En réduisant au même dénominateur :
\begin{itemize}
    \item $\frac{1}{90} = \frac{712979}{64168110}$
    \item $\frac{1}{8011} = \frac{8010}{64168110}$
    \item $\frac{1}{64168110} = \frac{1}{64168110}$
\end{itemize}
La somme des numérateurs est :
$$ \frac{1}{90} + \frac{1}{8011} + \frac{1}{64168110} = \frac{712979 + 8010 + 1}{64168110} = \frac{720990}{64168110} $$
Le PGCD du numérateur et du dénominateur est $\text{PGCD}(720990, 64168110) = 720990$.
La fraction irréductible est :
$$ \frac{720990}{64168110} = \frac{720990 \div 720990}{64168110 \div 720990} = \frac{1}{89} $$
Cette fraction est bien égale à $\frac{5}{445}$.

\subsection{Démonstration pour $n = 446$}
Soit $n = 446$. Nous cherchons $x, y, z \in \mathbb{Z}^{+}$ tels que $\frac{5}{446} = \frac{1}{x} + \frac{1}{y} + \frac{1}{z}$.
Posons $x = 90$, $y = 10036$, $z = 100711260$.
Les conditions $x > 0$, $y > 0$ et $z > 0$ sont satisfaites car $x, y, z \geq 1$.
Le PPCM des dénominateurs est $\text{PPCM}(90, 10036, 100711260) = 100711260$.
En réduisant au même dénominateur :
\begin{itemize}
    \item $\frac{1}{90} = \frac{1119014}{100711260}$
    \item $\frac{1}{10036} = \frac{10035}{100711260}$
    \item $\frac{1}{100711260} = \frac{1}{100711260}$
\end{itemize}
La somme des numérateurs est :
$$ \frac{1}{90} + \frac{1}{10036} + \frac{1}{100711260} = \frac{1119014 + 10035 + 1}{100711260} = \frac{1129050}{100711260} $$
Le PGCD du numérateur et du dénominateur est $\text{PGCD}(1129050, 100711260) = 225810$.
La fraction irréductible est :
$$ \frac{1129050}{100711260} = \frac{1129050 \div 225810}{100711260 \div 225810} = \frac{5}{446} $$
Cette fraction est bien égale à $\frac{5}{446}$.

\subsection{Démonstration pour $n = 447$}
Soit $n = 447$. Nous cherchons $x, y, z \in \mathbb{Z}^{+}$ tels que $\frac{5}{447} = \frac{1}{x} + \frac{1}{y} + \frac{1}{z}$.
Posons $x = 90$, $y = 13411$, $z = 179841510$.
Les conditions $x > 0$, $y > 0$ et $z > 0$ sont satisfaites car $x, y, z \geq 1$.
Le PPCM des dénominateurs est $\text{PPCM}(90, 13411, 179841510) = 179841510$.
En réduisant au même dénominateur :
\begin{itemize}
    \item $\frac{1}{90} = \frac{1998239}{179841510}$
    \item $\frac{1}{13411} = \frac{13410}{179841510}$
    \item $\frac{1}{179841510} = \frac{1}{179841510}$
\end{itemize}
La somme des numérateurs est :
$$ \frac{1}{90} + \frac{1}{13411} + \frac{1}{179841510} = \frac{1998239 + 13410 + 1}{179841510} = \frac{2011650}{179841510} $$
Le PGCD du numérateur et du dénominateur est $\text{PGCD}(2011650, 179841510) = 402330$.
La fraction irréductible est :
$$ \frac{2011650}{179841510} = \frac{2011650 \div 402330}{179841510 \div 402330} = \frac{5}{447} $$
Cette fraction est bien égale à $\frac{5}{447}$.

\subsection{Démonstration pour $n = 448$}
Soit $n = 448$. Nous cherchons $x, y, z \in \mathbb{Z}^{+}$ tels que $\frac{5}{448} = \frac{1}{x} + \frac{1}{y} + \frac{1}{z}$.
Posons $x = 90$, $y = 20161$, $z = 406445760$.
Les conditions $x > 0$, $y > 0$ et $z > 0$ sont satisfaites car $x, y, z \geq 1$.
Le PPCM des dénominateurs est $\text{PPCM}(90, 20161, 406445760) = 406445760$.
En réduisant au même dénominateur :
\begin{itemize}
    \item $\frac{1}{90} = \frac{4516064}{406445760}$
    \item $\frac{1}{20161} = \frac{20160}{406445760}$
    \item $\frac{1}{406445760} = \frac{1}{406445760}$
\end{itemize}
La somme des numérateurs est :
$$ \frac{1}{90} + \frac{1}{20161} + \frac{1}{406445760} = \frac{4516064 + 20160 + 1}{406445760} = \frac{4536225}{406445760} $$
Le PGCD du numérateur et du dénominateur est $\text{PGCD}(4536225, 406445760) = 907245$.
La fraction irréductible est :
$$ \frac{4536225}{406445760} = \frac{4536225 \div 907245}{406445760 \div 907245} = \frac{5}{448} $$
Cette fraction est bien égale à $\frac{5}{448}$.

\subsection{Démonstration pour $n = 449$}
Soit $n = 449$. Nous cherchons $x, y, z \in \mathbb{Z}^{+}$ tels que $\frac{5}{449} = \frac{1}{x} + \frac{1}{y} + \frac{1}{z}$.
Posons $x = 90$, $y = 40411$, $z = 1633008510$.
Les conditions $x > 0$, $y > 0$ et $z > 0$ sont satisfaites car $x, y, z \geq 1$.
Le PPCM des dénominateurs est $\text{PPCM}(90, 40411, 1633008510) = 1633008510$.
En réduisant au même dénominateur :
\begin{itemize}
    \item $\frac{1}{90} = \frac{18144539}{1633008510}$
    \item $\frac{1}{40411} = \frac{40410}{1633008510}$
    \item $\frac{1}{1633008510} = \frac{1}{1633008510}$
\end{itemize}
La somme des numérateurs est :
$$ \frac{1}{90} + \frac{1}{40411} + \frac{1}{1633008510} = \frac{18144539 + 40410 + 1}{1633008510} = \frac{18184950}{1633008510} $$
Le PGCD du numérateur et du dénominateur est $\text{PGCD}(18184950, 1633008510) = 3636990$.
La fraction irréductible est :
$$ \frac{18184950}{1633008510} = \frac{18184950 \div 3636990}{1633008510 \div 3636990} = \frac{5}{449} $$
Cette fraction est bien égale à $\frac{5}{449}$.

\subsection{Démonstration pour $n = 450$}
Soit $n = 450$. Nous cherchons $x, y, z \in \mathbb{Z}^{+}$ tels que $\frac{5}{450} = \frac{1}{x} + \frac{1}{y} + \frac{1}{z}$.
Posons $x = 91$, $y = 8191$, $z = 67084290$.
Les conditions $x > 0$, $y > 0$ et $z > 0$ sont satisfaites car $x, y, z \geq 1$.
Le PPCM des dénominateurs est $\text{PPCM}(91, 8191, 67084290) = 67084290$.
En réduisant au même dénominateur :
\begin{itemize}
    \item $\frac{1}{91} = \frac{737190}{67084290}$
    \item $\frac{1}{8191} = \frac{8190}{67084290}$
    \item $\frac{1}{67084290} = \frac{1}{67084290}$
\end{itemize}
La somme des numérateurs est :
$$ \frac{1}{91} + \frac{1}{8191} + \frac{1}{67084290} = \frac{737190 + 8190 + 1}{67084290} = \frac{745381}{67084290} $$
Le PGCD du numérateur et du dénominateur est $\text{PGCD}(745381, 67084290) = 745381$.
La fraction irréductible est :
$$ \frac{745381}{67084290} = \frac{745381 \div 745381}{67084290 \div 745381} = \frac{1}{90} $$
Cette fraction est bien égale à $\frac{5}{450}$.

\subsection{Démonstration pour $n = 451$}
Soit $n = 451$. Nous cherchons $x, y, z \in \mathbb{Z}^{+}$ tels que $\frac{5}{451} = \frac{1}{x} + \frac{1}{y} + \frac{1}{z}$.
Posons $x = 91$, $y = 10262$, $z = 60166106$.
Les conditions $x > 0$, $y > 0$ et $z > 0$ sont satisfaites car $x, y, z \geq 1$.
Le PPCM des dénominateurs est $\text{PPCM}(91, 10262, 60166106) = 60166106$.
En réduisant au même dénominateur :
\begin{itemize}
    \item $\frac{1}{91} = \frac{661166}{60166106}$
    \item $\frac{1}{10262} = \frac{5863}{60166106}$
    \item $\frac{1}{60166106} = \frac{1}{60166106}$
\end{itemize}
La somme des numérateurs est :
$$ \frac{1}{91} + \frac{1}{10262} + \frac{1}{60166106} = \frac{661166 + 5863 + 1}{60166106} = \frac{667030}{60166106} $$
Le PGCD du numérateur et du dénominateur est $\text{PGCD}(667030, 60166106) = 133406$.
La fraction irréductible est :
$$ \frac{667030}{60166106} = \frac{667030 \div 133406}{60166106 \div 133406} = \frac{5}{451} $$
Cette fraction est bien égale à $\frac{5}{451}$.

\subsection{Démonstration pour $n = 452$}
Soit $n = 452$. Nous cherchons $x, y, z \in \mathbb{Z}^{+}$ tels que $\frac{5}{452} = \frac{1}{x} + \frac{1}{y} + \frac{1}{z}$.
Posons $x = 91$, $y = 13711$, $z = 563960852$.
Les conditions $x > 0$, $y > 0$ et $z > 0$ sont satisfaites car $x, y, z \geq 1$.
Le PPCM des dénominateurs est $\text{PPCM}(91, 13711, 563960852) = 563960852$.
En réduisant au même dénominateur :
\begin{itemize}
    \item $\frac{1}{91} = \frac{6197372}{563960852}$
    \item $\frac{1}{13711} = \frac{41132}{563960852}$
    \item $\frac{1}{563960852} = \frac{1}{563960852}$
\end{itemize}
La somme des numérateurs est :
$$ \frac{1}{91} + \frac{1}{13711} + \frac{1}{563960852} = \frac{6197372 + 41132 + 1}{563960852} = \frac{6238505}{563960852} $$
Le PGCD du numérateur et du dénominateur est $\text{PGCD}(6238505, 563960852) = 1247701$.
La fraction irréductible est :
$$ \frac{6238505}{563960852} = \frac{6238505 \div 1247701}{563960852 \div 1247701} = \frac{5}{452} $$
Cette fraction est bien égale à $\frac{5}{452}$.

\subsection{Démonstration pour $n = 453$}
Soit $n = 453$. Nous cherchons $x, y, z \in \mathbb{Z}^{+}$ tels que $\frac{5}{453} = \frac{1}{x} + \frac{1}{y} + \frac{1}{z}$.
Posons $x = 91$, $y = 20612$, $z = 849688476$.
Les conditions $x > 0$, $y > 0$ et $z > 0$ sont satisfaites car $x, y, z \geq 1$.
Le PPCM des dénominateurs est $\text{PPCM}(91, 20612, 849688476) = 849688476$.
En réduisant au même dénominateur :
\begin{itemize}
    \item $\frac{1}{91} = \frac{9337236}{849688476}$
    \item $\frac{1}{20612} = \frac{41223}{849688476}$
    \item $\frac{1}{849688476} = \frac{1}{849688476}$
\end{itemize}
La somme des numérateurs est :
$$ \frac{1}{91} + \frac{1}{20612} + \frac{1}{849688476} = \frac{9337236 + 41223 + 1}{849688476} = \frac{9378460}{849688476} $$
Le PGCD du numérateur et du dénominateur est $\text{PGCD}(9378460, 849688476) = 1875692$.
La fraction irréductible est :
$$ \frac{9378460}{849688476} = \frac{9378460 \div 1875692}{849688476 \div 1875692} = \frac{5}{453} $$
Cette fraction est bien égale à $\frac{5}{453}$.

\subsection{Démonstration pour $n = 454$}
Soit $n = 454$. Nous cherchons $x, y, z \in \mathbb{Z}^{+}$ tels que $\frac{5}{454} = \frac{1}{x} + \frac{1}{y} + \frac{1}{z}$.
Posons $x = 91$, $y = 41315$, $z = 1706887910$.
Les conditions $x > 0$, $y > 0$ et $z > 0$ sont satisfaites car $x, y, z \geq 1$.
Le PPCM des dénominateurs est $\text{PPCM}(91, 41315, 1706887910) = 1706887910$.
En réduisant au même dénominateur :
\begin{itemize}
    \item $\frac{1}{91} = \frac{18757010}{1706887910}$
    \item $\frac{1}{41315} = \frac{41314}{1706887910}$
    \item $\frac{1}{1706887910} = \frac{1}{1706887910}$
\end{itemize}
La somme des numérateurs est :
$$ \frac{1}{91} + \frac{1}{41315} + \frac{1}{1706887910} = \frac{18757010 + 41314 + 1}{1706887910} = \frac{18798325}{1706887910} $$
Le PGCD du numérateur et du dénominateur est $\text{PGCD}(18798325, 1706887910) = 3759665$.
La fraction irréductible est :
$$ \frac{18798325}{1706887910} = \frac{18798325 \div 3759665}{1706887910 \div 3759665} = \frac{5}{454} $$
Cette fraction est bien égale à $\frac{5}{454}$.

\subsection{Démonstration pour $n = 455$}
Soit $n = 455$. Nous cherchons $x, y, z \in \mathbb{Z}^{+}$ tels que $\frac{5}{455} = \frac{1}{x} + \frac{1}{y} + \frac{1}{z}$.
Posons $x = 92$, $y = 8373$, $z = 70098756$.
Les conditions $x > 0$, $y > 0$ et $z > 0$ sont satisfaites car $x, y, z \geq 1$.
Le PPCM des dénominateurs est $\text{PPCM}(92, 8373, 70098756) = 70098756$.
En réduisant au même dénominateur :
\begin{itemize}
    \item $\frac{1}{92} = \frac{761943}{70098756}$
    \item $\frac{1}{8373} = \frac{8372}{70098756}$
    \item $\frac{1}{70098756} = \frac{1}{70098756}$
\end{itemize}
La somme des numérateurs est :
$$ \frac{1}{92} + \frac{1}{8373} + \frac{1}{70098756} = \frac{761943 + 8372 + 1}{70098756} = \frac{770316}{70098756} $$
Le PGCD du numérateur et du dénominateur est $\text{PGCD}(770316, 70098756) = 770316$.
La fraction irréductible est :
$$ \frac{770316}{70098756} = \frac{770316 \div 770316}{70098756 \div 770316} = \frac{1}{91} $$
Cette fraction est bien égale à $\frac{5}{455}$.

\subsection{Démonstration pour $n = 456$}
Soit $n = 456$. Nous cherchons $x, y, z \in \mathbb{Z}^{+}$ tels que $\frac{5}{456} = \frac{1}{x} + \frac{1}{y} + \frac{1}{z}$.
Posons $x = 92$, $y = 10489$, $z = 110008632$.
Les conditions $x > 0$, $y > 0$ et $z > 0$ sont satisfaites car $x, y, z \geq 1$.
Le PPCM des dénominateurs est $\text{PPCM}(92, 10489, 110008632) = 110008632$.
En réduisant au même dénominateur :
\begin{itemize}
    \item $\frac{1}{92} = \frac{1195746}{110008632}$
    \item $\frac{1}{10489} = \frac{10488}{110008632}$
    \item $\frac{1}{110008632} = \frac{1}{110008632}$
\end{itemize}
La somme des numérateurs est :
$$ \frac{1}{92} + \frac{1}{10489} + \frac{1}{110008632} = \frac{1195746 + 10488 + 1}{110008632} = \frac{1206235}{110008632} $$
Le PGCD du numérateur et du dénominateur est $\text{PGCD}(1206235, 110008632) = 241247$.
La fraction irréductible est :
$$ \frac{1206235}{110008632} = \frac{1206235 \div 241247}{110008632 \div 241247} = \frac{5}{456} $$
Cette fraction est bien égale à $\frac{5}{456}$.

\subsection{Démonstration pour $n = 457$}
Soit $n = 457$. Nous cherchons $x, y, z \in \mathbb{Z}^{+}$ tels que $\frac{5}{457} = \frac{1}{x} + \frac{1}{y} + \frac{1}{z}$.
Posons $x = 92$, $y = 14015$, $z = 589246660$.
Les conditions $x > 0$, $y > 0$ et $z > 0$ sont satisfaites car $x, y, z \geq 1$.
Le PPCM des dénominateurs est $\text{PPCM}(92, 14015, 589246660) = 589246660$.
En réduisant au même dénominateur :
\begin{itemize}
    \item $\frac{1}{92} = \frac{6404855}{589246660}$
    \item $\frac{1}{14015} = \frac{42044}{589246660}$
    \item $\frac{1}{589246660} = \frac{1}{589246660}$
\end{itemize}
La somme des numérateurs est :
$$ \frac{1}{92} + \frac{1}{14015} + \frac{1}{589246660} = \frac{6404855 + 42044 + 1}{589246660} = \frac{6446900}{589246660} $$
Le PGCD du numérateur et du dénominateur est $\text{PGCD}(6446900, 589246660) = 1289380$.
La fraction irréductible est :
$$ \frac{6446900}{589246660} = \frac{6446900 \div 1289380}{589246660 \div 1289380} = \frac{5}{457} $$
Cette fraction est bien égale à $\frac{5}{457}$.

\subsection{Démonstration pour $n = 458$}
Soit $n = 458$. Nous cherchons $x, y, z \in \mathbb{Z}^{+}$ tels que $\frac{5}{458} = \frac{1}{x} + \frac{1}{y} + \frac{1}{z}$.
Posons $x = 92$, $y = 21069$, $z = 443881692$.
Les conditions $x > 0$, $y > 0$ et $z > 0$ sont satisfaites car $x, y, z \geq 1$.
Le PPCM des dénominateurs est $\text{PPCM}(92, 21069, 443881692) = 443881692$.
En réduisant au même dénominateur :
\begin{itemize}
    \item $\frac{1}{92} = \frac{4824801}{443881692}$
    \item $\frac{1}{21069} = \frac{21068}{443881692}$
    \item $\frac{1}{443881692} = \frac{1}{443881692}$
\end{itemize}
La somme des numérateurs est :
$$ \frac{1}{92} + \frac{1}{21069} + \frac{1}{443881692} = \frac{4824801 + 21068 + 1}{443881692} = \frac{4845870}{443881692} $$
Le PGCD du numérateur et du dénominateur est $\text{PGCD}(4845870, 443881692) = 969174$.
La fraction irréductible est :
$$ \frac{4845870}{443881692} = \frac{4845870 \div 969174}{443881692 \div 969174} = \frac{5}{458} $$
Cette fraction est bien égale à $\frac{5}{458}$.

\subsection{Démonstration pour $n = 459$}
Soit $n = 459$. Nous cherchons $x, y, z \in \mathbb{Z}^{+}$ tels que $\frac{5}{459} = \frac{1}{x} + \frac{1}{y} + \frac{1}{z}$.
Posons $x = 92$, $y = 42229$, $z = 1783246212$.
Les conditions $x > 0$, $y > 0$ et $z > 0$ sont satisfaites car $x, y, z \geq 1$.
Le PPCM des dénominateurs est $\text{PPCM}(92, 42229, 1783246212) = 1783246212$.
En réduisant au même dénominateur :
\begin{itemize}
    \item $\frac{1}{92} = \frac{19383111}{1783246212}$
    \item $\frac{1}{42229} = \frac{42228}{1783246212}$
    \item $\frac{1}{1783246212} = \frac{1}{1783246212}$
\end{itemize}
La somme des numérateurs est :
$$ \frac{1}{92} + \frac{1}{42229} + \frac{1}{1783246212} = \frac{19383111 + 42228 + 1}{1783246212} = \frac{19425340}{1783246212} $$
Le PGCD du numérateur et du dénominateur est $\text{PGCD}(19425340, 1783246212) = 3885068$.
La fraction irréductible est :
$$ \frac{19425340}{1783246212} = \frac{19425340 \div 3885068}{1783246212 \div 3885068} = \frac{5}{459} $$
Cette fraction est bien égale à $\frac{5}{459}$.

\subsection{Démonstration pour $n = 460$}
Soit $n = 460$. Nous cherchons $x, y, z \in \mathbb{Z}^{+}$ tels que $\frac{5}{460} = \frac{1}{x} + \frac{1}{y} + \frac{1}{z}$.
Posons $x = 93$, $y = 8557$, $z = 73213692$.
Les conditions $x > 0$, $y > 0$ et $z > 0$ sont satisfaites car $x, y, z \geq 1$.
Le PPCM des dénominateurs est $\text{PPCM}(93, 8557, 73213692) = 73213692$.
En réduisant au même dénominateur :
\begin{itemize}
    \item $\frac{1}{93} = \frac{787244}{73213692}$
    \item $\frac{1}{8557} = \frac{8556}{73213692}$
    \item $\frac{1}{73213692} = \frac{1}{73213692}$
\end{itemize}
La somme des numérateurs est :
$$ \frac{1}{93} + \frac{1}{8557} + \frac{1}{73213692} = \frac{787244 + 8556 + 1}{73213692} = \frac{795801}{73213692} $$
Le PGCD du numérateur et du dénominateur est $\text{PGCD}(795801, 73213692) = 795801$.
La fraction irréductible est :
$$ \frac{795801}{73213692} = \frac{795801 \div 795801}{73213692 \div 795801} = \frac{1}{92} $$
Cette fraction est bien égale à $\frac{5}{460}$.

\subsection{Démonstration pour $n = 461$}
Soit $n = 461$. Nous cherchons $x, y, z \in \mathbb{Z}^{+}$ tels que $\frac{5}{461} = \frac{1}{x} + \frac{1}{y} + \frac{1}{z}$.
Posons $x = 93$, $y = 10719$, $z = 153185229$.
Les conditions $x > 0$, $y > 0$ et $z > 0$ sont satisfaites car $x, y, z \geq 1$.
Le PPCM des dénominateurs est $\text{PPCM}(93, 10719, 153185229) = 153185229$.
En réduisant au même dénominateur :
\begin{itemize}
    \item $\frac{1}{93} = \frac{1647153}{153185229}$
    \item $\frac{1}{10719} = \frac{14291}{153185229}$
    \item $\frac{1}{153185229} = \frac{1}{153185229}$
\end{itemize}
La somme des numérateurs est :
$$ \frac{1}{93} + \frac{1}{10719} + \frac{1}{153185229} = \frac{1647153 + 14291 + 1}{153185229} = \frac{1661445}{153185229} $$
Le PGCD du numérateur et du dénominateur est $\text{PGCD}(1661445, 153185229) = 332289$.
La fraction irréductible est :
$$ \frac{1661445}{153185229} = \frac{1661445 \div 332289}{153185229 \div 332289} = \frac{5}{461} $$
Cette fraction est bien égale à $\frac{5}{461}$.

\subsection{Démonstration pour $n = 462$}
Soit $n = 462$. Nous cherchons $x, y, z \in \mathbb{Z}^{+}$ tels que $\frac{5}{462} = \frac{1}{x} + \frac{1}{y} + \frac{1}{z}$.
Posons $x = 93$, $y = 14323$, $z = 205134006$.
Les conditions $x > 0$, $y > 0$ et $z > 0$ sont satisfaites car $x, y, z \geq 1$.
Le PPCM des dénominateurs est $\text{PPCM}(93, 14323, 205134006) = 205134006$.
En réduisant au même dénominateur :
\begin{itemize}
    \item $\frac{1}{93} = \frac{2205742}{205134006}$
    \item $\frac{1}{14323} = \frac{14322}{205134006}$
    \item $\frac{1}{205134006} = \frac{1}{205134006}$
\end{itemize}
La somme des numérateurs est :
$$ \frac{1}{93} + \frac{1}{14323} + \frac{1}{205134006} = \frac{2205742 + 14322 + 1}{205134006} = \frac{2220065}{205134006} $$
Le PGCD du numérateur et du dénominateur est $\text{PGCD}(2220065, 205134006) = 444013$.
La fraction irréductible est :
$$ \frac{2220065}{205134006} = \frac{2220065 \div 444013}{205134006 \div 444013} = \frac{5}{462} $$
Cette fraction est bien égale à $\frac{5}{462}$.

\subsection{Démonstration pour $n = 463$}
Soit $n = 463$. Nous cherchons $x, y, z \in \mathbb{Z}^{+}$ tels que $\frac{5}{463} = \frac{1}{x} + \frac{1}{y} + \frac{1}{z}$.
Posons $x = 93$, $y = 21530$, $z = 927060270$.
Les conditions $x > 0$, $y > 0$ et $z > 0$ sont satisfaites car $x, y, z \geq 1$.
Le PPCM des dénominateurs est $\text{PPCM}(93, 21530, 927060270) = 927060270$.
En réduisant au même dénominateur :
\begin{itemize}
    \item $\frac{1}{93} = \frac{9968390}{927060270}$
    \item $\frac{1}{21530} = \frac{43059}{927060270}$
    \item $\frac{1}{927060270} = \frac{1}{927060270}$
\end{itemize}
La somme des numérateurs est :
$$ \frac{1}{93} + \frac{1}{21530} + \frac{1}{927060270} = \frac{9968390 + 43059 + 1}{927060270} = \frac{10011450}{927060270} $$
Le PGCD du numérateur et du dénominateur est $\text{PGCD}(10011450, 927060270) = 2002290$.
La fraction irréductible est :
$$ \frac{10011450}{927060270} = \frac{10011450 \div 2002290}{927060270 \div 2002290} = \frac{5}{463} $$
Cette fraction est bien égale à $\frac{5}{463}$.

\subsection{Démonstration pour $n = 464$}
Soit $n = 464$. Nous cherchons $x, y, z \in \mathbb{Z}^{+}$ tels que $\frac{5}{464} = \frac{1}{x} + \frac{1}{y} + \frac{1}{z}$.
Posons $x = 93$, $y = 43153$, $z = 1862138256$.
Les conditions $x > 0$, $y > 0$ et $z > 0$ sont satisfaites car $x, y, z \geq 1$.
Le PPCM des dénominateurs est $\text{PPCM}(93, 43153, 1862138256) = 1862138256$.
En réduisant au même dénominateur :
\begin{itemize}
    \item $\frac{1}{93} = \frac{20022992}{1862138256}$
    \item $\frac{1}{43153} = \frac{43152}{1862138256}$
    \item $\frac{1}{1862138256} = \frac{1}{1862138256}$
\end{itemize}
La somme des numérateurs est :
$$ \frac{1}{93} + \frac{1}{43153} + \frac{1}{1862138256} = \frac{20022992 + 43152 + 1}{1862138256} = \frac{20066145}{1862138256} $$
Le PGCD du numérateur et du dénominateur est $\text{PGCD}(20066145, 1862138256) = 4013229$.
La fraction irréductible est :
$$ \frac{20066145}{1862138256} = \frac{20066145 \div 4013229}{1862138256 \div 4013229} = \frac{5}{464} $$
Cette fraction est bien égale à $\frac{5}{464}$.

\subsection{Démonstration pour $n = 465$}
Soit $n = 465$. Nous cherchons $x, y, z \in \mathbb{Z}^{+}$ tels que $\frac{5}{465} = \frac{1}{x} + \frac{1}{y} + \frac{1}{z}$.
Posons $x = 94$, $y = 8743$, $z = 76431306$.
Les conditions $x > 0$, $y > 0$ et $z > 0$ sont satisfaites car $x, y, z \geq 1$.
Le PPCM des dénominateurs est $\text{PPCM}(94, 8743, 76431306) = 76431306$.
En réduisant au même dénominateur :
\begin{itemize}
    \item $\frac{1}{94} = \frac{813099}{76431306}$
    \item $\frac{1}{8743} = \frac{8742}{76431306}$
    \item $\frac{1}{76431306} = \frac{1}{76431306}$
\end{itemize}
La somme des numérateurs est :
$$ \frac{1}{94} + \frac{1}{8743} + \frac{1}{76431306} = \frac{813099 + 8742 + 1}{76431306} = \frac{821842}{76431306} $$
Le PGCD du numérateur et du dénominateur est $\text{PGCD}(821842, 76431306) = 821842$.
La fraction irréductible est :
$$ \frac{821842}{76431306} = \frac{821842 \div 821842}{76431306 \div 821842} = \frac{1}{93} $$
Cette fraction est bien égale à $\frac{5}{465}$.

\subsection{Démonstration pour $n = 466$}
Soit $n = 466$. Nous cherchons $x, y, z \in \mathbb{Z}^{+}$ tels que $\frac{5}{466} = \frac{1}{x} + \frac{1}{y} + \frac{1}{z}$.
Posons $x = 94$, $y = 10952$, $z = 119935352$.
Les conditions $x > 0$, $y > 0$ et $z > 0$ sont satisfaites car $x, y, z \geq 1$.
Le PPCM des dénominateurs est $\text{PPCM}(94, 10952, 119935352) = 119935352$.
En réduisant au même dénominateur :
\begin{itemize}
    \item $\frac{1}{94} = \frac{1275908}{119935352}$
    \item $\frac{1}{10952} = \frac{10951}{119935352}$
    \item $\frac{1}{119935352} = \frac{1}{119935352}$
\end{itemize}
La somme des numérateurs est :
$$ \frac{1}{94} + \frac{1}{10952} + \frac{1}{119935352} = \frac{1275908 + 10951 + 1}{119935352} = \frac{1286860}{119935352} $$
Le PGCD du numérateur et du dénominateur est $\text{PGCD}(1286860, 119935352) = 257372$.
La fraction irréductible est :
$$ \frac{1286860}{119935352} = \frac{1286860 \div 257372}{119935352 \div 257372} = \frac{5}{466} $$
Cette fraction est bien égale à $\frac{5}{466}$.

\subsection{Démonstration pour $n = 467$}
Soit $n = 467$. Nous cherchons $x, y, z \in \mathbb{Z}^{+}$ tels que $\frac{5}{467} = \frac{1}{x} + \frac{1}{y} + \frac{1}{z}$.
Posons $x = 94$, $y = 14633$, $z = 642359434$.
Les conditions $x > 0$, $y > 0$ et $z > 0$ sont satisfaites car $x, y, z \geq 1$.
Le PPCM des dénominateurs est $\text{PPCM}(94, 14633, 642359434) = 642359434$.
En réduisant au même dénominateur :
\begin{itemize}
    \item $\frac{1}{94} = \frac{6833611}{642359434}$
    \item $\frac{1}{14633} = \frac{43898}{642359434}$
    \item $\frac{1}{642359434} = \frac{1}{642359434}$
\end{itemize}
La somme des numérateurs est :
$$ \frac{1}{94} + \frac{1}{14633} + \frac{1}{642359434} = \frac{6833611 + 43898 + 1}{642359434} = \frac{6877510}{642359434} $$
Le PGCD du numérateur et du dénominateur est $\text{PGCD}(6877510, 642359434) = 1375502$.
La fraction irréductible est :
$$ \frac{6877510}{642359434} = \frac{6877510 \div 1375502}{642359434 \div 1375502} = \frac{5}{467} $$
Cette fraction est bien égale à $\frac{5}{467}$.

\subsection{Démonstration pour $n = 468$}
Soit $n = 468$. Nous cherchons $x, y, z \in \mathbb{Z}^{+}$ tels que $\frac{5}{468} = \frac{1}{x} + \frac{1}{y} + \frac{1}{z}$.
Posons $x = 94$, $y = 21997$, $z = 483846012$.
Les conditions $x > 0$, $y > 0$ et $z > 0$ sont satisfaites car $x, y, z \geq 1$.
Le PPCM des dénominateurs est $\text{PPCM}(94, 21997, 483846012) = 483846012$.
En réduisant au même dénominateur :
\begin{itemize}
    \item $\frac{1}{94} = \frac{5147298}{483846012}$
    \item $\frac{1}{21997} = \frac{21996}{483846012}$
    \item $\frac{1}{483846012} = \frac{1}{483846012}$
\end{itemize}
La somme des numérateurs est :
$$ \frac{1}{94} + \frac{1}{21997} + \frac{1}{483846012} = \frac{5147298 + 21996 + 1}{483846012} = \frac{5169295}{483846012} $$
Le PGCD du numérateur et du dénominateur est $\text{PGCD}(5169295, 483846012) = 1033859$.
La fraction irréductible est :
$$ \frac{5169295}{483846012} = \frac{5169295 \div 1033859}{483846012 \div 1033859} = \frac{5}{468} $$
Cette fraction est bien égale à $\frac{5}{468}$.

\subsection{Démonstration pour $n = 469$}
Soit $n = 469$. Nous cherchons $x, y, z \in \mathbb{Z}^{+}$ tels que $\frac{5}{469} = \frac{1}{x} + \frac{1}{y} + \frac{1}{z}$.
Posons $x = 94$, $y = 44087$, $z = 1943619482$.
Les conditions $x > 0$, $y > 0$ et $z > 0$ sont satisfaites car $x, y, z \geq 1$.
Le PPCM des dénominateurs est $\text{PPCM}(94, 44087, 1943619482) = 1943619482$.
En réduisant au même dénominateur :
\begin{itemize}
    \item $\frac{1}{94} = \frac{20676803}{1943619482}$
    \item $\frac{1}{44087} = \frac{44086}{1943619482}$
    \item $\frac{1}{1943619482} = \frac{1}{1943619482}$
\end{itemize}
La somme des numérateurs est :
$$ \frac{1}{94} + \frac{1}{44087} + \frac{1}{1943619482} = \frac{20676803 + 44086 + 1}{1943619482} = \frac{20720890}{1943619482} $$
Le PGCD du numérateur et du dénominateur est $\text{PGCD}(20720890, 1943619482) = 4144178$.
La fraction irréductible est :
$$ \frac{20720890}{1943619482} = \frac{20720890 \div 4144178}{1943619482 \div 4144178} = \frac{5}{469} $$
Cette fraction est bien égale à $\frac{5}{469}$.

\subsection{Démonstration pour $n = 470$}
Soit $n = 470$. Nous cherchons $x, y, z \in \mathbb{Z}^{+}$ tels que $\frac{5}{470} = \frac{1}{x} + \frac{1}{y} + \frac{1}{z}$.
Posons $x = 95$, $y = 8931$, $z = 79753830$.
Les conditions $x > 0$, $y > 0$ et $z > 0$ sont satisfaites car $x, y, z \geq 1$.
Le PPCM des dénominateurs est $\text{PPCM}(95, 8931, 79753830) = 79753830$.
En réduisant au même dénominateur :
\begin{itemize}
    \item $\frac{1}{95} = \frac{839514}{79753830}$
    \item $\frac{1}{8931} = \frac{8930}{79753830}$
    \item $\frac{1}{79753830} = \frac{1}{79753830}$
\end{itemize}
La somme des numérateurs est :
$$ \frac{1}{95} + \frac{1}{8931} + \frac{1}{79753830} = \frac{839514 + 8930 + 1}{79753830} = \frac{848445}{79753830} $$
Le PGCD du numérateur et du dénominateur est $\text{PGCD}(848445, 79753830) = 848445$.
La fraction irréductible est :
$$ \frac{848445}{79753830} = \frac{848445 \div 848445}{79753830 \div 848445} = \frac{1}{94} $$
Cette fraction est bien égale à $\frac{5}{470}$.

\subsection{Démonstration pour $n = 471$}
Soit $n = 471$. Nous cherchons $x, y, z \in \mathbb{Z}^{+}$ tels que $\frac{5}{471} = \frac{1}{x} + \frac{1}{y} + \frac{1}{z}$.
Posons $x = 95$, $y = 11187$, $z = 166854105$.
Les conditions $x > 0$, $y > 0$ et $z > 0$ sont satisfaites car $x, y, z \geq 1$.
Le PPCM des dénominateurs est $\text{PPCM}(95, 11187, 166854105) = 166854105$.
En réduisant au même dénominateur :
\begin{itemize}
    \item $\frac{1}{95} = \frac{1756359}{166854105}$
    \item $\frac{1}{11187} = \frac{14915}{166854105}$
    \item $\frac{1}{166854105} = \frac{1}{166854105}$
\end{itemize}
La somme des numérateurs est :
$$ \frac{1}{95} + \frac{1}{11187} + \frac{1}{166854105} = \frac{1756359 + 14915 + 1}{166854105} = \frac{1771275}{166854105} $$
Le PGCD du numérateur et du dénominateur est $\text{PGCD}(1771275, 166854105) = 354255$.
La fraction irréductible est :
$$ \frac{1771275}{166854105} = \frac{1771275 \div 354255}{166854105 \div 354255} = \frac{5}{471} $$
Cette fraction est bien égale à $\frac{5}{471}$.

\subsection{Démonstration pour $n = 472$}
Soit $n = 472$. Nous cherchons $x, y, z \in \mathbb{Z}^{+}$ tels que $\frac{5}{472} = \frac{1}{x} + \frac{1}{y} + \frac{1}{z}$.
Posons $x = 95$, $y = 14947$, $z = 670223480$.
Les conditions $x > 0$, $y > 0$ et $z > 0$ sont satisfaites car $x, y, z \geq 1$.
Le PPCM des dénominateurs est $\text{PPCM}(95, 14947, 670223480) = 670223480$.
En réduisant au même dénominateur :
\begin{itemize}
    \item $\frac{1}{95} = \frac{7054984}{670223480}$
    \item $\frac{1}{14947} = \frac{44840}{670223480}$
    \item $\frac{1}{670223480} = \frac{1}{670223480}$
\end{itemize}
La somme des numérateurs est :
$$ \frac{1}{95} + \frac{1}{14947} + \frac{1}{670223480} = \frac{7054984 + 44840 + 1}{670223480} = \frac{7099825}{670223480} $$
Le PGCD du numérateur et du dénominateur est $\text{PGCD}(7099825, 670223480) = 1419965$.
La fraction irréductible est :
$$ \frac{7099825}{670223480} = \frac{7099825 \div 1419965}{670223480 \div 1419965} = \frac{5}{472} $$
Cette fraction est bien égale à $\frac{5}{472}$.

\subsection{Démonstration pour $n = 473$}
Soit $n = 473$. Nous cherchons $x, y, z \in \mathbb{Z}^{+}$ tels que $\frac{5}{473} = \frac{1}{x} + \frac{1}{y} + \frac{1}{z}$.
Posons $x = 95$, $y = 22468$, $z = 1009599580$.
Les conditions $x > 0$, $y > 0$ et $z > 0$ sont satisfaites car $x, y, z \geq 1$.
Le PPCM des dénominateurs est $\text{PPCM}(95, 22468, 1009599580) = 1009599580$.
En réduisant au même dénominateur :
\begin{itemize}
    \item $\frac{1}{95} = \frac{10627364}{1009599580}$
    \item $\frac{1}{22468} = \frac{44935}{1009599580}$
    \item $\frac{1}{1009599580} = \frac{1}{1009599580}$
\end{itemize}
La somme des numérateurs est :
$$ \frac{1}{95} + \frac{1}{22468} + \frac{1}{1009599580} = \frac{10627364 + 44935 + 1}{1009599580} = \frac{10672300}{1009599580} $$
Le PGCD du numérateur et du dénominateur est $\text{PGCD}(10672300, 1009599580) = 2134460$.
La fraction irréductible est :
$$ \frac{10672300}{1009599580} = \frac{10672300 \div 2134460}{1009599580 \div 2134460} = \frac{5}{473} $$
Cette fraction est bien égale à $\frac{5}{473}$.

\subsection{Démonstration pour $n = 474$}
Soit $n = 474$. Nous cherchons $x, y, z \in \mathbb{Z}^{+}$ tels que $\frac{5}{474} = \frac{1}{x} + \frac{1}{y} + \frac{1}{z}$.
Posons $x = 95$, $y = 45031$, $z = 2027745930$.
Les conditions $x > 0$, $y > 0$ et $z > 0$ sont satisfaites car $x, y, z \geq 1$.
Le PPCM des dénominateurs est $\text{PPCM}(95, 45031, 2027745930) = 2027745930$.
En réduisant au même dénominateur :
\begin{itemize}
    \item $\frac{1}{95} = \frac{21344694}{2027745930}$
    \item $\frac{1}{45031} = \frac{45030}{2027745930}$
    \item $\frac{1}{2027745930} = \frac{1}{2027745930}$
\end{itemize}
La somme des numérateurs est :
$$ \frac{1}{95} + \frac{1}{45031} + \frac{1}{2027745930} = \frac{21344694 + 45030 + 1}{2027745930} = \frac{21389725}{2027745930} $$
Le PGCD du numérateur et du dénominateur est $\text{PGCD}(21389725, 2027745930) = 4277945$.
La fraction irréductible est :
$$ \frac{21389725}{2027745930} = \frac{21389725 \div 4277945}{2027745930 \div 4277945} = \frac{5}{474} $$
Cette fraction est bien égale à $\frac{5}{474}$.

\subsection{Démonstration pour $n = 475$}
Soit $n = 475$. Nous cherchons $x, y, z \in \mathbb{Z}^{+}$ tels que $\frac{5}{475} = \frac{1}{x} + \frac{1}{y} + \frac{1}{z}$.
Posons $x = 96$, $y = 9121$, $z = 83183520$.
Les conditions $x > 0$, $y > 0$ et $z > 0$ sont satisfaites car $x, y, z \geq 1$.
Le PPCM des dénominateurs est $\text{PPCM}(96, 9121, 83183520) = 83183520$.
En réduisant au même dénominateur :
\begin{itemize}
    \item $\frac{1}{96} = \frac{866495}{83183520}$
    \item $\frac{1}{9121} = \frac{9120}{83183520}$
    \item $\frac{1}{83183520} = \frac{1}{83183520}$
\end{itemize}
La somme des numérateurs est :
$$ \frac{1}{96} + \frac{1}{9121} + \frac{1}{83183520} = \frac{866495 + 9120 + 1}{83183520} = \frac{875616}{83183520} $$
Le PGCD du numérateur et du dénominateur est $\text{PGCD}(875616, 83183520) = 875616$.
La fraction irréductible est :
$$ \frac{875616}{83183520} = \frac{875616 \div 875616}{83183520 \div 875616} = \frac{1}{95} $$
Cette fraction est bien égale à $\frac{5}{475}$.

\subsection{Démonstration pour $n = 476$}
Soit $n = 476$. Nous cherchons $x, y, z \in \mathbb{Z}^{+}$ tels que $\frac{5}{476} = \frac{1}{x} + \frac{1}{y} + \frac{1}{z}$.
Posons $x = 96$, $y = 11425$, $z = 130519200$.
Les conditions $x > 0$, $y > 0$ et $z > 0$ sont satisfaites car $x, y, z \geq 1$.
Le PPCM des dénominateurs est $\text{PPCM}(96, 11425, 130519200) = 130519200$.
En réduisant au même dénominateur :
\begin{itemize}
    \item $\frac{1}{96} = \frac{1359575}{130519200}$
    \item $\frac{1}{11425} = \frac{11424}{130519200}$
    \item $\frac{1}{130519200} = \frac{1}{130519200}$
\end{itemize}
La somme des numérateurs est :
$$ \frac{1}{96} + \frac{1}{11425} + \frac{1}{130519200} = \frac{1359575 + 11424 + 1}{130519200} = \frac{1371000}{130519200} $$
Le PGCD du numérateur et du dénominateur est $\text{PGCD}(1371000, 130519200) = 274200$.
La fraction irréductible est :
$$ \frac{1371000}{130519200} = \frac{1371000 \div 274200}{130519200 \div 274200} = \frac{5}{476} $$
Cette fraction est bien égale à $\frac{5}{476}$.

\subsection{Démonstration pour $n = 477$}
Soit $n = 477$. Nous cherchons $x, y, z \in \mathbb{Z}^{+}$ tels que $\frac{5}{477} = \frac{1}{x} + \frac{1}{y} + \frac{1}{z}$.
Posons $x = 96$, $y = 15265$, $z = 233004960$.
Les conditions $x > 0$, $y > 0$ et $z > 0$ sont satisfaites car $x, y, z \geq 1$.
Le PPCM des dénominateurs est $\text{PPCM}(96, 15265, 233004960) = 233004960$.
En réduisant au même dénominateur :
\begin{itemize}
    \item $\frac{1}{96} = \frac{2427135}{233004960}$
    \item $\frac{1}{15265} = \frac{15264}{233004960}$
    \item $\frac{1}{233004960} = \frac{1}{233004960}$
\end{itemize}
La somme des numérateurs est :
$$ \frac{1}{96} + \frac{1}{15265} + \frac{1}{233004960} = \frac{2427135 + 15264 + 1}{233004960} = \frac{2442400}{233004960} $$
Le PGCD du numérateur et du dénominateur est $\text{PGCD}(2442400, 233004960) = 488480$.
La fraction irréductible est :
$$ \frac{2442400}{233004960} = \frac{2442400 \div 488480}{233004960 \div 488480} = \frac{5}{477} $$
Cette fraction est bien égale à $\frac{5}{477}$.

\subsection{Démonstration pour $n = 478$}
Soit $n = 478$. Nous cherchons $x, y, z \in \mathbb{Z}^{+}$ tels que $\frac{5}{478} = \frac{1}{x} + \frac{1}{y} + \frac{1}{z}$.
Posons $x = 96$, $y = 22945$, $z = 526450080$.
Les conditions $x > 0$, $y > 0$ et $z > 0$ sont satisfaites car $x, y, z \geq 1$.
Le PPCM des dénominateurs est $\text{PPCM}(96, 22945, 526450080) = 526450080$.
En réduisant au même dénominateur :
\begin{itemize}
    \item $\frac{1}{96} = \frac{5483855}{526450080}$
    \item $\frac{1}{22945} = \frac{22944}{526450080}$
    \item $\frac{1}{526450080} = \frac{1}{526450080}$
\end{itemize}
La somme des numérateurs est :
$$ \frac{1}{96} + \frac{1}{22945} + \frac{1}{526450080} = \frac{5483855 + 22944 + 1}{526450080} = \frac{5506800}{526450080} $$
Le PGCD du numérateur et du dénominateur est $\text{PGCD}(5506800, 526450080) = 1101360$.
La fraction irréductible est :
$$ \frac{5506800}{526450080} = \frac{5506800 \div 1101360}{526450080 \div 1101360} = \frac{5}{478} $$
Cette fraction est bien égale à $\frac{5}{478}$.

\subsection{Démonstration pour $n = 479$}
Soit $n = 479$. Nous cherchons $x, y, z \in \mathbb{Z}^{+}$ tels que $\frac{5}{479} = \frac{1}{x} + \frac{1}{y} + \frac{1}{z}$.
Posons $x = 96$, $y = 45985$, $z = 2114574240$.
Les conditions $x > 0$, $y > 0$ et $z > 0$ sont satisfaites car $x, y, z \geq 1$.
Le PPCM des dénominateurs est $\text{PPCM}(96, 45985, 2114574240) = 2114574240$.
En réduisant au même dénominateur :
\begin{itemize}
    \item $\frac{1}{96} = \frac{22026815}{2114574240}$
    \item $\frac{1}{45985} = \frac{45984}{2114574240}$
    \item $\frac{1}{2114574240} = \frac{1}{2114574240}$
\end{itemize}
La somme des numérateurs est :
$$ \frac{1}{96} + \frac{1}{45985} + \frac{1}{2114574240} = \frac{22026815 + 45984 + 1}{2114574240} = \frac{22072800}{2114574240} $$
Le PGCD du numérateur et du dénominateur est $\text{PGCD}(22072800, 2114574240) = 4414560$.
La fraction irréductible est :
$$ \frac{22072800}{2114574240} = \frac{22072800 \div 4414560}{2114574240 \div 4414560} = \frac{5}{479} $$
Cette fraction est bien égale à $\frac{5}{479}$.

\subsection{Démonstration pour $n = 480$}
Soit $n = 480$. Nous cherchons $x, y, z \in \mathbb{Z}^{+}$ tels que $\frac{5}{480} = \frac{1}{x} + \frac{1}{y} + \frac{1}{z}$.
Posons $x = 97$, $y = 9313$, $z = 86722656$.
Les conditions $x > 0$, $y > 0$ et $z > 0$ sont satisfaites car $x, y, z \geq 1$.
Le PPCM des dénominateurs est $\text{PPCM}(97, 9313, 86722656) = 86722656$.
En réduisant au même dénominateur :
\begin{itemize}
    \item $\frac{1}{97} = \frac{894048}{86722656}$
    \item $\frac{1}{9313} = \frac{9312}{86722656}$
    \item $\frac{1}{86722656} = \frac{1}{86722656}$
\end{itemize}
La somme des numérateurs est :
$$ \frac{1}{97} + \frac{1}{9313} + \frac{1}{86722656} = \frac{894048 + 9312 + 1}{86722656} = \frac{903361}{86722656} $$
Le PGCD du numérateur et du dénominateur est $\text{PGCD}(903361, 86722656) = 903361$.
La fraction irréductible est :
$$ \frac{903361}{86722656} = \frac{903361 \div 903361}{86722656 \div 903361} = \frac{1}{96} $$
Cette fraction est bien égale à $\frac{5}{480}$.

\subsection{Démonstration pour $n = 481$}
Soit $n = 481$. Nous cherchons $x, y, z \in \mathbb{Z}^{+}$ tels que $\frac{5}{481} = \frac{1}{x} + \frac{1}{y} + \frac{1}{z}$.
Posons $x = 98$, $y = 5238$, $z = 61727211$.
Les conditions $x > 0$, $y > 0$ et $z > 0$ sont satisfaites car $x, y, z \geq 1$.
Le PPCM des dénominateurs est $\text{PPCM}(98, 5238, 61727211) = 123454422$.
En réduisant au même dénominateur :
\begin{itemize}
    \item $\frac{1}{98} = \frac{1259739}{123454422}$
    \item $\frac{1}{5238} = \frac{23569}{123454422}$
    \item $\frac{1}{61727211} = \frac{2}{123454422}$
\end{itemize}
La somme des numérateurs est :
$$ \frac{1}{98} + \frac{1}{5238} + \frac{1}{61727211} = \frac{1259739 + 23569 + 2}{123454422} = \frac{1283310}{123454422} $$
Le PGCD du numérateur et du dénominateur est $\text{PGCD}(1283310, 123454422) = 256662$.
La fraction irréductible est :
$$ \frac{1283310}{123454422} = \frac{1283310 \div 256662}{123454422 \div 256662} = \frac{5}{481} $$
Cette fraction est bien égale à $\frac{5}{481}$.

\subsection{Démonstration pour $n = 482$}
Soit $n = 482$. Nous cherchons $x, y, z \in \mathbb{Z}^{+}$ tels que $\frac{5}{482} = \frac{1}{x} + \frac{1}{y} + \frac{1}{z}$.
Posons $x = 97$, $y = 15585$, $z = 728661090$.
Les conditions $x > 0$, $y > 0$ et $z > 0$ sont satisfaites car $x, y, z \geq 1$.
Le PPCM des dénominateurs est $\text{PPCM}(97, 15585, 728661090) = 728661090$.
En réduisant au même dénominateur :
\begin{itemize}
    \item $\frac{1}{97} = \frac{7511970}{728661090}$
    \item $\frac{1}{15585} = \frac{46754}{728661090}$
    \item $\frac{1}{728661090} = \frac{1}{728661090}$
\end{itemize}
La somme des numérateurs est :
$$ \frac{1}{97} + \frac{1}{15585} + \frac{1}{728661090} = \frac{7511970 + 46754 + 1}{728661090} = \frac{7558725}{728661090} $$
Le PGCD du numérateur et du dénominateur est $\text{PGCD}(7558725, 728661090) = 1511745$.
La fraction irréductible est :
$$ \frac{7558725}{728661090} = \frac{7558725 \div 1511745}{728661090 \div 1511745} = \frac{5}{482} $$
Cette fraction est bien égale à $\frac{5}{482}$.

\subsection{Démonstration pour $n = 483$}
Soit $n = 483$. Nous cherchons $x, y, z \in \mathbb{Z}^{+}$ tels que $\frac{5}{483} = \frac{1}{x} + \frac{1}{y} + \frac{1}{z}$.
Posons $x = 97$, $y = 23426$, $z = 1097531526$.
Les conditions $x > 0$, $y > 0$ et $z > 0$ sont satisfaites car $x, y, z \geq 1$.
Le PPCM des dénominateurs est $\text{PPCM}(97, 23426, 1097531526) = 1097531526$.
En réduisant au même dénominateur :
\begin{itemize}
    \item $\frac{1}{97} = \frac{11314758}{1097531526}$
    \item $\frac{1}{23426} = \frac{46851}{1097531526}$
    \item $\frac{1}{1097531526} = \frac{1}{1097531526}$
\end{itemize}
La somme des numérateurs est :
$$ \frac{1}{97} + \frac{1}{23426} + \frac{1}{1097531526} = \frac{11314758 + 46851 + 1}{1097531526} = \frac{11361610}{1097531526} $$
Le PGCD du numérateur et du dénominateur est $\text{PGCD}(11361610, 1097531526) = 2272322$.
La fraction irréductible est :
$$ \frac{11361610}{1097531526} = \frac{11361610 \div 2272322}{1097531526 \div 2272322} = \frac{5}{483} $$
Cette fraction est bien égale à $\frac{5}{483}$.

\subsection{Démonstration pour $n = 484$}
Soit $n = 484$. Nous cherchons $x, y, z \in \mathbb{Z}^{+}$ tels que $\frac{5}{484} = \frac{1}{x} + \frac{1}{y} + \frac{1}{z}$.
Posons $x = 97$, $y = 46949$, $z = 2204161652$.
Les conditions $x > 0$, $y > 0$ et $z > 0$ sont satisfaites car $x, y, z \geq 1$.
Le PPCM des dénominateurs est $\text{PPCM}(97, 46949, 2204161652) = 2204161652$.
En réduisant au même dénominateur :
\begin{itemize}
    \item $\frac{1}{97} = \frac{22723316}{2204161652}$
    \item $\frac{1}{46949} = \frac{46948}{2204161652}$
    \item $\frac{1}{2204161652} = \frac{1}{2204161652}$
\end{itemize}
La somme des numérateurs est :
$$ \frac{1}{97} + \frac{1}{46949} + \frac{1}{2204161652} = \frac{22723316 + 46948 + 1}{2204161652} = \frac{22770265}{2204161652} $$
Le PGCD du numérateur et du dénominateur est $\text{PGCD}(22770265, 2204161652) = 4554053$.
La fraction irréductible est :
$$ \frac{22770265}{2204161652} = \frac{22770265 \div 4554053}{2204161652 \div 4554053} = \frac{5}{484} $$
Cette fraction est bien égale à $\frac{5}{484}$.

\subsection{Démonstration pour $n = 485$}
Soit $n = 485$. Nous cherchons $x, y, z \in \mathbb{Z}^{+}$ tels que $\frac{5}{485} = \frac{1}{x} + \frac{1}{y} + \frac{1}{z}$.
Posons $x = 98$, $y = 9507$, $z = 90373542$.
Les conditions $x > 0$, $y > 0$ et $z > 0$ sont satisfaites car $x, y, z \geq 1$.
Le PPCM des dénominateurs est $\text{PPCM}(98, 9507, 90373542) = 90373542$.
En réduisant au même dénominateur :
\begin{itemize}
    \item $\frac{1}{98} = \frac{922179}{90373542}$
    \item $\frac{1}{9507} = \frac{9506}{90373542}$
    \item $\frac{1}{90373542} = \frac{1}{90373542}$
\end{itemize}
La somme des numérateurs est :
$$ \frac{1}{98} + \frac{1}{9507} + \frac{1}{90373542} = \frac{922179 + 9506 + 1}{90373542} = \frac{931686}{90373542} $$
Le PGCD du numérateur et du dénominateur est $\text{PGCD}(931686, 90373542) = 931686$.
La fraction irréductible est :
$$ \frac{931686}{90373542} = \frac{931686 \div 931686}{90373542 \div 931686} = \frac{1}{97} $$
Cette fraction est bien égale à $\frac{5}{485}$.

\subsection{Démonstration pour $n = 486$}
Soit $n = 486$. Nous cherchons $x, y, z \in \mathbb{Z}^{+}$ tels que $\frac{5}{486} = \frac{1}{x} + \frac{1}{y} + \frac{1}{z}$.
Posons $x = 98$, $y = 11908$, $z = 141788556$.
Les conditions $x > 0$, $y > 0$ et $z > 0$ sont satisfaites car $x, y, z \geq 1$.
Le PPCM des dénominateurs est $\text{PPCM}(98, 11908, 141788556) = 141788556$.
En réduisant au même dénominateur :
\begin{itemize}
    \item $\frac{1}{98} = \frac{1446822}{141788556}$
    \item $\frac{1}{11908} = \frac{11907}{141788556}$
    \item $\frac{1}{141788556} = \frac{1}{141788556}$
\end{itemize}
La somme des numérateurs est :
$$ \frac{1}{98} + \frac{1}{11908} + \frac{1}{141788556} = \frac{1446822 + 11907 + 1}{141788556} = \frac{1458730}{141788556} $$
Le PGCD du numérateur et du dénominateur est $\text{PGCD}(1458730, 141788556) = 291746$.
La fraction irréductible est :
$$ \frac{1458730}{141788556} = \frac{1458730 \div 291746}{141788556 \div 291746} = \frac{5}{486} $$
Cette fraction est bien égale à $\frac{5}{486}$.

\subsection{Démonstration pour $n = 487$}
Soit $n = 487$. Nous cherchons $x, y, z \in \mathbb{Z}^{+}$ tels que $\frac{5}{487} = \frac{1}{x} + \frac{1}{y} + \frac{1}{z}$.
Posons $x = 98$, $y = 15909$, $z = 759272934$.
Les conditions $x > 0$, $y > 0$ et $z > 0$ sont satisfaites car $x, y, z \geq 1$.
Le PPCM des dénominateurs est $\text{PPCM}(98, 15909, 759272934) = 759272934$.
En réduisant au même dénominateur :
\begin{itemize}
    \item $\frac{1}{98} = \frac{7747683}{759272934}$
    \item $\frac{1}{15909} = \frac{47726}{759272934}$
    \item $\frac{1}{759272934} = \frac{1}{759272934}$
\end{itemize}
La somme des numérateurs est :
$$ \frac{1}{98} + \frac{1}{15909} + \frac{1}{759272934} = \frac{7747683 + 47726 + 1}{759272934} = \frac{7795410}{759272934} $$
Le PGCD du numérateur et du dénominateur est $\text{PGCD}(7795410, 759272934) = 1559082$.
La fraction irréductible est :
$$ \frac{7795410}{759272934} = \frac{7795410 \div 1559082}{759272934 \div 1559082} = \frac{5}{487} $$
Cette fraction est bien égale à $\frac{5}{487}$.

\subsection{Démonstration pour $n = 488$}
Soit $n = 488$. Nous cherchons $x, y, z \in \mathbb{Z}^{+}$ tels que $\frac{5}{488} = \frac{1}{x} + \frac{1}{y} + \frac{1}{z}$.
Posons $x = 98$, $y = 23913$, $z = 571807656$.
Les conditions $x > 0$, $y > 0$ et $z > 0$ sont satisfaites car $x, y, z \geq 1$.
Le PPCM des dénominateurs est $\text{PPCM}(98, 23913, 571807656) = 571807656$.
En réduisant au même dénominateur :
\begin{itemize}
    \item $\frac{1}{98} = \frac{5834772}{571807656}$
    \item $\frac{1}{23913} = \frac{23912}{571807656}$
    \item $\frac{1}{571807656} = \frac{1}{571807656}$
\end{itemize}
La somme des numérateurs est :
$$ \frac{1}{98} + \frac{1}{23913} + \frac{1}{571807656} = \frac{5834772 + 23912 + 1}{571807656} = \frac{5858685}{571807656} $$
Le PGCD du numérateur et du dénominateur est $\text{PGCD}(5858685, 571807656) = 1171737$.
La fraction irréductible est :
$$ \frac{5858685}{571807656} = \frac{5858685 \div 1171737}{571807656 \div 1171737} = \frac{5}{488} $$
Cette fraction est bien égale à $\frac{5}{488}$.

\subsection{Démonstration pour $n = 489$}
Soit $n = 489$. Nous cherchons $x, y, z \in \mathbb{Z}^{+}$ tels que $\frac{5}{489} = \frac{1}{x} + \frac{1}{y} + \frac{1}{z}$.
Posons $x = 98$, $y = 47923$, $z = 2296566006$.
Les conditions $x > 0$, $y > 0$ et $z > 0$ sont satisfaites car $x, y, z \geq 1$.
Le PPCM des dénominateurs est $\text{PPCM}(98, 47923, 2296566006) = 2296566006$.
En réduisant au même dénominateur :
\begin{itemize}
    \item $\frac{1}{98} = \frac{23434347}{2296566006}$
    \item $\frac{1}{47923} = \frac{47922}{2296566006}$
    \item $\frac{1}{2296566006} = \frac{1}{2296566006}$
\end{itemize}
La somme des numérateurs est :
$$ \frac{1}{98} + \frac{1}{47923} + \frac{1}{2296566006} = \frac{23434347 + 47922 + 1}{2296566006} = \frac{23482270}{2296566006} $$
Le PGCD du numérateur et du dénominateur est $\text{PGCD}(23482270, 2296566006) = 4696454$.
La fraction irréductible est :
$$ \frac{23482270}{2296566006} = \frac{23482270 \div 4696454}{2296566006 \div 4696454} = \frac{5}{489} $$
Cette fraction est bien égale à $\frac{5}{489}$.

\subsection{Démonstration pour $n = 490$}
Soit $n = 490$. Nous cherchons $x, y, z \in \mathbb{Z}^{+}$ tels que $\frac{5}{490} = \frac{1}{x} + \frac{1}{y} + \frac{1}{z}$.
Posons $x = 99$, $y = 9703$, $z = 94138506$.
Les conditions $x > 0$, $y > 0$ et $z > 0$ sont satisfaites car $x, y, z \geq 1$.
Le PPCM des dénominateurs est $\text{PPCM}(99, 9703, 94138506) = 94138506$.
En réduisant au même dénominateur :
\begin{itemize}
    \item $\frac{1}{99} = \frac{950894}{94138506}$
    \item $\frac{1}{9703} = \frac{9702}{94138506}$
    \item $\frac{1}{94138506} = \frac{1}{94138506}$
\end{itemize}
La somme des numérateurs est :
$$ \frac{1}{99} + \frac{1}{9703} + \frac{1}{94138506} = \frac{950894 + 9702 + 1}{94138506} = \frac{960597}{94138506} $$
Le PGCD du numérateur et du dénominateur est $\text{PGCD}(960597, 94138506) = 960597$.
La fraction irréductible est :
$$ \frac{960597}{94138506} = \frac{960597 \div 960597}{94138506 \div 960597} = \frac{1}{98} $$
Cette fraction est bien égale à $\frac{5}{490}$.

\subsection{Démonstration pour $n = 491$}
Soit $n = 491$. Nous cherchons $x, y, z \in \mathbb{Z}^{+}$ tels que $\frac{5}{491} = \frac{1}{x} + \frac{1}{y} + \frac{1}{z}$.
Posons $x = 99$, $y = 12153$, $z = 196915059$.
Les conditions $x > 0$, $y > 0$ et $z > 0$ sont satisfaites car $x, y, z \geq 1$.
Le PPCM des dénominateurs est $\text{PPCM}(99, 12153, 196915059) = 196915059$.
En réduisant au même dénominateur :
\begin{itemize}
    \item $\frac{1}{99} = \frac{1989041}{196915059}$
    \item $\frac{1}{12153} = \frac{16203}{196915059}$
    \item $\frac{1}{196915059} = \frac{1}{196915059}$
\end{itemize}
La somme des numérateurs est :
$$ \frac{1}{99} + \frac{1}{12153} + \frac{1}{196915059} = \frac{1989041 + 16203 + 1}{196915059} = \frac{2005245}{196915059} $$
Le PGCD du numérateur et du dénominateur est $\text{PGCD}(2005245, 196915059) = 401049$.
La fraction irréductible est :
$$ \frac{2005245}{196915059} = \frac{2005245 \div 401049}{196915059 \div 401049} = \frac{5}{491} $$
Cette fraction est bien égale à $\frac{5}{491}$.

\subsection{Démonstration pour $n = 492$}
Soit $n = 492$. Nous cherchons $x, y, z \in \mathbb{Z}^{+}$ tels que $\frac{5}{492} = \frac{1}{x} + \frac{1}{y} + \frac{1}{z}$.
Posons $x = 99$, $y = 16237$, $z = 263623932$.
Les conditions $x > 0$, $y > 0$ et $z > 0$ sont satisfaites car $x, y, z \geq 1$.
Le PPCM des dénominateurs est $\text{PPCM}(99, 16237, 263623932) = 263623932$.
En réduisant au même dénominateur :
\begin{itemize}
    \item $\frac{1}{99} = \frac{2662868}{263623932}$
    \item $\frac{1}{16237} = \frac{16236}{263623932}$
    \item $\frac{1}{263623932} = \frac{1}{263623932}$
\end{itemize}
La somme des numérateurs est :
$$ \frac{1}{99} + \frac{1}{16237} + \frac{1}{263623932} = \frac{2662868 + 16236 + 1}{263623932} = \frac{2679105}{263623932} $$
Le PGCD du numérateur et du dénominateur est $\text{PGCD}(2679105, 263623932) = 535821$.
La fraction irréductible est :
$$ \frac{2679105}{263623932} = \frac{2679105 \div 535821}{263623932 \div 535821} = \frac{5}{492} $$
Cette fraction est bien égale à $\frac{5}{492}$.

\subsection{Démonstration pour $n = 493$}
Soit $n = 493$. Nous cherchons $x, y, z \in \mathbb{Z}^{+}$ tels que $\frac{5}{493} = \frac{1}{x} + \frac{1}{y} + \frac{1}{z}$.
Posons $x = 99$, $y = 24404$, $z = 1191086028$.
Les conditions $x > 0$, $y > 0$ et $z > 0$ sont satisfaites car $x, y, z \geq 1$.
Le PPCM des dénominateurs est $\text{PPCM}(99, 24404, 1191086028) = 1191086028$.
En réduisant au même dénominateur :
\begin{itemize}
    \item $\frac{1}{99} = \frac{12031172}{1191086028}$
    \item $\frac{1}{24404} = \frac{48807}{1191086028}$
    \item $\frac{1}{1191086028} = \frac{1}{1191086028}$
\end{itemize}
La somme des numérateurs est :
$$ \frac{1}{99} + \frac{1}{24404} + \frac{1}{1191086028} = \frac{12031172 + 48807 + 1}{1191086028} = \frac{12079980}{1191086028} $$
Le PGCD du numérateur et du dénominateur est $\text{PGCD}(12079980, 1191086028) = 2415996$.
La fraction irréductible est :
$$ \frac{12079980}{1191086028} = \frac{12079980 \div 2415996}{1191086028 \div 2415996} = \frac{5}{493} $$
Cette fraction est bien égale à $\frac{5}{493}$.

\subsection{Démonstration pour $n = 494$}
Soit $n = 494$. Nous cherchons $x, y, z \in \mathbb{Z}^{+}$ tels que $\frac{5}{494} = \frac{1}{x} + \frac{1}{y} + \frac{1}{z}$.
Posons $x = 99$, $y = 48907$, $z = 2391845742$.
Les conditions $x > 0$, $y > 0$ et $z > 0$ sont satisfaites car $x, y, z \geq 1$.
Le PPCM des dénominateurs est $\text{PPCM}(99, 48907, 2391845742) = 2391845742$.
En réduisant au même dénominateur :
\begin{itemize}
    \item $\frac{1}{99} = \frac{24160058}{2391845742}$
    \item $\frac{1}{48907} = \frac{48906}{2391845742}$
    \item $\frac{1}{2391845742} = \frac{1}{2391845742}$
\end{itemize}
La somme des numérateurs est :
$$ \frac{1}{99} + \frac{1}{48907} + \frac{1}{2391845742} = \frac{24160058 + 48906 + 1}{2391845742} = \frac{24208965}{2391845742} $$
Le PGCD du numérateur et du dénominateur est $\text{PGCD}(24208965, 2391845742) = 4841793$.
La fraction irréductible est :
$$ \frac{24208965}{2391845742} = \frac{24208965 \div 4841793}{2391845742 \div 4841793} = \frac{5}{494} $$
Cette fraction est bien égale à $\frac{5}{494}$.

\subsection{Démonstration pour $n = 495$}
Soit $n = 495$. Nous cherchons $x, y, z \in \mathbb{Z}^{+}$ tels que $\frac{5}{495} = \frac{1}{x} + \frac{1}{y} + \frac{1}{z}$.
Posons $x = 100$, $y = 9901$, $z = 98019900$.
Les conditions $x > 0$, $y > 0$ et $z > 0$ sont satisfaites car $x, y, z \geq 1$.
Le PPCM des dénominateurs est $\text{PPCM}(100, 9901, 98019900) = 98019900$.
En réduisant au même dénominateur :
\begin{itemize}
    \item $\frac{1}{100} = \frac{980199}{98019900}$
    \item $\frac{1}{9901} = \frac{9900}{98019900}$
    \item $\frac{1}{98019900} = \frac{1}{98019900}$
\end{itemize}
La somme des numérateurs est :
$$ \frac{1}{100} + \frac{1}{9901} + \frac{1}{98019900} = \frac{980199 + 9900 + 1}{98019900} = \frac{990100}{98019900} $$
Le PGCD du numérateur et du dénominateur est $\text{PGCD}(990100, 98019900) = 990100$.
La fraction irréductible est :
$$ \frac{990100}{98019900} = \frac{990100 \div 990100}{98019900 \div 990100} = \frac{1}{99} $$
Cette fraction est bien égale à $\frac{5}{495}$.

\subsection{Démonstration pour $n = 496$}
Soit $n = 496$. Nous cherchons $x, y, z \in \mathbb{Z}^{+}$ tels que $\frac{5}{496} = \frac{1}{x} + \frac{1}{y} + \frac{1}{z}$.
Posons $x = 100$, $y = 12401$, $z = 153772400$.
Les conditions $x > 0$, $y > 0$ et $z > 0$ sont satisfaites car $x, y, z \geq 1$.
Le PPCM des dénominateurs est $\text{PPCM}(100, 12401, 153772400) = 153772400$.
En réduisant au même dénominateur :
\begin{itemize}
    \item $\frac{1}{100} = \frac{1537724}{153772400}$
    \item $\frac{1}{12401} = \frac{12400}{153772400}$
    \item $\frac{1}{153772400} = \frac{1}{153772400}$
\end{itemize}
La somme des numérateurs est :
$$ \frac{1}{100} + \frac{1}{12401} + \frac{1}{153772400} = \frac{1537724 + 12400 + 1}{153772400} = \frac{1550125}{153772400} $$
Le PGCD du numérateur et du dénominateur est $\text{PGCD}(1550125, 153772400) = 310025$.
La fraction irréductible est :
$$ \frac{1550125}{153772400} = \frac{1550125 \div 310025}{153772400 \div 310025} = \frac{5}{496} $$
Cette fraction est bien égale à $\frac{5}{496}$.

\subsection{Démonstration pour $n = 497$}
Soit $n = 497$. Nous cherchons $x, y, z \in \mathbb{Z}^{+}$ tels que $\frac{5}{497} = \frac{1}{x} + \frac{1}{y} + \frac{1}{z}$.
Posons $x = 100$, $y = 16567$, $z = 823379900$.
Les conditions $x > 0$, $y > 0$ et $z > 0$ sont satisfaites car $x, y, z \geq 1$.
Le PPCM des dénominateurs est $\text{PPCM}(100, 16567, 823379900) = 823379900$.
En réduisant au même dénominateur :
\begin{itemize}
    \item $\frac{1}{100} = \frac{8233799}{823379900}$
    \item $\frac{1}{16567} = \frac{49700}{823379900}$
    \item $\frac{1}{823379900} = \frac{1}{823379900}$
\end{itemize}
La somme des numérateurs est :
$$ \frac{1}{100} + \frac{1}{16567} + \frac{1}{823379900} = \frac{8233799 + 49700 + 1}{823379900} = \frac{8283500}{823379900} $$
Le PGCD du numérateur et du dénominateur est $\text{PGCD}(8283500, 823379900) = 1656700$.
La fraction irréductible est :
$$ \frac{8283500}{823379900} = \frac{8283500 \div 1656700}{823379900 \div 1656700} = \frac{5}{497} $$
Cette fraction est bien égale à $\frac{5}{497}$.

\subsection{Démonstration pour $n = 498$}
Soit $n = 498$. Nous cherchons $x, y, z \in \mathbb{Z}^{+}$ tels que $\frac{5}{498} = \frac{1}{x} + \frac{1}{y} + \frac{1}{z}$.
Posons $x = 100$, $y = 24901$, $z = 620034900$.
Les conditions $x > 0$, $y > 0$ et $z > 0$ sont satisfaites car $x, y, z \geq 1$.
Le PPCM des dénominateurs est $\text{PPCM}(100, 24901, 620034900) = 620034900$.
En réduisant au même dénominateur :
\begin{itemize}
    \item $\frac{1}{100} = \frac{6200349}{620034900}$
    \item $\frac{1}{24901} = \frac{24900}{620034900}$
    \item $\frac{1}{620034900} = \frac{1}{620034900}$
\end{itemize}
La somme des numérateurs est :
$$ \frac{1}{100} + \frac{1}{24901} + \frac{1}{620034900} = \frac{6200349 + 24900 + 1}{620034900} = \frac{6225250}{620034900} $$
Le PGCD du numérateur et du dénominateur est $\text{PGCD}(6225250, 620034900) = 1245050$.
La fraction irréductible est :
$$ \frac{6225250}{620034900} = \frac{6225250 \div 1245050}{620034900 \div 1245050} = \frac{5}{498} $$
Cette fraction est bien égale à $\frac{5}{498}$.

\subsection{Démonstration pour $n = 499$}
Soit $n = 499$. Nous cherchons $x, y, z \in \mathbb{Z}^{+}$ tels que $\frac{5}{499} = \frac{1}{x} + \frac{1}{y} + \frac{1}{z}$.
Posons $x = 100$, $y = 49901$, $z = 2490059900$.
Les conditions $x > 0$, $y > 0$ et $z > 0$ sont satisfaites car $x, y, z \geq 1$.
Le PPCM des dénominateurs est $\text{PPCM}(100, 49901, 2490059900) = 2490059900$.
En réduisant au même dénominateur :
\begin{itemize}
    \item $\frac{1}{100} = \frac{24900599}{2490059900}$
    \item $\frac{1}{49901} = \frac{49900}{2490059900}$
    \item $\frac{1}{2490059900} = \frac{1}{2490059900}$
\end{itemize}
La somme des numérateurs est :
$$ \frac{1}{100} + \frac{1}{49901} + \frac{1}{2490059900} = \frac{24900599 + 49900 + 1}{2490059900} = \frac{24950500}{2490059900} $$
Le PGCD du numérateur et du dénominateur est $\text{PGCD}(24950500, 2490059900) = 4990100$.
La fraction irréductible est :
$$ \frac{24950500}{2490059900} = \frac{24950500 \div 4990100}{2490059900 \div 4990100} = \frac{5}{499} $$
Cette fraction est bien égale à $\frac{5}{499}$.

\subsection{Démonstration pour $n = 500$}
Soit $n = 500$. Nous cherchons $x, y, z \in \mathbb{Z}^{+}$ tels que $\frac{5}{500} = \frac{1}{x} + \frac{1}{y} + \frac{1}{z}$.
Posons $x = 101$, $y = 10101$, $z = 102020100$.
Les conditions $x > 0$, $y > 0$ et $z > 0$ sont satisfaites car $x, y, z \geq 1$.
Le PPCM des dénominateurs est $\text{PPCM}(101, 10101, 102020100) = 102020100$.
En réduisant au même dénominateur :
\begin{itemize}
    \item $\frac{1}{101} = \frac{1010100}{102020100}$
    \item $\frac{1}{10101} = \frac{10100}{102020100}$
    \item $\frac{1}{102020100} = \frac{1}{102020100}$
\end{itemize}
La somme des numérateurs est :
$$ \frac{1}{101} + \frac{1}{10101} + \frac{1}{102020100} = \frac{1010100 + 10100 + 1}{102020100} = \frac{1020201}{102020100} $$
Le PGCD du numérateur et du dénominateur est $\text{PGCD}(1020201, 102020100) = 1020201$.
La fraction irréductible est :
$$ \frac{1020201}{102020100} = \frac{1020201 \div 1020201}{102020100 \div 1020201} = \frac{1}{100} $$
Cette fraction est bien égale à $\frac{5}{500}$.

\end{document}
